\documentclass[UTF8,11pt,twoside,openany]{ctexbook}
\usepackage[paperwidth=176mm,paperheight=250mm,inner=21mm,outer=18mm,top=20mm,bottom=20mm,headheight=16pt,headsep=6mm,footskip=10mm]{geometry}
\usepackage{amsmath,amssymb,amsthm,bm,mathtools,booktabs,array,enumitem,float}
\usepackage{tikz,fancyhdr,hyperref}
\usetikzlibrary{arrows.meta,calc,patterns,angles,quotes,decorations.pathmorphing}
\hypersetup{hidelinks,pdftitle={光学进阶：物理竞赛二轮讲义},pdfauthor={}}
\setlength{\parindent}{2em}\setlength{\parskip}{3pt}\setlength{\emergencystretch}{2em}
\linespread{1.18}\raggedbottom\let\cleardoublepage\clearpage\setcounter{tocdepth}{1}
\setlength{\intextsep}{8pt plus 1pt minus 1pt}
\ctexset{chapter={name={第,讲},number=\chinese{chapter},format=\LARGE\bfseries\centering,beforeskip=16pt,afterskip=22pt},section={format=\Large\bfseries},subsection={format=\large\bfseries}}
\pagestyle{fancy}\fancyhf{}\fancyhead[LE,RO]{\small\thepage}\fancyhead[LO]{\small 光学进阶}\fancyhead[RE]{\small\leftmark}
\renewcommand{\headrulewidth}{0pt}\fancypagestyle{plain}{\fancyhf{}\fancyfoot[C]{\small\thepage}\renewcommand{\headrulewidth}{0pt}}
\theoremstyle{definition}\newtheorem{example}{例}[chapter]\newtheorem{exercise}{习题}[chapter]
\newenvironment{solution}{\par\medskip\noindent\textbf{解}\quad\ignorespaces}{\par\medskip}
\numberwithin{equation}{chapter}
\newcommand{\dd}{\mathop{}\!\mathrm d}\newcommand{\ee}{\mathrm e}\newcommand{\kb}{k_{\mathrm B}}\newcommand{\vct}[1]{\bm{#1}}
\tikzset{axis/.style={->,>=Stealth,line width=.45pt},curve/.style={line width=.7pt},guide/.style={dashed,line width=.35pt},every node/.style={font=\small},particle/.style={circle,fill,inner sep=1.3pt}}
\setlist[enumerate]{itemsep=3pt,topsep=4pt,leftmargin=2em}
\let\LectureContents\tableofcontents\renewcommand{\tableofcontents}{\begingroup\small\linespread{1}\selectfont\LectureContents\endgroup}
\newcommand{\Diagram}[3]{\begin{figure}[H]\centering\begin{tikzpicture}#3\end{tikzpicture}\caption{#2}\label{#1}\end{figure}}
\newcommand{\OptAxis}{\draw[axis](-.5,0)--(7,0)node[right]{$z$};}
\providecommand{\ExpansionDiagram}[3]{\begin{figure}[H]\centering\begin{tikzpicture}#3\end{tikzpicture}\caption{#2}\label{#1}\end{figure}}


% 分册与十编合订本共用的定位页码；局部章、图、题号保持原有对应。
\newcommand{\CoursePartNumber}{10}
\geometry{inner=19.5mm,outer=19.5mm,headheight=16pt}
\AtBeginDocument{%
  \fancyhead{}\fancyhead[C]{\small 光学进阶}%
  \fancyfoot[C]{\small\CoursePartNumber-\thepage}%
  \fancypagestyle{plain}{\fancyhf{}\fancyfoot[C]{\small\CoursePartNumber-\thepage}\renewcommand{\headrulewidth}{0pt}}%
}

\begin{document}
\frontmatter
\begin{titlepage}\centering\vspace*{28mm}
{\Large 物理竞赛二轮讲义\quad 第十编\par}\vspace{8mm}
{\Huge\bfseries 光学进阶\par}\vspace{12mm}{\Large 物理竞赛二轮讲义\par}\vfill\end{titlepage}
\chapter*{编写说明}
成像问题从光线和参考面列式，干涉、衍射问题从光程和复振幅列式，偏振问题则要同时保留分量方向与相对相位。本编按这条线索安排四讲，先讨论近轴系统和渐变介质，再推导相干叠加、孔径积分及界面电磁波。

每个模型先说明介质、边界和近似条件，再求光线、场或能流。例题和习题的题设图与解答图相互接续，所用符号与计算保持一致。相位约定、角度方向和强度归一化在各讲中统一，便于比较不同装置的结果。
\tableofcontents
\mainmatter
\chapter{近轴成像、渐变介质与光程}
\section{符号与单球面光焦度}\label{th:o1power}
取光沿$z$正向传播，高度$y$向上，近轴倾角$\theta\simeq\dd y/\dd z$有正负。球面顶点为坐标原点，曲率中心在$+z$侧时$R>0$，否则$R<0$。物方、像方折射率为$n_1,n_2$，定义球面光焦度
\[
 \Phi=\frac{n_2-n_1}{R}.
\]
球面高度$y$处法线倾角近似$-y/R$，斯涅尔定律线性化给$n_1(\theta_1+y/R)=n_2(\theta_2+y/R)$，所以
\begin{equation}\label{eq:o1surface}
 n_2\theta_2=n_1\theta_1-\Phi y.
\end{equation}
若另把朝物方的几何半径记为$r>0$，该面$R=-r$，同式就写为$\Phi=(n_1-n_2)/r$。只要曲率和角度约定一致，两种写法给出同一折射方向。

取实物距$u>0$为物点到顶点沿传播方向的距离，实像距$v>0$为顶点到像点的距离，虚物、虚像用负值。轴上物点发出的近轴光有$\theta_1\simeq y/u$，会聚到轴上像点需$\theta_2\simeq-y/v$。代入式\eqref{eq:o1surface}得
\begin{equation}\label{eq:o1gauss}
 \frac{n_1}u+\frac{n_2}v=\Phi.
\end{equation}
前焦距$f_a=n_1/\Phi$，后焦距$f_b=n_2/\Phi$。它们一般不等，二者符号都由$\Phi$确定。将$n_1=\Phi f_a$、$n_2=\Phi f_b$代入成像式并乘$uv/\Phi$，得$f_av+f_bu=uv$，即$(u-f_a)(v-f_b)=f_af_b$。令$x_u=u-f_a$、$x_v=v-f_b$，就有牛顿式$x_ux_v=f_af_b$。这里$x_u,x_v$是相对焦点的偏移，已不是顶点物像距。

\Diagram{fig:o1surface}{球面折射的近轴光线和有向曲率}{\OptAxis\draw[curve](3,-1.6)arc(205:155:3.785);\fill(0,0)circle(1.3pt)node[below]{物点};\fill(6.4,0)circle(1.3pt)node[below]{$C$};\draw[guide](2.7,1)--(6.4,0);\draw[->](0,0)--(2.7,1);\draw[->](2.7,1)--(5.5,.7);\node at(1.2,1.6){$n_1$};\node at(4.8,1.6){$n_2$};}

\begin{example}\label{ex:o1surface}
折射率$n_1=1$的空气中有轴上点光源，位于顶点左方$u=30\,\mathrm{cm}$。右侧玻璃折射率$n_2=1.5$，界面曲率中心在右侧，$R=10\,\mathrm{cm}$。求像距、焦距，并求物距改为10 cm时的像。
\Diagram{fig:o1surfaceq}{空气射入凸球面，曲率中心在玻璃侧}{\OptAxis\draw[curve](3,-1.4)..controls(2.3,-.5)and(2.3,.5)..(3,1.4);\fill(0,0)circle(1.4pt)node[below]{$S$};\fill(4,0)circle(1.4pt)node[below]{$C$};\node at(1.2,1.2){空气};\node at(4.9,1.2){玻璃};}
\end{example}
\begin{solution}
采用图\ref{fig:o1surfaceq}的有向半径：曲率中心在右，故$R=+10\,\mathrm{cm}$。球面光焦度及前后主焦距大小为
\[
 \Phi=\frac{n_2-n_1}R=.05\,\mathrm{cm}^{-1},\quad
 f_a=\frac{n_1}\Phi=20\,\mathrm{cm},\quad
 f_b=\frac{n_2}\Phi=30\,\mathrm{cm}.
\]
其中前焦点在顶点左侧20 cm，后焦点在右侧30 cm，两侧折射率不同，两个焦距大小一般不等。物距为30 cm时，
\[
 \frac{n_2}v=\Phi-\frac{n_1}u
 =\frac1{20}-\frac1{30}=\frac1{60}\,\mathrm{cm}^{-1},
\]
所以$v=90\,\mathrm{cm}>0$，折射光在玻璃中会聚。
\Diagram{fig:o1surfacea}{物点移至前焦点内，折射后仍发散，反向延长得到虚像}{\OptAxis\draw[curve](3,-1.5)..controls(2.4,-.5)and(2.4,.5)..(3,1.5);\fill(1.5,0)circle(1.4pt)node[below]{$S$};\draw[->](1.5,0)--(2.55,.8);\draw[->](2.55,.8)--(5.8,1.65);\draw[guide](2.55,.8)--(-.5,0);\node[below]at(-.5,0){虚像方向};}
物距改为10 cm时，同一成像方程给$1.5/v=.05-.1=-.05\,\mathrm{cm}^{-1}$，故$v=-30\,\mathrm{cm}$。图\ref{fig:o1surfacea}中的虚像由折射光反向延长而成，其坐标在物侧，但整束折射后光线仍属于$n_2$一侧的传播解，成像方程仍使用像方折射率$n_2$。

可用任意近轴高度$y$核对：入射倾角近似为$\theta_1=y/u$，出射约化倾角为$n_2\theta_2=n_1y/u-\Phi y=-n_2y/v$。对实像$v>0$，出射角与$y$反号，向轴靠拢；对虚像$v<0$，出射角与$y$同号，向外发散。这个符号判断直接连接图中的光线走向与成像方程。
\end{solution}

\section{约化光线矩阵与复合系统}\label{th:o1matrix}
用列向量$\vct q=(y,n\theta)^{\mathsf T}$表示近轴光线。厚度$d$、折射率$n$的均匀层只改变高度；球面只改变约化倾角，所以
\[
 T(d,n)=\begin{pmatrix}1&d/n\\0&1\end{pmatrix},\quad
 S(\Phi)=\begin{pmatrix}1&0\\-\Phi&1\end{pmatrix}.
\]
光先经历的矩阵写在最右侧。各矩阵行列式均为1。两折射面相隔$d$、中间折射率$n_3$时，
\[
 M=S(\Phi_2)T(d,n_3)S(\Phi_1)
 =\begin{pmatrix}1-d\Phi_1/n_3&d/n_3\\-\Phi&1-d\Phi_2/n_3\end{pmatrix},
\]
其中$\Phi=\Phi_1+\Phi_2-d\Phi_1\Phi_2/n_3$。厚度项来自第一面改变倾角，进而改变第二面入射高度；简单相加光焦度只在$d=0$或该项可忽略时成立。

\subsection{主平面的位置}
一般近轴系统写成$M=\bigl(\begin{smallmatrix}A&B\\C&D\end{smallmatrix}\bigr)$。若$C\ne0$，等效光焦度为$\Phi=-C$。设第一主面在第一顶点右方$h_1$，第二主面在最后顶点右方$h_2$，则
\[
 M=T(-h_2,n_2)S(\Phi)T(h_1,n_1).
\]
比较$A,D$得$h_1=n_1(1-D)/\Phi$、$h_2=-n_2(1-A)/\Phi$。主面物像距满足$n_1/u_H+n_2/v_H=\Phi$，而不是从镜筒边缘量距离。

若$C=0$，系统无焦，不能用无穷大焦距加上发散的主面位置来计算实际光线。此时应直接使用有限的$A,B,D$矩阵，研究角放大率、束径与光束传播。无焦系统仍由束高、倾角之间的有限矩阵关系确定。

\begin{example}\label{ex:o1pair}
空气中两个薄透镜焦距为$f_1,f_2$，间隔$d=f_1+f_2+\Delta$。求等效焦距、两主面位置及其有向间距；讨论$\Delta=0$。
\Diagram{fig:o1pairq}{两薄透镜间隔$d$及各自焦点的间隔$\Delta$}{\OptAxis\draw[<->](1,-1.3)--(1,1.3);\draw[<->](5,-1.3)--(5,1.3);\node[above]at(1,1.3){$f_1$};\node[above]at(5,1.3){$f_2$};\draw[<->](1,-1.8)--(5,-1.8)node[midway,below]{$d$};\fill(2.7,0)circle(1.3pt)node[above]{$F'_1$};\fill(3.8,0)circle(1.3pt)node[above]{$F_2$};}
\end{example}
\begin{solution}
将图\ref{fig:o1pairq}的两透镜作为$\Phi_i=1/f_i$代入，
\[
 \Phi=\frac{f_1+f_2-d}{f_1f_2}=-\frac\Delta{f_1f_2},\quad f=-\frac{f_1f_2}\Delta.
\]
若$\Delta\ne0$，$A=1-d/f_1$、$D=1-d/f_2$，故从第一透镜位置$z=0$计，$z_{H_1}=fd/f_2$、$z_{H_2}=d-fd/f_1$。两主面的有向间距为
\begin{align*}
 z_{H_2}-z_{H_1}
 &=d-fd(1/f_1+1/f_2)\\
 &=d\left[1-\frac{f_1+f_2}{f_1+f_2-d}\right]
 =\frac{d^2}\Delta.
\end{align*}
\Diagram{fig:o1paira}{主面可能在实体透镜外侧；示例取$f_1=f_2=1,d=3$}{\draw[axis](-3.7,0)--(6.7,0)node[right]{$z$};\foreach\x/\t in{0/L_1,3/L_2}{\draw[<->](\x,-.8)--(\x,.8);\node[above]at(\x,.8){$\t$};}\foreach\x/\t in{-3/H_1,6/H_2}{\draw[guide](\x,-1.2)--(\x,1.2);\node[above]at(\x,1.2){$\t$};}}
图\ref{fig:o1paira}取$\Delta=1$，于是$f=-1$、$H_1$在$-3$、$H_2$在6，主面间隔9。主面是等效光学基准，允许在系统外甚至交叉。若$\Delta=0$，$\Phi=0$，矩阵仍有限，形成无焦系统；主面公式分母为零，已失去适用条件。
\end{solution}

\begin{example}\label{ex:o1ball}
空气中有半径$r$、折射率$n>1$的玻璃球。求等效焦距、主面及背焦距。轴上物点位于球心左方$nr$，求其像的位置并讨论$n=3/2$两侧的情况。
\Diagram{fig:o1ballq}{玻璃球的两个顶点与球心}{\OptAxis\draw[curve](3,0)circle(1.3);\fill(3,0)circle(1.2pt)node[below]{$O$};\fill(.2,0)circle(1.4pt)node[below]{$S$};\node[above]at(3,1.3){$n$};\node[below]at(1.7,0){$V_1$};\node[below]at(4.3,0){$V_2$};}
\end{example}
\begin{solution}
图\ref{fig:o1ballq}的第一面$R_1=r$、$n_1=1,n_2=n$，第二面$R_2=-r$、$n_1=n,n_2=1$，两面光焦度都是$(n-1)/r$。厚度$2r$、中间折射率$n$，所以
\[
 \Phi=\frac{2(n-1)}r-\frac{2r}n\frac{(n-1)^2}{r^2}
 =\frac{2(n-1)}{nr},\quad f=\frac{nr}{2(n-1)}.
\]
$A=D=(2-n)/n$，主面位置$h_1=r,h_2=-r$，两主面均经过球心。背焦距从后顶点起算为$f-r=r(2-n)/[2(n-1)]$。
\Diagram{fig:o1balla}{从球心主面量物像距，图示$3/2<n<3$时的球外实像}{\OptAxis\draw[curve](2.8,0)circle(1);\draw[guide](2.8,-1.4)--(2.8,1.5);\fill(0,0)circle(1.4pt)node[below]{$S$};\fill(6.1,0)circle(1.4pt)node[below]{$S'$};\draw[<->](0,-1.8)--(2.8,-1.8)node[midway,below]{$nr$};\draw[<->](2.8,-1.8)--(6.1,-1.8)node[midway,below]{$nr/(2n-3)$};}
图\ref{fig:o1balla}的主面物距$u_H=nr$。由$1/u_H+1/v_H=1/f$，得$v_H=nr/(2n-3)$。主面位于球心，判断出射束的实虚像还须改从后顶点量距：
\[
 v_{V_2}=v_H-r=\frac{(3-n)r}{2n-3}.
\]
当$3/2<n<3$，$v_{V_2}>0$，折射光在球外会聚成实像。$n=3/2$时出射平行、像在无穷远；$1<n<3/2$时像在球心左侧，为虚像。$n=3$时近轴像落在后顶点。若$n>3$，$0<v_H<r$，像虽在球心右侧，却位于后顶点之前，出射束反向延长相交于球内，仍为虚像。主面像距的正负不能单独代替出射束的会聚、发散判断。
\end{solution}

\section{费马原理与力学类比}\label{th:o1analogy}
在静态、各向同性介质中，几何光线使光程$\int n\dd s$取驻值。变分得到
\begin{equation}\label{eq:o1ray}
 \frac{\dd}{\dd s}(n\vct t)=\nabla n,
\end{equation}
其中$\vct t=\dd\vct r/\dd s$是单位切向量。对分层界面切向积分，由平移对称方向的“广义动量”连续，得$n\sin\theta$守恒，正是折射定律。

为了求轨迹，可以引入虚拟力学粒子，势能$U=-mv_0^2n^2/2$、总能量为零。它的虚拟速度为$|\dot{\vct r}|=v_0n$，运动方程$\ddot{\vct r}=v_0^2\nabla(n^2)/2$。零能量轨道的缩减作用量$\int p\dd s=mv_0\int n\dd s$，所以轨迹与光线相同。但虚拟速度与真实光速$c/n$成反比，不能把虚拟时间当传播时间。若虚拟参数为$q$，
\[
 \dd t_{\rm phys}=\frac{n\dd s}{c}=
 \frac{v_0n^2}{c}\dd q.
\]
对于水波等其他波，$n$可定义为参考波速除以局部相速度；真实时间应按当地群速度计算，浅水非色散模型中群速度等于相速度。

\begin{example}\label{ex:o1graded}
介质折射率满足$n^2(y)=n_0^2(1-y^2/a^2)$，只讨论$|y|\ll a$的近轴光线。一束光在$z=0$处高度$y_0$、倾角$\theta_0$。求轨迹和周期；若改为加号，轨迹有什么变化？
\Diagram{fig:o1gradedq}{中心折射率较大，光线弯向轴线}{\OptAxis\draw[guide](0,1.3)--(6,1.3);\draw[guide](0,-1.3)--(6,-1.3);\fill(0,.8)circle(1.4pt)node[left]{$y_0$};\draw[axis](0,.8)--(1,1.1)node[above]{$\theta_0$};\node[above]at(3,1.3){$|y|$增大时$n$减小};}
\end{example}
\begin{solution}
图\ref{fig:o1gradedq}的虚拟参数取$q$，其方程为$\ddot z=0$、$\ddot y=-v_0^2n_0^2y/a^2$。轴向速度为常数$\dot z=v_0n_i\cos\theta_0$，其中$n_i=n(y_0)$。链式法则给$\ddot y=(\dot z)^2y''(z)$，所以有限初角的精确射线方程为
\[
 y''+g^2y=0,\qquad g=\frac{n_0}{a n_i\cos\theta_0}.
\]
方程的通解为$y=A\cos gz+B\sin gz$，其导数为$y'=-gA\sin gz+gB\cos gz$。初值$y(0)=y_0$、$y'(0)=\tan\theta_0$给$A=y_0$、$B=\tan\theta_0/g$。取$|y_0|/a\ll1$、$|\theta_0|\ll1$，有$g\simeq1/a$、$\tan\theta_0\simeq\theta_0$，才得到题设近轴解
\[
 y=y_0\cos(z/a)+a\theta_0\sin(z/a),\qquad\Lambda_z=2\pi a.
\]
\Diagram{fig:o1gradeda}{负二次项产生周期聚焦，正二次项产生离轴增长}{\draw[axis](0,0)--(6.7,0)node[right]{$z/a$};\draw[axis](0,-1.5)--(0,2.5)node[above]{$y$};\draw[curve,domain=0:6.28,samples=100]plot(\x,{cos(\x r)});\draw[guide,domain=0:1.5,samples=40]plot(\x,{cosh(\x)});\node[right]at(1.6,2.3){加号};}
图\ref{fig:o1gradeda}把折射率平方改为加号。$y$方向虚拟加速度随之变号，近轴方程为$y''-y/a^2=0$。特征根为$\pm1/a$，通解为$C\ee^{z/a}+D\ee^{-z/a}$。初值给$C+D=y_0$、$(C-D)/a=\theta_0$，故$C=(y_0+a\theta_0)/2$、$D=(y_0-a\theta_0)/2$，合并为
\[
 y(z)=y_0\cosh(z/a)+a\theta_0\sinh(z/a).
\]
负号分布中折射率由轴向外减小，射线受回复作用；加号分布则把光线推向轴外。对负号分布，近轴振幅为$\sqrt{y_0^2+a^2\theta_0^2}$，须始终远小于$a$，因此不仅要检查初始高度，还要约束初始倾角引起的振幅。对加号分布，即使初始条件近轴，振幅也可能随传播增长而失效；应在实际截面和近轴允许长度内使用双曲解。
\end{solution}

\subsection{径向介质的角动量}\label{th:o1central}
若$n=n(r)$，虚拟力为中心力，$\vct r\times m\dot{\vct r}$守恒。除以$mv_0$得到
\[
 b=n(r)r\sin\psi,
\]
其中$\psi$是光线切向与径向的夹角，选取$\varphi$递增的绕行方向使$b>0$。弧长参数沿传播方向增加，有
\[
 \frac{\dd r}{\dd s}=\cos\psi
 =\sigma\sqrt{1-\frac{b^2}{n^2r^2}},\qquad
 r\frac{\dd\varphi}{\dd s}=\sin\psi=\frac b{nr},
\]
其中向外支$\sigma=+1$、向内支$\sigma=-1$。在$\dd r/\dd s\ne0$处，两式相除得到
\[
 \frac{\dd\varphi}{\dd r}
 =\sigma\frac{b}{r\sqrt{n^2(r)r^2-b^2}}.
\]
若有孤立最近点，则该点$\cos\psi=0$，故$b=n(r_t)r_t$；转点处不能继续用$r$作单调参数，应接续向内、向外两支。这一等式也可能对应圆轨道等退化情形，须代回射线方程判断。$b=0$为径向光线，直接使用$\dd\varphi/\dd s=0$。若根号内为负，则该区域没有此几何射线；波的穿透须由波动方程处理。

\begin{example}\label{ex:o1wake}
浅水区水深$h(y)=4h_0/(5-y^2/y_0^2)$，岸线$y=0$，船沿$y=y_0$向$+x$以$V_b=2\sqrt{gh_0}$运动。船在$t=0$经过$A=(0,y_0)$。由该处向岸传播的稳态船行波射线与岸线法向夹角$\theta_0$满足$\sin\theta_0=c_A/V_b$，$c_A=\sqrt{gh_0}$。求射线首先到岸的位置、传播时间及届时船的位置。
\Diagram{fig:o1wakeq}{船道、水深变化与向岸初始射线}{\draw[axis](-.3,0)--(6,0)node[right]{$x$};\draw[axis](0,-.2)--(0,3.2)node[above]{$y$};\draw[guide](0,2)--(5.5,2);\fill(0,2)circle(1.6pt)node[left]{$A$};\draw[axis](.2,2.4)--(1.7,2.4)node[above]{$V_b$};\draw[axis](0,2)--(.9,.45);\node[left]at(0,1){$y_0$};\node[below]at(3,0){岸线};}
\end{example}
\begin{solution}
局部浅水波速$c_w(y)=\sqrt{gh(y)}$。取相对折射率$n(y)=c_A/c_w(y)=\tfrac12\sqrt{5-y^2/y_0^2}$，在$A$处$n_A=1$。图\ref{fig:o1wakeq}初始$\sin\theta_0=1/2$，所以$n\sin\theta=1/2$守恒，射线向$+x,-y$运动。

虚拟运动满足$\ddot x=0$、$\ddot y=-v_0^2y/(4y_0^2)$，且$\dot x=v_0/2$。消去虚拟时间得$y''(x)=-y/y_0^2$。正弦型通解为$A\cos(x/y_0)+B\sin(x/y_0)$；由$y(0)=y_0$得$A=y_0$，由$y'(0)=B/y_0=-\cot\theta_0=-\sqrt3$得$B=-\sqrt3y_0$，故
\[
 y(x)=y_0\cos(x/y_0)-\sqrt3y_0\sin(x/y_0)
 =2y_0\cos(x/y_0+\pi/3).
\]
首个$y=0$在$x_B=\pi y_0/6$。
\Diagram{fig:o1wakea}{射线先到岸，船继续前进；横轴按$y_0$归一化}{\draw[axis](0,0)--(6.2,0)node[right]{$x/y_0$};\draw[axis](0,-.1)--(0,2.7)node[above]{$y/y_0$};\draw[guide](0,2)--(5.4,2);\draw[curve,domain=0:.5235988,samples=50]plot({2*\x},{4*cos((\x+pi/3) r)});\fill(1.0472,0)circle(1.5pt)node[below]{$B:\pi/6$};\fill(4.8736,2)circle(1.5pt)node[above]{$A':(\pi+\sqrt3)/2$};\node[left]at(0,2){$1$};}
图\ref{fig:o1wakea}的真实传播时间不可用虚拟周期直接代替。由$\dd s/\dd x=1/\sin\theta=2n$、$c_w=c_A/n$，得
\begin{align*}
 T&=\int_0^{x_B}\frac{\dd s}{c_w}
 =\frac2{c_A}\int_0^{x_B}n^2(x)\dd x\\
 &=\frac1{2c_A}\int_0^{\pi y_0/6}
 [5-4\cos^2(x/y_0+\pi/3)]\dd x\\
 &=\frac1{2c_A}\left[3x-y_0\sin(2x/y_0+2\pi/3)\right]_0^{\pi y_0/6}\\
 &=\frac1{2c_A}\left(\frac{\pi y_0}{2}+\frac{\sqrt3y_0}{2}\right)
 =\frac{y_0}{4c_A}(\pi+\sqrt3).
\end{align*}
最后一步用$4\cos^2u=2+2\cos2u$逐项积分。船位置$x_{A'}=V_bT=y_0(\pi+\sqrt3)/2$。图为清楚显示船位置，横纵采用同一缩放因子2；$B$的横坐标与船到岸时的横坐标不能互换。
\end{solution}

\section{复振幅与光程相位}\label{th:o1phase}
固定时间约定为$\ee^{-i\omega t}$。实电场是$\operatorname{Re}[\mathcal E(\vct r)\ee^{-i\omega t}]$，复振幅$\mathcal E=A\ee^{i\phi}$同时存储幅度和空间相位。同频同偏振光相加时，先加复振幅再求模平方：$I\propto|\sum_j\mathcal E_j|^2$。不同频率或不相关随机相位的交叉项，经足够长时间平均后可以消失。

几何光学中沿射线传播产生相位$k_0\int n\dd s$，$k_0=2\pi/\lambda_0$。若两光束起始相位相同，传播部分相位差为$2\pi\Delta\mathcal L/\lambda_0$；反射、焦散或其他元件的相位须另加。薄片正常插入一束空气光路，以厚$h$的折射率$n$介质替换同厚空气，附加光程$(n-1)h$。斜入射时路径和出射点都改变，不能继续把物理厚度直接代入正常入射公式。

\begin{example}\label{ex:o1plate}
两等振幅相干光原本在某点相长，一束中垂直插入折射率$n$、厚度$h$的平行薄片。求该点光强及使其第一次变暗的最小正厚度。忽略反射损失。
\Diagram{fig:o1plateq}{一束光中以薄片替换同厚空气}{\draw[->](0,1)--(5,1);\draw[->](0,0)--(5,0);\draw[curve](2,.6)rectangle(2.7,1.4);\node[above]at(2.35,1.4){$n,h$};\node[right]at(5,.5){合束};}
\end{example}
\begin{solution}
图\ref{fig:o1plateq}中两复振幅可取$A$、$A\ee^{i\delta}$，$\delta=2\pi(n-1)h/\lambda_0$。合振幅$2A\ee^{i\delta/2}\cos(\delta/2)$，故$I=4I_1\cos^2[\pi(n-1)h/\lambda_0]$。
\Diagram{fig:o1platea}{两相量夹角为$\delta$，合相量沿角平分线}{\draw[axis](0,0)--(3,0)node[right]{$A$};\draw[axis](0,0)--(1.5,2.598076)node[above left]{$A\ee^{i\delta}$};\draw[very thick,->](0,0)--(4.5,2.598076)node[right]{$\mathcal E$};\draw[guide](3,0)--(4.5,2.598076)--(1.5,2.598076);\node at(.9,.4){$\delta$};}
图\ref{fig:o1platea}的两相量首次反向时$\delta=\pi$，所以$h_{\min}=\lambda_0/[2(n-1)]$，$n>1$。若两束振幅不等，反向时仍有残余强度，不能称完全变暗。若频率不相同，静态薄片也不能把随时间移动的拍频图样变成永久暗点。
\end{solution}
\section{费马原理的变分推导}
\subsection{驻值不必是最短路程}
静态各向同性介质中的光线使光程$\mathcal L=\int_A^B n(\vct r)\dd s$在固定端点的路径变分下取驻值。高折射率区的单位路程耗时较长，所以几何短路未必是光程短路。费马原理只要求一阶变分为零，允许局部最小、鞍值及退化驻值。

用任意单调参数$q$表示路径，拉格朗日量为$\mathcal G=n(\vct r)|\dot{\vct r}|$。变分导数
\[
 \frac{\partial\mathcal G}{\partial\dot{\vct r}}=n\vct t,
 \qquad \frac{\partial\mathcal G}{\partial\vct r}=|\dot{\vct r}|\nabla n.
\]
欧拉—拉格朗日方程除以$\dd s/\dd q=|\dot{\vct r}|$，得到$\dd(n\vct t)/\dd s=\nabla n$。把左侧展开为$(\dd n/\dd s)\vct t+n\dd\vct t/\dd s$，再以$\dd n/\dd s=\vct t\cdot\nabla n$，得曲率方程
\[
 n\frac{\dd\vct t}{\dd s}
 =\nabla n-(\vct t\cdot\nabla n)\vct t=\nabla_\perp n.
\]
只有折射率梯度的横向部分使射线转弯；沿射线方向的变化只改变相位积累和局部速度。弯曲方向朝横向折射率较大处，这与分层界面逐层折射的极限一致。

\subsection{界面的切向条件}
在两均匀介质间界面选交点$Q$，端点$A,B$固定。光程为$n_1|A-Q|+n_2|Q-B|$。让$Q$沿界面切向移动$\delta\vct q$，一阶变化为$(n_1\vct t_1-n_2\vct t_2)\cdot\delta\vct q$。任意切向变化都取驻值，要求切向分量$n_1\vct t_{1\parallel}=n_2\vct t_{2\parallel}$，于是$n_1\sin i=n_2\sin r$。

\Diagram{fig:o1fermat}{界面交点的切向变分给出折射定律}{\draw[curve](-3,0)--(3,0);\fill(-2,1.6)circle(1.5pt)node[left]{$A$};\fill(2,-1.4)circle(1.5pt)node[right]{$B$};\fill(0,0)circle(1.5pt)node[above]{$Q$};\draw[curve](-2,1.6)--(0,0)--(2,-1.4);\draw[guide](-2,1.6)--(.6,0)--(2,-1.4);\draw[->](0,.3)--(.6,.3)node[midway,above]{$\delta q$};}
图\ref{fig:o1fermat}的驻值路径在真实界面法向上不要求约化动量连续，因为折射率在该处跳变。光沿边界的相位必须连续，给出相同切向条件。反射则可由镜像法或在同一介质中改变法向传播方向得到，入射角等于反射角。

\subsection{全反射和几何模型的边界}
从$n_1>n_2$射向低折射率介质时，切向条件可能要求$\sin r>1$，此时没有实传播方向。几何光学把它描述为全反射，但并未说明低折射率侧的场严格为零。波动边界解会出现虚法向波矢、倏逝场及偏振相关反射相位。这些内容在第四讲推导，说明同一切向相位约束可以继续使用，而实射线概念已有适用边界。

\section{主平面如何连接实际光线}
主平面是在把系统矩阵分解为“物侧平移—等效折射—像侧平移”时选出的两个参考面。在相同物像介质中，一条假想入射光从实际入口沿直线延长到第一主面，再在等效薄系统中改变倾角，最后从第二主面延长回实际出口；这个构造保持系统外光线关系，不声称系统内部光线按该折线传播。

把矩阵乘法完整写出，设$h_1$从第一顶点向右、$h_2$从最后顶点向右，
\[
 T(-h_2,n_2)S(\Phi)T(h_1,n_1)
 =\begin{pmatrix}
 1+h_2\Phi/n_2&h_1/n_1-h_2/n_2+h_1h_2\Phi/(n_1n_2)\\
 -\Phi&1-h_1\Phi/n_1
 \end{pmatrix}.
\]
由$C=-\Phi$先定总光焦度，再由$A,D$定两个主面，$B$因行列式条件自动一致。若$C=0$，这套分解失去有限主面解，但原$M$依然完整；无焦望远系统、扩束器应直接用矩阵分析。

\subsection{焦距和顶点焦距的差别}
系统外平行光与后焦点的关系，从出射参考面的$y_o=Ay_i$、$n_2\theta_o=Cy_i$可得后焦点距该面的距离$v_F=-n_2A/C=n_2A/\Phi$。另一方面，主面后焦距为$n_2/\Phi$，第二主面相对出口位置为$h_2=n_2(A-1)/\Phi$，相加正好也是$n_2A/\Phi$。这证明主面公式与实际外光线相交位置一致。

玻璃球的两主面在球心，主面后焦距为$f=nr/[2(n-1)]$，从后表面量起的背焦距则为$f-r$。当$n=2$，平行近轴光的焦点在后顶点；当$n>2$，背焦距为负，总光焦度$\Phi$仍为正。实虚像应由后表面外光线的会聚或发散判断，并注明距离的参考面。

\section{渐变介质力学类比的参数选择}
由力学零能量条件$|\dot{\vct r}|=v_0n$，切向单位矢量$\vct t=\dot{\vct r}/(v_0n)$。对$n\vct t=\dot{\vct r}/v_0$按弧长求导，利用$\dd/\dd s=(v_0n)^{-1}\dd/\dd q$，得到
\[
 \frac{\dd(n\vct t)}{\dd s}
 =\frac{\ddot{\vct r}}{v_0^2n}.
\]
令它等于$\nabla n$，便得$\ddot{\vct r}=v_0^2n\nabla n=\tfrac12v_0^2\nabla(n^2)$。这一等式证明零能量力学轨迹满足费马射线方程。任意正值$v_0$只改变虚拟参数的快慢，不改变空间轨迹。

真实传播时间中的权重由$\dd t=n\dd s/c$和$\dd s=v_0n\dd q$共同给出，即$\dd t=v_0n^2\dd q/c$。船行波模型的当地波速为$c_A/n$，相同计算须把分母$c$换成$c_A$。有色散时，波包到达时间应按群速度计算，相位轨迹未必就是能量传播轨迹。

\subsection{对称性给出的积分常数}
若$n$不依$z$，虚拟势具有轴向平移对称性，$\dot z$及$n\cos\alpha$守恒，其中$\alpha$相对$z$轴量取。若$n$只依$r$，转动对称性给$nr\sin\psi$守恒，其中$\psi$相对径向量取。使用守恒量时，须同时保留角度的几何定义。

\Diagram{fig:o1symmetries}{平移对称守恒轴向分量，转动对称守恒角动量}{\begin{scope}[xshift=-2.6cm]\draw[axis](0,0)--(2.5,0)node[right]{$z$};\draw[curve](0,.2)..controls(1,1.3)and(1.6,1.4)..(2.2,1);\draw[guide](0,1.7)--(2.3,1.7);\node[below]at(1.2,-.5){$n\cos\alpha$};\end{scope}\begin{scope}[xshift=1.5cm]\draw[curve](0,0)circle(1);\draw[->](0,0)--(1.6,.8);\draw[->](1.6,.8)--(1.9,1.8);\node[below]at(.5,-1.5){$nr\sin\psi$};\end{scope}}
图\ref{fig:o1symmetries}中的守恒量限定方向和可达区域，继而用于积分求轨迹。若下一界面破坏相应对称性，应按界面相位匹配确定新区域的积分常数。

\subsection{同一共轭关系的横向与纵向变化}
物距和像距必须从各自主面量取。设物、像侧折射率分别为$n_O,n_I$，系统光焦度为$\Phi$，则$n_O/u+n_I/v=\Phi$。在系统不变时微分，
\[
 -\frac{n_O}{u^2}\dd u-\frac{n_I}{v^2}\dd v=0,
 \qquad \frac{\dd v}{\dd u}=-\frac{n_Ov^2}{n_Iu^2}.
\]
物点轴向坐标取$z_O=-u$、像点取$z_I=v$，所以同向轴坐标的纵向放大率为$\dd z_I/\dd z_O=n_Ov^2/(n_Iu^2)>0$。由横向放大率$m_y=-n_Ov/(n_Iu)$，得
\[
 m_z=\frac{\dd z_I}{\dd z_O}=\frac{n_I}{n_O}m_y^2.
\]
在相同外部介质中，轴向微小长度的放大率是横向放大率的平方。横向倒像的负号不使纵向关系变负；物点向右微移，实像在相应共轭分支上也向右移。
\Diagram{fig:o1longmagnify}{两个相近轴上物点与其共轭像点；箭头表示微小轴向位移}{\draw[axis](-3.3,0)--(3.7,0)node[right]{$z$};\draw[thick](0,-1.2)--(0,1.2)node[above]{系统};\fill(-2.8,0)circle(1.5pt)node[below]{$O$};\fill(-2.3,0)circle(1.5pt)node[above]{$O+\dd z_O$};\fill(2,0)circle(1.5pt)node[below]{$I$};\fill(3,0)circle(1.5pt)node[above]{$I+\dd z_I$};\draw[->](-2.8,-.6)--(-2.3,-.6);\draw[->](2,-.6)--(3,-.6);}
图\ref{fig:o1longmagnify}只表示微分共轭，不把有限厚度物体的所有截面当作具有同一个放大率。对有限景深，应分别求前后物面的像距；当物面跨越焦面时，实虚像分支还会改变。矩阵成像式能直接处理这些情况，而只报一个等效焦距不足以定位物像。

当$|m_y|$较大时，物面的小定位误差会被放大为像面位移。几何离焦产生的弥散直径还依赖像方孔径角；衍射斑则由有限孔径决定。用纵向放大率估算焦面位置时，应保持近轴、小位移和同一成像分支三个条件。

\section{成像矩阵与横向放大率}
设物面到像面的总矩阵为$M_{OI}=\bigl(\begin{smallmatrix}A&B\\C&D\end{smallmatrix}\bigr)$。同一物点高度$y_O$发出不同倾角的光，像面高度为$y_I=Ay_O+B n_O\theta_O$。所有出射方向都汇于同一像点，要求$B=0$，横向放大率就是$A$。因为行列式为1，此时$AD=1$，角放大与横向放大有互为倒数的约化关系。

对物像面两条光线，双线性量$y_1n\theta_2-y_2n\theta_1$在每个行列式为1的矩阵下不变，这就是近轴光学不变量。它约束像高与孔径角的乘积：同一系统不可能既把像无代价放大，又保持全部出射角不变。有限孔径、亮度和通光范围需在这个几何约束之外再结合能流研究。

\begin{example}\label{ex:o1matriximage}
近轴系统两端都在空气，总矩阵为$\bigl(\begin{smallmatrix}A&B\\C&D\end{smallmatrix}\bigr)$。物面在入射参考面前$u$，像面在出射参考面后$v$。推导成像条件及放大率，说明它比只给焦距多保留了什么信息。
\Diagram{fig:o1matrixq}{物、像距离分别从系统两参考面量起}{\OptAxis\draw[curve](1.8,-1)rectangle(3.8,1);\node at(2.8,.5){$M$};\draw[->](0,0)--(0,.8);\draw[guide](5.8,-1.2)--(5.8,1.2);\node[below]at(.9,-.5){$u$};\node[below]at(4.8,-.5){$v$};}
\end{example}
\begin{solution}
图\ref{fig:o1matrixq}的光线变量取$(y,p)$，空气中$p=\theta$。由物面先向右传播$u$，再经过系统$M$，最后传播$v$，因此从右向左依次相乘：
\[
 T(v)MT(u)=
 \begin{pmatrix}
 A+vC&Au+B+v(Cu+D)\\ C&Cu+D
 \end{pmatrix}.
\]
同一物点发出的光线具有相同$y_O$、不同$p_O$。要使它们到像面有同一高度，输出高度中$p_O$的系数必须为零，即$Au+B+v(Cu+D)=0$。若$Cu+D\ne0$，
\[
 v=-\frac{Au+B}{Cu+D}.
\]
\Diagram{fig:o1matrixa}{同一物点的两条不同初始倾角光线须落在同一像点}{\OptAxis\draw[curve](2,-1.2)rectangle(3,1.2);\draw[->](0,.8)--(2,1);\draw[->](0,.8)--(2,-.2);\draw[->](3,.5)--(5.5,-.6);\draw[->](3,-.8)--(5.5,-.6);\fill(5.5,-.6)circle(1.4pt);}
图\ref{fig:o1matrixa}中两条入射光线在像面汇于同点，横向放大率为矩阵左上元：
\begin{align*}
 m_y&=A+vC
 =\frac{A(Cu+D)-C(Au+B)}{Cu+D}\\
 &=\frac{AD-BC}{Cu+D}=\frac1{Cu+D}.
\end{align*}
最后使用约化光线矩阵的$\det M=1$。若$Cu+D=0$，行列式又保证$Au+B$不能同时为零，否则矩阵两列线性相关；此时出射为平行束，有限$v$处没有像面。$C=0$的无焦系统仍可用完整矩阵求共轭，无须引入发散的主面坐标。$C$只决定总光焦度，$A,D$还保留焦点相对系统外壳的位置，$B$则影响有限物距的共轭关系。
\end{solution}

\begin{example}\label{ex:o1thicklens}
空气中厚透镜折射率$n$，轴向厚度$d$，前后曲率半径$R_1,R_2$均按本讲有向约定。推导厚透镜制作者公式，并求对称双凸$R_1=R,R_2=-R$的焦距与主面偏移。
\Diagram{fig:o1thickq}{前后曲率中心的方向不同，$R_1>0,R_2<0$}{\OptAxis\draw[curve](2,-1.5)..controls(1.4,-.5)and(1.4,.5)..(2,1.5);\draw[curve](2,-1.5)..controls(3,-.5)and(3,.5)..(2,1.5);\node[above]at(2.2,1.7){$n,d$};}
\end{example}
\begin{solution}
图\ref{fig:o1thickq}中两界面光焦度为$\Phi_1=(n-1)/R_1$、$\Phi_2=(1-n)/R_2$。折射矩阵$S_j=\bigl(\begin{smallmatrix}1&0\\-\Phi_j&1\end{smallmatrix}\bigr)$，玻璃内传播矩阵$T=\bigl(\begin{smallmatrix}1&d/n\\0&1\end{smallmatrix}\bigr)$，故
\[
 S_2TS_1=\begin{pmatrix}
 1-d\Phi_1/n&d/n\\
 -\Phi_1-\Phi_2+d\Phi_1\Phi_2/n&1-d\Phi_2/n
 \end{pmatrix}.
\]
由$C=-\Phi$读出厚透镜制作者公式
\[
 \Phi=(n-1)\left[\frac1{R_1}-\frac1{R_2}
 +\frac{(n-1)d}{nR_1R_2}\right].
\]
\Diagram{fig:o1thicka}{主面向透镜内部移动，两个偏移相反}{\OptAxis\draw[guide](2,-1.5)--(2,1.5);\draw[guide](3,-1.5)--(3,1.5);\draw[very thick](2.2,-1.1)--(2.2,1.1);\draw[very thick](2.8,-1.1)--(2.8,1.1);\node[above]at(2,1.6){$V_1$};\node[above]at(3,1.6){$V_2$};\node[below]at(2.2,-1.1){$H_1$};\node[below]at(2.8,-1.1){$H_2$};}
图\ref{fig:o1thicka}采用$R_1=R$、$R_2=-R$，所以$\Phi_1=\Phi_2=\Phi_s=(n-1)/R$，$\Phi=2\Phi_s-d\Phi_s^2/n$、$f=1/\Phi$。以两个顶点分别为基准的主面有向偏移为
\[
 h_1=\frac{1-D}\Phi=\frac{d\Phi_s}{n\Phi},\qquad
 h_2=\frac{A-1}\Phi=-\frac{d\Phi_s}{n\Phi}.
\]
常规会聚情形$\Phi>0$，前主面向右偏、后主面向左偏；背焦距是从后顶点量到后焦点，故为$f+h_2$，而非$f$。令$d\to0$，两主面均回到薄透镜平面，$\Phi\to2(n-1)/R$。若参数使$\Phi=0$，总矩阵仍可使用，但这些除以$\Phi$的焦距、主面公式不适用，应按无焦系统处理。
\end{solution}

\begin{example}\label{ex:o1apparentdepth}
水中点物体距平界面深$h$，水折射率$n$，空气中观察。求近轴表观深度，并用球面成像的$R\to\infty$极限和射线几何两种方法核对。
\Diagram{fig:o1depthq}{折射光的反向延长给出水中的虚像}{\draw[curve](-2.5,0)--(2.5,0);\fill(0,-2)circle(1.5pt)node[below]{物点};\draw[->](0,-2)--(1,0)--(2.3,1.4);\draw[guide](0,-2)--(0,1.7);\node[left]at(-2.5,-.8){$n$};\node[left]at(-2.5,.8){空气};}
\end{example}
\begin{solution}
沿向上方向取光轴，图\ref{fig:o1depthq}物方$n_1=n$、像方$n_2=1$，平面$\Phi=0$。式\eqref{eq:o1gauss}给$n/h+1/v=0$，所以$v=-h/n$，表观深度为$h/n$。
\Diagram{fig:o1deptha}{近轴时折射角变大，延长线交于较浅位置}{\draw[curve](-2.5,0)--(2.5,0);\fill(0,-2)circle(1.5pt)node[left]{$h$};\fill(0,-1.2)circle(1.5pt)node[left]{$h/n$};\draw[->](0,-2)--(1,0);\draw[->](1,0)--(2.5,1.8);\draw[guide](1,0)--(0,-1.2);}
射线在水中角$i$、空气角$r$，$n\sin i=\sin r$。若界面交点距轴$y$，真实深度$h=y/\tan i$，延长交点深度$h_a=y/\tan r$。近轴$r\simeq ni$，所以$h_a\simeq h/n$。图\ref{fig:o1deptha}画的是近轴虚像；大角度时$h_a=h\tan i/\tan r$依角度变，平面折射并不对所有孔径形成严格同一点像。
\end{solution}

\section{渐变介质的传像与精确光线}
\begin{example}\label{ex:o1grinrelay}
近轴渐变介质满足$y''+g^2y=0$，轴上折射率近似$n_0$。求长度$l$的约化矩阵，分别讨论四分之一节距和半节距的用途。
\Diagram{fig:o1grinq}{渐变杆内的周期光线}{\OptAxis\draw[curve](0,-1.2)rectangle(6,1.2);\draw[curve,domain=0:6,samples=80]plot(\x,{.8*cos(60*\x)});\node[above]at(3,1.2){$y''+g^2y=0$};}
\end{example}
\begin{solution}
由初值解$y(l)=y(0)\cos gl+y'(0)\sin gl/g$、$y'(l)=-gy(0)\sin gl+y'(0)\cos gl$，图\ref{fig:o1grinq}的约化矩阵为
\[
 M(l)=\begin{pmatrix}\cos gl&\sin gl/(n_0g)\\-n_0g\sin gl&\cos gl\end{pmatrix}.
\]
四分之一节距$gl=\pi/2$，输出高度由输入倾角决定、输出倾角由输入高度决定；轴上点物发出的光变为轴向平行光。半节距$gl=\pi$，$M=-I$，输入面被倒立等大地传到输出面。
\Diagram{fig:o1grina}{半节距中，来自同一输入点的不同初始倾角光线在反高处重聚}{\draw[axis](0,0)--(6.6,0);\draw[curve,domain=0:6,samples=80]plot(\x,{.7*cos(30*\x)});\draw[guide,domain=0:6,samples=80]plot(\x,{.7*cos(30*\x)+.6*sin(30*\x)});\fill(0,.7)circle(1.4pt);\fill(6,-.7)circle(1.4pt);\node[below]at(3,-1.1){$gl=\pi$};}
图\ref{fig:o1grina}的两光线初始斜率不同，末高却同为$-y_0$，体现$B=0$成像条件。四分之一节距不能同时又说输出面是输入面的等大像；它交换了位置和倾角的信息。若杆端与空气折射，还需接上两个平面界面的角度变换，本约化变量已自动保持$n\theta$连续。
\end{solution}

\begin{example}\label{ex:o1linearn2}
折射率平方在某区域满足$n^2(y)=n_A^2+2ay$，$a$为常数。光从原点沿$+z$水平进入，求虚拟力学轨迹和真实传播时间，限定在折射率为实的区间。
\Diagram{fig:o1linearq}{折射率平方的梯度使光线弯向高折射率侧}{\OptAxis\draw[axis](0,0)--(0,2.7)node[above]{$y$};\draw[->](0,0)--(1.7,0);\node at(3,1.4){$n^2=n_A^2+2ay$};}
\end{example}
\begin{solution}
图\ref{fig:o1linearq}中，虚拟参数$q$满足$\ddot{\vct r}=(v_0^2/2)\nabla n^2$，故$\ddot z=0$、$\ddot y=av_0^2$。光从原点水平进入，$z(0)=y(0)=0$、$\dot z(0)=n_Av_0$、$\dot y(0)=0$。先积分加速度，$\dot y=\int_0^qav_0^2\dd s=av_0^2q$、$\dot z=n_Av_0$；再积分速度，
\[
 z=n_Av_0q,\qquad y=\tfrac12av_0^2q^2
 =\frac{az^2}{2n_A^2}.
\]
此处的速度是对虚拟参数求导所得；真实时间仍由光程积分给出。沿轨迹$y'=az/n_A^2$，且$n^2=n_A^2+a^2z^2/n_A^2=n_A^2(1+y'^2)$。
\Diagram{fig:o1lineara}{轨迹为抛物线，但真实传播时间不是匀加速粒子时间}{\OptAxis\draw[axis](0,0)--(0,2.9)node[above]{$y$};\draw[curve,domain=0:4.5,samples=70]plot(\x,{.12*\x*\x});\draw[guide](3,0)--(3,1.08);\node[below]at(3,0){$Z$};}
真实传播时间为$\dd t=n\dd s/c$，图\ref{fig:o1lineara}中的弧长满足$\dd s=\sqrt{1+y'^2}\dd z=(n/n_A)\dd z$，因此
\begin{align*}
 t(Z)&=\frac1{cn_A}\int_0^Z n^2(z)\dd z\\
 &=\frac1{cn_A}\left[n_A^2z+\frac{a^2z^3}{3n_A^2}\right]_0^Z
 =\frac{n_AZ}c+\frac{a^2Z^3}{3cn_A^3}.
\end{align*}
几何路程的增加和沿途折射率的增加都已包含在$n\dd s$中。若$a<0$，光线向负$y$偏折，但$a y>0$，沿射线仍是折射率增大；时间修正含$a^2$，保持正值。$a$的量纲为长度倒数，$a^2Z^3/c$为时间。以上轨迹只适用于该分布覆盖的区域，遇到界面后须重新匹配出射方向。
\end{solution}

\begin{example}\label{ex:o1luneburg}
半径$R$的球形渐变透镜在球内$n^2(r)=2-r^2/R^2$，球外$n=1$。证明平行于$x$、冲量参数满足$|b|<R$的入射光线会聚于球面右端$(R,0)$，并求从左侧平面$x=-R$到焦点的总光程。
\Diagram{fig:o1luneq}{球边$n=1$，平行光以不同冲量参数$b$进入}{\draw[axis](-3,0)--(3,0)node[right]{$x$};\draw[axis](0,-2.2)--(0,2.4)node[above]{$y$};\draw[curve](0,0)circle(1.8);\foreach\y in{-.9,0,.9}{\draw[->](-2.8,\y)--({-sqrt(3.24-\y*\y)},\y);}\node[right]at(2,1.4){$n^2=2-r^2/R^2$};}
\end{example}
\begin{solution}
取图\ref{fig:o1luneq}二维截面，入点$\vct r_0=(-a,b)$，$a=\sqrt{R^2-b^2}$。球边折射率连续且等于1，入射方向不跳变。虚拟力为$\ddot{\vct r}=-(v_0^2/R^2)\vct r$，初始速度$v_0\hat{\vct x}$。令$\eta=v_0q/R$，精确解
\[
 x=-a\cos\eta+R\sin\eta,\qquad y=b\cos\eta.
\]
在$0<\eta<\pi/2$，$r^2=R^2-2aR\sin\eta\cos\eta<R^2$，光线确在球内；首个再次到球面时$\eta=\pi/2$，得到$x=R,y=0$，与$b$无关。
\Diagram{fig:o1lunea}{精确光线在球面同一点相交；这里无需近轴条件}{\draw[axis](-2.8,0)--(2.6,0);\draw[curve](0,0)circle(1.8);\foreach\b in{-.9,0,.9}{\pgfmathsetmacro{\a}{sqrt(3.24-\b*\b)}\draw[curve,domain=0:90,samples=70]plot({-\a*cos(\x)+1.8*sin(\x)},{\b*cos(\x)});}\fill(1.8,0)circle(1.7pt)node[right]{$F$};}
沿轨迹$n^2=1+(a/R)\sin2\eta$，真实光程为
\[
 ct=R\int_0^{\pi/2}\left(1+\frac aR\sin2\eta\right)\dd\eta
 =R\left[\eta-\frac{a}{2R}\cos2\eta\right]_0^{\pi/2}
 =\frac{\pi R}{2}+a.
\]
从左平面到入点的空气程为$R-a$，相加得总光程$R(1+\pi/2)$，不依冲量参数。图\ref{fig:o1lunea}的几何聚焦同时满足等光程。严格掠射的$|b|=R$不属于穿入球内的光线，以上结论可取其从内侧逼近的极限。
\end{solution}

\section{光程、相位与短波近似}
设标量场满足$\nabla^2U+k_0^2n^2U=0$，写$U=A\ee^{ik_0S}$，$A,S$缓变且$S$为实。第一次求导给$\nabla U=\ee^{ik_0S}(\nabla A+ik_0A\nabla S)$；再用乘积法则求散度，
\[
 \nabla^2U=\ee^{ik_0S}\left[
 \nabla^2A+2ik_0\nabla A\cdot\nabla S
 +ik_0A\nabla^2S-k_0^2A|\nabla S|^2\right].
\]
代回波动方程并除去指数，得$\nabla^2A+ik_0(2\nabla A\cdot\nabla S+A\nabla^2S)+k_0^2A(n^2-|\nabla S|^2)=0$。短波渐近展开中，先令$k_0^2$阶系数为零，得到程函方程$|\nabla S|^2=n^2$；再令$k_0$阶系数为零，得到输运方程$2\nabla A\cdot\nabla S+A\nabla^2S=0$。后式乘$A$，利用$\nabla(A^2)=2A\nabla A$，可写成$\nabla\cdot(A^2\nabla S)=0$。第一式决定相位面的法线方向和光程，第二式决定光线束收缩时的振幅变化；$\nabla^2A$留在更低阶，焦区等振幅急变区域不能将它忽略。

\begin{example}\label{ex:o1eikonal}
在均匀介质中，分别求平面波和球面出射波的程函$S$及振幅随传播距离的变化，解释焦散附近为何几何光学振幅会失效。
\Diagram{fig:o1eikonalq}{平行光管截面积不变，球面光管面积随$r^2$增大}{\begin{scope}[xshift=-2cm]\draw[->](0,-.6)--(2,-.6);\draw[->](0,.6)--(2,.6);\node[above]at(1,1){平行光管};\end{scope}\draw[->](2,0)--(4.5,1.2);\draw[->](2,0)--(4.5,-1.2);\node[above]at(3.5,1.7){发散光管};}
\end{example}
\begin{solution}
图\ref{fig:o1eikonalq}的平面波可取$S=nz$，$\nabla^2S=0$。输运方程给$\partial_z A=0$，积分为$A(z)-A(z_0)=\int_{z_0}^z0\dd q=0$，故振幅不随$z$衰减。球面波取$S=nr$，径向拉普拉斯为$\nabla^2S=r^{-2}\partial_r(r^2\partial_rS)=r^{-2}\partial_r(nr^2)=2n/r$。输运方程变为$A'=-A/r$，在$r>0$分离变量，
\[
 \int_{A(r_0)}^{A(r)}\frac{\dd a}{a}
 =-\int_{r_0}^{r}\frac{\dd q}{q},\qquad
 \ln\frac{A(r)}{A(r_0)}=-\ln\frac r{r_0}.
\]
因此$A(r)=A(r_0)r_0/r$，强度随$1/r^2$衰减。
\Diagram{fig:o1eikonala}{会聚光管在焦点截面积趋零，短波振幅公式不能无限使用}{\draw[->,thick](-2,1.2)--(0,0)--(2,-1.2);\draw[->,thick](-2,-1.2)--(0,0)--(2,1.2);\draw[guide](-1,-.6)--(-1,.6);\fill(0,0)circle(1.5pt)node[above]{$F$};\node[below]at(0,-1.7){焦区需波动解};}
球面能流乘球面积保持不变。图\ref{fig:o1eikonala}会聚时光管面积趋零，几何输运式形式上给出发散振幅，表明缓变假设在焦区失效。焦区应改用衍射解，得到有限场强及焦点相移。
\end{solution}

\providecommand{\ExpansionDiagram}[3]{\begin{figure}[H]\centering\begin{tikzpicture}#3\end{tikzpicture}\caption{#2}\label{#1}\end{figure}}
\section{有限孔径、光瞳与可成像的视场}
\subsection{成像条件并不保证光线能够通过}
矩阵给出光线的交点，镜框、光阑和透镜边缘则限制哪些光线能够到达该处。求出像面后，还需逐面计算光线高度，与各面的通光口径比较。

取物面光线坐标为$(y,p)^{\mathsf T}$，其中$p=n_O\theta$。从物面到第$j$个孔径平面的矩阵第一行为$(A_j,B_j)$，该孔径半宽为$a_j$，则通过条件是
\[
 |A_jy+B_jp|\leq a_j.
\]
给定物点高度$y$，每个孔径在$p$轴上限定一个区间，通过全部孔径的光线方向由这些区间的交集确定。$B_j=0$时，该孔径只限制物高；$B_j\ne0$时，它限制该物点可用的光线方向。交集缩小会造成渐晕，交集为空则该物点没有光到达像面。采用同一套有向光线坐标时，这一判定也适用于虚物和无焦系统。

对轴上物点$y=0$，最先限制入射角的孔径是孔径光阑。它通过前方光学元件所成的像称为入瞳，通过后方元件所成的像称为出瞳。光瞳可以是虚像，不能把入瞳一律画在真实光阑处。主光线通过孔径光阑中心；边缘光线通过光阑边缘。前者便于追踪视场，后者便于判断数值孔径。

\ExpansionDiagram{fig:dep-o1-interval}{固定物高时，不同孔径对入射方向的限制形成区间交集}{
\draw[->](0,0)--(8,0)node[right]{$p$};
\draw[very thick](1,.7)--(5,.7);\draw(1,.55)--(1,.85);\draw(5,.55)--(5,.85);\node[left]at(.9,.7){孔径1};
\draw[very thick](3,1.5)--(7,1.5);\draw(3,1.35)--(3,1.65);\draw(7,1.35)--(7,1.65);\node[left]at(2.9,1.5){孔径2};
\fill[gray!20](3,-.25)rectangle(5,.25);\node[below]at(4,-.35){通过区间};
\draw[dashed](3,-.25)--(3,1.6);\draw[dashed](5,-.25)--(5,1.6);}
图\ref{fig:dep-o1-interval}中的阴影是两个允许区间的交集。物高变化时，区间通常发生相对平移，所以边缘视场的有效通光量可能低于轴上通光量，即使两个物点都严格成像。

\subsection{光学不变量与通光量的几何限制}
设同一系统中的两条近轴光线为$(y_1,p_1)$、$(y_2,p_2)$。约化矩阵的行列式为1，故
\[
 H=y_1p_2-y_2p_1
\]
在近轴传播和折射中保持不变。直接代入$y_j'=Ay_j+Bp_j$、$p_j'=Cy_j+Dp_j$，展开两乘积，$ACy_1y_2$及$BDp_1p_2$项相消，剩下$H'=(AD-BC)(y_1p_2-y_2p_1)=H$。它是二维光线相空间中的有向面积。选一条从轴上物点出发的边缘光线和一条从轴外物点出发的主光线，在物面得到$|H|=n_O|y_O\theta_O|$，在像面则为$n_I|y_I\theta_I|$。因此
\[
 n_O|y_O\theta_O|=n_I|y_I\theta_I|.
\]
放大像高以后，像方角孔径相应减小。孔径遮挡能删去一部分光线，却不能在保留全部入射光线的同时任意压缩这块相空间面积。这里使用近轴角度；大角度系统应以正弦形式及相应光学条件重新讨论。

\begin{example}两孔径引起的渐晕
空气中薄透镜焦距$f>0$、半口径$a$，物面在透镜前$2f$。透镜后$f$处有半宽同为$a$的光阑。求像面位置、轴上角孔径，以及物高$h$变化时可通过的方向区间和无光视场边界。
\ExpansionDiagram{fig:dep-o1-aperture-q}{透镜与后置光阑分别限制光线高度}{
\draw[->](-3.5,0)--(4,0)node[right]{$z$};
\draw[<->,thick](0,-1)--(0,1);\node[above]at(0,1.2){透镜};\draw[<->](-.45,0)--(-.45,1)node[midway,left]{$a$};
\draw[very thick](1.5,1)--(1.5,1.7);\draw[very thick](1.5,-1)--(1.5,-1.7);\node[above]at(1.5,1.8){光阑};
\draw[->](-3,0)--(-3,.8)node[above]{$h$};\draw[dashed](3,-1.4)--(3,1.4)node[above]{像面};
\draw(-3,-.5)--(0,-.5)node[midway,below]{$2f$};\draw(0,-1.9)--(1.5,-1.9)node[midway,below]{$f$};}
\end{example}
\begin{solution}
由$1/(2f)+1/v=1/f$得$v=2f$，横向放大率为$-1$。图\ref{fig:dep-o1-aperture-q}中像面在光阑之后，光阑本身不改变成像位置。

设物点发射角为$\theta$。到达透镜时高度$y_L=h+2f\theta$，折射后方向为
\[
 \theta_+=\theta-\frac{y_L}{f}=-\frac hf-\theta.
\]
再传播$f$到光阑，高度变为$y_S=y_L+f\theta_+=f\theta$。两个通过条件分别给出
\[
 \frac{-a-h}{2f}\leq\theta\leq\frac{a-h}{2f},
 \qquad -\frac af\leq\theta\leq\frac af.
\]
轴上第一段较窄，所以透镜边缘是孔径光阑，轴上角半宽$a/(2f)$。
\ExpansionDiagram{fig:dep-o1-aperture-a}{方向区间随物高平移；临界时交集收缩为一点}{
\draw[->](0,0)--(7,0)node[right]{$f\theta/a$};
\draw[very thick](1,.5)--(5,.5);\node[left]at(.9,.5){后光阑};
\draw[very thick](2,1.2)--(4,1.2);\node[right]at(4.1,1.2){$h=0$};
\draw[very thick](0,1.9)--(2,1.9);\node[right]at(2.1,1.9){$h=2a$};
\draw[very thick](-1,2.6)--(1,2.6);\node[right]at(1.1,2.6){$h=3a$};
\node[below]at(1,0){$-1$};\node[below]at(3,0){$0$};\node[below]at(5,0){$1$};\draw[dashed](1,.2)--(1,2.8);}
在图\ref{fig:dep-o1-aperture-a}中，透镜允许区间的中心为$-h/(2f)$，半宽$a/(2f)$；后光阑的区间固定。由两个端点比较可得方向区间长度
\[
 \Delta\theta=\begin{cases}
 a/f,& |h|\leq a,\\
 (3a-|h|)/(2f),&a<|h|<3a,\\
 0,&|h|\geq3a.
 \end{cases}
\]
$|h|=3a$时仅剩一条几何边界光线，其连续光束通量为零；$|h|>3a$时完全遮挡。这里求的是一个子午截面的方向宽度，不能直接把它当作三维圆孔系统的光通量比例；三维情况还需积分允许方向的面积。
\end{solution}

\begin{example}像面偏离与弥散斑
仍取焦距$f$、半口径$a$的薄透镜，轴上物点在$2f$处。屏从理想像面沿$+z$移开$\delta$。求几何弥散斑半径；若物点还发生小轴向位移$\Delta z_O$，求一阶合并后的离焦量。忽略衍射，要求$|\delta|,|\Delta z_O|\ll f$。
\ExpansionDiagram{fig:dep-o1-defocus-q}{像点两侧的截面都可能形成弥散斑}{
\draw[->](-.5,0)--(6,0)node[right]{$z$};\draw[<->,thick](0,-1.3)--(0,1.3);
\draw(0,1.2)--(4,0)--(5,-.3);\draw(0,-1.2)--(4,0)--(5,.3);
\fill(4,0)circle(1.4pt)node[below]{$I$};\draw[thick](5,-1)--(5,1)node[above]{屏};\draw[<->](4,.65)--(5,.65)node[midway,above]{$\delta$};}
\end{example}
\begin{solution}
图\ref{fig:dep-o1-defocus-q}中任一到达透镜高度$y_L$的光线，在透镜后为$y(z)=y_L(1-z/(2f))$。屏在$z=2f+\delta$，故$y=-y_L\delta/(2f)$。取$|y_L|\leq a$即得斑半径
\[
 r_g=\frac{a|\delta|}{2f}.
\]
物点坐标向右移动$\Delta z_O$意味着物距由$2f$变为$2f-\Delta z_O$。对$1/u+1/v=1/f$求微分，$-\dd u/u^2-\dd v/v^2=0$，故$\dd v/\dd u=-v^2/u^2$。在$u=v=2f$代入得$\dd v=-\dd u$，所以像点向右移动$\Delta z_I=\Delta z_O$。屏相对新像面的离焦量为$\delta-\Delta z_O$。
\ExpansionDiagram{fig:dep-o1-defocus-a}{物面移动与屏面移动共同决定实际离焦}{
\draw[->](0,0)--(7,0)node[right]{$z$};\fill(2,0)circle(1.4pt)node[below]{$I$};\fill(3,0)circle(1.4pt)node[below]{$I'$};
\draw[thick](5,-1)--(5,1)node[above]{屏};\draw[<->](2,.4)--(3,.4)node[midway,above]{$\Delta z_O$};\draw[<->](3,-.7)--(5,-.7)node[midway,below]{$\delta-\Delta z_O$};}
因此一阶结果为$r_g\simeq a|\delta-\Delta z_O|/(2f)$。图\ref{fig:dep-o1-defocus-a}中两种位移同向且相等时，一阶离焦消失。若需要更高精度，应使用新的像距与新的像方角度；若斑半径已接近衍射尺度，则几何斑不能单独作为实际点像宽度。
\end{solution}

\section{矩阵成像与高斯光束的连接}
\subsection{高斯积分的归一化与带相位形式}
光束传播与相干平均都会用到
\[
 \mathcal G(a,b)=\int_{-\infty}^{\infty}\exp(-ax^2+bx)\dd x,
 \qquad \operatorname{Re}a>0.
\]
先取实数$a>0,b=0$。将积分平方写成二维积分，再取极坐标，
\[
 \mathcal G(a,0)^2=\int_0^{2\pi}\dd\varphi\int_0^\infty
 r\ee^{-ar^2}\dd r
 =2\pi\left[-\frac{\ee^{-ar^2}}{2a}\right]_0^\infty=\frac\pi a.
\]
原积分为正，故$\mathcal G(a,0)=\sqrt{\pi/a}$。对$b$求导并分部积分，
\begin{align*}
 \partial_b\mathcal G
 &=\int_{-\infty}^{\infty}x\ee^{-ax^2+bx}\dd x\\
 &=-\frac1{2a}[\ee^{-ax^2+bx}]_{-\infty}^{\infty}
 +\frac b{2a}\int_{-\infty}^{\infty}\ee^{-ax^2+bx}\dd x
 =\frac b{2a}\mathcal G.
\end{align*}
边界项由$\operatorname{Re}a>0$保证为零。乘积分因子$\exp[-b^2/(4a)]$，有$\partial_b\{\ee^{-b^2/(4a)}\mathcal G\}=0$，以$b=0$的值确定常数，得到
\begin{equation}\label{eq:o1gaussintegral}
 \mathcal G(a,b)=\sqrt{\frac\pi a}\exp\frac{b^2}{4a}.
\end{equation}
对复数$a$取在正实轴上为正的连续平方根。积分与右侧在$\operatorname{Re}a>0$均解析，上式由实正$a$连续延拓到该半平面；它适用于后面含二次相位的收敛积分。

\subsection{光束宽度与波前曲率需要同时传播}
束腰附近的有限光斑会发生衍射扩展，需要同时计算场的振幅与相位。对沿$+z$传播、时间因子为$\ee^{-\mathrm i\omega t}$的标量场，写作$E=U(x,y,z)\ee^{\mathrm ikz}$。在$|\partial_z^2U|\ll|k\partial_zU|$的近轴条件下，Helmholtz方程化为
\[
 2\mathrm ik\partial_zU+\nabla_\perp^2U=0.
\]
均匀介质中试取轴对称解$U=Cq^{-1}\exp[\mathrm ikr^2/(2q)]$。对$z,r$分别求导，
\begin{align*}
 \frac{\partial_zU}{U}&=-\frac{q'}q-\frac{\mathrm ikr^2q'}{2q^2},\\
 \frac{\partial_rU}{U}&=\frac{\mathrm ikr}q,\qquad
 \frac{\partial_r^2U}{U}=\frac{\mathrm ik}q-\frac{k^2r^2}{q^2},\\
 \frac{\nabla_\perp^2U}{U}&=\frac{2\mathrm ik}q-\frac{k^2r^2}{q^2}.
\end{align*}
代入近轴方程，剩余$(1-q')[2\mathrm ik/q-k^2r^2/q^2]=0$，要求$q'=1$。以束腰处$q(z_w)=-\mathrm iz_R$为初值，$q(z)-q(z_w)=\int_{z_w}^z\dd s=z-z_w$，故
\[
 q=z-z_w-\mathrm iz_R,\qquad z_R>0,
 \quad\frac1q=\frac1{R_w}+\mathrm i\frac2{kw^2}.
\]
这里$z_w$为束腰位置，$w$是场振幅降至轴上$1/\ee$的半径，也即强度降至$1/\ee^2$的半径；$R_w$为波前有向曲率半径。这个$q$的虚部为负，与所选传播相位一致，不能与另一套相位约定的$q=z+\mathrm iz_R$混用。

由复数倒数直接比较实虚部，得到
\[
 w^2(z)=w_0^2\left[1+\frac{(z-z_w)^2}{z_R^2}\right],
 \quad z_R=\frac{kw_0^2}{2}=\frac{\pi n w_0^2}{\lambda_0},
\]
\[
 R_w=(z-z_w)+\frac{z_R^2}{z-z_w}.
\]
束腰处$R_w$无限大，只表示该处波前近乎平面，并不意味着光束以后不发散。远处的发散半角为$w_0/z_R=\lambda_0/(\pi n w_0)$。
\ExpansionDiagram{fig:dep-o1-gaussian}{束腰附近的光束边界；虚线为远区渐近线}{
\draw[->](-3.4,0)--(3.5,0)node[right]{$z-z_w$};\draw[thick,domain=-3:3,samples=90]plot(\x,{.55*sqrt(1+(\x)^2)});\draw[thick,domain=-3:3,samples=90]plot(\x,{-.55*sqrt(1+(\x)^2)});\draw[dashed](-3,1.65)--(3,-1.65);\draw[dashed](-3,-1.65)--(3,1.65);\draw[<->](0,-.55)--(0,.55);\node[right,fill=white,inner sep=1pt]at(.12,.23){$2w_0$};}
图\ref{fig:dep-o1-gaussian}的上下曲线分别为$\pm w(z)$，不是硬边界；边界外仍有指数尾部。计算遮挡功率时要积分实际强度，不能把$\pi w^2$内外简单当作有光与无光。

\subsection{同一ABCD矩阵怎样作用于复参数}
在同一均匀介质的光线坐标$(y,\theta)$下，自由传播使$q\mapsto q+d$。薄透镜附加相位$\exp[-\mathrm ikr^2/(2f)]$，使$1/q\mapsto1/q-1/f$。两者与光线矩阵一致，因此由这些元件组成的系统满足
\[
 q_{\rm out}=\frac{Aq_{\rm in}+B}{Cq_{\rm in}+D}.
\]
可逐个乘矩阵验证分式变换的复合规律。若使用约化坐标$(y,n\theta)$或跨越不同外部介质，应同步改变复参数定义，不能在矩阵中保留折射率而在$q$中忽略它。本节例题两侧均为空气。

\begin{example}透镜后的束腰未必在几何焦平面
一高斯光束恰在薄正透镜平面有束腰，腰半径$w_0$、Rayleigh长度$z_R$，透镜焦距$f>0$。求透镜后的新束腰位置和腰半径，比较$z_R\gg f$与$z_R=f$两种情况。
\ExpansionDiagram{fig:dep-o1-gaussfocus-q}{透镜位于原束腰面，入射波前平坦但束宽有限}{
\draw[->](-3,0)--(4.5,0)node[right]{$z$};\draw[<->,thick](0,-1.4)--(0,1.4)node[above]{透镜};\draw[thick,domain=-2.5:0,samples=70]plot(\x,{.7*sqrt(1+(\x/3)^2)});\draw[thick,domain=-2.5:0,samples=70]plot(\x,{-.7*sqrt(1+(\x/3)^2)});\draw[dashed](3,-1.2)--(3,1.2)node[above]{$f$};\node[left]at(0,-1.6){原束腰};}
\end{example}
\begin{solution}
图\ref{fig:dep-o1-gaussfocus-q}中入射$q_-=-\mathrm iz_R$。薄透镜后
\[
 \frac1{q_+}=-\frac1f+\frac{\mathrm i}{z_R},\qquad
 q_+=-\frac{fz_R^2}{f^2+z_R^2}
 -\mathrm i\frac{f^2z_R}{f^2+z_R^2}.
\]
继续传播距离$z$后，$q(z)=q_++z$。新束腰要求实部为零，所以
\[
 z_w'=\frac{fz_R^2}{f^2+z_R^2},\qquad
 z_R'=\frac{f^2z_R}{f^2+z_R^2},\qquad
 \frac{w_0'}{w_0}=\sqrt{\frac{z_R'}{z_R}}
 =\frac f{\sqrt{f^2+z_R^2}}.
\]
\ExpansionDiagram{fig:dep-o1-gaussfocus-a}{新束腰在几何焦点之前；图示$z_R=f$时位于$f/2$}{
\draw[->](0,0)--(6,0)node[right]{$z$};\draw[thick,domain=0:5,samples=90]plot(\x,{.55*sqrt(1+(\x-2)^2/4)});\draw[thick,domain=0:5,samples=90]plot(\x,{-.55*sqrt(1+(\x-2)^2/4)});\draw[<->](0,-1.5)--(0,1.5);\draw[dashed](2,-1.2)--(2,1.2)node[above]{$f/2$};\draw[dashed](4,-1.5)--(4,1.5)node[above]{$f$};}
当$z_R\gg f$，$z_w'\simeq f$、$w_0'\simeq\lambda_0f/(\pi w_0)$，几何焦平面近似成立。若$z_R=f$，图\ref{fig:dep-o1-gaussfocus-a}中$z_w'=f/2,z_R'=f/2,w_0'=w_0/\sqrt2$。原波前虽平坦，有限腰半径意味着角谱本来就有宽度，不能把它等同于所有光线严格平行的无限平面波。

该结果假定透镜孔径足够大、无明显像差，且新束腰仍满足近轴条件$\lambda_0/(\pi w_0')\ll1$。若透镜把光束截断，单一高斯$q$参数已不能表示完整衍射场，应回到带实际孔径的积分。
\end{solution}

\section{习题与解答}
\begin{exercise}\label{ex:o1oval}
要用一个折射界面把介质$n_1$中的点$A$严格成像到介质$n_2$中的点$B$。忽略反射，证明等光程曲面$n_1|Q-A|+n_2|B-Q|=C$的适当光滑部分满足折射定律，说明与球面近轴成像的区别。
\Diagram{fig:o1ovalq}{笛卡尔卵形面的一个分支，图取$n_1=1,n_2=1.5,C=7.5$}{\fill(0,0)circle(1.5pt)node[below]{$A$};\fill(5.5,0)circle(1.5pt)node[below]{$B$};\draw[curve]plot coordinates {(3.732186,-1.500000) (3.532345,-1.462025) (3.375574,-1.424051) (3.242276,-1.386076) (3.124437,-1.348101) (3.017837,-1.310127) (2.919932,-1.272152) (2.829045,-1.234177) (2.744007,-1.196203) (2.663963,-1.158228) (2.588269,-1.120253) (2.516428,-1.082278) (2.448047,-1.044304) (2.382811,-1.006329) (2.320466,-0.968354) (2.260808,-0.930380) (2.203667,-0.892405) (2.148908,-0.854430) (2.096420,-0.816456) (2.046116,-0.778481) (1.997930,-0.740506) (1.951812,-0.702532) (1.907731,-0.664557) (1.865669,-0.626582) (1.825624,-0.588608) (1.787608,-0.550633) (1.751645,-0.512658) (1.717772,-0.474684) (1.686040,-0.436709) (1.656510,-0.398734) (1.629252,-0.360759) (1.604347,-0.322785) (1.581883,-0.284810) (1.561951,-0.246835) (1.544645,-0.208861) (1.530055,-0.170886) (1.518266,-0.132911) (1.509353,-0.094937) (1.503375,-0.056962) (1.500375,-0.018987) (1.500375,0.018987) (1.503375,0.056962) (1.509353,0.094937) (1.518266,0.132911) (1.530055,0.170886) (1.544645,0.208861) (1.561951,0.246835) (1.581883,0.284810) (1.604347,0.322785) (1.629252,0.360759) (1.656510,0.398734) (1.686040,0.436709) (1.717772,0.474684) (1.751645,0.512658) (1.787608,0.550633) (1.825624,0.588608) (1.865669,0.626582) (1.907731,0.664557) (1.951812,0.702532) (1.997930,0.740506) (2.046116,0.778481) (2.096420,0.816456) (2.148908,0.854430) (2.203667,0.892405) (2.260808,0.930380) (2.320466,0.968354) (2.382811,1.006329) (2.448047,1.044304) (2.516428,1.082278) (2.588269,1.120253) (2.663963,1.158228) (2.744007,1.196203) (2.829045,1.234177) (2.919932,1.272152) (3.017837,1.310127) (3.124437,1.348101) (3.242276,1.386076) (3.375574,1.424051) (3.532345,1.462025) (3.732186,1.500000)};\draw[curve](0,0)--(2.372225,1)--(5.5,0);\draw[guide](0,0)--(2.372225,-1)--(5.5,0);\node[above]at(2.372225,1){$Q$};}
\end{exercise}
\begin{solution}
定义$F(Q)=n_1r_1+n_2r_2$，沿入射传播方向$\vct t_1=(Q-A)/r_1$、出射方向$\vct t_2=(B-Q)/r_2$。梯度为$\nabla F=n_1\vct t_1-n_2\vct t_2$。图\ref{fig:o1ovalq}的等值面法线沿$\nabla F$，任何切向位移$\delta\vct q$满足$\nabla F\cdot\delta\vct q=0$，所以$n_1\vct t_{1\parallel}=n_2\vct t_{2\parallel}$。
\Diagram{fig:o1ovala}{等光程面的法线由入射、出射约化方向矢量之差确定}{\draw[axis](0,0)--(2.4,1.2)node[above left]{$n_1\vct t_1$};\draw[axis](0,0)--(2.1,-.5)node[below]{$n_2\vct t_2$};\draw[->,thick](2.1,-.5)--(2.4,1.2);\node[right]at(2.6,.35){差矢量};\draw[guide](0,0)--(.3,1.7);}
上式为斯涅尔切向条件，法线的矢量构造见图\ref{fig:o1ovala}。须选择使光由介质1进入介质2、无全反射且不被其他界面截断的分支。球面一般只能在轴上小孔径处近似满足固定物像等光程；这个等光程面通常不是球面，可以为指定物像对校正球差，却不保证别的物点也严格成像。
\end{solution}

\begin{exercise}\label{ex:o1galilean}
空气中正物镜焦距$f_1>0$，负目镜$f_2<0$，间隔$d=f_1+f_2>0$。求无焦系统角放大率、轴向平行束宽度比和像的正倒，与双正透镜系统比较。
\Diagram{fig:o1galileoq}{负目镜截取尚未到物镜焦点的会聚光束}{\OptAxis\draw[<->](1,-1.3)--(1,1.3);\draw[>-<](4,-1)--(4,1);\draw[guide](5.5,-1.5)--(5.5,1.5);\node[above]at(5.5,1.5){物镜焦点};\node[above]at(4,1){负目镜};}
\end{exercise}
\begin{solution}
图\ref{fig:o1galileoq}仍用两透镜矩阵，把$d=f_1+f_2$代入得到
\[
 M=\begin{pmatrix}-f_2/f_1&f_1+f_2\\0&-f_1/f_2\end{pmatrix}.
\]
因$f_2<0$，两个对角元都正。轴向平行光的宽度比$|f_2|/f_1$，小视角放大率$\theta_o/\theta_i=f_1/|f_2|>0$，所以角方向不反转。
\Diagram{fig:o1galileoa}{光束在到达几何焦点之前被负透镜重新准直}{\OptAxis\draw[<->](1,-1.3)--(1,1.3);\draw[>-<](4,-1)--(4,1);\draw[->](0,1)--(1,1)--(4,.333333)--(6,.333333);\draw[->](0,-1)--(1,-1)--(4,-.333333)--(6,-.333333);\draw[guide](4,.333333)--(5.5,0);\draw[guide](4,-.333333)--(5.5,0);}
图\ref{fig:o1galileoa}按$f_1=4.5,f_2=-1.5,d=3$作图，束宽缩为1/3、角放大率3，内部没有经过实焦点。双正系统角放大率为负，会产生倒向角像。两系统都满足$\Phi=0$，但角放大率符号和内部光路不同，对应正立、倒立两种观察效果。
\end{solution}

\begin{exercise}\label{ex:o1layers}
平行分层透明介质的各层厚度为$h_j$、折射率$n_j$。底部点物体从顶面外的介质$n_0$中近轴观察。求表观深度，说明为何各层的真实光程厚度不能直接作为表观深度相加。
\Diagram{fig:o1layersq}{平行分层结构保持每条光线的切向约化分量}{\foreach\y in{0,.7,1.4,2.1}{\draw[curve](-2,\y)--(2,\y);}\fill(0,0)circle(1.5pt)node[below]{物点};\draw[->](0,0)--(.35,.7)--(.6,1.4)--(1.1,2.1)--(1.8,3);\node[right]at(2,.35){$n_3,h_3$};\node[right]at(2,1.05){$n_2,h_2$};\node[right]at(2,1.75){$n_1,h_1$};}
\end{exercise}
\begin{solution}
图\ref{fig:o1layersq}所有平面界面无光焦度，$p=n\theta$在近轴下保持。光从底部穿过第$j$层的横向位移为$h_j\theta_j=h_jp/n_j$，到顶面总横移$y=p\sum_jh_j/n_j$。外部出射斜率为$\theta_0=p/n_0$。
\Diagram{fig:o1layersa}{外部光线反向延长至虚像，深度由横移除外部斜率决定}{\draw[curve](-2,1.6)--(2,1.6);\fill(0,0)circle(1.5pt)node[below]{真实物};\fill(0,.6)circle(1.5pt)node[left]{虚像};\draw[guide](0,.6)--(1,1.6);\draw[->](1,1.6)--(2.3,2.9);\draw[guide](0,0)--(0,2.6);}
图\ref{fig:o1layersa}表观深度$h_a=y/\theta_0=n_0\sum_jh_j/n_j$。真实传播光程为$\sum_jn_jh_j$，一个含$n_j$倒数、另一个含$n_j$本身，物理意义不同。若各层折射率相同$n$，便恢复$n_0h/n$；若$n_0=n$，深度不变。
\end{solution}

\begin{exercise}\label{ex:o1immersion}
同一薄双凸透镜折射率$n_l=1.50$，在空气中焦距$f_a$。完全浸入折射率$n_0=1.33$的液体，求主面焦距与$f_a$之比，讨论液体折射率超过玻璃时的光焦度符号。
\Diagram{fig:o1immersionq}{同一曲率透镜置于不同外部介质}{\OptAxis\draw[curve](2,-1.3)..controls(1.4,-.4)and(1.4,.4)..(2,1.3);\draw[curve](2,-1.3)..controls(2.6,-.4)and(2.6,.4)..(2,1.3);\node[above]at(2,1.5){$n_l$};\node at(4,1){$n_0$};}
\end{exercise}
\begin{solution}
令曲率组合$K=1/R_1-1/R_2>0$。图\ref{fig:o1immersionq}在外部介质中约化光焦度$\Phi=(n_l-n_0)K$，两侧相同折射率的主面焦距$f_0=n_0/\Phi$。空气中$f_a=1/[(n_l-1)K]$，所以
\[
 \frac{f_0}{f_a}=\frac{n_0(n_l-1)}{n_l-n_0}
 =\frac{1.33\times.50}{.17}\simeq3.912.
\]
\Diagram{fig:o1immersiona}{外介质接近透镜折射率，折射偏折变弱、焦距增大}{\OptAxis\draw[<->](1.5,-1.2)--(1.5,1.2);\draw[->](0,.8)--(1.5,.8)--(5.8,0);\draw[guide](1.5,.8)--(2.6,0);\node[below]at(2.6,0){空气焦点};\node[below]at(5.8,0){液体焦点};}
图\ref{fig:o1immersiona}焦距增大是相对折射率差减小的结果。$n_0=n_l$时界面不可光学分辨，薄透镜光焦度为零；$n_0>n_l$时$\Phi<0$，相同双凸形状转为发散。透镜的会聚、发散性质由相对折射率和几何形状共同决定。
\end{solution}

\begin{exercise}\label{ex:o1mirage}
平移不变介质中$n(y)=n_A-gy$，$g>0$，只取$n>0$区域。光从$(z,y)=(0,0)$向上方以相对$+z$角$\alpha_0\in(0,\pi/2)$进入。求精确轨迹、最高点和回到$y=0$的水平距离。
\Diagram{fig:o1mirageq}{折射率向上减小，光线弯回高折射率区}{\OptAxis\draw[axis](0,0)--(0,3)node[above]{$y$};\draw[axis](0,0)--(1.4,1)node[above]{$\alpha_0$};\draw[guide](0,2.5)--(5.5,2.5);\node[right]at(5.5,2.5){$n$较小};}
\end{exercise}
\begin{solution}
图\ref{fig:o1mirageq}的守恒量$b=n\cos\alpha=n_A\cos\alpha_0$。轨迹斜率$y'=\tan\alpha$，所以$n^2=b^2(1+y'^2)$。在上升支$y'>0$，令$n=b\cosh u$，则$y'=\sinh u$；由$n=n_A-gy$，$\dd y=-(b/g)\sinh u\dd u$，于是$\dd z=-(b/g)\dd u$。
\Diagram{fig:o1miragea}{轨迹是倒置悬链线的一段，对最高点左右对称}{\OptAxis\draw[axis](0,0)--(0,2.7)node[above]{$y$};\draw[curve,domain=0:2.8,samples=80]plot(\x,{1.4*(cosh(1)-cosh((\x-1.4)/1.4))});\draw[guide](1.4,0)--(1.4,.760313);\node[above]at(1.4,.85){最高点};\node[below]at(2.8,0){返回};}
以最高点$u=0,z=z_t$为参考点，上升支满足
\[
 \int_{z_t}^{z}\dd Z=-\frac bg\int_0^u\dd q,
 \qquad z-z_t=-\frac bg u.
\]
故$u=g(z_t-z)/b$。由初始条件$n(0)=n_A=b\cosh(gz_t/b)$求出$z_t$，再以$\cosh$的偶性接续下降支，得到
\[
 y(z)=\frac{n_A-b\cosh[g(z-z_t)/b]}g,\quad
 z_t=\frac bg\operatorname{arcosh}\frac{n_A}b.
\]
最高点$y_t=(n_A-b)/g$，回到$y=0$时$z=2z_t$。图\ref{fig:o1miragea}按一组比例画出同一悬链线结构。这个解适用于有限角度。前面的$n^2$抛物分布也有精确的正弦型轨迹，但其轴向周期依赖初始方向；仅在近轴极限下才化为固定节距$2\pi a$。两种分布的轨迹都须限制在实际介质覆盖的区域内。
\end{solution}

\begin{exercise}\label{ex:o1chromatic}
薄透镜几何曲率固定，空气中折射率随波长变化。已知$n(\lambda_0)=n_0$，一小段波长变化对应$\Delta n$。求焦距一阶变化及平行白光的轴向色差方向。
\Diagram{fig:o1chromq}{同一透镜对不同波长的光焦度不同}{\OptAxis\draw[<->](1.8,-1.3)--(1.8,1.3);\draw[->](0,1)--(1.8,1);\draw[curve](1.8,1)--(4,0);\draw[guide](1.8,1)--(5.6,0);\node[below]at(4,0){短波};\node[below]at(5.6,0){长波};}
\end{exercise}
\begin{solution}
图\ref{fig:o1chromq}的$f=K^{-1}(n-1)^{-1}$。固定$K$求导，$\dd f/\dd n=-K^{-1}(n-1)^{-2}=-f/(n-1)$，所以$\Delta f\simeq-f_0\Delta n/(n_0-1)$。对题图中的正透镜，正常色散使短波折射率较大，蓝光因而比红光更早会聚；发散透镜也服从同一有向焦距微分式。
\Diagram{fig:o1chroma}{像面固定时，不同焦点在屏上形成不同半径模糊斑}{\OptAxis\draw[<->](1,-1.2)--(1,1.2);\draw[guide](4,-1.6)--(4,1.6);\draw[curve](1,.9)--(3,0)--(5,-.9);\draw[curve](1,-.9)--(3,0)--(5,.9);\draw[guide](1,.9)--(4.5,0);\draw[guide](1,-.9)--(4.5,0);}
图\ref{fig:o1chroma}说明固定屏幕不可能同时处在所有波长的几何焦面。色差是$n(\lambda)$改变光焦度的几何效应，单色衍射则即使焦距相同仍形成有限光斑；两种宽度不能不分机制地直接相加。本一阶式还要求$|\Delta n|\ll n_0-1$。
\end{solution}

\begin{exercise}\label{ex:o1prismminimum}
折射率$n$、顶角$A$的棱镜在空气中，光由一个斜面入射、另一个出射。推导对称路径下的最小偏向角公式，说明该状态需满足的传播条件。
\Diagram{fig:o1prismq}{一般穿透光路；法线与界面垂直，入射、折射角分别对法线量取}{
\draw[curve](0,0)--(2,3.4641)--(4,0)--cycle;
\draw[->](-.7,.395855)--(.692820,1.2);
\draw[curve](.692820,1.2)--(3.175128,1.428721);
\draw[->](3.175128,1.428721)--(4.7,1.190565);
\draw[guide](0,1.6)--(1.385641,.8);
\draw[guide](2.482307,1.028721)--(3.867948,1.828721);
\node at(2,2.75){$A$};\node at(2,.55){$n>1$};
\node[above left]at(.28,1.23){$i_1$};\node[below right]at(.98,1.2){$r_1$};
}
\end{exercise}
\begin{solution}
图\ref{fig:o1prismq}中折射角满足$r_1+r_2=A$，空气侧入、出射角满足$\sin i_j=n\sin r_j$，总偏向为$D=i_1+i_2-A$。令$f(r)=\arcsin(n\sin r)$，则$D=f(r_1)+f(A-r_1)-A$。在$n>1$、$0\leq r_j<\arcsin(1/n)$的通常穿透分支，
\begin{align*}
 f'(r)&=\frac{n\cos r}{\sqrt{1-n^2\sin^2r}},\\
 f''(r)&=-\frac{n\sin r}{\sqrt{1-n^2\sin^2r}}
 +\frac{n^3\sin r\cos^2r}{(1-n^2\sin^2r)^{3/2}}\\
 &=\frac{n(n^2-1)\sin r}{(1-n^2\sin^2r)^{3/2}}\geq0.
\end{align*}
凸性保证$r_1=r_2=A/2$为极小点，光路可逆性则给出对称驻点。
\Diagram{fig:o1prisma}{最小偏向的对称光路；图取顶角60度、折射率1.5}{
\draw[curve](0,0)--(2,3.4641)--(4,0)--cycle;
\draw[->](-.7,.8929)--(.8083,1.4);
\draw[->](.8083,1.4)--(3.1917,1.4);
\draw[->](3.1917,1.4)--(4.7,.8929);
\draw[guide](2,-.2)--(2,3.7);
\draw[guide](.1155,1.8)--(1.5011,1);
\draw[guide](2.4989,1)--(3.8845,1.8);
\node at(1.3,1.1){$r_1$};\node at(2.7,1.1){$r_2$};
\node[below]at(2,-.3){$r_1=r_2=A/2$};
\node at(2,2.85){$A$};
}
图\ref{fig:o1prisma}的对称路径有$i_1=i_2=i$，因此
\[
 D_{\min}=2i-A,
 \qquad n=\frac{\sin[(A+D_{\min})/2]}{\sin(A/2)}.
\]
可传播条件为$n\sin(A/2)\leq1$；等号对应掠入、掠出极限，有限通光通常要求严格小于1。若不满足，无法让光沿这个两面穿透分支通过。薄棱镜$A\ll1$且$nA\ll1$时$i\simeq nA/2$，于是$D_{\min}\simeq(n-1)A$。近似要求两个空气侧角度也小，不能只检查棱镜顶角而忽略很大的折射率。
\end{solution}

\begin{exercise}\label{ex:o1obliqueplate}
平面波以角$i$穿过厚$h$、折射率$n$的平行板，板外为空气。比较板后同一平面上的相位与未插板情况，推导附加相位，避免把折线实际路长误当程差。
\Diagram{fig:o1obliqueplateq}{比较同一输出平面上的波相位}{\draw[curve](-2,0)--(2,0);\draw[curve](-2,1)--(2,1);\draw[->](-1.5,-1)--(-.5,0)--(.1,1)--(1.1,2);\draw[guide](-.5,0)--(.5,1);\node[right]at(2,.5){$h$};}
\end{exercise}
\begin{solution}
图\ref{fig:o1obliqueplateq}中平行界面保持切向波矢，$\sin i=n\sin r$。在固定横坐标比较上下参考面的相位增长，板内法向波矢$k_{zg}=k_0n\cos r$，原空气中为$k_{za}=k_0\cos i$。所以附加相位$\Delta\phi=k_0h(n\cos r-\cos i)$。
\Diagram{fig:o1obliqueplatea}{沿不同几何射线积累的光程须先归到同一参考波面}{\draw[curve](0,0)--(4,0);\draw[curve](0,1)--(4,1);\draw[guide](1,-.5)--(1,1.7);\draw[->](.3,-.6)--(1,0)--(1.5,1)--(2.2,1.6);\node[above]at(1,1.8){固定横坐标};}
图\ref{fig:o1obliqueplatea}的计算自动扣除了横向位移相关相位。正常入射恢复$k_0(n-1)h$。若题目改为固定两个点之间的封闭干涉臂，需同时满足光线连接、元件位置和合束条件，不能只把本平面波相位式套给任意折线路径。
\end{solution}

\begin{exercise}\label{ex:o1matrixinverse}
空气近轴系统矩阵为$M$且$\det M=1$。写出逆矩阵，解释“数学逆传播”与把整套透镜实际前后翻转为何须先统一坐标基底。
\Diagram{fig:o1inverseq}{同一对参考面之间的前向、反向关系}{\draw[curve](1,-1)rectangle(3.5,1);\node at(2.25,0){$M$};\draw[->](-.5,.5)--(.8,.5);\draw[->](3.7,-.5)--(5,-.5);}
\end{exercise}
\begin{solution}
图\ref{fig:o1inverseq}用同一固定$z$和有向倾角描述时，数学逆为$M^{-1}=\bigl(\begin{smallmatrix}D&-B\\-C&A\end{smallmatrix}\bigr)$，因为与原矩阵相乘给单位阵。逆关系把末端的$y,n\theta$恢复到起点，并不自动改变倾角的参考轴。
\Diagram{fig:o1inversea}{反向实际传播若重新令光轴向前，倾角符号也改变}{\draw[axis](0,.7)--(4,.7)node[right]{$+z$};\draw[axis](4,-.7)--(0,-.7)node[left]{$+\tilde z$};\node at(2,0){同一光线需换基};}
若反向传播时重新把轴指向左方而高度仍向上，约化倾角需变号，换基矩阵$K=\operatorname{diag}(1,-1)$，反向系统矩阵为$KM^{-1}K=\bigl(\begin{smallmatrix}D&B\\C&A\end{smallmatrix}\bigr)$。图\ref{fig:o1inversea}说明相同光学系统反向后的$B,C$正负取决于坐标约定，不能只看一处符号判定互易性失效。
\end{solution}

\begin{exercise}\label{ex:o1focusplate}
物点到固定屏距离$D$，中间放焦距$f$的薄透镜，空气中近轴成像。求允许的透镜位置数及放大率乘积，讨论$D<4f$。
\Diagram{fig:o1focusq}{固定物屏间距，透镜位置待定}{\OptAxis\draw[->](0,0)--(0,1);\draw[<->](2,-1.2)--(2,1.2);\draw[curve](6,-1.5)--(6,1.5);\draw[<->](0,-1.8)--(6,-1.8)node[midway,below]{$D$};}
\end{exercise}
\begin{solution}
图\ref{fig:o1focusq}中取物体在0、屏在$D$，透镜位置为$u$，于是像距$v=D-u$。由$1/u+1/v=1/f$，相乘得到$uv=f(u+v)=fD$，即
\[
 u^2-Du+fD=0,\qquad
 u_\pm=\frac{D\pm\sqrt{D^2-4fD}}2.
\]
题中的实物、实屏位于透镜两侧，故要求$0<u<D$；当$D>4f$、$f>0$时两个根均满足此条件。
\Diagram{fig:o1focusa}{两次成像透镜位置关于物屏中点对称}{\OptAxis\draw[<->](1.7,-1.2)--(1.7,1.2);\draw[guide](4.3,-1.2)--(4.3,1.2);\draw[curve](6,-1.5)--(6,1.5);\draw[guide](3,-1.7)--(3,1.7);\node[above]at(1.7,1.2){位置一};\node[above]at(4.3,1.2){位置二};}
图\ref{fig:o1focusa}的两个位置将$u,v$互换，放大率分别为$-v/u$和$-u/v$，乘积1。一幅放大像与一幅缩小像配对出现。两位置间距
\[
 l=u_+-u_- =\sqrt{D^2-4fD},\qquad
 f=\frac{D^2-l^2}{4D}.
\]
这给出了用固定物屏间距与两次清晰成像位置测焦距的方法。$D=4f$时判别式为零，唯一位置$u=v=2f$、放大率$-1$；$D<4f$时没有满足实物实像要求的透镜位置。也可由$u+v\geq2\sqrt{uv}=2\sqrt{fD}$直接得到$D\geq4f$。实际厚透镜要以主面测距，若机械测的是两顶点位置，须先修正主面偏移。
\end{solution}

\begin{exercise}\label{ex:o1etendue}
理想近轴成像系统横向放大率为$m_y$，物像介质折射率分别$n_O,n_I$。求同一轴上物点光束的角放大率，并解释放大像时为何常需减小像方数值孔径。
\Diagram{fig:o1etendueq}{同一物点的两条光线定义孔径角}{\draw[axis](-.3,0)--(6,0);\draw[curve](2,-1.1)rectangle(3,1.1);\draw[curve](0,0)--(2,1);\draw[curve](0,0)--(2,-1);\draw[curve](3,.6)--(5,0);\draw[curve](3,-.6)--(5,0);}
\end{exercise}
\begin{solution}
图\ref{fig:o1etendueq}中，物像面之间使用约化变量$(y,p=n\theta)$。成像要求$B=0$，横向放大率为$A=m_y$。矩阵行列式恒为1，故$AD=1$、$D=1/m_y$，有
\[
 \begin{pmatrix}y_I\\n_I\theta_I\end{pmatrix}
 =\begin{pmatrix}m_y&0\\ C&1/m_y\end{pmatrix}
 \begin{pmatrix}y_O\\n_O\theta_O\end{pmatrix}.
\]
对轴上同一物点$y_O=0$，下行式化为$n_I\theta_I=n_O\theta_O/m_y$，所以角放大率为$n_O/(n_Im_y)$。
\Diagram{fig:o1etenduea}{像高和约化孔径角的乘积保持}{\draw[curve](0,0)rectangle(1.8,.7);\node at(.9,.35){$h_O$};\draw[curve](3,0)rectangle(5.7,1.1);\node at(4.35,.55){$h_I=m_yh_O$};\node[below]at(3,-.8){$n_Oh_O\theta_O=n_Ih_I\theta_I$};}
图\ref{fig:o1etenduea}所画两束使用同一组通光光线，其数值孔径满足$\mathrm{NA}_I\simeq\mathrm{NA}_O/|m_y|$。负放大率还带来倾角反号，孔径大小则取绝对值。对轴上附近小面元，面积放大$|m_y|^2$，小圆锥立体角约为$\pi\theta^2$，于是
\[
 n_I^2\dd A_I\dd\Omega_I
 =n_O^2\dd A_O\dd\Omega_O.
\]
该关系给出近轴通光相空间体积的不变性。无损时同一光束的功率也守恒，因此辐亮度除以$n^2$保持。扩大像面积会减小相应像方角锥；若另行增大物方收集孔径，所比较的已是另一组光线，不能把两种条件下的亮度混为一谈。
\end{solution}

\begin{exercise}\label{ex:o1magnification}
薄透镜在空气中$f>0$，物高$h$，物距$u>f$。求像距与横向放大率，并由两条光线画出$u=5f/3$的像。
\Diagram{fig:o1magnifyq}{已知物高、物距与焦距}{\OptAxis\draw[<->](2,-1.8)--(2,1.8);\draw[->,thick](0,0)--(0,1);\node[left]at(0,1){$h$};\fill(3.2,0)circle(1.3pt)node[below]{$F'$};\draw[<->](0,-.6)--(2,-.6)node[midway,below]{$u$};}
\end{exercise}
\begin{solution}
图\ref{fig:o1magnifyq}中$1/u+1/v=1/f$，解得$v=fu/(u-f)$。通过光心的近轴光线保持方向，由相似三角形$h'/v=-h/u$，所以横向放大率$m_y=h'/h=-v/u=-f/(u-f)$。
\Diagram{fig:o1magnifya}{平行轴光线经后焦点，过光心光线不偏折，两者交于倒立实像}{\OptAxis\draw[<->](2,-1.9)--(2,1.9);\draw[->,thick](0,0)--(0,1);\draw[->,thick](5,0)--(5,-1.5);\draw[->](0,1)--(2,1);\draw[->](2,1)--(5,-1.5);\draw[->](0,1)--(5,-1.5);\fill(3.2,0)circle(1.2pt)node[above]{$F'$};}
取$u=5f/3$，则$v=5f/2$、$m_y=-3/2$。图\ref{fig:o1magnifya}按$f=1.2$、$u=2$作图，实像在$z=5$、像高$-1.5h$，交点与计算吻合。真实透镜在大孔径或大物高时需考虑像差，本图只表达近轴模型。
\end{solution}

\begin{exercise}\label{ex:o1afocal}
正透镜$f_1>0$、$f_2>0$相隔$f_1+f_2$。求平行细光束通过后的束径比与小角偏转比，并说明无焦系统是否“没有作用”。
\Diagram{fig:o1afocalq}{共焦双正透镜系统}{\OptAxis\draw[<->](1,-1.4)--(1,1.4);\draw[<->](5,-1.1)--(5,1.1);\draw[guide](3.5,-1.6)--(3.5,1.6);\node[above]at(3.5,1.6){$F'_1=F_2$};}
\end{exercise}
\begin{solution}
图\ref{fig:o1afocalq}中$d=f_1+f_2$，直接乘矩阵得到
\[
 M=\begin{pmatrix}-f_2/f_1&f_1+f_2\\0&-f_1/f_2\end{pmatrix}.
\]
对于轴向平行光$\theta_i=0$，$y_o=-(f_2/f_1)y_i$，输出仍平行，束径绝对比$f_2/f_1$。
\Diagram{fig:o1afocala}{轴向光束先会聚再重新平行，图取两焦距不同}{\OptAxis\draw[<->](1,-1.4)--(1,1.4);\draw[<->](5,-1.4)--(5,1.4);\draw[->](0,1)--(1,1);\draw[->](1,1)--(5,-.6);\draw[->](5,-.6)--(6.5,-.6);\draw[->](0,-1)--(1,-1);\draw[->](1,-1)--(5,.6);\draw[->](5,.6)--(6.5,.6);}
对整体倾斜的小角平行光，$\theta_o=-(f_1/f_2)\theta_i$，与高度无关。图\ref{fig:o1afocala}取$f_1/f_2=5/3$，中间焦点位于第一透镜后2.5图示单位。无焦仅表示$C=0$、平行束仍为平行束；系统仍可扩束、缩束和改变角度，绝非恒等映射。
\end{solution}

\begin{exercise}\label{ex:o1glassplate}
一块厚度$d$、折射率$n$的平行玻璃板插入原本将于空气中会聚的近轴光束，板面垂直轴。求焦点沿轴的位移，忽略反射。
\Diagram{fig:o1glassq}{会聚光束在平行板内传播的斜率减小}{\OptAxis\draw[curve](2,-1.1)rectangle(3,1.1);\draw[guide](0,1.3)--(4.8,0);\draw[guide](0,-1.3)--(4.8,0);\node[above]at(2.5,1.1){$n,d$};\fill(4.8,0)circle(1.3pt)node[below]{$F_0$};}
\end{exercise}
\begin{solution}
图\ref{fig:o1glassq}中原光线以斜率$\theta<0$会聚。进入平面板后$\theta_g=\theta/n$，传播$d$导致高度变化$d\theta/n$；若仍为空气则变化$d\theta$。出板后倾角恢复$\theta$，故高度比原光线高出$\Delta y=d\theta(1/n-1)>0$。
\Diagram{fig:o1glassa}{出板后两光线平行，新的交轴点向后移动}{\OptAxis\draw[curve](2,-1.1)rectangle(3,1.1);\draw[->](0,1.2)--(2,.7)--(3,.575)--(5.3,0);\draw[guide](2,.7)--(4.8,0);\fill(4.8,0)circle(1.2pt)node[below]{$F_0$};\fill(5.3,0)circle(1.2pt)node[above right]{$F$};}
图\ref{fig:o1glassa}中的轴向位移由$\theta\Delta z=-\Delta y$得$\Delta z=d(1-1/n)>0$。矩阵角度看，是把空气传播长度$d$的约化距离$d$替换为$d/n$，而折射后的外部倾角不变。大孔径时不同光线位移并不完全相同，会引入球差。
\end{solution}

\begin{exercise}\label{ex:o1turning}
水平分层介质$n(y)=n_0\sqrt{1-y^2/a^2}$。光线从$y=0$以与$+z$轴成$\alpha_0$的方向进入。用守恒量求最大高度，不使用近轴条件。
\Diagram{fig:o1turnq}{光线在折射率下降区弯折至水平}{\OptAxis\draw[axis](0,0)--(0,2.7)node[above]{$y$};\draw[curve](0,0)..controls(1,1)and(2,1.8)..(3.2,1.8);\draw[guide](0,1.8)--(4,1.8);\node[left]at(0,1.8){$y_t$};\node at(.8,.3){$\alpha_0$};}
\end{exercise}
\begin{solution}
图\ref{fig:o1turnq}中，介质沿$z$平移不变，$b=n\cos\alpha=n_0\cos\alpha_0$守恒。用$y'=\tan\alpha$表示方向，得到
\[
 n^2(y)=b^2(1+y'^2),\qquad
 y'^2=\frac{\sin^2\alpha_0-y^2/a^2}{\cos^2\alpha_0}.
\]
真实射线要求右侧非负，允许区域为$|y|\leq a|\sin\alpha_0|$。最高点$y'=0$，所以正向上射线的最大高度为$y_t=a\sin\alpha_0$；若只讨论幅度，则为$a|\sin\alpha_0|$。
\Diagram{fig:o1turna}{允许高度由$n(y)\geq n_0\cos\alpha_0$确定}{\draw[axis](0,0)--(4.2,0)node[right]{$y/a$};\draw[axis](0,0)--(0,2.8)node[above]{$n/n_0$};\draw[curve,domain=0:.99,samples=70]plot({3.5*\x},{2.3*sqrt(1-\x*\x)});\draw[guide](0,1.7)--(3.7,1.7);\node[left]at(0,1.7){$\cos\alpha_0$};}
图\ref{fig:o1turna}中的折射率曲线与水平$b$线相交处就是转点。还可将方程积分，以核对近轴极限。对$0<\alpha_0<\pi/2$的上升支，令$A=a\sin\alpha_0$，则
\[
 \int_0^{y}\frac{\dd q}{\sqrt{A^2-q^2}}
 =\frac1{a\cos\alpha_0}\int_0^z\dd s,\qquad
 [\arcsin(q/A)]_0^{y}=\frac{z}{a\cos\alpha_0}.
\]
这里用$q=A\sin u$，在上升支$0\leq u\leq\pi/2$有$\dd q=A\cos u\dd u$、$\sqrt{A^2-q^2}=A\cos u$。反解得$y=A\sin[z/(a\cos\alpha_0)]$，在转点处接续下降支。
精确轴向周期是$2\pi a\cos\alpha_0$，最高点在$z_t=\pi a\cos\alpha_0/2$。小角时振幅$A\simeq a\alpha_0$、周期趋$2\pi a$，恢复近轴解。有限角度下转点高度与周期都应使用轴向守恒量求出，不能仅保留振幅修正而仍把近轴周期称为精确周期。
\end{solution}

\begin{exercise}\label{ex:o1centralturn}
圆对称介质在某区域$n(r)=n_a a/r$。求一条非径向光线能否有孤立最近点，推导其轨迹形状。讨论守恒量$b=n_a a$的特殊情形。
\Diagram{fig:o1centralq}{径向折射率场中的光线与径向夹角}{\draw[curve](0,0)circle(1);\draw[curve](0,0)circle(2);\draw[axis](0,0)--(2,1)node[right]{$r$};\draw[axis](2,1)--(2.5,2.2)node[above]{切向};\node[below left]at(0,0){$O$};}
\end{exercise}
\begin{solution}
图\ref{fig:o1centralq}的$n(r)r=n_aa$为常数，所以$\sin\psi=b/(n_aa)$也为常数。若$0<b<n_aa$，夹角固定且不等于直角，不存在孤立转点。由$\dd r/\dd s=\cos\psi$、$r\dd\varphi/\dd s=\sin\psi$，得$\dd\varphi/\dd r=\tan\psi/r$。
\Diagram{fig:o1centrala}{固定夹角对应对数螺线，图只画有限介质区域}{\draw[axis](-2.3,0)--(2.5,0);\draw[axis](0,-2.3)--(0,2.5);\draw[curve,domain=0:300,samples=100]plot({.35*exp(.006*\x)*cos(\x)},{.35*exp(.006*\x)*sin(\x)});\node[below left]at(0,0){$O$};}
以$(r_0,\varphi_0)$为初始点积分，$\int_{\varphi_0}^{\varphi}\dd\phi=\tan\psi\int_{r_0}^{r}\dd q/q=\tan\psi[\ln q]_{r_0}^{r}$，得$\varphi-\varphi_0=\tan\psi\ln(r/r_0)$，即图\ref{fig:o1centrala}的对数螺线。向内支的$\cos\psi<0$，须保留$\tan\psi$的符号。若$b=n_aa$，则$\psi=\pi/2$，光线可保持任意给定半径的圆周轨道；此时不能再除以$\dd r/\dd s$，应直接使用切向关系。该折射率在$r=0$奇异，轨迹仅适用于题定介质区域。
\end{solution}

\begin{exercise}\label{ex:o1phasors}
同频同偏振两束光复振幅为$A$和$2A\ee^{i\delta}$。求合振幅的大小、相位以及最大最小强度，并说明相位求值应怎样处理象限。
\Diagram{fig:o1phasorq}{两相量大小之比为1:2}{\draw[axis](0,0)--(1.5,0)node[right]{$A$};\draw[axis](0,0)--(1.5,2.598076)node[right]{$2A\ee^{i\delta}$};\draw(.6,0)arc(0:60:.6);\node at(.8,.4){$\delta$};}
\end{exercise}
\begin{solution}
以第一束的相位为零，图\ref{fig:o1phasorq}中合复振幅为
\[
 \frac{\mathcal E}A=1+2\ee^{i\delta}
 =1+2\cos\delta+2i\sin\delta.
\]
因此$|\mathcal E|^2/A^2=(1+2\cos\delta)^2+4\sin^2\delta=5+4\cos\delta$，合振幅大小为$A\sqrt{5+4\cos\delta}$，相位为
\[
 \phi=\operatorname{atan2}(2\sin\delta,1+2\cos\delta).
\]
该函数由实、虚两分量共同确定象限；单独取反正切商只能给出相差整数$\pi$的一组可能值。
\Diagram{fig:o1phasora}{合相量实部可能为负，单值反正切会丢失象限}{\begin{scope}[scale=.7]\draw[axis](-2.2,0)--(3.4,0)node[right]{实部};\draw[axis](0,-2.3)--(0,2.8)node[above]{虚部};\draw[curve](1,0)circle(2);\draw[->,thick](0,0)--(-.5,1.322876);\node[left]at(-.5,1.322876){$\mathcal E$};\end{scope}}
图\ref{fig:o1phasora}的合相量末端以$(A,0)$为圆心、$2A$为半径。原点位于圆内，合相量会经过左半平面，但不会为零，因为两束振幅不等。若$A$对应单束强度$I_1$，则
\[
 I_{\max}=9I_1\quad(\delta=2m\pi),\qquad
 I_{\min}=I_1\quad(\delta=(2m+1)\pi),
\]
可见度为$(9-1)/(9+1)=4/5$。合相量越过虚轴、实部变号时，相位连续经过$\pm\pi/2$；越过负实轴时，取值区间为$(-\pi,\pi]$的主值相位才会发生$2\pi$跳变。本题在$\delta=\pi$通过负实轴，可按连续相位消去这种记号上的跳变。场本身始终连续且不为零，相位也始终有定义。
\end{solution}


\chapter{双束干涉、薄膜与多束相干叠加}
\section{杨氏干涉与光程差}\label{th:o2young}
两狭缝位于$x=\pm d/2$，屏距$L$，观测点横坐标$x$。定义下标$+$对应$+d/2$的缝，则传播距离$r_\pm=\sqrt{L^2+(x\mp d/2)^2}$。取程差$\Delta=r_--r_+$，当$L\gg d,|x|$时，用$\sqrt{L^2+a^2}\simeq L+a^2/(2L)$展开，得到
\[
 \Delta\simeq\frac{(x+d/2)^2-(x-d/2)^2}{2L}=\frac{dx}{L}.
\]
如果两缝初相位相同、同偏振，则亮纹$x_m=m\lambda L/d$，暗纹$x_{m+1/2}=(m+1/2)\lambda L/d$，相邻亮纹间距$\lambda L/d$。这里$\lambda$是在缝屏间介质中的波长；若用真空波长，则需把几何程差乘介质折射率。

一般双束强度为
\begin{equation}\label{eq:o2two}
 I=I_1+I_2+2\sqrt{I_1I_2}\,|\gamma_{12}|\cos\delta,
\end{equation}
其中$\gamma_{12}$表示互相干程度，已把其相位并入$\delta$。完全相干且偏振相同时$|\gamma_{12}|=1$，可见度$\mathcal V=2\sqrt{I_1I_2}/(I_1+I_2)$。若偏振单位矢量不同，还要把交叉项乘它们的内积；正交偏振即使同频同源也不出现普通强度干涉条纹。

\Diagram{fig:o2young}{双缝到屏上$P$点的两条路径}{\draw[curve](0,-1.6)--(0,1.6);\draw[curve](5,-2)--(5,2.3);\fill(0,.7)circle(1.5pt)node[left]{$S_+$};\fill(0,-.7)circle(1.5pt)node[left]{$S_-$};\fill(5,1.3)circle(1.5pt)node[right]{$P(x)$};\draw[curve](0,.7)--(5,1.3);\draw[curve](0,-.7)--(5,1.3);\draw[guide](0,0)--(5,0);\node[below]at(2.5,0){$L$};}

\begin{example}\label{ex:o2shift}
点光源位于双缝前方距离$a$、横坐标$x_s$处，屏在后方距离$L$。在$S_+$前插入正常入射薄片，附加光程$\ell=(n-1)h$。求条纹位置与中央亮纹的移动方向。
\Diagram{fig:o2shiftq}{离轴源和单臂薄片共同改变到达屏幕的总光程差}{\draw[axis](-2,0)--(5,0)node[right]{$z$};\fill(-1.6,1)circle(1.5pt)node[left]{$x_s$};\fill(0,.6)circle(1.5pt)node[above]{$S_+$};\fill(0,-.6)circle(1.5pt)node[below]{$S_-$};\draw[curve](-1.6,1)--(0,.6)--(4,1.1);\draw[curve](-1.6,1)--(0,-.6)--(4,1.1);\draw[curve](.5,.45)rectangle(.8,.95);\draw[curve](4,-1.5)--(4,1.8);}
\end{example}
\begin{solution}
图\ref{fig:o2shiftq}中，取两缝横坐标为$x_\pm=\pm d/2$，程差统一定义为“$-$臂减$+$臂”。源到两缝的距离为$\sqrt{a^2+(x_s\mp d/2)^2}$。用$\sqrt{a^2+q^2}\simeq a+q^2/(2a)$展开，相减得
\[
 r_{s-}-r_{s+}\simeq
 \frac{(x_s+d/2)^2-(x_s-d/2)^2}{2a}=\frac{dx_s}a.
\]
同理缝到屏点$x$的程差为$dx/L$。薄片加在$+$臂，使该臂增加$\ell=(n-1)h$，故总程差为
\[
 \Delta=d(x_s/a+x/L)-\ell.
\]
\Diagram{fig:o2shifta}{中央纹位置由入射波面倾斜和薄片补偿相位共同确定}{\draw[axis](-2.6,0)--(3.7,0)node[right]{$x$};\foreach\x in{-2,-1,0,1,2,3}{\draw[curve](\x,-.3)--(\x,.3);}\draw[->,thick](0,.7)--(1.4,.7);\node[above]at(.7,.7){薄片导致$+L\ell/d$};\draw[->](-.2,-.8)--(-1.4,-.8);\node[below]at(-.8,-.8){$x_s>0$导致负向};}
图\ref{fig:o2shifta}中相长条件$\Delta=m\lambda$给
\[
 x_m=\frac Ld(m\lambda+\ell)-\frac{Lx_s}a,
 \qquad x_{m+1}-x_m=\frac{\lambda L}d.
\]
插片引起位移$+L\ell/d$，中央亮纹向$+$缝侧移动；源向正向移则引起$-Lx_s/a$，方向相反。两种位移方向均由所选程差的符号确定。若希望插片后原中央亮纹仍位于$x=0$，可令$x_s=a\ell/d$补偿。该式要求源、缝、屏之间均近轴；大离轴时先写四段准确距离，再加同一臂中的附加光程，不能只保留一个线性位移项。
\end{solution}

\begin{example}\label{ex:o2lens}
双缝在薄透镜前方，后焦面记录远场干涉。证明条纹间距为$\lambda f/d$，并解释为什么不能再随意给两条光线各加一段“透镜厚度光程”。
\Diagram{fig:o2lensq}{透镜把同一出射方向的平行光会聚到后焦面}{\draw[axis](-.3,0)--(6,0);\fill(0,.8)circle(1.4pt)node[left]{$S_+$};\fill(0,-.8)circle(1.4pt)node[left]{$S_-$};\draw[<->](2.5,-1.6)--(2.5,1.6);\draw[curve](5,-1.8)--(5,1.8);\node[above]at(5,1.8){后焦面};\draw[guide](2.5,0)--(5,0);\node[below]at(3.75,0){$f$};}
\end{example}
\begin{solution}
在图\ref{fig:o2lensq}中取两缝沿同一角度$\theta$发出的平行光，垂直该方向的参考波面上的程差为$d\sin\theta$。理想透镜将此平面波聚焦到$x=f\tan\theta\simeq f\theta$。等相位波面到其理想像点的总光程相等，透镜相位已包括在这个成像性质中。因此总相位差仍是$2\pi d\sin\theta/\lambda$，近轴亮纹$x_m=m\lambda f/d$。
\Diagram{fig:o2lensa}{在同一个参考波面比较相位，再由透镜等光程传播到焦点}{\draw[axis](-.3,0)--(6,0);\draw[<->](2.5,-1.7)--(2.5,1.7);\draw[->](0,.6)--(2.5,1.1)--(5,.7);\draw[->](0,-.8)--(2.5,-.3)--(5,.7);\draw[guide](1.6,-1.1)--(1.16,1.1);\fill(5,.7)circle(1.4pt)node[right]{$P$};\node[above]at(1.2,1.2){参考波面};}
图\ref{fig:o2lensa}的光线在透镜前具有相同倾角，到透镜后的倾角不同。若先使用理想透镜聚焦的等光程性质，再额外添加某条光线的玻璃厚度，就重复计算了透镜相位。另一常见配置用两束夹角$2\theta$的平面波叠加，$\Delta=2x\sin\theta$；当$2\theta\simeq d/f$时，同样得到$\lambda f/d$。
\end{solution}

\subsection{焦点相移的使用条件}\label{th:o2focus}
几何光程在焦点附近不足以描述波的相位。近轴高斯束的轴上相位含$-\arctan(z/z_R)$，光束由远负轴穿过焦点到远正轴，这一项总变化为$-\pi$；二维横向均受限的普通点聚焦由此产生总$\pi$焦点相移。只在一个横向方向聚焦的柱面束对应总$\pi/2$。符号还随时间指数约定改变，但两束是否多跨一个焦散的相对相移可直接比较。

该相移来自波的会聚—发散传播，不是介质界面反射的半波损失。若两条干涉光路经历相同焦点过程，相移抵消；只有两路焦散数或维数不同才需保留。不能在每经过一片透镜时都机械加$\pi$。

\begin{example}\label{ex:o2mirror}
点光源到两平面镜交线附近的距离为$l$，两镜夹角为很小的$\alpha$。两反射光在距虚源平面$L$的屏上重叠。求虚光源间距与干涉间距，并说明反射相位的作用。
\Diagram{fig:o2mirrorq}{双镜产生同一光源的两个虚像}{\draw[curve](0,0)--(3,.12);\draw[curve](0,0)--(-3,.12);\fill(0,1.8)circle(1.5pt)node[above]{$S$};\draw[guide](0,1.8)--(1,.04);\draw[guide](0,1.8)--(-1,.04);\node at(1.2,-.4){两镜法线夹角$\alpha$};}
\end{example}
\begin{solution}
反射像相对镜面与原物点对称。转动镜面$\alpha$使像的方向转动$2\alpha$，所以图\ref{fig:o2mirrorq}对应的两个虚源在以交线为中心、半径$l$的圆上夹角$2\alpha$，其距离$d=2l\sin\alpha\simeq2l\alpha$。
\Diagram{fig:o2mirrora}{等效双虚源的路径构成普通双源干涉}{\fill(0,.5)circle(1.5pt)node[left]{$S_1$};\fill(0,-.5)circle(1.5pt)node[left]{$S_2$};\draw[curve](5,-1.5)--(5,1.7);\draw[curve](0,.5)--(5,.8);\draw[curve](0,-.5)--(5,.8);\draw[guide](0,0)--(5,0);\node[below]at(2.5,0){$L$};}
图\ref{fig:o2mirrora}给近轴间距$\Delta x=\lambda L/d\simeq\lambda L/(2l\alpha)$。若两反射使用相同材料、相近入射角、同一偏振，两路反射相移近似相同而抵消。若两镜的反射系数不同，应把它们的相位差加入总相位，不能只凭“都反射一次”保证中央纹一定亮。
\end{solution}

\begin{example}\label{ex:o2biprism}
菲涅耳双棱镜的两个小棱角均为$\alpha$，折射率$n$，点光源位于棱镜前$l$。求两虚源间距；屏距虚源平面$L$时求条纹间距。
\Diagram{fig:o2biprismq}{薄双棱镜把两部分波面向中间偏折}{\fill(0,0)circle(1.4pt)node[left]{$S$};\draw[curve](2,-1.6)--(2.5,0)--(2,1.6)--cycle;\draw[->](0,0)--(2,.8)--(4.6,.35);\draw[->](0,0)--(2,-.8)--(4.6,-.35);\node[above]at(2.3,1.6){$n,\alpha$};}
\end{example}
\begin{solution}
小角棱镜偏向角$\epsilon\simeq(n-1)\alpha$，可由两个界面的近轴折射定律相减得到：内部折射角满足$r_1+r_2=\alpha$，总偏向$i_1+i_2-\alpha\simeq n(r_1+r_2)-\alpha$。
\Diagram{fig:o2biprisma}{向后延长偏折光线，得到上下两个虚源}{\fill(0,1.2)circle(1.4pt)node[left]{$S_1$};\fill(0,-1.2)circle(1.4pt)node[left]{$S_2$};\draw[guide](0,1.2)--(2,.8);\draw[guide](0,-1.2)--(2,-.8);\draw[->](2,.8)--(5,.2);\draw[->](2,-.8)--(5,-.2);\draw[<->](0,-1.2)--(0,1.2);\node[below]at(1,-1.5){$d\simeq2l\epsilon$};}
按图\ref{fig:o2biprisma}，每个虚源相对原源横移约$l\epsilon$，故$d\simeq2l(n-1)\alpha$。双源干涉间距$\lambda L/[2l(n-1)\alpha]$。这个结果假定两束在屏幕上实际有重叠区域，且棱镜孔径不截断所需光束；没有重叠时，虚源形式存在也不会产生可见干涉。
\end{solution}

\begin{example}\label{ex:o2lloyd}
点光源在平面镜上方高度$h$，屏距沿镜方向为$L\gg h$。直接光和一次反射光在屏上重叠。若反射产生$\pi$相移，求暗亮纹位置及镜面延长线附近条纹的性质。
\Diagram{fig:o2lloydq}{劳埃德镜的直接光与反射光}{\draw[curve](-.5,0)--(5,0);\fill(0,.8)circle(1.5pt)node[left]{$S$};\draw[curve](5,0)--(5,2.5);\fill(5,1.3)circle(1.5pt)node[right]{$P(y)$};\draw[->](0,.8)--(5,1.3);\draw[->](0,.8)--(1.904762,0)--(5,1.3);}
\end{example}
\begin{solution}
将图\ref{fig:o2lloydq}的反射光反向延长，等效虚源$S'$在$-h$，两源间距$2h$。几何程差$r_{S'P}-r_{SP}\simeq2hy/L$，总相位差为$2\pi(2hy/L)/\lambda+\pi$。
\Diagram{fig:o2lloyda}{镜像源使几何程差为零处因反射相移而暗}{\draw[guide](-.3,0)--(5,0);\fill(0,.8)circle(1.4pt)node[left]{$S$};\fill(0,-.8)circle(1.4pt)node[left]{$S'$};\draw[curve](0,.8)--(5,1.3);\draw[guide](0,-.8)--(5,1.3);\draw[curve](5,0)--(5,2);\node[below]at(2.5,0){镜面延长线};}
因此暗纹$y_m=m\lambda L/(2h)$，亮纹$y_m=(m+1/2)\lambda L/(2h)$。图\ref{fig:o2lloyda}中$y=0$几何程差为零却是暗纹。公式适用于两束充分重叠的近轴区域；镜边遮挡及掠射反射系数变化会使边缘条纹变形。
\end{solution}

\section{薄膜的斜入射程差与反射相位}\label{th:o2film}
厚度$h$、折射率$n$的平行薄膜内折射角为$r$。相邻两束反射光一束在上面反射，另一束进入膜内往返一次再出射。比较必须取出射方向垂直的同一波面。膜内光程为$2nh/\cos r$，还要减去外部波面截取引起的相位，化简得到
\begin{equation}\label{eq:o2film}
 \Delta_{\rm geom}=2nh\cos r.
\end{equation}
也可用膜内法向波矢$k_z=2\pi n\cos r/\lambda_0$，一次往返相位为$2k_zh$，立即得到同式。真实膜内路长$2h/\cos r$与干涉等效程差$2h\cos r$不是同一量。

若两路反射相位差为$\pi$，反射暗纹满足$2nh\cos r=m\lambda_0$，亮纹满足$(m+1/2)\lambda_0$；若相位差为0则互换。非全反射、非吸收介质的反射系数为实数，符号负给$\pi$相移，但$p$偏振越过布儒斯特角时符号也可能改变。亮暗条件应采用相应偏振、介质排列和入射角的反射系数。

\begin{example}\label{ex:o2wedge}
两块玻璃间有折射率$n$的薄劈尖，夹角$\alpha\ll1$，厚度$h(x)=h_0+\alpha x$。近乎垂直照射，反射的两束有一次相对半波损失。求纹距及由条纹反求夹角的方法。
\Diagram{fig:o2wedgeq}{劈尖厚度随$x$线性增长}{\draw[curve](0,0)--(5,0);\draw[curve](0,.2)--(5,1);\draw[axis](0,-.5)--(5.6,-.5)node[right]{$x$};\draw[<->](4,0)--(4,.84)node[midway,right]{$h(x)$};\draw[->](2,2)--(2,.6);}
\end{example}
\begin{solution}
图\ref{fig:o2wedgeq}的暗纹$2n(h_0+\alpha x_m)=m\lambda_0$，相邻两式相减给$2n\alpha(x_{m+1}-x_m)=\lambda_0$，所以纹距$\Delta x=\lambda_0/(2n\alpha)$。
\Diagram{fig:o2wedgea}{同一厚度的横截线对应一条干涉纹}{\draw[curve](0,0)--(5,0)--(5,1)--(0,.2)--cycle;\foreach\x in{.6,1.5,2.4,3.3,4.2}{\draw[guide](\x,0)--(\x,{.2+.16*\x});}\draw[<->](1.5,-.4)--(2.4,-.4)node[midway,below]{$\Delta x$};}
若在长度$D$内跨过$N$个完整纹距，则$\alpha=N\lambda_0/(2nD)$。图\ref{fig:o2wedgea}的纹序受$h_0$影响，纹距不受影响。若接触端$h_0=0$，此端反射为暗；若端点混入胶层或尘粒，就不能直接把可见第一条暗纹当$m=0$。
\end{solution}

\begin{example}\label{ex:o2rings}
曲率半径$R$的平凸透镜凸面置于平板上，间隙折射率$n$，中心残余间隙$h_0$。垂直单色光形成牛顿环。求反射暗环半径及用两条暗环测$R$的方法。
\Diagram{fig:o2ringsq}{球面和平板构成厚度随半径增加的薄膜}{\draw[curve](-2.5,0)--(2.5,0);\draw[curve,domain=-2.2:2.2,samples=70]plot(\x,{.12+.16*\x*\x});\draw[guide](0,0)--(0,1.8);\draw[<->](1.8,0)--(1.8,.638)node[right]{$h(r)$};}
\end{example}
\begin{solution}
图\ref{fig:o2ringsq}中球面几何给$h=h_0+R-\sqrt{R^2-r^2}\simeq h_0+r^2/(2R)$。一次相对半波损失使反射暗纹$2nh=m\lambda_0$，所以$r_m^2=R(m\lambda_0/n-2h_0)$，只取右侧非负的级次。
\Diagram{fig:o2ringsa}{等厚圆环的半径平方等间隔}{\foreach\r in{.6,.95,1.2,1.4,1.58}{\draw[curve](0,0)circle(\r);}\draw[axis](0,0)--(2.6,0)node[right]{$r$};\node[above]at(0,1.7){$r_{m+p}^2-r_m^2=pR\lambda_0/n$};}
取相隔$p$级的两环，图\ref{fig:o2ringsa}给$R=n(r_{m+p}^2-r_m^2)/(p\lambda_0)$，消除了未知$h_0$。若中心真正接触，$m=0$对应中心暗点；透射图样在理想无损、相同孔径收集条件下互补，中心为亮。半径不等间隔，等间隔的是半径平方。
\end{solution}

\begin{example}\label{ex:o2michelson}
迈克耳孙干涉仪的两等效反射面平行，轴向间隔$d$。求等倾条纹条件、近轴环半径和一面镜移动$\Delta d$时中心跨过的条纹数。
\Diagram{fig:o2michelsonq}{将两臂展开为两平行虚反射面}{\draw[curve](0,-1)--(0,1.6);\draw[curve](1.5,-1)--(1.5,1.6);\draw[axis](2.1,.3)--(5,1.4)node[right]{观察方向};\draw[guide](2.1,.3)--(5,.3);\node at(3.2,.6){$\theta$};\draw[<->](0,-1.4)--(1.5,-1.4)node[midway,below]{$d$};}
\end{example}
\begin{solution}
图\ref{fig:o2michelsonq}中相位差传播部分为$2kd\cos\theta$，若两臂固有附加相位已补偿，亮纹满足$2d\cos\theta=m\lambda$。中心$\theta=0$的级次最大，离轴$\cos\theta$下降，级次依次减小。
\Diagram{fig:o2michelsona}{后焦面将等倾角映为圆环，中心级次最高}{\foreach\r in{.45,.85,1.15,1.4}{\draw[curve](0,0)circle(\r);}\fill(0,0)circle(1.5pt);\draw[axis](0,0)--(2.4,0)node[right]{$\rho$};\node[right]at(2.5,1){$\rho\simeq f\theta$};}
若$2d/\lambda=m_0$为整数，近轴$\cos\theta\simeq1-\theta^2/2$，第$p$个离轴同类亮环满足$\rho_p^2\simeq pf^2\lambda/d$，如图\ref{fig:o2michelsona}。镜移改变往返程差$2\Delta d$，中心通过条纹数$N=2|\Delta d|/\lambda$。实际非整数中心相位使所有环的半径平方有相同偏移；环间平方差仍遵守该式。
\end{solution}

\section{多束几何级数与分辨能力}\label{th:o2multiple}
等振幅、相邻相位差$\delta$的$N$束光（以下取$N\geq2$）有
\[
 \mathcal E=A\sum_{j=0}^{N-1}\ee^{ij\delta}
 =A\ee^{i(N-1)\delta/2}\frac{\sin(N\delta/2)}{\sin(\delta/2)},
\]
因此$I=I_1[\sin(N\delta/2)/\sin(\delta/2)]^2$。在$\delta=2m\pi$分子分母同时为零，须取极限，合振幅$NA$，主极大强度$N^2I_1$。两个主极大之间有$N-1$个零点$\delta=2\pi j/N$，$j=1,\ldots,N-1$，其间有$N-2$个次极大。

主峰邻近令$\delta=2m\pi+\epsilon$，首零点$|\epsilon|=2\pi/N$，这是主极大到首零点的半宽，不是半高宽。大$N$时强度归一化曲线近似$(\sin x/x)^2$，$x=N\epsilon/2$。对它求导，
\[
 \frac{\dd}{\dd x}\left(\frac{\sin x}{x}\right)^2
 =2\frac{\sin x}{x}\frac{x\cos x-\sin x}{x^2}.
\]
旁瓣内部$\sin x\ne0$，故极值条件为$x\cos x-\sin x=0$，即$\tan x=x$。第一正根在$\pi$与$2\pi$之间，约4.493，强度比约0.0472。将旁瓣位置粗略取$3\pi/2$得到$4/(9\pi^2)\simeq0.0450$，只能作估计。

\subsection{光栅式相位与瑞利判据}\label{th:o2resolution}
若相邻光束间几何间距为$d$，$\delta=2\pi d\sin\theta/\lambda$，主峰$d\sin\theta_m=m\lambda$。在非掠射主峰附近，对$\delta$求角导数，$\dd\delta/\dd\theta=2\pi d\cos\theta_m/\lambda$；将首零相位偏移$2\pi/N$除以它，得首零角半宽$\Delta\theta_0\simeq\lambda/(Nd\cos\theta_m)$。固定非零级次$m$，对主峰方程微分，$d\cos\theta_m\,\dd\theta=m\,\dd\lambda$，故两近邻波长的角分离为$|\Delta\theta|\simeq|m|\Delta\lambda/(d\cos\theta_m)$。按瑞利条件令它等于首零半宽，得到
\[
 \mathcal R=\frac\lambda{\Delta\lambda}=|m|N.
\]
零级$m=0$没有这种角色散，不能用它作上述分光。两谱线需在测量中互不相干地相加强度；瑞利判据是一种分辨约定，实际双峰边界还取决于线形、强度比与观测条件。
\begin{example}\label{ex:o2three}
三束等振幅、等相邻相位差光叠加，求主极大、零点与次极大，并画出一个周期内强度。
\Diagram{fig:o2threeq}{三个等长相量相邻夹角相同}{\draw[axis](0,0)--(2,0);\draw[axis](0,0)--(1.285575,1.532089);\draw[axis](0,0)--(-.347296,1.969616);\node[right]at(2,0){$A$};\node[above right]at(1.285575,1.532089){$A\ee^{i\delta}$};\node[left]at(-.347296,1.969616){$A\ee^{2i\delta}$};}
\end{example}
\begin{solution}
图\ref{fig:o2threeq}的三相量依次为$A,A\ee^{i\delta},A\ee^{2i\delta}$。提出公共相位得
\[
 \mathcal E=A\ee^{i\delta}(\ee^{-i\delta}+1+\ee^{i\delta})
 =A\ee^{i\delta}(1+2\cos\delta),
\]
所以$I/I_1=(1+2\cos\delta)^2$。零场条件为$\cos\delta=-1/2$，在$0\leq\delta\leq2\pi$内给$2\pi/3$、$4\pi/3$。
\Diagram{fig:o2threea}{三束干涉的两个零点和中央次极大}{\draw[axis](0,0)--(6.7,0)node[right]{$\delta$};\draw[axis](0,0)--(0,3.3)node[above]{$I/I_1$};\draw[curve,domain=0:6.283,samples=120]plot(\x,{(1+2*cos(\x r))^2/3});\foreach\x/\t in{2.094/{2\pi/3},3.142/\pi,4.189/{4\pi/3},6.283/{2\pi}}{\node[below]at(\x,0){$\t$};}\node[left]at(0,3){$9$};}
为区别极大与零点，对强度求导：
\[
 \frac1{I_1}\frac{\dd I}{\dd\delta}
 =-4(1+2\cos\delta)\sin\delta.
\]
除零场点外，驻点为$0,\pi,2\pi$。在$0,2\pi$三相量同向，强度为9；$\pi$处相量依次为$A,-A,A$，余下$A$，强度为1。由$\delta=\pi+\epsilon$时$1+2\cos\delta=-1+\epsilon^2+O(\epsilon^4)$，得到$I/I_1=1-2\epsilon^2+\cdots$，因此$\pi$确为次极大。

图\ref{fig:o2threea}的一个周期含一个次极大和两个完全暗点。周期平均强度为$\langle I\rangle/I_1=1+4\langle\cos\delta\rangle+4\langle\cos^2\delta\rangle=3$，等于三束单独强度之和；峰值9来自相干重排，不能解释为总能量增为三倍。
\end{solution}

\begin{example}\label{ex:o2unequal}
三束光的振幅依次为$A,2A,3A$，相邻相位差$\delta$。求强度、全局极值和可见度，判断是否有完全暗点。
\Diagram{fig:o2unequalq}{三个相量的长度为1、2、3，方向按$\delta$递增}{\draw[axis](0,0)--(1,0)node[right]{$1$};\draw[axis](0,0)--(1.285575,1.532089)node[right]{$2$};\draw[axis](0,0)--(-.520945,2.954423)node[left]{$3$};}
\end{example}
\begin{solution}
图\ref{fig:o2unequalq}给$\mathcal E/A=1+2\ee^{i\delta}+3\ee^{2i\delta}$。乘共轭时，每一对相量只产生一项余弦：
\begin{align*}
 I/I_1&=1+4+9+4\cos\delta+12\cos\delta+6\cos2\delta\\
 &=12\cos^2\delta+16\cos\delta+8.
\end{align*}
令$x=\cos\delta\in[-1,1]$，这是开口向上的二次式，顶点$x=-2/3$，最小值$8/3$。端点$x=1$强度36，$x=-1$强度4，所以全局最大36。
\Diagram{fig:o2unequala}{不等振幅的强度曲线不接触零轴}{\draw[axis](0,0)--(6.8,0)node[right]{$\delta$};\draw[axis](0,0)--(0,3.4)node[above]{$I/I_1$};\draw[curve,domain=0:6.283,samples=140]plot(\x,{(12*cos(\x r)^2+16*cos(\x r)+8)/12});\node[left]at(0,3){$36$};\draw[guide](0,.2222)--(6.3,.2222);\node[right]at(6.3,.35){$8/3$};}
图\ref{fig:o2unequala}没有完全暗点。全局可见度$(36-8/3)/(36+8/3)=25/29$。在$\delta=\pi$虽有局部次极大4，但它不决定全局可见度；求导时要区分关于$x$的端点和关于$\delta$的驻点。
\end{solution}

\begin{example}\label{ex:o2spacing}
三等振幅缝位置为$0,d,5d$。令$\delta=2\pi d\sin\theta/\lambda$，求强度和轴旁第一个零点。
\Diagram{fig:o2spacingq}{三缝的相邻间距为$d$和$4d$}{\draw[axis](-.3,0)--(6.1,0)node[right]{$x/d$};\foreach\x/\t in{0/0,1/1,5/5}{\draw[very thick](\x,-.4)--(\x,.5);\node[above]at(\x,.5){$\t$};}\draw[<->](1,-.8)--(5,-.8)node[midway,below]{$4d$};}
\end{example}
\begin{solution}
图\ref{fig:o2spacingq}给$\mathcal E/A=1+\ee^{i\delta}+\ee^{i5\delta}$，乘共轭得到$I/I_1=3+2\cos\delta+2\cos4\delta+2\cos5\delta$。完全相长在$\delta=2k\pi$，强度9。
\Diagram{fig:o2spacinga}{完全相消时三个等长相量构成正三角形}{\draw[->,thick](0,0)--(2,0)node[midway,below]{$1$};\draw[->,thick](2,0)--(1,1.732)node[midway,right]{$\ee^{i\delta}$};\draw[->,thick](1,1.732)--(0,0)node[midway,left]{$\ee^{i5\delta}$};}
三个等长矢量和为零必须彼此相差$2\pi/3$。由图\ref{fig:o2spacinga}，含有相位0和$\delta$的一对须满足$\delta=\pm2\pi/3$模$2\pi$；此时$5\delta$确为另一方向。因此首个正零点$\delta=2\pi/3$，$\sin\theta_0=\lambda/(3d)$。这一零点只有在$\lambda\leq3d$时属于传播角域；等号时位于掠射方向。不等间距不能直接把$N=3$代入等间距几何级数来解释全部旁瓣，但本题首零点恰好相同。
\end{solution}

\section{相干叠加的时间平均}
\subsection{交叉项由什么决定}
设两复包络为$\mathcal E_1(t),\mathcal E_2(t)$，共同载频为$\omega_0$。探测器先平均光学快振荡，再按曝光时间平均包络，强度含$\langle|\mathcal E_1|^2\rangle+\langle|\mathcal E_2|^2\rangle+2\operatorname{Re}\langle\mathcal E_1\mathcal E_2^*\rangle$。定义归一化互相干度
\[
 \gamma_{12}=\frac{\langle\mathcal E_1\mathcal E_2^*\rangle}
 {\sqrt{\langle|\mathcal E_1|^2\rangle\langle|\mathcal E_2|^2\rangle}},
 \quad |\gamma_{12}|\leq1.
\]
最后不等式由柯西不等式给出，保证总强度非负。相干度模控制对比度，其幅角控制条纹偏移；整体初相位虽可任取，两臂相对相位不能删去。

两个独立光源即使名义频率相同，相位也可能随机变化，使长曝光交叉项消失。稍异频而相位关系稳定的两束光，则可在短曝光中形成移动条纹；曝光足够长时，交叉项又会平均为零。是否记录到条纹，取决于频谱、相位变化与探测时间的共同作用。

\subsection{时间延迟与光程差}
同一稳定光源分成两路，一路相对另一路延迟$\tau$。取包含载频的复场$E(t)$，其互相关为$\langle E(t)E^*(t-\tau)\rangle$；不同频率分量互不相干时，这个量就是功率谱乘$\ee^{-i\omega\tau}$的频率积分。若使用前文除去载频的包络$\mathcal E(t)$，谱的变量应改为偏移频率$\Omega=\omega-\omega_0$，两种相关函数只差载频相位因子$\ee^{-i\omega_0\tau}$，模相同。矩形谱给sinc包络，双线谱给拍频包络，高斯谱给高斯包络。计算时先求每一频率的同频两路干涉，再将不同频率的强度相加。

在真空静态路径中，传播延迟为光程差除以$c$；色散材料中相位差为$\phi(\omega)$，一阶频率展开的群延迟是$\dd\phi/\dd\omega$。两臂在中心频率相位补偿为零，并不保证宽带光的群延迟也为零。白光干涉需要各频率的交叉相位在谱宽内不过度散开，因此一般要同时补偿中心相位和主要色散。

\subsection{空间延伸与相干面积}
扩展非相干光源的不同源元之间没有稳定相位关系，每个源元却可同时照亮两个观察点，产生同频相关项。设窄带光的波长为$\lambda$，源方向的近轴角坐标为$\boldsymbol\theta=(\theta_x,\theta_y)$，观察基线为$\vct d=(d_x,d_y)$。各源元的亮度权重为$B(\boldsymbol\theta)$，且两点到源的距离远大于源宽和基线。去掉与源元无关的共同相位，归一化互相干度为
\[
 \gamma(\vct d)=
 \frac{\iint B(\boldsymbol\theta)
 \ee^{ik(d_x\theta_x+d_y\theta_y)}\dd\theta_x\dd\theta_y}
 {\iint B(\boldsymbol\theta)\dd\theta_x\dd\theta_y},
 \qquad k=\frac{2\pi}{\lambda}.
\]
这里积分的是每个源元的两路相关项；独立源元之间的交叉项在平均后为零。它不是把独立源元的确定电场直接相加。

取图\ref{fig:o2rectcoherence}中的均匀矩形角源，两个方向的全角宽分别为$\Theta_x,\Theta_y$。积分可分离。对$i=x,y$，当$d_i\ne0$时，
\[
 \frac1{\Theta_i}\int_{-\Theta_i/2}^{\Theta_i/2}
 \ee^{ikd_i\theta_i}\dd\theta_i
 =\left[\frac{\ee^{ikd_i\theta_i}}{ikd_i\Theta_i}\right]_{-\Theta_i/2}^{\Theta_i/2}
 =\operatorname{sinc}\left(\frac{\pi\Theta_i d_i}{\lambda}\right);
\]
$d_i=0$时直接积分为1。因此
\[
 \gamma(d_x,d_y)=
 \operatorname{sinc}\left(\frac{\pi\Theta_xd_x}{\lambda}\right)
 \operatorname{sinc}\left(\frac{\pi\Theta_yd_y}{\lambda}\right).
\]
\Diagram{fig:o2rectcoherence}{矩形角源与基线平面上的首零范围；两个方向的宽度分别决定相干尺度}{
\begin{scope}
\fill[gray!15](-1.2,-.7)rectangle(1.2,.7);
\draw(-1.2,-.7)rectangle(1.2,.7);
\draw[axis](-1.9,0)--(2,0)node[right]{$\theta_x$};
\draw[axis](0,-1.3)--(0,1.8)node[above]{$\theta_y$};
\draw[<->](-1.2,-1)--(1.2,-1)node[midway,below]{$\Theta_x$};
\draw[<->](-1.5,-.7)--(-1.5,.7)node[midway,left,fill=white,inner sep=1pt]{$\Theta_y$};
\end{scope}
\begin{scope}[xshift=5.7cm]
\fill[gray!15](-.95,-.85)rectangle(.95,.85);
\draw[dashed](-.95,-.85)rectangle(.95,.85);
\draw[axis](-1.7,0)--(1.9,0)node[right]{$d_x$};
\draw[axis](0,-1.4)--(0,1.8)node[above]{$d_y$};
\node[above right]at(.95,0){$\lambda/\Theta_x$};
\node[above right]at(.1,.85){$\lambda/\Theta_y$};
\end{scope}}
两路等强时，可见度为$|\gamma|$。沿坐标轴的首零基线分别为$\lambda/\Theta_x$、$\lambda/\Theta_y$；首零围成的中央矩形面积为$4\lambda^2/(\Theta_x\Theta_y)$。若源实际宽度为$w_x,w_y$、距离为$a$，近轴有$\Theta_i\simeq w_i/a$，线源结果见例\ref{ex:o2extendedsource}。

以近轴相关函数的全平面延拓定义积分相干面积$A_c=\iint|\gamma(\vct d)|^2\dd d_x\dd d_y$。分别令$u_i=\pi\Theta_id_i/\lambda$，$\dd d_i=\lambda\dd u_i/(\pi\Theta_i)$，上下限仍为整个实轴，得到
\[
 A_c=\frac{\lambda^2}{\pi^2\Theta_x\Theta_y}
 \left(\int_{-\infty}^{\infty}\operatorname{sinc}^2u\,\dd u\right)^2
 =\frac{\lambda^2}{\Theta_x\Theta_y}.
\]
最后用了习题\ref{ex:o3energy}由帕塞瓦尔关系证明的积分$\int\operatorname{sinc}^2u\,\dd u=\pi$。这个积分面积与首零中央矩形面积采用不同定义；若按指定可见度阈值划区，系数还会改变。源角宽越小，相干基线范围越大；某方向角宽趋零时，该理想模型在对应基线方向上不衰减。

\section{薄膜程差的同波面比较}
膜内折射角为$r$，外部入射角为$i$，膜厚$h$。入射点$A$、底面反射点$B$、再出射点$C$的水平间距$AC=2h\tan r$。第二束在膜内累计光程$2nh/\cos r$，但两束要比较到同一个出射参考波面，第一束从$A$还走了一段$AD=AC\sin i$。若外部折射率$n_0$，相减得到
\begin{align*}
 \Delta&=\frac{2nh}{\cos r}-n_0(2h\tan r)\sin i\\
 &=\frac{2nh}{\cos r}-\frac{2nh\sin^2r}{\cos r}
 =2nh\cos r.
\end{align*}
第二步用$n_0\sin i=n\sin r$。扣除同一出射波面上的参考光程后，程差中便出现$\cos r$因子。

\Diagram{fig:o2filmphase}{$CD$垂直出射方向，$A\to D$光程须从膜内往返光程中扣除}{\draw[curve](-.7,0)--(3.2,0);\draw[curve](-.7,-1)--(3.2,-1);\fill(0,0)circle(1.3pt)node[left]{$A$};\fill(.57735,-1)circle(1.3pt)node[below]{$B$};\fill(1.1547,0)circle(1.3pt)node[below right]{$C$};\draw[->](0,0)--(.57735,-1)--(1.1547,0)--(2.8,1.6453);\draw[->](0,0)--(2,2);\draw[guide](.57735,.57735)--(1.1547,0);\fill(.57735,.57735)circle(1.3pt)node[above left]{$D$};}
图\ref{fig:o2filmphase}中两束出射光平行，$AD$沿第一束方向，$CD$为参考波面。若膜厚、折射率、入射角在观察区域变化，等相位条件相应形成等厚纹或等倾纹。实际观测透镜把哪一组光线汇到同一点，决定条纹的定域位置，不能只凭“是薄膜”就断言条纹一定在膜面。

总干涉相位还含反射相位。采用一致的偏振基时，两路实反射系数同号给相对相位0，异号给$\pi$；全反射的复系数则须取各自辐角。将它与传播相位相加后，才能确定亮暗级次。

\section{焦散相移的波动来由}
在腰面$z=0$取圆对称高斯场$U(r,0)=U_0\exp(-r^2/w_0^2)$。对于$z\ne0$，代入近轴菲涅耳传播积分，并展开$(x-\xi)^2+(y-\eta)^2$。记
\[
 a=\frac1{w_0^2}-\frac{ik}{2z},\qquad
 b_x=-\frac{ikx}{z},\qquad b_y=-\frac{iky}{z}.
\]
两个横向积分分别具有式\eqref{eq:o1gaussintegral}的形式，故
\begin{align*}
 U(r,z)&=\frac{U_0\ee^{ikz}}{i\lambda z}\ee^{ikr^2/(2z)}
 \mathcal G(a,b_x)\mathcal G(a,b_y)\\
 &=\frac{\pi U_0\ee^{ikz}}{i\lambda za}
 \exp\left[\frac{ikr^2}{2z}-\frac{k^2r^2}{4az^2}\right].
\end{align*}
取$z_R=kw_0^2/2$、$\lambda=2\pi/k$，记$\zeta=z/z_R$。则
\[
 aw_0^2=1-\frac{i}{\zeta},\qquad
 \frac{i\lambda za}{\pi}=1+i\zeta,\qquad
 \frac1{1-i/\zeta}=\frac{i\zeta}{1+i\zeta}.
\]
前因子因而为$U_0\ee^{ikz}/(1+i\zeta)$。指数中的两项先提取$r^2/w_0^2$，再通分：
\begin{align*}
 \frac{ikr^2}{2z}-\frac{k^2r^2}{4az^2}
 &=\frac{r^2}{w_0^2}
 \left[\frac i\zeta-\frac1{\zeta^2(1-i/\zeta)}\right]\\
 &=\frac{r^2}{w_0^2}
 \left[\frac i\zeta-\frac{i}{\zeta(1+i\zeta)}\right]
 =-\frac{r^2}{w_0^2(1+i\zeta)}.
\end{align*}
代回得到
\[
 U(r,z)=\frac{U_0\ee^{ikz}}{1+iz/z_R}
 \exp\left[-\frac{r^2}{w_0^2(1+iz/z_R)}\right].
\]
在轴上$r=0$，除去平面传播相位$kz$，剩余相位为$-\arctan(z/z_R)$。从负远处经过腰面到正远处，其总变化为$-\pi$。一维柱面高斯场只有一个高斯积分，前因子为$(1+iz/z_R)^{-1/2}$，额外相位总变化为$-\pi/2$。

\Diagram{fig:o2gouy}{高斯束腰附近，波动相位连续经过焦区}{\draw[axis](-3,0)--(3.2,0)node[right]{$z$};\draw[curve,domain=-2.7:2.7,samples=80]plot(\x,{.45*sqrt(1+\x*\x)});\draw[curve,domain=-2.7:2.7,samples=80]plot(\x,{-.45*sqrt(1+\x*\x)});\draw[guide](0,-1.6)--(0,1.6);\node[above]at(0,1.7){束腰};}
图\ref{fig:o2gouy}中的场在焦区有限、相位连续，几何光学常把整个焦区压成一个点，才把总变化记为一次焦点相移。焦点相移的数值取决于横向受限维数，因此点聚焦与线聚焦不同。两条干涉臂如果经历相同束腰和相同方向传播，这些相移通常相消，不应在每个透镜后重复添加。

\section{有限束数与多次反射腔的比较}
有限$N$束等幅的合场是有限几何级数，零点来自相量正多边形闭合；多次反射腔是无限几何级数，但每次场幅乘$R<1$，闭合多边形的零点结构不再相同。反射率越高，后续束衰减越慢，相干参与的有效往返数大致增多，透射共振更尖。

对正常入射腔，$2nl=m\lambda$，频率共振间隔$c/(2nl)$。若精细度为$\mathcal F_c$，线宽约$c/(2nl\mathcal F_c)$，中心频率$\nu_m=mc/(2nl)$，所以局部频率分辨本领$\nu_m/\delta\nu\simeq m\mathcal F_c$。它与有限光栅的$mN$形式类似，但$\mathcal F_c$由反射损耗和腔镜决定，不能把它严格解释成某个整数条数。

当腔镜不对称、存在吸收或散射，共振峰值一般不再为1，最佳入射耦合也依赖两镜透射率和内部损耗。前述单位峰值只属于理想无损对称腔。分析实测峰高时，应分别计入相位失配、孔径损失和吸收。

\subsection{扫描相位与扫描角度的极值区别}
对两束确定相干场，写$I(x)=I_1(x)+I_2(x)+2\sqrt{I_1(x)I_2(x)}|\gamma(x)|\cos\delta(x)$。若在固定空间点只改变一臂的均匀相位，前面的系数不变，最大、最小强度严格为$I_1+I_2\pm2\sqrt{I_1I_2}|\gamma|$，可见度公式可直接使用。

若沿屏幕扫描，包络、偏振投影和相干度也可能随$x$变化，实际极值应对整个$I(x)$求导。局部条纹可见度要求这些系数在一个纹距内近似不变。衍射包络零点、窄谱共振边缘或强渐变照明区需要保留系数的空间变化。可先在同一$x$固定位置扫描相位，区分包络变化与干涉对比度变化。
\Diagram{fig:o2localvisibility}{缓变包络内的局部明暗纹；靠近包络边缘时需对完整强度求极值}{\draw[axis](-.2,0)--(6.3,0)node[right]{$x$};\draw[axis](0,0)--(0,2.8)node[above]{$I$};\draw[guide,domain=0:5.8,samples=70]plot(\x,{2.4*exp(-.11*(\x-2.4)^2)});\draw[curve,domain=0:5.8,samples=200]plot(\x,{1.2*exp(-.11*(\x-2.4)^2)*(1+.8*cos(300*\x))});}
图\ref{fig:o2localvisibility}的上包络不等于每个实际亮纹都恰好落在$\delta=2m\pi$；相位极值与包络斜率共同决定峰位。这一点在有限$N$束加权、法布里—珀罗有吸收以及不等宽双缝中都会出现。

多束完全相干时还应区别场相加与功率归一化。若场权重为$a_j$，主峰场为$\sum a_j$，而输入功率按$\sum|a_j|^2$计。改变权重后比较峰高，必须保持同一功率条件；否则峰高的改变可能仅来自总照明功率改变。旁瓣降低、主峰变宽和峰值变化应分别从同一个加权和推导，不能用其中一个指标代替全部光学性能。

\section{有限带宽、有限光源与曝光时间}
\begin{example}\label{ex:o2bandwidth}
双束干涉的功率谱在频率$\nu_0-\Delta\nu/2$到$\nu_0+\Delta\nu/2$内恒定，其他处为零。不同频率互不相干，同频两路完全相干且等强。两臂时间差为$\tau$，求干涉可见度包络及首个零点。
\Diagram{fig:o2bandq}{矩形功率谱的每个频率分别产生自己的干涉项}{\draw[axis](0,0)--(5,0)node[right]{$\nu$};\draw[axis](0,0)--(0,2.8)node[above]{功率谱};\draw[curve](1.2,0)--(1.2,2)--(3.8,2)--(3.8,0);\node[below]at(2.5,0){$\nu_0$};\node[above]at(2.5,2){$\Delta\nu$};}
\end{example}
\begin{solution}
图\ref{fig:o2bandq}总干涉项需对频率积分。先计算长度为$L$的均匀区间平均：当$b\ne0$时，
\begin{align}
 \frac1L\int_{-L/2}^{L/2}\cos(a+bs)\dd s
 &=\left[\frac{\sin(a+bs)}{bL}\right]_{-L/2}^{L/2}\notag\\
 &=\frac{2\cos a\sin(bL/2)}{bL}
 =\cos a\operatorname{sinc}(bL/2).\label{eq:o2rectaverage}
\end{align}
当$b=0$时直接积分得$\cos a$，与连续极限相同。对正弦作同一计算，得$\sin a\operatorname{sinc}(bL/2)$，因此复指数的平均为$\ee^{ia}\operatorname{sinc}(bL/2)$。本题以$\nu=\nu_0+s$换元，取$L=\Delta\nu,a=2\pi\nu_0\tau,b=2\pi\tau$，得到归一化干涉项$\cos(2\pi\nu_0\tau)\operatorname{sinc}(\pi\Delta\nu\tau)$。
\Diagram{fig:o2banda}{可见度是谱积分的归一化模，有限带宽使远离零光程差的条纹变弱}{\draw[axis](-3.2,0)--(3.4,0)node[right]{$\Delta\nu\tau$};\draw[axis](0,0)--(0,2.7)node[above]{$\mathcal V$};\draw[curve,domain=.01:3,samples=80]plot(\x,{2.2*abs(sin(180*\x)/(pi*\x))});\draw[curve,domain=-3:-.01,samples=80]plot(\x,{2.2*abs(sin(180*\x)/(pi*\x))});}
故图\ref{fig:o2banda}的可见度$\mathcal V=|\operatorname{sinc}(\pi\Delta\nu\tau)|$，首零$|\tau|=1/\Delta\nu$，对应光程差尺度$c/\Delta\nu$。这一定义依赖所选的矩形谱及“首零”标准；高斯谱没有有限首零，不能把相干长度当所有光源共有的硬截止距离。
\end{solution}

\begin{example}\label{ex:o2extendedsource}
宽$w$的均匀线光源平行于双缝连线，位于缝前距离$a$。源上不同点互不相干，每个点各自照明两缝。双缝距$d$。求近轴干涉可见度及第一消失源宽。
\Diagram{fig:o2sourceq}{扩展光源各点产生相位不同的条纹，相加强度}{\draw[very thick](0,-1.3)--(0,1.3);\node[left]at(0,0){$w$};\fill(3,.6)circle(1.5pt)node[right]{$S_+$};\fill(3,-.6)circle(1.5pt)node[right]{$S_-$};\draw[guide](0,1)--(3,.6);\draw[guide](0,1)--(3,-.6);\node[below]at(1.5,-1.5){$a$};}
\end{example}
\begin{solution}
图\ref{fig:o2sourceq}横坐标$x_s$的源元附加双缝相位$2\pi dx_s/(\lambda a)$。屏上该点的共同相位记为$\delta_0$。在式\eqref{eq:o2rectaverage}已展开的定积分中，取区间长度$w$、相位斜率$2\pi d/(\lambda a)$，均匀源强度平均的交叉项比例为
\[
 \frac1w\int_{-w/2}^{w/2}\cos\left(\delta_0+\frac{2\pi dx_s}{\lambda a}\right)\dd x_s
 =\cos\delta_0\,\operatorname{sinc}\left(\frac{\pi dw}{\lambda a}\right).
\]
\Diagram{fig:o2sourcea}{源角宽越大，不同源点的中央纹偏移越难重合}{\draw[axis](0,0)--(5.3,0)node[right]{屏坐标};\draw[curve,domain=0:4.7,samples=100]plot(\x,{1+.5*cos(300*\x)});\draw[guide,domain=0:4.7,samples=100]plot(\x,{1+.5*cos(300*\x+120)});\node[above]at(2.5,2.1){不同源点条纹示意};}
所以$\mathcal V=|\operatorname{sinc}(\pi dw/(\lambda a))|$，首零$w=\lambda a/d$。图\ref{fig:o2sourcea}表达的是不同条纹强度相加，不能先把独立源点的电场相干相加。源尺寸沿另一方向延伸时，其投影到双缝基线的宽度才进入本式。
\end{solution}

\begin{example}\label{ex:o2movingfringes}
两平面波频率分别为$\omega_1,\omega_2$，波矢分别$\vct k_1,\vct k_2$，偏振相同。求等相位差条纹的运动速度，并讨论曝光时间$T$对可见度的影响。
\Diagram{fig:o2movingq}{两个不同频率和方向的平面波形成移动干涉纹}{\draw[axis](0,0)--(3.2,1.2)node[above]{$\vct k_1$};\draw[axis](0,0)--(3.2,-1.2)node[below]{$\vct k_2$};\foreach\y in{-1,0,1}{\draw[curve](4,\y)--(6,\y);}\draw[axis](6.3,-1)--(6.3,1)node[right]{$\Delta\vct k$};}
\end{example}
\begin{solution}
图\ref{fig:o2movingq}交叉相位为$\Delta\vct k\cdot\vct r-\Delta\omega t+\phi$。沿条纹随动点求时间导数，$\Delta\vct k\cdot\vct v_f-\Delta\omega=0$。取速度沿条纹法线，且$\Delta\vct k\ne0$，得到$\vct v_f=(\Delta\omega/|\Delta\vct k|^2)\Delta\vct k$。若$\Delta\vct k=0$，干涉项在空间中均匀；$\Delta\omega\ne0$时只有时间拍动，二者同时为零时则相位差恒定，均无可定义的空间条纹移动速度。纹速描述等相位面的移动，光能和信息仍按各束光的传播规律传递。
\Diagram{fig:o2movinga}{固定像素在有限曝光内平均正弦变化，长曝光削弱干涉项}{\draw[axis](0,0)--(5.5,0)node[right]{$t$};\draw[axis](0,0)--(0,2.7)node[above]{光强};\draw[curve,domain=0:5,samples=130]plot(\x,{1.3+cos(360*\x)});\draw[guide](1,0)--(1,2.5);\draw[guide](3.5,0)--(3.5,2.5);\node[above]at(2.25,2.5){曝光$T$};}
取曝光变量$s=t-t_0$，则交叉相位为$a-\Delta\omega s$，其中$a=\Delta\vct k\cdot\vct r-\Delta\omega t_0+\phi$。在式\eqref{eq:o2rectaverage}中取$L=T,b=-\Delta\omega$，交叉项便乘$\operatorname{sinc}(\Delta\omega T/2)$，等强可见度为其绝对值。图\ref{fig:o2movinga}在$T\gg1/|\Delta\omega|$时条纹变弱，短曝光则可记录移动条纹；“不同频率不干涉”包含了长时间平均的条件。
\end{solution}

\section{多次反射与干涉测量}
\begin{example}\label{ex:o2fabry}
无损对称法布里—珀罗平行腔两镜的能量反射率均为$R$，腔内单程相位$\phi$。取镜的振幅系数使腔往返倍率为$R\ee^{2i\phi}$。求透射强度并推导高反射率时的精细度。
\Diagram{fig:o2fabryq}{腔内每多一次往返，场乘$R\ee^{2i\phi}$}{\draw[very thick](1,-1.3)--(1,1.3);\draw[very thick](4,-1.3)--(4,1.3);\draw[->](0,.8)--(4,.8);\draw[->](4,.5)--(1,.5);\draw[->](1,.2)--(4,.2);\draw[->](4,-.1)--(5.5,-.1);\node[below]at(2.5,-1.4){$\phi$};}
\end{example}
\begin{solution}
图\ref{fig:o2fabryq}中，每个无损镜的能量透射率为$1-R$，首束透射场模为$(1-R)A$，包含单程相位$\ee^{i\phi}$。后续每束再乘一次往返倍率$R\ee^{2i\phi}$，几何级数给
\[
 A_t=(1-R)A\ee^{i\phi}\sum_{j=0}^\infty R^j\ee^{2ij\phi}
 =\frac{(1-R)A\ee^{i\phi}}{1-R\ee^{2i\phi}}.
\]
分母模平方为$1+R^2-2R\cos2\phi=(1-R)^2+4R\sin^2\phi$，因此
\[
 \frac{I_t}{I_i}=\frac1{1+\mathcal A\sin^2\phi},\qquad
 \mathcal A=\frac{4R}{(1-R)^2}.
\]
\Diagram{fig:o2fabrya}{高反射率使透射峰变窄，理想共振峰值仍为1}{\draw[axis](0,0)--(6.7,0)node[right]{$\phi$};\draw[axis](0,0)--(0,2.7)node[above]{$I_t/I_i$};\draw[curve,domain=0:6.283,samples=220]plot(\x,{2.3/(1+80*sin(\x r)^2)});\node[below]at(3.142,0){$\pi$};}
图\ref{fig:o2fabrya}的峰在$\phi=m\pi$，峰高为1。若$\mathcal A\geq1$，峰两侧半高点满足$\sin^2\delta\phi=1/\mathcal A$，故准确半高全宽为$2\arcsin(1/\sqrt{\mathcal A})$，精细度为
\[
 \mathcal F_c=\frac{\pi}{2\arcsin[(1-R)/(2\sqrt R)]}.
\]
高反射率时$(1-R)/(2\sqrt R)\ll1$，用$\arcsin x\simeq x$得到$\mathcal F_c\simeq\pi\sqrt R/(1-R)$。若$\mathcal A<1$，整条曲线的最低值仍高于峰的一半，按这个基线为零的半高定义没有半高点，不应使用窄峰线宽近似。腔内各束振幅逐次衰减，属于无限几何级数。无损对称腔在共振处可完全透射，与有限等幅多束的$N^2$峰高采用不同归一化。
\end{solution}

\begin{example}\label{ex:o2gascell}
迈克耳孙一臂放入长度$l$的气室，光在其中往返。抽真空后充入折射率$n$的气体，记录中心通过$N$条纹。求$n$；若只改变另一臂镜子的位置补偿，需移动多少？
\Diagram{fig:o2gasq}{气室在同一干涉臂内被光往返通过}{\draw[axis](0,0)--(6.3,0);\draw[curve](2,-.8)rectangle(4,.8);\draw[very thick](5.5,-1.2)--(5.5,1.2);\node[above]at(3,.8){气室$n,l$};\draw[->](1,.3)--(5.3,.3);\draw[->](5.3,-.3)--(1,-.3);}
\end{example}
\begin{solution}
图\ref{fig:o2gasq}以抽空状态为参考，每一程附加光程为$(n-1)l$，往返总改变量$2(n-1)l=N\lambda_0$，所以$n=1+N\lambda_0/(2l)$。
\Diagram{fig:o2gasa}{另一臂镜移一段，往返补偿两倍该长度}{\draw[very thick](2,-1)--(2,1);\draw[guide](3.3,-1)--(3.3,1);\draw[<->](2,-1.5)--(3.3,-1.5)node[midway,below]{$\Delta d$};\node[right]at(3.6,0){$2\Delta d=2(n-1)l$};}
图\ref{fig:o2gasa}需要镜移大小$|\Delta d|=(n-1)l$，方向选为增加另一臂光程来补偿气室臂的增加。若气室窗在充气前后保持不变，窗片的静态相位相消；若充气使窗形变或气室长度变化，测得的不再单是折射率项，需把额外光程纳入。
\end{solution}

\begin{example}\label{ex:o2ringmotion}
牛顿环中心原间隙$h_0$，将透镜整体抬高$\Delta h>0$且不改变曲率。跟踪固定级次$m$的反射暗环，求半径平方变化和中心每暗一次所需的抬升量。
\Diagram{fig:o2ringmotionq}{整体抬高改变所有半径处的共同间隙}{\draw[curve](-2.3,0)--(2.3,0);\draw[curve,domain=-2:2,samples=60]plot(\x,{.2+.15*\x*\x});\draw[guide,domain=-2:2,samples=60]plot(\x,{.7+.15*\x*\x});\draw[->](0,.25)--(0,.65)node[right]{$\Delta h$};}
\end{example}
\begin{solution}
图\ref{fig:o2ringmotionq}的暗环$r_m^2=R[m\lambda_0/n-2h_0]$，抬高后$r_m'^2=r_m^2-2R\Delta h$。所以同一级次向内收缩，直到$r_m'=0$后消失。
\Diagram{fig:o2ringmotiona}{跟踪同一级次时圆环向中心收缩}{\draw[curve](0,0)circle(1.6);\draw[guide](0,0)circle(1.05);\draw[->](1.55,0)--(1.1,0);\draw[->](-1.55,0)--(-1.1,0);\fill(0,0)circle(1pt);}
中心相位的暗纹条件$2nh_0=m\lambda_0$，相邻暗态间隔$\Delta h=\lambda_0/(2n)$。图\ref{fig:o2ringmotiona}所说向内是固定$m$的跟踪；如果每次都把“从中心数的第一环”重新编号，编号跳变会掩盖运动方向，必须先说明跟踪对象。
\end{solution}

\begin{example}\label{ex:o2weightedbinomial}
四束等间距光的振幅依次为$1,3,3,1$倍$A$。求干涉强度，比较它与四束等振幅时的零点和旁瓣，讨论权重改变的代价。
\Diagram{fig:o2binomialq}{四束相邻相位差为$\delta$，振幅取二项式权重}{\draw[axis](-.3,0)--(4.3,0)node[right]{束序};\foreach\x/\y in{0/1,1/3,2/3,3/1}{\draw[very thick](\x,0)--(\x,{.7*\y});\node[above]at(\x,{.7*\y}){$\y$};}}
\end{example}
\begin{solution}
图\ref{fig:o2binomialq}复振幅$\mathcal E/A=1+3\ee^{i\delta}+3\ee^{2i\delta}+\ee^{3i\delta}=(1+\ee^{i\delta})^3$，所以$I/I_A=64\cos^6(\delta/2)$。
\Diagram{fig:o2binomiala}{二项式加权消去周期内旁瓣，但主峰更宽}{\draw[axis](0,0)--(6.8,0)node[right]{$\delta$};\draw[axis](0,0)--(0,2.7)node[above]{归一化强度};\draw[curve,domain=0:6.283,samples=100]plot(\x,{2.3*cos((\x/2) r)^6});\node[below]at(3.142,0){$\pi$};}
图\ref{fig:o2binomiala}一个周期内只有$\delta=\pi$的零点，没有次极大。四束等幅时有三个零点和两个次极大、主峰首零在$\pi/2$。加权减少旁瓣却加宽主峰；若总入射功率固定，还应把权重整体归一化，不能直接拿峰值64与等幅16比较“效率”。
\end{solution}

\section{腔内损耗、耦合匹配与相位测量}
\subsection{把振幅损耗与能量损耗分别保留}
腔的透射峰取决于两端耦合和内部损耗。设两镜能量反射率为$0\leq R_1,R_2\leq1$，镜本身无吸收，腔内单程的能量存活率为$0\leq\eta\leq1$。一次往返的振幅模为$q=\eta\sqrt{R_1R_2}$，以下稳态级数要求$q<1$。用$\delta$表示包括镜相移在内的往返相位，略去各束共有的单程相位后，透射场为
\[
 E_t=\sqrt{(1-R_1)(1-R_2)\eta}\,E_{
 \rm in}\sum_{j=0}^{\infty}q^j\ee^{\mathrm i j\delta}.
\]
因此
\[
 T(\delta)=\frac{(1-R_1)(1-R_2)\eta}
 {(1-q)^2+4q\sin^2(\delta/2)}.
\]
此处$\sqrt\eta$是单程振幅因子，$\eta$才是往返振幅因子。两者混用会使峰高和线宽同时错误。

在$\delta=2m\pi$附近取$\delta=2m\pi+\epsilon$，若$q$接近1，则半高点满足$q\epsilon^2=(1-q)^2$，故往返相位全宽为$2(1-q)/\sqrt q$。精确半高点要求$\sin^2(\epsilon/2)=(1-q)^2/(4q)$；右边超过1时，一个自由光谱区内没有下降到峰值一半的位置，便不宜继续套用高精细度线宽公式。

\ExpansionDiagram{fig:dep-o2-roundtrip}{每次新增透射束比前一束多一次往返及其损耗}{
\draw[very thick](0,-1)--(0,1);\draw[very thick](4,-1)--(4,1);\node[above]at(0,1){$R_1$};\node[above]at(4,1){$R_2$};
\draw[->](-1,.6)--(4,.6);\draw[->](4,.2)--(0,.2);\draw[->](0,-.2)--(5,-.2);\node[below]at(2,-.6){单程能量因子$\eta$};}
图\ref{fig:dep-o2-roundtrip}中的一次往返含两次单程衰减，但振幅每次只乘$\sqrt\eta$。把它们与两镜反射振幅相乘，才得到$q$。

\begin{example}有内损耗的腔怎样消除入射端反射
采用上述模型，入射端的振幅反射系数取
\[
 r_{\rm cav}(\delta)=\frac{\sqrt{R_1}-\eta\sqrt{R_2}\ee^{\mathrm i\delta}}
 {1-\eta\sqrt{R_1R_2}\ee^{\mathrm i\delta}}.
\]
求共振时反射为零的条件、此时透射率与腔内吸收率。解释反射为零是否意味着全部透射。
\ExpansionDiagram{fig:dep-o2-match-q}{入射端反射由直接反射和腔内返回光相干组成}{
\draw[very thick](1,-1)--(1,1);\draw[very thick](5,-1)--(5,1);
\draw[->](-1,.5)--(1,.5);\draw[->](1,.1)--(-1,.1)node[left]{直接反射};\draw[->](5,-.3)--(-1,-.3)node[left]{返回光};\node[above]at(3,1){有损腔};}
\end{example}
\begin{solution}
图\ref{fig:dep-o2-match-q}中的两项在共振时反相，其模相等要求$\sqrt{R_1}=\eta\sqrt{R_2}$，即
\[
 R_1=\eta^2R_2.
\]
这就是给定输出镜与内损耗时的入射耦合匹配条件。代入透射公式，注意$q=\eta^2R_2$，有
\[
 T_{\rm res}=\frac{\eta(1-R_2)}{1-\eta^2R_2},\qquad
 A_{\rm res}=1-T_{\rm res}
 =\frac{(1-\eta)(1+\eta R_2)}{1-\eta^2R_2}.
\]
这里镜无吸收，$A_{\rm res}$全是介质内部损耗。$\eta=1$时匹配变为$R_1=R_2$，此时全部透射；$R_2=1$且$\eta<1$时没有透射口，共振入射能量全部在腔内耗散。
\ExpansionDiagram{fig:dep-o2-match-a}{共振匹配时，外部能量分成透射和内部损耗}{
\draw[->,thick](0,0)--(2,0)node[midway,above]{$P$};\draw(2,-.6)rectangle(4,.6)node[midway]{腔};\draw[->,thick](4,0)--(6,0)node[midway,above]{$TP$};\draw[->,thick](3,-.6)--(3,-1.6)node[below]{$AP$};\node[above]at(3,.8){反射为零};}
例如$\eta=0.9,R_2=0.8$，所需$R_1=0.648$，得到$T_{\rm res}=0.18/0.352\simeq0.5114$、$A_{\rm res}\simeq0.4886$。图\ref{fig:dep-o2-match-a}说明反射消失是相干抵消条件，能量的最终去向还须由透射和吸收分别计算。若镜面也吸收，应重新区分镜的透射率与$1-R_j$，不能直接使用本题的分子。
\end{solution}

\subsection{从四次强度测量还原复干涉项}
只测一次强度不能同时确定振幅和相位。给一臂施加已知附加相位$\alpha$，强度写成
\[
 I(\alpha)=B+C\cos(\phi+\alpha),\qquad B=I_1+I_2,
 \quad C=2\sqrt{I_1I_2}|\gamma|.
\]
$\gamma$是两路的复相干度，其辐角已并入$\phi$。在四个等间隔相位$0,\pi/2,\pi,3\pi/2$测量，余弦项和正弦项可以分别提取；这与把两幅条纹的位置目测相减不同。

\begin{example}四步移相与相位的象限
某点四次强度依次为$14,6,2,10$。已知附加相位顺序为$0,\pi/2,\pi,3\pi/2$。求$B,C,\phi$与可见度；若单独遮挡两臂测得$I_1=I_2=4$，求$|\gamma|$。
\ExpansionDiagram{fig:dep-o2-step-q}{参考臂的已知移相使同一测量点获得四个强度样本}{
\draw(0,0)--(1,1)--(4,1)--(5,0);\draw(0,0)--(1,-1)--(4,-1)--(5,0);\draw[thick](3,-1.35)--(3,-.65);\node[below right]at(3,-1.35){$\alpha$};\fill(5,0)circle(1.4pt)node[right]{$P$};\node[above]at(2,1){物臂};\node[below]at(1,-1){参考臂};}
\end{example}
\begin{solution}
图\ref{fig:dep-o2-step-q}中两臂功率不变，只有参考臂相位变化。用相差$\pi$的样本相加，有$B=(I_0+I_\pi)/2=8$，另一对也给$(6+10)/2=8$。相减得到
\[
 X=\frac{I_0-I_\pi}{2}=C\cos\phi=6,\qquad
 Y=\frac{I_{3\pi/2}-I_{\pi/2}}2=C\sin\phi=2.
\]
于是$C=\sqrt{X^2+Y^2}=2\sqrt{10}$，$\phi=\operatorname{atan2}(2,6)\simeq18.435^\circ$，而不是仅由一个反余弦随意选分支。
\ExpansionDiagram{fig:dep-o2-step-a}{强度差给出干涉项在复平面上的两个分量}{
\draw[->](0,0)--(4,0)node[right]{$X$};\draw[->](0,0)--(0,2.4)node[above]{$Y$};\draw[->,thick](0,0)--(3,1)node[above right]{$6+2\mathrm i$};\draw[dashed](3,0)--(3,1)--(0,1);\node[below]at(3,0){$6$};\node[left]at(0,1){$2$};}
由图\ref{fig:dep-o2-step-a}的模得到可见度$\mathcal V=C/B=\sqrt{10}/4\simeq0.7906$。单臂测量给$2\sqrt{I_1I_2}=8$，所以$|\gamma|=C/8=\sqrt{10}/4$。若不知道单臂强度，仅凭四步数据只能得$B,C$，不能区分可见度降低来自功率不等还是相干度不足。两对样本和相等是该模型的一个数据相容条件；漂移、移相误差或探测非线性会破坏它。
\end{solution}

\section{谱线形状、相位扩散与色散包络}
\subsection{相位扩散对应Lorentz线形}
在$D\ll\omega_0$的窄带近似下，用$E(t)=E_0\exp[-\mathrm i\omega_0t+\mathrm i\varphi(t)]$表示幅度近乎恒定、相位缓慢扩散的场。设相位增量为均值零的高斯变量，$\langle[\varphi(t+\tau)-\varphi(t)]^2\rangle=2D|\tau|$，初相位在整周均匀分布。高斯平均给
\[
 \langle\ee^{\mathrm i[\varphi(t+\tau)-\varphi(t)]}\rangle
 =\ee^{-D|\tau|}.
\]
高斯平均可由习题\ref{ex:o2phaseerror}的特征函数推导核对：在其中取方差$s^2=2D|\tau|$、$k=1$。于是归一化场相关函数含$\ee^{-D|\tau|}\ee^{-\mathrm i\omega_0\tau}$。令$\Omega=\omega-\omega_0$，其傅里叶积分按正负时间分开：
\begin{align*}
 S(\omega)&=\frac1{2\pi}\left[
 \int_{-\infty}^0\ee^{(D+\mathrm i\Omega)\tau}\dd\tau
 +\int_0^\infty\ee^{(-D+\mathrm i\Omega)\tau}\dd\tau\right]\\
 &=\frac1{2\pi}\left[
 \frac{\ee^{(D+\mathrm i\Omega)\tau}}{D+\mathrm i\Omega}\right]_{-\infty}^0
 +\frac1{2\pi}\left[
 \frac{\ee^{(-D+\mathrm i\Omega)\tau}}{-D+\mathrm i\Omega}\right]_0^\infty\\
 &=\frac1{2\pi}\left(\frac1{D+\mathrm i\Omega}
 +\frac1{D-\mathrm i\Omega}\right).
\end{align*}
通分后得到归一化角频率谱
\[
 S(\omega)=\frac{D}{\pi[(\omega-\omega_0)^2+D^2]}.
\]
角频率半高全宽为$2D$，普通频率半高全宽为$\Delta\nu=D/\pi$。等强两臂的可见度为$\ee^{-D|\tau|}$，不存在矩形谱那样的有限首零点。若将相干长度定义为可见度降至$1/\ee$的光程差，则$L_c=c/D=c/(\pi\Delta\nu)$。定义换为半高宽时常数会改变，应连同定义一起报告。

\begin{example}光源线宽与扫描曝光共同降低可见度
等强双臂干涉仪采用上述相位扩散光源，中心波长$600\,\mathrm{nm}$，谱线全宽$\Delta\nu=3.0\,\mathrm{MHz}$。两臂光程差$10\,\mathrm m$。扫描镜以$30\,\mu\mathrm{m/s}$匀速运动，使往返光程差以$2v$改变，曝光时间为$5.0\,\mathrm{ms}$。在曝光期间相关包络近似不变，求记录的条纹可见度。
\ExpansionDiagram{fig:dep-o2-diffusion-q}{扫描镜改变载波相位，光源线宽另行决定延迟相关包络}{
\draw[->](-1.8,0)--(-.25,0);\draw[thick](-.2,-.2)--(.2,.2);\draw[->](0,.2)--(0,2);\draw[->](.12,2)--(.12,.2);\draw[very thick](-.4,2)--(.5,2)node[right]{固定镜};\draw[->](.25,0)--(4,0);\draw[->](4,-.12)--(.2,-.12);\draw[very thick](4,-.6)--(4,.6);\draw[->](4,.9)--(5.2,.9)node[right]{$v$};\draw[->](0,-.3)--(0,-1.5)node[below]{记录点};}
\end{example}
\begin{solution}
图\ref{fig:dep-o2-diffusion-q}中静止光程差给$\tau=10/c$，所以包络因子为$\exp[-\pi\Delta\nu\,10/c]=\exp(-\pi/10)\simeq0.7304$，取$c=3.0\times10^8\,\mathrm{m/s}$。

镜面移动使相位角速度$\Omega_s=4\pi v/\lambda$。在以曝光中点为原点的时间区间内平均余弦，
\[
 \frac1{T_e}\int_{-T_e/2}^{T_e/2}\cos(\phi_0+\Omega_st)\dd t
 =\left[\frac{\sin(\phi_0+\Omega_st)}{\Omega_sT_e}\right]_{-T_e/2}^{T_e/2}
 =\operatorname{sinc}(\Omega_sT_e/2)\cos\phi_0.
\]
本题$\Omega_sT_e/2=2\pi vT_e/\lambda=\pi/2$，所以记录可见度为
\[
 \mathcal V=\ee^{-\pi/10}\frac2\pi\simeq0.4650.
\]
\ExpansionDiagram{fig:dep-o2-diffusion-a}{可见度的两个独立因子：延迟包络与曝光平均}{
\node[draw]at(0,0){$0.7304$};\node at(1.5,0){$\times$};\node[draw]at(3,0){$2/\pi$};\node at(4.5,0){$=$};\node[draw]at(6,0){$0.4650$};\node[below]at(0,-.4){谱线宽度};\node[below]at(3,-.4){扫描曝光};}
图\ref{fig:dep-o2-diffusion-a}中的两项分别来自不同平均。本题曝光内包络变化参数$D(2vT_e/c)$远小于1，可以近似相乘；若这个量不可略，应把随时间改变的相关包络一并放入曝光积分。减小曝光时间只能改善后一项，不能恢复已经因长光程差而失去的时间相干性。
\end{solution}

\subsection{群光程决定包络，中心相位决定细条纹}
白光或短脉冲干涉中，两臂相位差通常依赖频率。设中心频率附近$\Omega=\omega-\omega_0$，相位展开为
\[
 \phi(\omega)=\phi_0+\tau_g\Omega+\tfrac12D_2\Omega^2,
\]
入射功率谱的归一化形状为$\exp[-\Omega^2/(2\sigma_\omega^2)]/(\sqrt{2\pi}\sigma_\omega)$。取$\sigma_\omega\ll\omega_0$，将偏移频率积分延伸到整个实轴。复干涉项是$\langle\ee^{\mathrm i\phi(\omega)}\rangle$。记$a=(1-\mathrm iD_2\sigma_\omega^2)/(2\sigma_\omega^2)$、$b=\mathrm i\tau_g$，则$\operatorname{Re}a>0$，并有
\[
 \Gamma=\frac{\ee^{\mathrm i\phi_0}}{\sqrt{2\pi}\sigma_\omega}
 \int_{-\infty}^{\infty}\ee^{-a\Omega^2+b\Omega}\dd\Omega
 =\frac{\ee^{\mathrm i\phi_0}}{\sqrt{2\pi}\sigma_\omega}
 \sqrt{\frac\pi a}\ee^{b^2/(4a)}.
\]
最后一步使用式\eqref{eq:o1gaussintegral}已推导的积分。代入$a,b$，得
\[
 \Gamma=\frac{\ee^{\mathrm i\phi_0}}{\sqrt{1-\mathrm iD_2\sigma_\omega^2}}
 \exp\!\left[-\frac{\sigma_\omega^2\tau_g^2}
 {2(1-\mathrm iD_2\sigma_\omega^2)}\right].
\]
平方根取在$D_2=0$处连续为1的分支。模为
\[
 |\Gamma|=[1+(D_2\sigma_\omega^2)^2]^{-1/4}
 \exp\!\left[-\frac{\sigma_\omega^2\tau_g^2}
 {2[1+(D_2\sigma_\omega^2)^2]}\right].
\]
二阶色散使包络展宽、峰值降低。包络最大值位于$\tau_g=0$，细条纹相位却还含有$\phi_0$和积分产生的附加相位，因此包络峰的位置一般不等于中心频率的零程差位置。

\begin{example}插入色散片后应补偿哪一种光程
单程双臂干涉仪在一臂垂直插入厚度$l$、折射率$n(\omega)$的透明片，另一臂增加空气光程$x$。忽略反射，近中心频率相位差为$\phi(\omega)=\omega[l(n(\omega)-1)-x]/c$。求使相关包络最大的一阶补偿量，与使中心相位为零的补偿量比较。
\ExpansionDiagram{fig:dep-o2-group-q}{一臂色散片、另一臂空气光程补偿}{
\draw(0,0)--(1,1.2)--(4,1.2)--(5,0);\draw(0,0)--(1,-1.2)--(4,-1.2)--(5,0);\draw[fill=gray!15](2.2,.9)rectangle(2.6,1.5);\node[above]at(2.4,1.6){$l,n(\omega)$};\node[below]at(2.5,-1.3){空气补偿$x$};}
\end{example}
\begin{solution}
图\ref{fig:dep-o2-group-q}的中心相位为$\phi_0=\omega_0[l(n_0-1)-x]/c$。对完整的$\omega n(\omega)$求导，
\[
 \tau_g=\left.\frac{\dd\phi}{\dd\omega}\right|_{\omega_0}
 =\left.\frac{l[n(\omega)-1]-x+\omega l n'(\omega)}c\right|_{\omega_0}
 =\frac{l[n_0+\omega_0n'_0-1]-x}{c}.
\]
因此包络最大要求$x_g=l(n_g-1)$，其中$n_g=n_0+\omega_0n'_0$；中心频率相位恰为零则要求$x_p=l(n_0-1)$。两者差为$x_g-x_p=l\omega_0n'_0$。
\ExpansionDiagram{fig:dep-o2-group-a}{细条纹位于相关包络之内，零载波相位不必位于包络中心}{
\draw[->](0,0)--(7,0)node[right]{$x$};\draw[thick,domain=.1:6.6,samples=130]plot(\x,{1.2*exp(-.35*(\x-3.8)^2)});\draw[dashed,domain=.1:6.6,samples=200]plot(\x,{1.2*exp(-.35*(\x-3.8)^2)*cos(600*(\x-2.5))});\draw[dotted](2.5,-1.2)--(2.5,1.6)node[above]{$x_p$};\draw[dotted](3.8,-1.2)--(3.8,1.6)node[above]{$x_g$};}
图\ref{fig:dep-o2-group-a}示意两种中心的位置，细条纹的精确位置还由谱积分相位决定。忽略色散$n'_0=0$时，两种补偿相同。二阶色散显著时，改变$x$只调整常数项和一次频率项，$D_2$仍然保留；恢复原有峰值需要另一种材料或其他色散补偿，增加采样密度对这一展宽没有作用。
\end{solution}

\section{习题与解答}
\begin{exercise}\label{ex:o2bragg}
正常入射的无损多层反射镜由$N$组高、低折射率层$H,L$组成，入射侧先遇$H$，每层在设计波长的光学厚度均为$\lambda_0/4$。入射介质、基底导纳分别为$n_0,n_s$。推导设计波长处的反射率。
\Diagram{fig:o2braggq}{交替四分之一波厚层叠加界面反射}{\foreach\x in{0,1,2,3}{\draw[curve](\x,-1)rectangle({\x+.6},1);\node[above]at({\x+.3},1){$H$};\draw[guide]({\x+.6},-1)rectangle({\x+1},1);\node[below]at({\x+.8},-1){$L$};}\draw[axis](-1,0)--(5,0);}
\end{exercise}
\begin{solution}
把磁场归一化为$\mathcal H=Z_0H$，其中$Z_0$为真空波阻抗。正常入射层中前、后行波幅为$A,B$，在某参考面有$E=A+B$、$\mathcal H=n(A-B)$。厚度为$l$的层内，前行波向右乘$\ee^{i\phi}$、后行波向右乘$\ee^{-i\phi}$，$\phi=k_0nl$。因此左面振幅可由右面振幅写为$A_L=A_R\ee^{-i\phi}$、$B_L=B_R\ee^{i\phi}$。代入场的和、差，得到从层右面到左面的特征矩阵
\[
 \begin{pmatrix}E_L\\\mathcal H_L\end{pmatrix}
 =C(n,\phi)\begin{pmatrix}E_R\\\mathcal H_R\end{pmatrix},\qquad
 C(n,\phi)=\begin{pmatrix}\cos\phi&-i\sin\phi/n\\
 -in\sin\phi&\cos\phi\end{pmatrix}.
\]
界面处切向$E,\mathcal H$连续，可以直接接续相邻层。图\ref{fig:o2braggq}在设计波长$\phi_H=\phi_L=\pi/2$，一组矩阵
\[
 C_HC_L=\begin{pmatrix}-n_L/n_H&0\\0&-n_H/n_L\end{pmatrix}.
\]
设$q=n_H/n_L$，$N$组矩阵为$(-1)^N\operatorname{diag}(q^{-N},q^N)$。基底中只有向右的波，故$\mathcal H_s=n_sE_s$，于是$E_{\rm in}=(-1)^Nq^{-N}E_s$、$\mathcal H_{\rm in}=(-1)^Nq^Nn_sE_s$，等效输入导纳为$Y=\mathcal H_{\rm in}/E_{\rm in}=n_sq^{2N}$。
\Diagram{fig:o2bragga}{层组把基底导纳变换为越来越大的输入导纳}{\draw[curve](0,0)rectangle(2.8,1);\node at(1.4,.5){$n_s$};\draw[->](3,.5)--(4,.5);\node[right]at(4.2,.5){$Y=n_s(n_H/n_L)^{2N}$};\node[below]at(2.3,-.7){$r=(n_0-Y)/(n_0+Y)$};}
入射和反射叠加给$E=1+r$、$\mathcal H=n_0(1-r)$，与$\mathcal H=YE$联立得到图\ref{fig:o2bragga}的$r$，反射率$R=[(n_0-Y)/(n_0+Y)]^2$。当$n_H>n_L$且$N$增大，$R\to1$。该结果只在设计波长和给定层序精确成立；交换层序会倒置导纳倍率，层数和层序均应保留。
\end{solution}

\begin{exercise}\label{ex:o2jointcoherence}
均匀线源宽$w$、距双缝$a$，双缝间距$d$，光谱为中心$\nu_0$、宽$\Delta\nu$的矩形谱。源位置与频率统计独立，两路等强、偏振相同，中心源点的两路时间差为$\tau$。在窄带近似下求可见度，并说明何时不能简单把时间与空间可见度相乘。
\Diagram{fig:o2jointq}{源位置和频率各自造成一类相位散布}{\draw[curve](0,0)rectangle(2.2,1.2);\node at(1.1,.6){位置$x_s$};\draw[curve](3.6,0)rectangle(5.8,1.2);\node at(4.7,.6){频率$\nu$};\node[below]at(2.9,-.7){两变量分布独立};}
\end{exercise}
\begin{solution}
图\ref{fig:o2jointq}中各源元、各频率的两路相位差为$2\pi\nu\tau+2\pi\nu d x_s/(ca)$。对均匀位置分布和矩形谱，归一化干涉因子为
\[
 \Gamma=\frac1{w\Delta\nu}\int_{-w/2}^{w/2}\dd x_s
 \int_{\nu_0-\Delta\nu/2}^{\nu_0+\Delta\nu/2}
 \ee^{i2\pi\nu\tau+i2\pi\nu d x_s/(ca)}\dd\nu.
\]
窄带时在位置项中以$\nu_0$代替$\nu$，积分核分解。频率积分给$\ee^{i2\pi\nu_0\tau}\operatorname{sinc}(\pi\Delta\nu\tau)$，位置积分给$\operatorname{sinc}[\pi\nu_0dw/(ca)]$。假定两路等强且偏振相同，可见度为$|\Gamma|$，因此
\[
 \mathcal V\simeq\left|\operatorname{sinc}(\pi\Delta\nu\tau)
 \operatorname{sinc}\left(\frac{\pi dw}{\lambda_0a}\right)\right|.
\]
频率平均控制不同延迟的同点相干，位置平均控制两缝的照明相干，两者都要满足才有高对比条纹。
\Diagram{fig:o2jointa}{两个独立平均因子相乘，任一个接近零都会洗去条纹}{\node at(0,.6){时间因子};\node at(0,-.6){空间因子};\draw[->](1.4,.6)--(3,0);\draw[->](1.4,-.6)--(3,0);\node[right]at(3.2,0){总可见度};}
被略去的交叉相位尺度为$2\pi\Delta\nu d w/(ca)$，须远小于1。若谱很宽、源很大，或不同频率来自不同源位置，图\ref{fig:o2jointa}的分解不成立，应保留联合分布做二重积分。两个单独可见度接近1只是窄带独立模型中的充分描述，不能对任意源结构机械相乘。
\end{solution}

\begin{exercise}\label{ex:o2coherencenum}
中心波长600 nm、窄矩形谱宽度0.2 nm的光源用于双束干涉。估计首零光程差及迈克耳孙单镜相对零光程差允许的首零移动量。
\Diagram{fig:o2cohnumq}{镜移通过往返光程差扫描相干包络}{\draw[very thick](1,-1)--(1,1);\draw[guide](3,-1)--(3,1);\draw[<->](1,-1.5)--(3,-1.5)node[midway,below]{$d$};\draw[->](0,.4)--(2.7,.4);\draw[->](2.7,-.3)--(0,-.3);}
\end{exercise}
\begin{solution}
由$\nu=c/\lambda$求导，$\dd\nu=-c\dd\lambda/\lambda^2$；谱宽取正值，窄谱时$\Delta\nu\simeq c\Delta\lambda/\lambda_0^2$。图\ref{fig:o2cohnumq}的首零光程差为$\mathcal D_0=c/\Delta\nu\simeq\lambda_0^2/\Delta\lambda=(600\,\mathrm{nm})^2/(.2\,\mathrm{nm})=1.8$ mm。
\Diagram{fig:o2cohnuma}{单镜移动量是扫描光程差的一半}{\draw[axis](0,0)--(5,0)node[right]{毫米};\draw[curve](1,-.4)--(1,.4);\draw[curve](2.8,-.4)--(2.8,.4);\node[below]at(1,-.4){$d=.9$};\node[below]at(2.8,-.4){$\mathcal D=1.8$};}
图\ref{fig:o2cohnuma}单镜首零移动量$d_0=.9$ mm。此数以光程差零点为起点；若初始已有臂长差，应把它计入，只有从零光程差开始移动，才在.9 mm处首次消失。矩形谱理想模型在首零以外仍有较小旁瓣，可见度不是永久为零。
\end{solution}

\begin{exercise}\label{ex:o2sourceorientation}
均匀线光源长度$w$与双缝连线成角$\psi$，距缝$a$。求空间相干可见度，比较光源沿缝长方向和垂直缝长方向的情况。
\Diagram{fig:o2orientationq}{光源的基线投影决定两缝之间的入射相位散布}{\draw[axis](-.3,0)--(3.7,0)node[right]{双缝基线};\draw[very thick](0,0)--(2.2,1.7);\draw[guide](2.2,0)--(2.2,1.7);\node at(.8,.3){$\psi$};\node[right]at(2.2,1.7){$w$};}
\end{exercise}
\begin{solution}
图\ref{fig:o2orientationq}线源参数$s\in[-w/2,w/2]$，其在双缝基线上的投影为$s\cos\psi$，源元相位为$2\pi d s\cos\psi/(\lambda a)$。将式\eqref{eq:o2rectaverage}中的相位斜率改为$2\pi d\cos\psi/(\lambda a)$，区间长度仍取$w$，得到$\mathcal V=|\operatorname{sinc}(\pi dw\cos\psi/(\lambda a))|$；$\cos\psi=0$时用该定积分的常相位分支，结果为1。
\Diagram{fig:o2orientationa}{垂直基线延伸时，近轴模型中两缝相位差不随源点改变}{\fill(-1,0)circle(1.5pt)node[below]{$S_-$};\fill(1,0)circle(1.5pt)node[below]{$S_+$};\draw[very thick](0,.7)--(0,2.5);\draw[guide](0,2)--(-1,0);\draw[guide](0,2)--(1,0);\node[right]at(.2,2.5){源的投影示意};}
沿缝长方向通常垂直双缝基线，$\cos\psi=0$，图\ref{fig:o2orientationa}的源点不引起一阶基线相位差，可见度保持1。沿基线延伸最容易洗掉条纹。实际有限孔径和非近轴高阶项仍可限制可见度，不能把此一阶结论外推为任意大光源都完全相干。
\end{solution}

\begin{exercise}\label{ex:o2phaseerror}
两束等强光相位差含随机误差$\epsilon$，服从均值0、方差$s^2$的高斯分布。求长曝光可见度，解释相位抖动对条纹位置和对比度的不同影响。
\Diagram{fig:o2errorq}{相位随时间在平均位置附近抖动}{\draw[axis](0,0)--(5,0)node[right]{$t$};\draw[axis](0,-1.4)--(0,1.7)node[above]{$\epsilon$};\draw[curve](0,.2)--(.5,-.6)--(1,.8)--(1.5,.2)--(2,-.9)--(2.5,.6)--(3,-.2)--(3.5,1)--(4,.1)--(4.5,-.4);}
\end{exercise}
\begin{solution}
图\ref{fig:o2errorq}中，瞬时强度可写为$2I_1[1+\cos(\delta_0+\epsilon)]$。长曝光的交叉项由$G=\langle\ee^{i\epsilon}\rangle$控制，其中
\[
 G=\frac1{\sqrt{2\pi}s}\int_{-\infty}^\infty
 \exp[-\epsilon^2/(2s^2)]\ee^{i\epsilon}\dd\epsilon.
\]
引入$G(k)=\langle\ee^{ik\epsilon}\rangle$。概率密度的导数为$p'(\epsilon)=-(\epsilon/s^2)p(\epsilon)$；对$k$求导，再分部积分，
\begin{align*}
 G'(k)&=i\int_{-\infty}^{\infty}\epsilon p(\epsilon)\ee^{ik\epsilon}\dd\epsilon
 =-is^2\int_{-\infty}^{\infty}p'(\epsilon)\ee^{ik\epsilon}\dd\epsilon\\
 &=-is^2[p(\epsilon)\ee^{ik\epsilon}]_{-\infty}^{\infty}
 +is^2\int_{-\infty}^{\infty}ikp(\epsilon)\ee^{ik\epsilon}\dd\epsilon
 =-s^2kG(k).
\end{align*}
边界项因高斯衰减而为零。以$G(0)=1$分离变量积分，$\int_1^{G(k)}\dd g/g=-s^2\int_0^kq\dd q=-s^2k^2/2$，因此$G(k)=\ee^{-s^2k^2/2}$。
\Diagram{fig:o2errora}{零均值抖动不移动平均条纹，却降低可见度}{\draw[axis](0,0)--(4.5,0)node[right]{$s$};\draw[axis](0,0)--(0,2.7)node[above]{$\mathcal V$};\draw[curve,domain=0:4,samples=60]plot(\x,{2.3*exp(-\x*\x/2)});}
取$k=1$得到图\ref{fig:o2errora}的平均强度$2I_1[1+\ee^{-s^2/2}\cos\delta_0]$，因此$\mathcal V=\ee^{-s^2/2}$。平均相位误差为零时亮暗位置不变；若均值为$\bar\epsilon$，则$G=\ee^{i\bar\epsilon}\ee^{-s^2/2}$，条纹另有整体位移。

这里$s$以弧度计，曝光须覆盖足够多次抖动，且误差统计与光强无关。若机械位移$z$使迈克耳孙一臂镜往返程差改变$2z$，则相位标准差为$s=4\pi\sigma_z/\lambda$。同一位置抖动对短波长的对比度损失更大，不能只用$\sigma_z$而不换成相位来比较。
\end{solution}

\begin{exercise}\label{ex:o2filmangle}
一薄膜反射具有一次相对半波损失。正常入射某波长$\lambda_0$在膜厚$h$处形成$m$级暗纹。入射角改为$i$，求对应同级暗纹波长，介质色散忽略。
\Diagram{fig:o2angleq}{斜入射改变膜内法向相位积累}{\draw[curve](-2,0)--(2,0);\draw[curve](-2,-.7)--(2,-.7);\draw[->](-1.6,1.6)--(0,0)--(.4,-.7)--(.8,0)--(2.4,1.6);\node[left]at(-2,-.35){$n$};}
\end{exercise}
\begin{solution}
图\ref{fig:o2angleq}正常暗纹条件为$2nh=m\lambda_0$。对正膜厚$h>0$，取$m\geq1$；斜入射时$2nh\cos r=m\lambda_i$，且$\sin r=\sin i/n$，两暗纹条件相除，得$\lambda_i=\lambda_0\sqrt{1-\sin^2i/n^2}$。若$h=0,m=0$，接触处在该反射相位模型下对所有波长都暗，不能由零级条件选出唯一的$\lambda_i$。
\Diagram{fig:o2anglea}{同级暗纹波长随角度增加向短波移动}{\draw[axis](0,0)--(4.4,0)node[right]{$i$};\draw[axis](0,0)--(0,2.7)node[above]{$\lambda_i/\lambda_0$};\draw[curve,domain=0:85,samples=60]plot({\x/22},{2.3*sqrt(1-sin(\x)^2/2.25)});\node[left]at(0,2.3){$1$};}
图\ref{fig:o2anglea}对应$n=1.5$。角度增大使法向波矢减小，暗纹波长蓝移。若两界面反射相位在所扫角度内因$p$偏振过布儒斯特点而改变，原来的暗纹编号规则需要重新讨论，不能只连续追踪$\cos r$。
\end{solution}

\begin{exercise}\label{ex:o2freespectral}
平行腔长度$l$、折射率$n$无色散，正常入射。求相邻透射共振的频率间隔，并由精细度$\mathcal F_c$求线宽。
\Diagram{fig:o2fsrq}{共振要求往返相位为$2m\pi$}{\draw[very thick](0,-1)--(0,1);\draw[very thick](4,-1)--(4,1);\draw[->](.2,.4)--(3.8,.4);\draw[->](3.8,-.4)--(.2,-.4);\node at(2,0){$n,l$};}
\end{exercise}
\begin{solution}
图\ref{fig:o2fsrq}的单程相位为$\phi=2\pi nl\nu/c$，共振要求$2\phi=2\pi m$，因此
\[
 \nu_m=\frac{mc}{2nl},\qquad
 \Delta\nu_{\rm FSR}=\nu_{m+1}-\nu_m=\frac c{2nl}.
\]
精细度按相邻峰距除以单峰半高全宽定义，故$\delta\nu=\Delta\nu_{\rm FSR}/\mathcal F_c$。高反射率时再代入$\mathcal F_c\simeq\pi\sqrt R/(1-R)$。
\Diagram{fig:o2fsra}{峰间距与峰宽是不同的两个频率尺度}{\draw[axis](0,0)--(5.5,0)node[right]{$\nu$};\foreach\x in{1,3,5}{\draw[curve]({\x-.3},0)..controls({\x-.1},2.8)and({\x+.1},2.8)..({\x+.3},0);}\draw[<->](1,-.5)--(3,-.5)node[midway,below]{FSR};}
图\ref{fig:o2fsra}中增大$R$使单峰变窄，而相邻峰距仍由光学长度决定。若$n$有色散，往返相位应写$\Psi(\nu)=4\pi l\nu n(\nu)/c$，近邻共振满足
\[
 2\pi\simeq\Psi'(\nu)\Delta\nu,
 \quad \Delta\nu\simeq\frac{c}{2l[n+\nu\dd n/\dd\nu]}
 =\frac c{2ln_g}.
\]
其中$n_g=n-\lambda\dd n/\dd\lambda$是群折射率。若镜的反射相位也随频率变化，还须把两个镜面的反射相位导数加进$\Psi'$；此时有效往返群延迟不再只是$2nl/c$。只要相邻峰之间参数变化足够小，上述导数法仍适用。
\end{solution}

\begin{exercise}\label{ex:o2Nphasechange}
$N$束等振幅光相邻相位差由$\delta$均匀增到$\delta+\epsilon$。求主峰中心的相位移动与峰高变化；若只在最后一束附加小相位$\epsilon$，轴上峰高的二阶变化如何？
\Diagram{fig:o2Nphaseq}{线性相位斜坡与单束相位误差有不同后果}{\draw[axis](0,0)--(4.5,0)node[right]{$j$};\draw[axis](0,0)--(0,2.7)node[above]{附加相位};\draw[curve](0,0)--(3.8,2);\draw[guide](0,0)--(3.8,0)--(3.8,1);}
\end{exercise}
\begin{solution}
图\ref{fig:o2Nphaseq}线性斜坡只把多束强度函数自变量换成$\delta+\epsilon$，主峰移到$\delta=2m\pi-\epsilon$，峰高仍$N^2I_1$。若只有最后一束出错，原轴上合场为$A[(N-1)+\ee^{i\epsilon}]$。
\Diagram{fig:o2Nphasea}{单束误差使一条相量偏离其余同向相量}{\draw[->,very thick](0,0)--(3.3,0)node[below]{$(N-1)A$};\draw[->,thick](3.3,0)--(4.1,1)node[right]{$A\ee^{i\epsilon}$};\draw[guide](3.3,0)--(4.5,0);}
图\ref{fig:o2Nphasea}给$I/I_1=(N-1)^2+1+2(N-1)\cos\epsilon\simeq N^2-(N-1)\epsilon^2$。这是原轴上强度，严格新峰位置还需对$\delta$求极值；不能把它直接当全部角度中已重新寻优的最大值。相位斜坡可用峰位补偿，局部误差一般同时改变旁瓣和峰形。
\end{solution}

\begin{exercise}\label{ex:o2visibility}
两相干光强比为9:1，偏振方向夹角$\eta$，求条纹可见度。若希望完全相消，需同时满足什么条件？
\Diagram{fig:o2visibilityq}{两束电场的偏振方向夹角}{\draw[axis](0,0)--(3,0)node[right]{$\vct e_1$};\draw[axis](0,0)--(1.5,2.2)node[right]{$\vct e_2$};\node at(.8,.4){$\eta$};}
\end{exercise}
\begin{solution}
以图\ref{fig:o2visibilityq}的第一束偏振为$x$轴，令两束实振幅为$3A,A$。第二束可分成$A\cos\eta$的平行分量和$A\sin\eta$的正交分量。合场分量为
\[
 E_x=3A+A\cos\eta\ee^{i\delta},\qquad
 E_y=A\sin\eta\ee^{i\delta}.
\]
探测器记录两正交分量强度之和，因此$I/I_A=|E_x/A|^2+|E_y/A|^2=10+6\cos\eta\cos\delta$。
\Diagram{fig:o2visibilitya}{正交分量不能被另一束光抵消}{\draw[axis](0,0)--(2.4,0)node[right]{$\parallel\vct e_1$};\draw[axis](0,0)--(0,2.2)node[above]{垂直分量};\draw[->,thick](0,0)--(1.4,1.6);\draw[guide](1.4,0)--(1.4,1.6)--(0,1.6);}
图\ref{fig:o2visibilitya}中的正交分量不能与第一束相消。相位变化时，$I_{\max}/I_A=10+6|\cos\eta|$、$I_{\min}/I_A=10-6|\cos\eta|$，所以$\mathcal V=3|\cos\eta|/5$。即使偏振平行，最暗仍为$4I_A$，因为$3A-A=2A$。

一般两束完全相干线偏振光的可见度为$2\sqrt{I_1I_2}|\cos\eta|/(I_1+I_2)$。由$2\sqrt{I_1I_2}\leq I_1+I_2$与$|\cos\eta|\leq1$，完全暗点要求两不等式同时取等号：两束等强、偏振共线，再选择相反的复场相位。若把偏振矢量方向反转，振幅相位也随之改变$\pi$，应以两束总电场反向为最终判据。部分相干时交叉项还要乘$|\gamma_{12}|$。
\end{solution}

\begin{exercise}\label{ex:o2filmindex}
空气—薄膜—基底的折射率依次为$1,n_f,n_s$，$1<n_f<n_s$。垂直照射，求使反射最弱的最小膜厚；进一步求两反射振幅可完全抵消的折射率条件。
\Diagram{fig:o2antirefq}{法向入射时两反射同向；分列画出各束以显示传播次序}{
\draw[curve](-2,0)--(2.4,0);\draw[curve](-2,-.9)--(2.4,-.9);
\node[left]at(-2,1.25){$1$};\node[left]at(-2,-.45){$n_f$};\node[left]at(-2,-1.3){$n_s$};
\draw[->](-1.2,1.5)--(-1.2,0)node[midway,left]{入射};
\draw[->](-.65,0)--(-.65,1.5)node[above]{$r_{01}$};
\draw[->](.35,.2)--(.35,-.9);
\draw[->](.9,-.9)--(.9,1.5)node[above]{$r_{12}\ee^{2i\phi}$};
\draw[<->](1.9,0)--(1.9,-.9)node[midway,right]{$h$};
}
\end{exercise}
\begin{solution}
图\ref{fig:o2antirefq}中两界面均由低折射率向高折射率反射，$r_{01},r_{12}$都为负；两反射的相对初相位为零。要使两束因传播反相，最小正厚度满足$2n_fh=\lambda_0/2$，故$h=\lambda_0/(4n_f)$。这给出相位条件，完全消反还需振幅匹配。

令单程膜内相位为$\phi$。纳入多次反射，首个由下界面返回的振幅为$t_{01}t_{10}r_{12}\ee^{2i\phi}$，后续倍率为$r_{10}r_{12}\ee^{2i\phi}$。因$r_{10}=-r_{01}$、$t_{01}t_{10}=1-r_{01}^2$，
\[
 r=r_{01}+\frac{(1-r_{01}^2)r_{12}\ee^{2i\phi}}
 {1+r_{01}r_{12}\ee^{2i\phi}}
 =\frac{r_{01}+r_{12}\ee^{2i\phi}}{1+r_{01}r_{12}\ee^{2i\phi}}.
\]
\Diagram{fig:o2antirefa}{四分之一波膜让往返传播相位成为$\pi$}{\draw[axis](0,0)--(2.7,0)node[right]{$r_{01}$};\draw[axis](0,0)--(-2.7,0)node[left]{$r_{12}\ee^{i\pi}$};\node[above]at(0,1){两反射相量等大反向时消失};}
在图\ref{fig:o2antirefa}的四分之一波厚度，$\ee^{2i\phi}=-1$，消反要求$r_{01}=r_{12}$。代入界面系数并交叉相乘，
\[
 (1-n_f)(n_f+n_s)=(n_f-n_s)(1+n_f)
 \quad\Longrightarrow\quad n_f^2=n_s.
\]
因此$n_f=\sqrt{n_s}$。若入射介质改为$n_0$，相同推导给$n_f^2=n_0n_s$。匹配只针对给定波长和法向入射；色散或斜入射会改变相位与两偏振的导纳。若折射率不匹配，四分之一波厚度仍使反射最小，但最低反射率为$[(r_{01}-r_{12})/(1-r_{01}r_{12})]^2$，一般不为零。
\end{solution}

\begin{exercise}\label{ex:o2ringliquid}
同一牛顿环装置在空气和未知液体中用同一波长测得相隔$p$级暗环的直径平方差分别为$D_a,D_l$。求液体折射率，说明如何消去中心小间隙影响。
\Diagram{fig:o2liquidq}{同一球面间隙先后由空气和液体填充}{\draw[curve](-2.3,0)--(2.3,0);\draw[curve,domain=-2:2,samples=50]plot(\x,{.15+.18*\x*\x});\node at(0,.6){$n$};\node[above]at(0,1.4){几何$R$不变};}
\end{exercise}
\begin{solution}
图\ref{fig:o2liquidq}中暗环直径$d_m=2r_m$，所以$d_{m+p}^2-d_m^2=4pR\lambda_0/n$。空气近似$n_a=1$，两次相除得$n=D_a/D_l$。
\Diagram{fig:o2liquida}{直径平方差消除所有环共同的厚度偏移}{\draw[axis](0,0)--(4.5,0)node[right]{级次$m$};\draw[axis](0,-.5)--(0,3)node[above]{$d_m^2$};\draw[curve](.7,0)--(3.8,2.6);\draw[curve](.7,0)--(4.2,1.6);\node[right]at(3.8,2.6){空气};\node[right]at(4.2,1.6){液体};}
中心间隙$h_0$只使$d_m^2$增加共同截距$-8Rh_0$，图\ref{fig:o2liquida}的斜率不受影响。如果充液引起透镜受力、曲率或接触状态改变到不能视为同一$R$，该消元就不成立；应使用同一机械装置、相同测量级差。
\end{solution}

\begin{exercise}\label{ex:o2doubleline}
迈克耳孙干涉仪受两条等强、互不相干的近邻谱线$\lambda$和$\lambda+\Delta\lambda$照射。求随光程差变化的条纹可见度包络，以及相邻可见度消失位置之间的镜移量。
\Diagram{fig:o2doubleq}{两频率各形成条纹，记录器相加强度}{\draw[curve](0,0)rectangle(2,1);\node at(1,.5){两条谱线};\draw[->](2,.5)--(3,.5);\draw[curve](3,0)rectangle(5.4,1);\node at(4.2,.5){光程差$\mathcal D$};}
\end{exercise}
\begin{solution}
取每条谱线的条纹强度为$I_0[1+\cos(k_i\mathcal D)]$。图\ref{fig:o2doubleq}的和为$2I_0[1+\cos((k_1-k_2)\mathcal D/2)\cos((k_1+k_2)\mathcal D/2)]$，所以包络可见度为$|\cos(\Delta k\mathcal D/2)|$。
\Diagram{fig:o2doublea}{快速干涉条纹受慢变化拍频包络调制}{\draw[axis](0,0)--(6.5,0)node[right]{$\mathcal D$};\draw[curve,domain=0:6,samples=170]plot(\x,{1+cos(60*\x)*cos(720*\x)});\draw[guide,domain=0:6,samples=80]plot(\x,{1+abs(cos(60*\x))});\draw[guide,domain=0:6,samples=80]plot(\x,{1-abs(cos(60*\x))});}
图\ref{fig:o2doublea}相邻包络零点相差$\Delta\mathcal D=2\pi/|\Delta k|$。由于$|\Delta k|\simeq2\pi|\Delta\lambda|/\lambda^2$，得到$\Delta\mathcal D\simeq\lambda^2/|\Delta\lambda|$。单镜位移是光程差改变的一半，所以$\Delta d\simeq\lambda^2/(2|\Delta\lambda|)$。不能把两条不同频率的复电场长期相干相加来计算静态条纹。
\end{solution}

\begin{exercise}\label{ex:o2Nwidth}
有$N=100$条等间距相干缝，周期$d=10\,\mu\mathrm m$，波长$500\,\mathrm{nm}$。求一级主极大角度、首零半宽及瑞利可分辨的最小波长差。
\Diagram{fig:o2Nq}{有限数目缝的主峰角宽}{\draw[curve](0,-1.3)--(0,1.3);\foreach\y in{-1,-.5,0,.5,1}{\draw[->](0,\y)--(3,{\y+.8});}\node[above]at(0,1.5){$N$缝};\draw[guide](0,0)--(4,0);\node[right]at(3,.8){$\theta_1$};}
\end{exercise}
\begin{solution}
图\ref{fig:o2Nq}的相邻缝相位差为$\delta=2\pi d\sin\theta/\lambda$。一级主极大对应$\delta=2\pi$，所以$\sin\theta_1=.5\,\mu\mathrm m/(10\,\mu\mathrm m)=.05$、$\theta_1=2.866^\circ$。在主峰附近令$\delta=2\pi+\epsilon$，阵列因子首零满足$N\epsilon/2=\pm\pi$，即$|\epsilon|=2\pi/N$。对角度线性展开，
\[
 \Delta\theta_0=\frac{\lambda}{Nd\cos\theta_1}
 =5.006\times10^{-4}\ \mathrm{rad}.
\]
\Diagram{fig:o2Na}{主峰到相邻首零的距离与另一谱线峰位比较}{\draw[axis](-3,0)--(3.7,0)node[right]{角度偏移};\draw[curve,domain=-2.7:-.01,samples=100]plot(\x,{2.3*(sin(180*\x)/(pi*\x))^2});\draw[curve,domain=.01:2.7,samples=100]plot(\x,{2.3*(sin(180*\x)/(pi*\x))^2});\draw[curve](-.01,2.29924)--(0,2.3)--(.01,2.29924);\draw[guide](1,0)--(1,2.3);\node[below]at(0,0){$\theta_1$};\node[below]at(1,0){首零};}
图\ref{fig:o2Na}中的瑞利条件要求一条谱线的主峰落在另一条首零处。同级角色散由$d\sin\theta=m\lambda$给$\Delta\theta\simeq m\Delta\lambda/(d\cos\theta)$。令它等于首零半宽，得到
\[
 \frac{m\Delta\lambda}{d\cos\theta}
 =\frac{\lambda}{Nd\cos\theta},\qquad
 \Delta\lambda_{\min}=\frac{\lambda}{mN}=5\,\mathrm{nm}.
\]
此处$m=1$，分辨本领100。首零半宽不是半高全宽；两者都随$1/N$变化，但系数不同。题设未给单缝宽度，因此只计算阵列峰的宽度和理想瑞利分辨率；该级能否被包络压低，还需单缝孔径信息。
\end{solution}


\chapter{衍射积分、孔径与光栅分辨率}
\section{电磁波、复数平均与光强}\label{th:o3wave}
在均匀、线性、各向同性、无自由电荷电流的介质中，$\vct D=\varepsilon\vct E$、$\vct B=\mu\vct H$，麦克斯韦方程给$\nabla\cdot\vct E=0$。对法拉第定律取旋度，利用$\nabla\times\nabla\times\vct E=\nabla(\nabla\cdot\vct E)-\nabla^2\vct E$，得到
\[
 \nabla^2\vct E-\varepsilon\mu\partial_t^2\vct E=0.
\]
代入$\vct E=\vct E_0\ee^{i(\vct k\cdot\vct r-\omega t)}$，得到$\vct k\cdot\vct E_0=0$、$k^2=\varepsilon\mu\omega^2$。相速度$v=1/\sqrt{\varepsilon\mu}=c/n$，磁场$\vct B_0=\vct k\times\vct E_0/\omega$，所以$\vct E,\vct H,\vct k$构成定向正交关系。若介质吸收或各向异性，这些实数和共线关系需重新讨论。

\subsection{时间平均中的共轭}\label{th:o3intensity}
同频复振幅$A,B$的实部乘积平均为$\langle\operatorname{Re}(A\ee^{-i\omega t})\operatorname{Re}(B\ee^{-i\omega t})\rangle=\tfrac12\operatorname{Re}(AB^*)$。可把每个实部写成两个指数之和，乘积中的$\ee^{\pm2i\omega t}$在整周期平均中消失，只留下互为共轭的两项。

因而平均坡印廷矢量$\langle\vct S\rangle=\tfrac12\operatorname{Re}(\vct E_0\times\vct H_0^*)$。无损平面波的光强为$I=E_0^2/(2\mu v)$。非磁性介质$\mu=\mu_0$时，$I=n\varepsilon_0cE_0^2/2$。这个强度沿传播方向定义；穿过倾斜表面的法向能流还要乘余弦。

复振幅叠加先做矢量求和，再作时间平均。若两同偏振场为$A_1\ee^{i\phi_1}$和$A_2\ee^{i\phi_2}$，合相位应由实虚部使用$\operatorname{atan2}$确定，不能把各自相位直接平均。只有振幅相同、相位差在适当支路内，角平分线解释才可直接使用。

\section{从惠更斯积分到远场傅里叶积分}\label{th:o3integral}
设孔径面为$z=0$，孔内复振幅为$U(Q)$。标量近似下，观察点$P$的场可写为
\[
 U(P)\simeq\frac1{i\lambda}\iint_\Sigma
 U(Q)\frac{\ee^{ikr}}r\frac{\cos\theta+\cos\theta_0}2\dd S,
\]
其中$r=|P-Q|$，$\theta,\theta_0$分别描述出射和入射方向与孔面法向的夹角。该式忽略孔边附近的严格矢量边界场，适于孔径远大于波长、观察方向非极端掠射的情形。

观察屏在$z=L$，孔径坐标$(\xi,\eta)$、观察坐标$(x,y)$。近轴距离展开
\[
 r\simeq L+\frac{(x-\xi)^2+(y-\eta)^2}{2L}.
\]
幅度因子$1/r$可近似$1/L$，但指数中的小距离会乘大波数，不能同时全部省略。于是
\begin{align}\label{eq:o3fresnel}
 U(x,y)&\simeq\frac{\ee^{ikL}}{i\lambda L}
 \ee^{ik(x^2+y^2)/(2L)}\\
 &\quad\times\iint U(\xi,\eta)
 \ee^{ik(\xi^2+\eta^2)/(2L)}
 \ee^{-ik(x\xi+y\eta)/L}\dd\xi\dd\eta.\notag
\end{align}
保留孔径二次相位是菲涅耳衍射。若孔径特征半宽$a$满足$a^2/(\lambda L)\ll1$，二次相位变化小，可略去，得到夫琅禾费积分。薄透镜后焦面在近轴模型中也可给出孔径的傅里叶分布，无须真的把屏移到无限远；其相位补偿已纳入透镜传播。

\Diagram{fig:o3integral}{孔径点$Q$到屏上点$P$的路径及坐标}{\draw[curve](0,-1.6)--(0,1.7);\draw[curve](5,-1.6)--(5,2.2);\draw[guide](0,0)--(5,0);\fill(0,.6)circle(1.4pt)node[left]{$Q(\xi,\eta)$};\fill(5,1.4)circle(1.4pt)node[right]{$P(x,y)$};\draw[curve](0,.6)--(5,1.4)node[midway,above]{$r$};\node[below]at(2.5,0){$L$};}

\section{单缝的积分与相量圆弧}\label{th:o3slit}
均匀照亮宽度$a$的长直缝，先只处理宽度方向。沿$x$方向观察角为$\theta$，孔内坐标$\xi\in[-a/2,a/2]$，忽略共同相位因子，
\[
 U(\theta)=C\int_{-a/2}^{a/2}\ee^{-ik\xi\sin\theta}\dd\xi
 =C\left[\frac{\ee^{-ik\xi\sin\theta}}{-ik\sin\theta}\right]_{-a/2}^{a/2}
 =Ca\frac{\sin\beta}{\beta},\quad \beta=\frac{\pi a\sin\theta}\lambda.
\]
因此$I(\theta)=I(0)(\sin\beta/\beta)^2$。$\beta=0$取极限1，不能因分子为零判作暗纹。暗纹$a\sin\theta_m=m\lambda$，$m=\pm1,\pm2,\ldots$。中央主瓣首零全角宽近轴为$2\lambda/a$。

若把缝坐标改为$[0,a]$，合场多一个$\ee^{-i\beta}$因子，但强度不变；对后续与其他孔径的相干叠加，这个位置相位则不能删掉。场振幅的实值sinc随旁瓣交替变号，强度始终非负。画图时需说明纵轴是场、场模还是强度。

\begin{example}\label{ex:o3phasor}
把单缝分为无穷多个等宽小元，以相量首尾相接求合场，解释第一个暗纹和第一旁瓣场的反号。
\Diagram{fig:o3phasorq}{缝上各小元到同一远场方向的相位线性变化}{\draw[curve](0,-1.6)--(0,1.6);\foreach\y in{-1.4,-.7,0,.7,1.4}{\draw[->](0,\y)--(3,{\y+.8});}\node[left]at(0,0){$a$};\node[right]at(3,.8){$\theta$};}
\end{example}
\begin{solution}
图\ref{fig:o3phasorq}中相邻小元振幅相等、相位递增相同，首尾相接形成圆弧。整缝相位变化为$ka\sin\theta=2\beta$，弧长与轴上场$Ca$成正比。若相量圆半径为$R_p$，弧长$2R_p\beta=Ca$，合场弦长为$2R_p\sin\beta=Ca\sin\beta/\beta$。
\Diagram{fig:o3phasora}{相量圆弧的端点连线是合场；一整圆对应首零}{\draw[curve](0,0)circle(1.4);\draw[->,thick](1.4,0)arc(0:150:1.4);\draw[->,very thick](1.4,0)--(-1.212,.7);\draw[guide](0,0)--(1.4,0);\draw[guide](0,0)--(-1.212,.7);\node[above]at(0,1.7){圆心角$2\beta$};\node[below]at(0,-1.7){整圆时端点重合};}
图\ref{fig:o3phasora}的一整圆对应$2\beta=2\pi$，相量闭合为零。超过整圆后，合弦相对原公共相位反号，故第一旁瓣的实场sinc为负；模平方仍为正。上下半缝的面元配对可解释首零，圆弧构造还能给出任意角度的振幅与相位。
\end{solution}

\begin{example}\label{ex:o3rectangle}
矩形孔径宽$a$、高$H$，均匀平面波正入射。求远场二维强度，并讨论长缝近似的适用方向。
\Diagram{fig:o3rectangleq}{以孔中心为坐标原点的矩形孔径}{
\draw[axis](-1.8,0)--(2,0)node[right]{$\xi$};
\draw[axis](0,-2.1)--(0,2.3)node[above]{$\eta$};
\draw[curve](-1.1,-1.6)rectangle(1.1,1.6);
\draw[<->](-1.1,-1.9)--(1.1,-1.9)node[midway,below]{$a$};
\draw[<->](1.5,-1.6)--(1.5,1.6);\node[right]at(1.5,.7){$H$};
\node[below left]at(0,0){$O$};
}
\end{example}
\begin{solution}
图\ref{fig:o3rectangleq}中取孔中心为原点，$-a/2<\xi<a/2$、$-H/2<\eta<H/2$。观察方向余弦为$l_x,l_y$，令$q_x=kl_x$、$q_y=kl_y$，场积分为
\[
 U=C\int_{-a/2}^{a/2}\ee^{-iq_x\xi}\dd\xi
 \int_{-H/2}^{H/2}\ee^{-iq_y\eta}\dd\eta.
\]
两个积分分别取原函数$\ee^{-iq_x\xi}/(-iq_x)$、$\ee^{-iq_y\eta}/(-iq_y)$，在$\pm a/2$和$\pm H/2$代入，得到$2\sin(q_xa/2)/q_x$和$2\sin(q_yH/2)/q_y$；零横向波矢处直接积分为$a,H$。于是
\[
 U=CaH\operatorname{sinc}(\pi a l_x/\lambda)
 \operatorname{sinc}(\pi H l_y/\lambda),
\]
其中$\operatorname{sinc}q=\sin q/q$，零点取连续极限1。
\Diagram{fig:o3rectanglea}{孔径越长的方向，首零边界越靠近中心}{\draw[axis](-2.4,0)--(2.6,0)node[right]{$l_x$};\draw[axis](0,-1.5)--(0,1.7)node[above]{$l_y$};\draw[curve](-1.8,-.45)rectangle(1.8,.45);\node[above]at(0,1){首零边界};\node[below]at(1.8,-.5){$\lambda/a$};\node[left]at(0,.6){$\lambda/H$};}
图\ref{fig:o3rectanglea}中$H>a$，所以沿$y$的角分布更窄：首零方向余弦分别为$|l_x|=\lambda/a$和$|l_y|=\lambda/H$。强度为两个sinc平方之积，中央矩形亮斑在宽孔对应方向上较窄。取$l_y=0$切面，第二因子等于1，才恢复一维长缝结果。

轴上场正比孔面积$aH$，轴上强度正比$(aH)^2$；但两方向角尺度分别含$1/a$和$1/H$，对整个远场积分后总功率仍只正比$aH$。若仅增大$H$，$x$切面的归一化形状不变，绝对峰强和$y$方向宽度却都改变。孔高不能在二维问题中被无条件省去。
\end{solution}

\section{双缝包络与有限宽光栅}\label{th:o3grating}
两条宽$a$的缝，其中心相隔$b$。利用平移对应相位因子，合场为$U=U_a[1+\ee^{-ikb\sin\theta}]$，所以
\[
 I=4I_a(0)\operatorname{sinc}^2\beta\cos^2\alpha,
 \quad\beta=\pi a\sin\theta/\lambda,\quad
 \alpha=\pi b\sin\theta/\lambda.
\]
对$N$条等宽等间距缝，同理
\begin{equation}\label{eq:o3grating}
 I=I_a(0)\operatorname{sinc}^2\beta
 \left[\frac{\sin N\alpha}{\sin\alpha}\right]^2.
\end{equation}
单缝包络决定每一级主峰获得多少光能，多缝因子决定主峰宽度和细结构。改变$a/b$会改变缺级和效率，改变被照亮缝数$N$才直接改变同级主峰宽度。这里所有缝照明振幅相同、相位线性递增；照明不均匀时须改成带权和。

\begin{example}\label{ex:o3missing}
双缝或光栅的周期$b=3a$。按干涉因子的主极大条件，求单缝中央包络内的级次，指出缺级，并画出主极大与包络的对应位置。另讨论传播条件$|\sin\theta|\le1$对这些级次的限制。
\Diagram{fig:o3missingq}{有限缝宽$a$和中心周期$b=3a$}{\draw[curve](0,-.5)rectangle(.5,1);\draw[curve](1.5,-.5)rectangle(2,1);\draw[<->](0,-.9)--(.5,-.9)node[midway,below]{$a$};\draw[<->](.25,1.4)--(1.75,1.4)node[midway,above]{$b$};}
\end{example}
\begin{solution}
干涉因子的主极大满足$b\sin\theta=m\lambda$，单缝零点满足$a\sin\theta=p\lambda$，$p\ne0$。图\ref{fig:o3missingq}取$b=3a$，两条件重合时$m=3p$，故在能够传播的级次中，$\pm3,\pm6,\ldots$缺级。
\Diagram{fig:o3missinga}{以$b\sin\theta/\lambda$为横轴，包络首零在$\pm3$}{\draw[axis](-3.7,0)--(3.9,0)node[right]{$m$};\draw[axis](0,0)--(0,2.7)node[above]{强度};\draw[curve,domain=.01:3.5,samples=70]plot(\x,{2.2*(sin(60*\x)/(pi*\x/3))^2});\draw[curve,domain=-3.5:-.01,samples=70]plot(\x,{2.2*(sin(60*\x)/(pi*\x/3))^2});\foreach\x in{-2,-1,1,2}{\draw[very thick](\x,0)--(\x,{2.2*(sin(60*\x)/(pi*\x/3))^2});}\draw[very thick](0,0)--(0,2.2);\node[below]at(3,0){$3$};\node[below]at(-3,0){$-3$};}
单缝中央包络严格内部为$|a\sin\theta/\lambda|<1$，代入主极大条件得到$|m|<3$，候选级次为$m=0,\pm1,\pm2$。还须逐个检验传播条件
\[
 |\sin\theta_m|=\frac{|m|\lambda}{b}\le1.
\]
当$b>2\lambda$时，五个候选级次都在非掠射方向传播；当$\lambda<b<2\lambda$时只剩$0,\pm1$；当$0<b<\lambda$时只剩零级。$b=\lambda$或$2\lambda$时，相应一级或二级位于掠射方向，是否被接收还取决于观察几何。一般应取整数集合$\{m:|m|<3,\ |m|\lambda/b\le1\}$。

图\ref{fig:o3missinga}的横轴是无量纲方向变量$b\sin\theta/\lambda$；图示范围需取$b>3\lambda$，才能同时画出两侧首零。竖线标干涉因子的主极大位置，平滑线标单缝包络；实际有限$N$的峰有有限宽。总强度是两因子的乘积，除对称中心等特殊位置外，其局部峰还可能因包络斜率而略有偏移。
\end{solution}

\section{分辨判据与孔径宽度}
\begin{example}\label{ex:o3sodium}
要在第$m$级分辨波长589.0 nm和589.6 nm两线，按瑞利判据至少需照亮多少条缝？若只照亮200条缝，最低可用级次为何？
\Diagram{fig:o3sodiumq}{两条谱线在同一级产生相邻主峰}{\draw[axis](0,0)--(5.5,0)node[right]{$\theta$};\draw[curve](1,0)..controls(1.5,2.5)and(1.8,2.5)..(2.4,0);\draw[curve](2,0)..controls(2.5,2.5)and(2.8,2.5)..(3.4,0);\node[above]at(1.7,2){589.0};\node[above]at(2.7,2.7){589.6};}
\end{example}
\begin{solution}
取平均波长$\bar\lambda=589.3$ nm、差$\Delta\lambda=.6$ nm，所需分辨本领约982.17。图\ref{fig:o3sodiumq}按瑞利比较要求$mN\geq982.17$，因此$N_{\min}=\lceil982.17/m\rceil$。
\Diagram{fig:o3sodiuma}{提高级次缩小所需缝数，但必须有该传播级次}{\draw[axis](0,0)--(4.8,0)node[right]{$m$};\draw[axis](0,0)--(0,3.2)node[above]{$N_{\min}$};\foreach\x/\y/\t in{1/2.95/983,2/1.476/492,3/.984/328,4/.738/246}{\fill(\x,\y)circle(1.5pt);\node[right]at(\x,\y){$\t$};}}
一级至少983条，二级492条。若$N=200$，则$m\geq\lceil982.17/200\rceil=5$。图\ref{fig:o3sodiuma}的提高级次并非无条件可行，还须$|m\lambda/b|\leq1$、该级不缺级且效率足够。只满足$mN$而传播角不存在时，不能实现分辨。
\end{solution}

\begin{example}\label{ex:o3sparrow}
两条等强、互不相干谱线产生相同的窄主峰，单线近似$F(x)=(\sin x/x)^2$。两峰中心为$\pm a$。以总曲线中点由极大转为极小作为双峰边界，推导所需间距，并与瑞利间距比较。
\Diagram{fig:o3sparrowq}{两条单线的强度相加，峰中心相距$2a$}{\draw[axis](-3,0)--(3.4,0)node[right]{$x$};\draw[curve,domain=-2.7:2.7,samples=90]plot(\x,{1.5*exp(-2*(\x-1)^2)});\draw[guide,domain=-2.7:2.7,samples=90]plot(\x,{1.5*exp(-2*(\x+1)^2)});\node[below]at(-1,0){$-a$};\node[below]at(1,0){$a$};}
\end{example}
\begin{solution}
图\ref{fig:o3sparrowq}只示意两中心位置，定量使用sinc平方。总强度$G(x)=F(x-a)+F(x+a)$为偶函数，$G'(0)=0$、$G''(0)=2F''(a)$。曲率边界为$F''(a)=0$。完整求导给
\[
 F'(x)=\frac{\sin2x}{x^2}-\frac{2\sin^2x}{x^3},
\]
\[
 F''(x)=\frac{2\cos2x}{x^2}-\frac{4\sin2x}{x^3}
 +\frac{6\sin^2x}{x^4}.
\]
在$x=a$处通分并用$2\sin^2a=1-\cos2a$，得
\[
 F''(a)=\frac{(2a^2-3)\cos2a-4a\sin2a+3}{a^4}.
\]
取第一非零正根，数值$a\simeq1.30308171$，所以两中心间距$2a\simeq2.60616$。
\Diagram{fig:o3sparrowa}{曲率判据的临界中点变平，超过临界后出现中央凹陷}{\draw[axis](-3,0)--(3.3,0)node[right]{$x$};\draw[axis](0,0)--(0,2.6)node[above]{$G$};\draw[curve,domain=-2.8:2.8,samples=120]plot(\x,{(sin((\x-1.30308) r)/(\x-1.30308))^2+(sin((\x+1.30308) r)/(\x+1.30308))^2});\draw[guide](-.6,1.096)--(.6,1.096);}
瑞利判据的中心间距是单峰中心到首零的$\pi$，故图\ref{fig:o3sparrowa}对应的曲率间距约为瑞利间距的$2.60616/\pi\simeq0.82957$。求导时$x^{-2}$的变化与$\sin^2x$的变化同阶，不能把分母固定后仅对分子求导。此判据还依赖两线等强、无噪声和相同形状；实际能否可靠辨认需结合测量条件。
\end{solution}

\begin{example}\label{ex:o3width}
光栅照明总宽度$W$，周期$b$，正常入射。证明任意可传播级次的瑞利分辨本领受$W/\lambda$限制，并解释为何无限缩小周期不能无限提高分辨率。
\Diagram{fig:o3widthq}{照明宽度决定实际参与的缝数}{\draw[curve](0,-.5)rectangle(5,.7);\foreach\x in{.2,.6,1,1.4,1.8,2.2,2.6,3,3.4,3.8,4.2,4.6}{\draw[curve](\x,-.5)--(\x,.7);}\draw[<->](0,-1)--(5,-1)node[midway,below]{$W\simeq Nb$};}
\end{example}
\begin{solution}
图\ref{fig:o3widthq}中$N\simeq W/b$。主极大$b\sin\theta=m\lambda$要求$|m|\leq b/\lambda$。相乘得到$|m|N\lesssim W/\lambda$。
\Diagram{fig:o3widtha}{周期越小虽增加缝数，却降低最高可用级次}{\draw[axis](0,0)--(4.6,0)node[right]{$b/\lambda$};\draw[axis](0,0)--(0,3)node[above]{最高$|m|$};\draw[very thick](0,0)--(1,0)--(1,.7)--(2,.7)--(2,1.4)--(3,1.4)--(3,2.1)--(4,2.1);\draw[guide](0,0)--(4,2.8);}
图\ref{fig:o3widtha}中的级次是整数，实际最大值为$\lfloor b/\lambda\rfloor$。当$b<\lambda$只剩零级，不能靠通常角色散分辨不同波长。上限在掠射方向逼近，近轴角宽公式及实际效率此时需要更细处理；这一上限由实际参与相干叠加的孔径宽度确定。
\end{solution}

\section{闪耀反射光栅}\label{th:o3blaze}
反射光栅每周期的小面可视为有限平面镜。设宏观法线沿$+z$，入射光沿$-z$，小面法线向$+x$倾斜$\alpha$，出射方向相对$+z$成$\theta$。小面切向单位矢量为$(\cos\alpha,-\sin\alpha)$，其上坐标$s$的相位梯度为
\[
 (\vct k_i-\vct k_o)\cdot\vct t
 =k[\sin\alpha-\sin(\theta-\alpha)].
\]
单小面包络在梯度为零时最强，即$\theta=2\alpha$，正是镜面反射方向。周期阵列的主峰仍由$b\sin\theta_m=m\lambda$确定；使某级落在包络峰上，闪耀波长为$\lambda_B=b\sin2\alpha/m$。阴影遮挡、偏振及真实沟槽形状会影响效率，简单标量面元积分只描述主要相位匹配。

\begin{example}\label{ex:o3blaze}
正常入射的闪耀光栅周期$b=2\,\mu\mathrm m$，小面倾角$\alpha=15^\circ$。求二级闪耀波长，并比较此处相邻周期相位和单小面内相位各自的作用。
\Diagram{fig:o3blazeq}{倾斜小面使镜面包络偏向某一衍射级}{\foreach\x in{0,1.5,3,4.5}{\draw[curve](\x,0)--({\x+1.1},-.3)--({\x+1.5},0);}\draw[->](2,2)--(2,-.14);\draw[->](2,-.14)--(3.3,2.11);\node[above]at(2,2){入射};\node[right]at(3.3,2.11){$2\alpha$};}
\end{example}
\begin{solution}
图\ref{fig:o3blazeq}的包络中心$\theta=30^\circ$，二级光栅条件$2\lambda_B=b\sin30^\circ$，所以$\lambda_B=.5\,\mu\mathrm m=500$ nm。
\Diagram{fig:o3blazea}{相同多缝窄峰乘以偏移包络，强度集中在选定级次}{\draw[axis](0,0)--(5.8,0)node[right]{衍射级次};\draw[curve,domain=0:5,samples=80]plot(\x,{2.5*exp(-1.1*(\x-2)^2)});\foreach\x in{0,1,2,3,4,5}{\draw[very thick](\x,0)--(\x,{2.5*exp(-1.1*(\x-2)^2)});}\node[below]at(2,0){$m=2$};}
小面内同相叠加使该角度单缝包络大，周期之间相差$4\pi$使二级主峰相长，图\ref{fig:o3blazea}表示两条件同时满足。闪耀主要提高选定级的能量效率；固定$N,m$下，瑞利分辨本领仍是$mN$，不会因包络更高而自动变成更窄的几何峰。
\end{solution}
\section{衍射积分中的尺度与归一化}
\subsection{几何小量不等于相位小量}
传播距离可展开为$\sqrt{L^2+\rho^2}=L+\rho^2/(2L)-\rho^4/(8L^3)+\cdots$。在幅度因子$1/r$中，只要$\rho/L\ll1$就可近似$r=L$；在相位$kr$中，误差还乘$k=2\pi/\lambda$，要求$k\rho^4/(8L^3)$足够小才能只保留二次项。把三角近轴条件与相位截断条件分开，才能判断菲涅耳近似在特定尺寸下的精度。

夫琅禾费近似还要求孔面二次相位$k(\xi^2+\eta^2)/(2L)$的变化足够小，或已由透镜补偿。增大屏距、缩小孔径都可减小该项；仅让观察角很小还不够。高斯束、直边及大孔径近场计算通常仍须保留二次相位。

在孔径前还有有限距离点光源时，入射到孔面已带一个二次球面相位。传播到屏的二次相位与它相加，等效距离由两段波前曲率共同决定。不能只看孔到屏的$L$来断定远场；若点光源很近，入射曲率项可能仍很大。理想透镜可用反号二次相位补偿相应曲率，形成清楚的傅里叶变换面。

\subsection{远场坐标的量纲}
一维远场场可写$U(x)=C F(q)$，其中$F(q)=\int U_0(\xi)\ee^{-iq\xi}\dd\xi$，$q=kx/L$。$q$量纲为长度倒数，$x$是屏幕长度。若使用空间频率$f_x=x/(\lambda L)$，指数应写$\ee^{-i2\pi f_x\xi}$。这两套约定只差$2\pi$，但单缝零点和高斯宽度的常数必须随之同步。

\Diagram{fig:o3scales}{孔面长度尺度与远场空间频率尺度互为倒数}{\draw[axis](0,1.4)--(5,1.4)node[right]{$\xi$};\draw[very thick](1.5,1.4)--(3.5,1.4);\node[above]at(2.5,1.4){宽$a$};\draw[axis](0,0)--(5,0)node[right]{$q$};\draw[very thick](2,0)--(3,0);\node[below]at(2.5,0){宽度尺度$1/a$};}
图\ref{fig:o3scales}说明孔越宽、角谱越窄，但比较绝对亮度还要保留孔面照明归一化。均匀单缝振幅正比$a$，所以中央强度随$a^2$增长。把每条曲线都除以各自峰值后只剩形状，不能再从归一化图上读取绝对功率变化。

\section{傅里叶积分中的能量守恒}
对平方可积孔场，采用上述无$2\pi$的傅里叶约定，帕塞瓦尔关系为
\[
 \int_{-\infty}^{\infty}|F(q)|^2\dd q
 =2\pi\int_{-\infty}^{\infty}|U_0(\xi)|^2\dd\xi.
\]
它可从双重积分展开$|F|^2$后，使用$\int\ee^{-iq(\xi-\eta)}\dd q=2\pi\delta(\xi-\eta)$得到。自由传播核的系数与屏面变量$\dd x=L\dd q/k$相配合，使积分能流等于孔面透过能流。二维情况沿两个方向各用一次同样关系。

单缝可直接核对功率归一化：孔面振幅恒定时，$F(q)=aU_0\operatorname{sinc}(qa/2)$。换元$u=qa/2$，得$\int|F|^2\dd q=2a|U_0|^2\int\operatorname{sinc}^2u\dd u$，积分正比于$a$。孔宽增大使中央强度升高、角分布变窄，二者共同满足总功率守恒。

\subsection{相位孔与振幅孔}
若孔面仅乘单位模相位$\ee^{i\phi(\xi)}$，总透过功率不变，但远场分布可明显改变。线性相位整体转向，二次相位改变聚焦或观察面，周期相位形成多个衍射级。若孔面乘振幅透过率$|t(\xi)|<1$，则总功率减少；它仍可通过压低边缘权重改善旁瓣，代价是峰宽或通光量。

多束加权的二项式系数和连续高斯照明都是这种思想的离散、连续例子。比较不同孔径时应固定哪一个条件：总输入功率、最大照明强度、有效孔径宽度，还是透过总能量。若约束未说明，峰高、旁瓣与“效率”的比较往往没有唯一答案。

\section{点像、物像与分辨判据}
\subsection{单点响应怎样组合成像}
理想无像差透镜的几何像点在波动理论中仍有有限衍射斑。孔函数的傅里叶变换给出点源的复振幅响应$h(x)$，其强度响应为$|h(x)|^2$。多个彼此不相干物点成像时，各点像强度相加；若物点相干，则应先把复振幅相加。两者即使几何点位完全相同，图样也可能因相对相位不同而改变。

\Diagram{fig:o3pointimages}{不相干物点的点像相加强度，相干物点则先加场}{\fill(0,.7)circle(1.5pt)node[left]{$O_1$};\fill(0,-.7)circle(1.5pt)node[left]{$O_2$};\draw[<->](2,-1.5)--(2,1.5);\draw[curve](4.5,.7)circle(.45);\draw[curve](4.5,-.7)circle(.45);\draw[guide](0,.7)--(4.5,-.7);\draw[guide](0,-.7)--(4.5,.7);}
图\ref{fig:o3pointimages}只标出点像中心与宽度。对两个等强不相干谱线或星点，瑞利判据使用主峰与首暗环重合；对两相干点，交叉项会使中间强度取决于相位，不能沿用同一张强度相加曲线。

\subsection{瑞利、曲率和测量误差}
瑞利判据以主峰与首零点重合作为分离尺度。曲率判据则考察中央由峰变谷的边界，在对称等强模型中令二阶导数为零。对sinc平方线形，后者间距约为瑞利尺度的.8296倍；高斯、圆孔等线形须另求这个系数。

实际谱线很弱时，另一个强峰的翼部与背景噪声可能比中心曲率更重要。高信噪比、已知线形的拟合有时能估计小于传统瑞利间隔的峰位，但这不等于光学点像变窄；估计不确定度取决于光子数、噪声、模型先验和参数相关。竞赛题若指定某种判据，应按该定义计算，不把不同定义混作唯一“最小可分辨距离”。

\subsection{光栅参数中的三个长度}
缝宽$a$、周期$b$、照明总宽$W\simeq Nb$决定不同结构。$a$给单缝包络和缺级；$b$给各级中心角和最高可传播级次；$W$给同级窄峰的宽度。若固定$W$而缩小$b$，虽然$N$增大，可用级次却下降；若固定$b$而增加$N$，主峰中心不动，宽度变窄。

闪耀小面的倾角把单元包络峰移向选定级次，是效率设计；增加总照明宽度才真正增加可参与的相干横向范围。真实光栅还受表面误差、局部周期误差和照明相位不均影响，实际窄峰可能比理想$mN$更宽。这些误差应按具体分布进入孔面场，再求实际线形。

\subsection{孔径、相干照明和点像强度}
夫琅禾费场是孔面复场的傅里叶变换，而非孔面强度的傅里叶变换。若孔面透过率为$t(\xi)$、照明复场为$A(\xi)$，实际参与积分的是$A(\xi)t(\xi)$。均匀照明、振幅加权和附加相位因而改变不同对象：加权改变面元振幅，相位片改变面元方向，遮挡则同时改变通光范围。

对若干点源成像，还要在像面区分相干与不相干叠加。若两个点源经同一系统形成复振幅点响应$h(x-x_1)$、$h(x-x_2)$，完全相干时强度为
\[
 |a_1h(x-x_1)+a_2\ee^{i\phi}h(x-x_2)|^2.
\]
互不相干时交叉项长期平均为零，才得到$|a_1|^2|h(x-x_1)|^2+|a_2|^2|h(x-x_2)|^2$。瑞利判据和本讲等强双峰曲率判据采用后一种模型；在孔径和点源间距固定时，相干性仍会改变中央峰谷。
\Diagram{fig:o3coherentpairs}{实线为同相合场，曲线虚线为反相合场；竖直虚线标出点像中心}{\draw[axis](-.3,0)--(5.7,0)node[right]{$x$};\draw[axis](0,-1)--(0,2.7)node[above]{场};\draw[guide](2,-.4)--(2,2.3);\draw[guide](3.2,-.4)--(3.2,2.3);\draw[curve,domain=.2:5.2,samples=100]plot(\x,{exp(-2*(\x-2)^2)+exp(-2*(\x-3.2)^2)});\draw[dashed,domain=.2:5.2,samples=100]plot(\x,{exp(-2*(\x-2)^2)-exp(-2*(\x-3.2)^2)});\node at(4.7,2){同相};\node at(4.7,-.7){反相};}
图\ref{fig:o3coherentpairs}为平滑实场的示意，实线同相相加，虚线反相相加；它不代替具体圆孔或单缝积分。使用测得的强度曲线讨论分辨时，应先明确各点响应、光源统计和探测平均方式，再选择判据。

远场强度还能检验算式归一化：对完整角域或相应近轴傅里叶平面积分，应与孔面透过功率相符。若只截取主瓣，则所得只是该区域功率；忽略旁瓣后不能继续把主瓣积分称为总功率。若孔径尺寸接近波长、观察角接近掠射或介质边界引入明显矢量效应，标量近轴模型的适用性也须重新判断。

\section{近场边缘与孔径的空间频率}
夫琅禾费近似把孔面场变为空间频率的傅里叶变换；菲涅耳近似则保留每个孔元的二次相位。二者使用同一传播积分，区别在于可省略哪些相位。描述近场时，不能把每条射线的几何阴影边界当成强度突然为零的位置。

直边衍射需要的振荡积分也可由第一讲的复高斯积分确定。先乘收敛因子$\ee^{-\epsilon t^2}$，$\epsilon>0$，在式\eqref{eq:o1gaussintegral}中取$a=\epsilon-i\pi/2$、$b=0$，
\[
 \int_{-\infty}^{\infty}\ee^{-\epsilon t^2+i\pi t^2/2}\dd t
 =\sqrt{\frac{\pi}{\epsilon-i\pi/2}}
 \xrightarrow{\epsilon\to0^+}\sqrt2\ee^{i\pi/4}=1+i.
\]
平方根沿$\operatorname{Re}a>0$连续取支。该收敛极限给出菲涅耳振荡积分；被积函数为偶函数，所以$\int_0^\infty\ee^{i\pi t^2/2}\dd t=(1+i)/2$。实、虚部分别给$C(\infty)=S(\infty)=1/2$，用于下面的无遮挡场归一化。

\begin{example}\label{ex:o3edge}
平面波照射遮住$x<0$的直边，光沿$+z$传播。距边$L$的屏上坐标为$x$，定义$v=x\sqrt{2/(\lambda L)}$。用菲涅耳积分表示归一化强度，求几何边界$x=0$的强度。
\Diagram{fig:o3edgeq}{半平面边缘及几何阴影线}{\draw[curve](0,-1.5)--(0,0);\foreach\y in{-1.4,-1.1,-.8,-.5,-.2}{\draw(0,\y)--(-.3,{\y-.2});}\draw[curve](5,-1.8)--(5,2);\draw[guide](0,0)--(5,0);\draw[->](-1,1)--(1,1);\node[right]at(5,-1){几何阴影};\node[right]at(5,1){照明区};}
\end{example}
\begin{solution}
图\ref{fig:o3edgeq}的透光半平面为$\xi\geq0$。对沿边方向的积分提出共同因子后，垂直边方向场正比$\int_0^\infty\exp[i\pi(x-\xi)^2/(\lambda L)]\dd\xi$。令$t=(x-\xi)\sqrt{2/(\lambda L)}$，$\xi=0$对应$t=v$，$\xi\to\infty$对应$t\to-\infty$，负的微分将上下限反转，故积分正比$\int_{-\infty}^v\ee^{i\pi t^2/2}\dd t$。
\[
 C(v)=\int_0^v\cos(\pi t^2/2)\dd t,\quad
 S(v)=\int_0^v\sin(\pi t^2/2)\dd t.
\]
用$C(\infty)=S(\infty)=1/2$和奇函数性质，分子为$(1/2+C)+i(1/2+S)$。
\Diagram{fig:o3edgea}{按菲涅耳积分计算的直边衍射强度，几何边界处为四分之一}{\draw[axis](-3.2,0)--(3.7,0)node[right]{$v$};\draw[axis](0,0)--(0,3.2)node[above]{$I/I_0$};\draw[curve]plot coordinates {(-3.00000,0.012310) (-2.96425,0.012605) (-2.92849,0.012910) (-2.89274,0.013227) (-2.85698,0.013556) (-2.82123,0.013897) (-2.78547,0.014250) (-2.74972,0.014618) (-2.71397,0.014999) (-2.67821,0.015395) (-2.64246,0.015807) (-2.60670,0.016235) (-2.57095,0.016681) (-2.53520,0.017144) (-2.49944,0.017627) (-2.46369,0.018131) (-2.42793,0.018655) (-2.39218,0.019203) (-2.35642,0.019774) (-2.32067,0.020370) (-2.28492,0.020993) (-2.24916,0.021645) (-2.21341,0.022326) (-2.17765,0.023039) (-2.14190,0.023785) (-2.10615,0.024568) (-2.07039,0.025388) (-2.03464,0.026248) (-1.99888,0.027151) (-1.96313,0.028100) (-1.92737,0.029097) (-1.89162,0.030146) (-1.85587,0.031250) (-1.82011,0.032412) (-1.78436,0.033638) (-1.74860,0.034930) (-1.71285,0.036294) (-1.67709,0.037735) (-1.64134,0.039258) (-1.60559,0.040869) (-1.56983,0.042575) (-1.53408,0.044382) (-1.49832,0.046297) (-1.46257,0.048330) (-1.42682,0.050488) (-1.39106,0.052782) (-1.35531,0.055221) (-1.31955,0.057817) (-1.28380,0.060583) (-1.24804,0.063531) (-1.21229,0.066676) (-1.17654,0.070034) (-1.14078,0.073622) (-1.10503,0.077458) (-1.06927,0.081564) (-1.03352,0.085960) (-0.99777,0.090672) (-0.96201,0.095725) (-0.92626,0.101147) (-0.89050,0.106971) (-0.85475,0.113229) (-0.81899,0.119958) (-0.78324,0.127198) (-0.74749,0.134992) (-0.71173,0.143388) (-0.67598,0.152435) (-0.64022,0.162189) (-0.60447,0.172709) (-0.56872,0.184060) (-0.53296,0.196312) (-0.49721,0.209538) (-0.46145,0.223821) (-0.42570,0.239246) (-0.38994,0.255906) (-0.35419,0.273900) (-0.31844,0.293335) (-0.28268,0.314321) (-0.24693,0.336979) (-0.21117,0.361434) (-0.17542,0.387818) (-0.13966,0.416269) (-0.10391,0.446932) (-0.06816,0.479956) (-0.03240,0.515493) (0.00335,0.553700) (0.03911,0.594733) (0.07486,0.638752) (0.11061,0.685910) (0.14637,0.736357) (0.18212,0.790236) (0.21788,0.847678) (0.25363,0.908798) (0.28939,0.973691) (0.32514,1.042429) (0.36089,1.115053) (0.39665,1.191567) (0.43240,1.271932) (0.46816,1.356059) (0.50391,1.443804) (0.53966,1.534957) (0.57542,1.629234) (0.61117,1.726277) (0.64693,1.825637) (0.68268,1.926775) (0.71844,2.029055) (0.75419,2.131740) (0.78994,2.233992) (0.82570,2.334871) (0.86145,2.433339) (0.89721,2.528269) (0.93296,2.618454) (0.96872,2.702622) (1.00447,2.779462) (1.04022,2.847639) (1.07598,2.905834) (1.11173,2.952776) (1.14749,2.987280) (1.18324,3.008296) (1.21899,3.014955) (1.25475,3.006620) (1.29050,2.982935) (1.32626,2.943878) (1.36201,2.889802) (1.39777,2.821479) (1.43352,2.740125) (1.46927,2.647416) (1.50503,2.545493) (1.54078,2.436936) (1.57654,2.324730) (1.61229,2.212195) (1.64804,2.102900) (1.68380,2.000547) (1.71955,1.908834) (1.75531,1.831298) (1.79106,1.771145) (1.82682,1.731064) (1.86257,1.713059) (1.89832,1.718273) (1.93408,1.746857) (1.96983,1.797854) (2.00559,1.869154) (2.04134,1.957496) (2.07709,2.058541) (2.11285,2.167031) (2.14860,2.277011) (2.18436,2.382138) (2.22011,2.476043) (2.25587,2.552741) (2.29162,2.607058) (2.32737,2.635058) (2.36313,2.634413) (2.39888,2.604699) (2.43464,2.547578) (2.47039,2.466829) (2.50615,2.368210) (2.54190,2.259137) (2.57765,2.148196) (2.61341,2.044493) (2.64916,1.956893) (2.68492,1.893216) (2.72067,1.859441) (2.75642,1.859019) (2.79218,1.892370) (2.82793,1.956642) (2.86369,2.045788) (2.89944,2.150984) (2.93520,2.261397) (2.97095,2.365235) (3.00670,2.451002) (3.04246,2.508830) (3.07821,2.531726) (3.11397,2.516577) (3.14972,2.464759) (3.18547,2.382227) (3.22123,2.279016) (3.25698,2.168165) (3.29274,2.064152) (3.32849,1.981003) (3.36425,1.930323) (3.40000,1.919516)};\fill(0,.55)circle(1.5pt)node[left]{$1/4$};\draw[guide](-3,2.2)--(3.4,2.2);\node[right]at(3.4,2.2){$1$};}
无遮挡平面对应积分$-\infty<t<\infty$，值为$1+i$。因此图\ref{fig:o3edgea}使用的准确归一化为
\[
 \frac{U}{U_0}=\frac{(1/2+C)+i(1/2+S)}{1+i},\quad
 \frac I{I_0}=\frac{(1/2+C)^2+(1/2+S)^2}{2}.
\]
在$v=0$，分子为$(1+i)/2$，所以场为一半、强度为四分之一。$v\to-\infty$时$C,S\to-1/2$，深阴影强度趋零；$v\to+\infty$时均趋$1/2$，强度趋1。边界附近的空间尺度由$|v|\sim1$给出，即$|x|\sim\sqrt{\lambda L/2}$，所以屏越远，几何边缘被衍射展开的范围越大。正照、直边及均匀平面波是本归一化的条件。
\end{solution}

\begin{example}\label{ex:o3shiftproperty}
孔面场$U_0(\xi)$整体横移$\xi_0$。求夫琅禾费场的变化；若孔未移动而给孔面加线性相位$\exp(iq_0\xi)$，又怎样改变远场？
\Diagram{fig:o3shiftq}{孔的位置平移与孔内相位斜坡是两种不同操作}{\draw[axis](-1,0)--(5,0)node[right]{$\xi$};\draw[curve](0,0)rectangle(1.2,1);\draw[guide](2.5,0)rectangle(3.7,1);\draw[->](.6,1.5)--(3.1,1.5)node[midway,above]{$\xi_0$};}
\end{example}
\begin{solution}
取远场变量$q=k\sin\theta$，$F(q)=\int U_0(\xi)\ee^{-iq\xi}\dd\xi$。图\ref{fig:o3shiftq}平移孔为$U_0(\xi-\xi_0)$，令$s=\xi-\xi_0$，得$F_{\rm shift}(q)=\ee^{-iq\xi_0}F(q)$。单孔强度分布不变，只改变方向相关相位。
\Diagram{fig:o3shifta}{线性孔面相位使远场分布平移，而整体孔位置只加远场相位}{\draw[axis](-2.3,0)--(3.4,0)node[right]{$q$};\draw[curve,domain=-2:3,samples=70]plot(\x,{2*exp(-2*\x*\x)});\draw[guide,domain=-2:3,samples=70]plot(\x,{2*exp(-2*(\x-1.3)^2)});\node[above]at(1.3,2.1){$q_0$};}
若孔内乘$\ee^{iq_0\xi}$，积分变为$F(q-q_0)$，强度图样向$q_0$平移，如图\ref{fig:o3shifta}。斜入射照明的切向波矢或相位片的线性相位均按此式改变远场中心。孔整体平移在有另一参考孔相干叠加时会改变条纹，不能因单孔强度不变就删去其相位。
\end{solution}

\begin{example}\label{ex:o3twowidths}
两条中心相隔$b$的缝宽分别为$a_1,a_2$，单位面积照明场相同且无初始相差。求远场强度，说明可见度为何随角度变化。
\Diagram{fig:o3twowidthq}{不等缝宽带来不同角度包络和轴上振幅}{\draw[curve](0,-.6)rectangle(.6,1);\draw[curve](2,-.6)rectangle(3.2,1);\node[below]at(.3,-.6){$a_1$};\node[below]at(2.6,-.6){$a_2$};\draw[<->](.3,1.4)--(2.6,1.4)node[midway,above]{$b$};}
\end{example}
\begin{solution}
图\ref{fig:o3twowidthq}中每缝中心对称积分给实包络$F_j=a_j\operatorname{sinc}(\pi a_j\sin\theta/\lambda)$。合场$C[F_1+F_2\ee^{-ikb\sin\theta}]$，故$I=|C|^2[F_1^2+F_2^2+2F_1F_2\cos(kb\sin\theta)]$。
\Diagram{fig:o3twowidtha}{取$a_2=2a_1$时的有符号场包络，$\beta=\pi a_1\sin\theta/\lambda$}{
\draw[axis](-3.2,0)--(3.4,0)node[right]{$\beta$};
\draw[axis](0,-.65)--(0,2.7)node[above]{$F_j/a_1$};
\draw[curve,domain=-6.7:6.7,samples=181]plot({\x/2.3},{abs(\x)<.001 ? 1 : sin(deg(\x))/\x});
\draw[guide,domain=-6.7:6.7,samples=181]plot({\x/2.3},{abs(\x)<.001 ? 2 : sin(deg(2*\x))/\x});
\node[left]at(0,1){$1$};\node[left]at(0,2){$2$};
\node[below]at(1.366,0){$\pi$};\node[below]at(2.732,0){$2\pi$};
\draw[curve](1.6,2.2)--(2.2,2.2)node[right]{$F_1$};
\draw[guide](1.6,1.65)--(2.2,1.65)node[right]{$F_2$};}
在包络随一个纹距缓变的局部区域，$\mathcal V=2|F_1F_2|/(F_1^2+F_2^2)$。图\ref{fig:o3twowidtha}中两包络的零点不同：某一包络过零而另一不为零时，可见度降至0；过零后的$F_j$反号又使干涉纹相位翻转。若两者同时为零，该方向无光，不定义可见度。宽度比只给出轴上振幅比，离轴计算须保留各自的sinc因子。
\end{solution}

\section{连续孔径、二维光栅与采样}
\begin{example}\label{ex:o3circular}
半径$a$的均匀圆孔正入射。将远场场写成一维贝塞尔积分并求归一化形式，说明第一暗环如何决定角分辨尺度。
\Diagram{fig:o3circleq}{圆孔的径向对称性使远场只依赖离轴角}{\draw[curve](0,0)circle(1.6);\draw[axis](0,0)--(2.6,0)node[right]{$\xi$};\draw[axis](0,0)--(0,2.3)node[above]{$\eta$};\draw[->](0,0)--(1.1,1.1)node[right]{$r$};}
\end{example}
\begin{solution}
图\ref{fig:o3circleq}中，取观察方向的横向投影为极角零方向，$q=k\sin\theta$。孔面点的远场相位为$-qr\cos\varphi$，故
\[
 U=C\int_0^a r\dd r\int_0^{2\pi}\ee^{-iqr\cos\varphi}\dd\varphi
 =2\pi C\int_0^a rJ_0(qr)\dd r.
\]
这里用角积分定义$J_0(z)=(2\pi)^{-1}\int_0^{2\pi}\ee^{-iz\cos\varphi}\dd\varphi$。为求径向原函数，将指数展开；奇次余弦积分为零，偶次积分记为$C_{2j}=\int_0^{2\pi}\cos^{2j}\varphi\dd\varphi$。取$u=\cos^{2j-1}\varphi$、$\dd v=\cos\varphi\dd\varphi$，分部积分为
\begin{align*}
 C_{2j}&=[\cos^{2j-1}\varphi\sin\varphi]_0^{2\pi}
 +(2j-1)\int_0^{2\pi}\cos^{2j-2}\varphi\sin^2\varphi\dd\varphi\\
 &=(2j-1)(C_{2j-2}-C_{2j}),\qquad j\geq1.
\end{align*}
移项得$C_{2j}=(2j-1)C_{2j-2}/(2j)$，从$C_0=2\pi$递推，得$C_{2j}=2\pi(2j)!/[2^{2j}(j!)^2]$。于是
\[
 J_0(z)=\sum_{j=0}^{\infty}\frac{(-1)^j(z/2)^{2j}}{(j!)^2},
 \qquad
 J_1(z)=\sum_{j=0}^{\infty}
 \frac{(-1)^j(z/2)^{2j+1}}{j!(j+1)!}.
\]
第二式定义此处所需的一阶贝塞尔函数。逐项求导时，$zJ_1(z)$的第$j$项乘出$2j+2$，约去$(j+1)!$，便得到$\dd[zJ_1(z)]/\dd z=zJ_0(z)$。令$z=qr$、$r\dd r=z\dd z/q^2$，上下限由$0,a$变为$0,qa$，得到
\[
 U=\frac{2\pi C}{q^2}[zJ_1(z)]_0^{qa}
 =\frac{2\pi Ca}{q}J_1(qa).
\]
\Diagram{fig:o3circlea}{圆孔归一化远场强度；首零由贝塞尔函数确定}{
\draw[axis](0,0)--(5.1,0)node[right]{$qa$};
\draw[axis](0,0)--(0,2.9)node[above]{$I/I(0)$};
\draw[curve]plot coordinates{(0.000000,2.4500000) (0.015000,2.4494488) (0.030000,2.4477958) (0.045000,2.4450429) (0.060000,2.4411932) (0.075000,2.4362510) (0.090000,2.4302219) (0.105000,2.4231125) (0.120000,2.4149310) (0.135000,2.4056864) (0.150000,2.3953891) (0.165000,2.3840506) (0.180000,2.3716836) (0.195000,2.3583018) (0.210000,2.3439200) (0.225000,2.3285544) (0.240000,2.3122218) (0.255000,2.2949402) (0.270000,2.2767288) (0.285000,2.2576076) (0.300000,2.2375974) (0.315000,2.2167202) (0.330000,2.1949986) (0.345000,2.1724563) (0.360000,2.1491175) (0.375000,2.1250075) (0.390000,2.1001521) (0.405000,2.0745777) (0.420000,2.0483117) (0.435000,2.0213819) (0.450000,1.9938165) (0.465000,1.9656444) (0.480000,1.9368951) (0.495000,1.9075983) (0.510000,1.8777842) (0.525000,1.8474834) (0.540000,1.8167266) (0.555000,1.7855450) (0.570000,1.7539699) (0.585000,1.7220326) (0.600000,1.6897649) (0.615000,1.6571983) (0.630000,1.6243646) (0.645000,1.5912954) (0.660000,1.5580223) (0.675000,1.5245769) (0.690000,1.4909905) (0.705000,1.4572943) (0.720000,1.4235191) (0.735000,1.3896958) (0.750000,1.3558546) (0.765000,1.3220255) (0.780000,1.2882382) (0.795000,1.2545218) (0.810000,1.2209050) (0.825000,1.1874160) (0.840000,1.1540825) (0.855000,1.1209315) (0.870000,1.0879896) (0.885000,1.0552825) (0.900000,1.0228356) (0.915000,0.9906731) (0.930000,0.9588190) (0.945000,0.9272962) (0.960000,0.8961270) (0.975000,0.8653328) (0.990000,0.8349343) (1.005000,0.8049512) (1.020000,0.7754025) (1.035000,0.7463063) (1.050000,0.7176799) (1.065000,0.6895395) (1.080000,0.6619004) (1.095000,0.6347773) (1.110000,0.6081837) (1.125000,0.5821322) (1.140000,0.5566346) (1.155000,0.5317015) (1.170000,0.5073428) (1.185000,0.4835675) (1.200000,0.4603834) (1.215000,0.4377977) (1.230000,0.4158163) (1.245000,0.3944445) (1.260000,0.3736865) (1.275000,0.3535458) (1.290000,0.3340247) (1.305000,0.3151248) (1.320000,0.2968469) (1.335000,0.2791909) (1.350000,0.2621558) (1.365000,0.2457397) (1.380000,0.2299401) (1.395000,0.2147537) (1.410000,0.2001762) (1.425000,0.1862030) (1.440000,0.1728282) (1.455000,0.1600458) (1.470000,0.1478488) (1.485000,0.1362295) (1.500000,0.1251797) (1.515000,0.1146908) (1.530000,0.1047533) (1.545000,0.0953573) (1.560000,0.0864925) (1.575000,0.0781479) (1.590000,0.0703123) (1.605000,0.0629739) (1.620000,0.0561205) (1.635000,0.0497397) (1.650000,0.0438187) (1.665000,0.0383442) (1.680000,0.0333029) (1.695000,0.0286811) (1.710000,0.0244650) (1.725000,0.0206406) (1.740000,0.0171935) (1.755000,0.0141095) (1.770000,0.0113742) (1.785000,0.0089731) (1.800000,0.0068915) (1.815000,0.0051150) (1.830000,0.0036291) (1.845000,0.0024191) (1.860000,0.0014707) (1.875000,0.0007695) (1.890000,0.0003013) (1.905000,0.0000519) (1.920000,0.0000074) (1.935000,0.0001540) (1.950000,0.0004782) (1.965000,0.0009666) (1.980000,0.0016061) (1.995000,0.0023838) (2.010000,0.0032871) (2.025000,0.0043038) (2.040000,0.0054219) (2.055000,0.0066296) (2.070000,0.0079157) (2.085000,0.0092692) (2.100000,0.0106794) (2.115000,0.0121362) (2.130000,0.0136295) (2.145000,0.0151499) (2.160000,0.0166883) (2.175000,0.0182360) (2.190000,0.0197847) (2.205000,0.0213265) (2.220000,0.0228539) (2.235000,0.0243598) (2.250000,0.0258376) (2.265000,0.0272810) (2.280000,0.0286843) (2.295000,0.0300419) (2.310000,0.0313489) (2.325000,0.0326008) (2.340000,0.0337932) (2.355000,0.0349225) (2.370000,0.0359853) (2.385000,0.0369784) (2.400000,0.0378994) (2.415000,0.0387459) (2.430000,0.0395160) (2.445000,0.0402083) (2.460000,0.0408215) (2.475000,0.0413548) (2.490000,0.0418077) (2.505000,0.0421799) (2.520000,0.0424715) (2.535000,0.0426830) (2.550000,0.0428150) (2.565000,0.0428684) (2.580000,0.0428444) (2.595000,0.0427444) (2.610000,0.0425702) (2.625000,0.0423234) (2.640000,0.0420063) (2.655000,0.0416211) (2.670000,0.0411701) (2.685000,0.0406560) (2.700000,0.0400815) (2.715000,0.0394495) (2.730000,0.0387630) (2.745000,0.0380250) (2.760000,0.0372387) (2.775000,0.0364073) (2.790000,0.0355343) (2.805000,0.0346229) (2.820000,0.0336766) (2.835000,0.0326988) (2.850000,0.0316929) (2.865000,0.0306625) (2.880000,0.0296109) (2.895000,0.0285416) (2.910000,0.0274579) (2.925000,0.0263634) (2.940000,0.0252611) (2.955000,0.0241545) (2.970000,0.0230467) (2.985000,0.0219407) (3.000000,0.0208397) (3.015000,0.0197465) (3.030000,0.0186640) (3.045000,0.0175949) (3.060000,0.0165419) (3.075000,0.0155073) (3.090000,0.0144937) (3.105000,0.0135032) (3.120000,0.0125381) (3.135000,0.0116002) (3.150000,0.0106914) (3.165000,0.0098135) (3.180000,0.0089680) (3.195000,0.0081564) (3.210000,0.0073800) (3.225000,0.0066398) (3.240000,0.0059370) (3.255000,0.0052724) (3.270000,0.0046467) (3.285000,0.0040605) (3.300000,0.0035142) (3.315000,0.0030081) (3.330000,0.0025424) (3.345000,0.0021172) (3.360000,0.0017324) (3.375000,0.0013877) (3.390000,0.0010828) (3.405000,0.0008173) (3.420000,0.0005906) (3.435000,0.0004020) (3.450000,0.0002507) (3.465000,0.0001360) (3.480000,0.0000567) (3.495000,0.0000119) (3.510000,0.0000003) (3.525000,0.0000209) (3.540000,0.0000723) (3.555000,0.0001531) (3.570000,0.0002620) (3.585000,0.0003974) (3.600000,0.0005580) (3.615000,0.0007420) (3.630000,0.0009479) (3.645000,0.0011742) (3.660000,0.0014191) (3.675000,0.0016810) (3.690000,0.0019582) (3.705000,0.0022491) (3.720000,0.0025519) (3.735000,0.0028651) (3.750000,0.0031869) (3.765000,0.0035157) (3.780000,0.0038498) (3.795000,0.0041877) (3.810000,0.0045278) (3.825000,0.0048685) (3.840000,0.0052084) (3.855000,0.0055459) (3.870000,0.0058798) (3.885000,0.0062085) (3.900000,0.0065309) (3.915000,0.0068455) (3.930000,0.0071514) (3.945000,0.0074473) (3.960000,0.0077321) (3.975000,0.0080048) (3.990000,0.0082646) (4.005000,0.0085106) (4.020000,0.0087419) (4.035000,0.0089578) (4.050000,0.0091577) (4.065000,0.0093410) (4.080000,0.0095071) (4.095000,0.0096558) (4.110000,0.0097865) (4.125000,0.0098991) (4.140000,0.0099933) (4.155000,0.0100689) (4.170000,0.0101259) (4.185000,0.0101642) (4.200000,0.0101841) (4.215000,0.0101854) (4.230000,0.0101686) (4.245000,0.0101337) (4.260000,0.0100812) (4.275000,0.0100114) (4.290000,0.0099246) (4.305000,0.0098215) (4.320000,0.0097025) (4.335000,0.0095682) (4.350000,0.0094191) (4.365000,0.0092561) (4.380000,0.0090796) (4.395000,0.0088906) (4.410000,0.0086897) (4.425000,0.0084777) (4.440000,0.0082555) (4.455000,0.0080238) (4.470000,0.0077836) (4.485000,0.0075356) (4.500000,0.0072808)};
\node[left]at(0,2.45){$1$};
\draw[guide](1.915853,0)--(1.915853,.55);
\node[below]at(1.915853,0){$3.8317$};
\node[right]at(2.15,.45){首暗环};
}
轴上直接积分为$U(0)=C\pi a^2$，因此图\ref{fig:o3circlea}的归一化场和强度为
\[
 \frac{U}{U(0)}=\frac{2J_1(qa)}{qa},\qquad
 \frac I{I(0)}=\left[\frac{2J_1(qa)}{qa}\right]^2.
\]
$J_1(z)=z/2+O(z^3)$使轴上极限为1，不存在奇点。首个非零根$qa=3.8317$给
\[
 \sin\theta_0=\frac{3.8317}{2\pi}\frac\lambda a
 =1.2197\frac\lambda D,\quad D=2a.
\]
在小角度下这也是瑞利角间隔。两个互不相干点源的像应把两套强度相加；若相干，交叉项依相位而变，不能直接引用同一个双点强度判据。系数1.22由均匀圆孔的贝塞尔首零决定，矩形孔和高斯孔径需重新计算。
\end{solution}

\begin{example}\label{ex:o3gaussian}
孔面场振幅为$U(\xi)=U_0\ee^{-\xi^2/w^2}$，孔径足够大不截断高斯尾部。求远场强度及其$1/\ee^2$角半宽，与均匀硬边孔径比较。
\Diagram{fig:o3gaussq}{高斯振幅没有硬边界}{\draw[axis](-2.5,0)--(2.8,0)node[right]{$\xi/w$};\draw[axis](0,0)--(0,2.7)node[above]{$U/U_0$};\draw[curve,domain=-2.3:2.3,samples=80]plot(\x,{2.2*exp(-\x*\x)});\node[left]at(0,2.2){$1$};\draw(1,0)--(1,-.08)node[below]{$1$};}
\end{example}
\begin{solution}
图\ref{fig:o3gaussq}的远场积分为$F(q)=U_0\int_{-\infty}^{\infty}\ee^{-\xi^2/w^2-iq\xi}\dd\xi$。它正是第一讲已用分部积分及初值方程证明的式\eqref{eq:o1gaussintegral}。代入该式的$a=w^{-2}>0$、$b=-iq$，有$\sqrt{\pi/a}=w\sqrt\pi$、$b^2/(4a)=-q^2w^2/4$，因此
\[
 F(q)=U_0w\sqrt\pi\ee^{-q^2w^2/4}.
\]
轴上$q=0$时恢复孔面场的积分$U_0w\sqrt\pi$；方向相关的衰减完全由指数中的$q^2w^2$决定。
\Diagram{fig:o3gaussa}{高斯远场无旁瓣，角宽与孔面场宽成反比}{\draw[axis](-2.6,0)--(2.8,0)node[right]{$\theta/\theta_d$};\draw[axis](0,0)--(0,2.7)node[above]{$I/I(0)$};\draw[curve,domain=-2.4:2.4,samples=80]plot(\x,{2.3*exp(-2*\x*\x)});\draw[guide](-2.4,.3113)--(2.4,.3113);\node[left]at(0,.3113){$\ee^{-2}$};\node[left]at(0,2.3){$1$};\draw[guide](1,0)--(1,.3113);\draw(1,0)--(1,-.08)node[below]{$1$};}
图\ref{fig:o3gaussa}中的强度比为$\exp(-q^2w^2/2)$。取$I/I(0)=\ee^{-2}$，则$|q|=2/w$；近轴$q\simeq k\theta$，所以$1/\ee^2$角半宽为
\[
 \theta_d=\frac2{kw}=\frac\lambda{\pi w}.
\]
题中的$w$是场下降至$1/\ee$、强度下降至$1/\ee^2$的孔面半宽。若给的是强度高斯式$\exp(-\xi^2/w_I^2)$，须先取场的平方根，不能沿用相同的$w$。平滑高斯孔面对应无有限零点、无旁瓣的高斯远场，宽度须按约定强度比定义。$1/\ee^2$半宽与硬孔首零半宽采用不同的宽度标准。
\end{solution}

\begin{example}\label{ex:o3twoDgrating}
二维矩形周期阵列沿$x,y$周期为$b_x,b_y$，分别有$N_x,N_y$个相同小孔，正常入射。求主极大的方向条件及两方向首零宽度。
\Diagram{fig:o3twoDq}{二维有限孔阵列，图只显示部分周期}{\foreach\x in{0,1,2,3}{\foreach\y in{0,1,2}{\draw[curve](\x,\y)circle(.13);}}\draw[axis](-.5,-.5)--(4,-.5)node[right]{$x$};\draw[axis](-.5,-.5)--(-.5,3)node[above]{$y$};}
\end{example}
\begin{solution}
令$k_0=2\pi/\lambda$。图\ref{fig:o3twoDq}的平移相位可分离为
\[
 \sum_{j=0}^{N_x-1}\ee^{-ik_0b_xl_xj}
 \sum_{\ell=0}^{N_y-1}\ee^{-ik_0b_yl_y\ell}.
\]
两方向同时相长要求$b_xl_x=m\lambda$、$b_yl_y=n\lambda$。
\Diagram{fig:o3twoDa}{允许主峰位于方向余弦平面内的二维格点}{\draw[axis](-2,0)--(2.4,0)node[right]{$l_x$};\draw[axis](0,-1.8)--(0,2.2)node[above]{$l_y$};\draw[guide](0,0)circle(1.65);\foreach\x in{-1,0,1}{\foreach\y in{-1,0,1}{\fill(\x,\y)circle(1.5pt);}}}
还须$l_x^2+l_y^2\leq1$，图\ref{fig:o3twoDa}圆外格点不对应传播波。固定另一方向，首零宽度为$\Delta l_x=\lambda/(N_xb_x)$、$\Delta l_y=\lambda/(N_yb_y)$。单孔包络再乘到阵列因子上，某格点可能因孔形而缺级。二维条件不能只用总孔数$N_xN_y$替代一个方向的有效周期数。
\end{solution}

\begin{example}\label{ex:o3unequalresolution}
两条互不相干谱线的单线分布为偶函数$F$，强度比为$1:r$，峰中心为$\pm a$。解释为何等强时中点曲率判据不能直接用于$r\ne1$，写出真正峰合并边界的方程。
\Diagram{fig:o3unequalq}{弱谱线可能被强谱线的旁瓣或翼部覆盖}{\draw[axis](-3,0)--(3.3,0)node[right]{$x$};\draw[curve,domain=-2.8:2.8,samples=90]plot(\x,{2.3*exp(-2*(\x+1)^2)});\draw[guide,domain=-2.8:2.8,samples=90]plot(\x,{.7*exp(-2*(\x-1)^2)});}
\end{example}
\begin{solution}
图\ref{fig:o3unequalq}总强度$G(x)=F(x+a)+rF(x-a)$。当$r\ne1$时$G$不再关于0对称，$G'(0)=(1-r)F'(a)$一般不为零。因此即使$G''(0)=0$，中点也不是极值，不能用它判定两峰是否分开。
\Diagram{fig:o3unequala}{不等强高斯线形在临界间距处的合强度，$r=0.3$}{\draw[axis](-3.6,0)--(3.7,0)node[right]{$x$};\draw[axis](0,0)--(0,2.9)node[above]{$G$};\draw[curve]plot coordinates {(-3.40000,0.008830) (-3.35976,0.010673) (-3.31953,0.012859) (-3.27929,0.015443) (-3.23905,0.018487) (-3.19882,0.022059) (-3.15858,0.026235) (-3.11834,0.031102) (-3.07811,0.036752) (-3.03787,0.043289) (-2.99763,0.050823) (-2.95740,0.059476) (-2.91716,0.069376) (-2.87692,0.080663) (-2.83669,0.093484) (-2.79645,0.107991) (-2.75621,0.124347) (-2.71598,0.142717) (-2.67574,0.163271) (-2.63550,0.186182) (-2.59527,0.211622) (-2.55503,0.239760) (-2.51479,0.270761) (-2.47456,0.304783) (-2.43432,0.341970) (-2.39408,0.382454) (-2.35385,0.426349) (-2.31361,0.473745) (-2.27337,0.524709) (-2.23314,0.579277) (-2.19290,0.637453) (-2.15266,0.699204) (-2.11243,0.764459) (-2.07219,0.833103) (-2.03195,0.904978) (-1.99172,0.979877) (-1.95148,1.057549) (-1.91124,1.137690) (-1.87101,1.219952) (-1.83077,1.303937) (-1.79053,1.389206) (-1.75030,1.475273) (-1.71006,1.561619) (-1.66982,1.647686) (-1.62959,1.732891) (-1.58935,1.816630) (-1.54911,1.898282) (-1.50888,1.977220) (-1.46864,2.052818) (-1.42840,2.124458) (-1.38817,2.191542) (-1.34793,2.253499) (-1.30769,2.309791) (-1.26746,2.359927) (-1.22722,2.403466) (-1.18698,2.440025) (-1.14675,2.469287) (-1.10651,2.491006) (-1.06627,2.505012) (-1.02604,2.511211) (-0.98580,2.509591) (-0.94556,2.500220) (-0.90533,2.483249) (-0.86509,2.458905) (-0.82485,2.427492) (-0.78462,2.389386) (-0.74438,2.345025) (-0.70414,2.294912) (-0.66391,2.239598) (-0.62367,2.179680) (-0.58343,2.115789) (-0.54320,2.048583) (-0.50296,1.978736) (-0.46272,1.906932) (-0.42249,1.833851) (-0.38225,1.760163) (-0.34201,1.686518) (-0.30178,1.613541) (-0.26154,1.541821) (-0.22130,1.471903) (-0.18107,1.404289) (-0.14083,1.339426) (-0.10059,1.277707) (-0.06036,1.219464) (-0.02012,1.164969) (0.02012,1.114433) (0.06036,1.068004) (0.10059,1.025767) (0.14083,0.987750) (0.18107,0.953924) (0.22130,0.924203) (0.26154,0.898455) (0.30178,0.876497) (0.34201,0.858111) (0.38225,0.843037) (0.42249,0.830988) (0.46272,0.821651) (0.50296,0.814695) (0.54320,0.809771) (0.58343,0.806528) (0.62367,0.804608) (0.66391,0.803657) (0.70414,0.803329) (0.74438,0.803290) (0.78462,0.803223) (0.82485,0.802831) (0.86509,0.801841) (0.90533,0.800006) (0.94556,0.797108) (0.98580,0.792959) (1.02604,0.787405) (1.06627,0.780321) (1.10651,0.771619) (1.14675,0.761239) (1.18698,0.749154) (1.22722,0.735369) (1.26746,0.719914) (1.30769,0.702847) (1.34793,0.684251) (1.38817,0.664228) (1.42840,0.642900) (1.46864,0.620404) (1.50888,0.596890) (1.54911,0.572517) (1.58935,0.547451) (1.62959,0.521859) (1.66982,0.495912) (1.71006,0.469776) (1.75030,0.443616) (1.79053,0.417588) (1.83077,0.391839) (1.87101,0.366507) (1.91124,0.341720) (1.95148,0.317590) (1.99172,0.294219) (2.03195,0.271694) (2.07219,0.250087) (2.11243,0.229459) (2.15266,0.209856) (2.19290,0.191309) (2.23314,0.173839) (2.27337,0.157456) (2.31361,0.142157) (2.35385,0.127930) (2.39408,0.114756) (2.43432,0.102606) (2.47456,0.091446) (2.51479,0.081237) (2.55503,0.071934) (2.59527,0.063491) (2.63550,0.055858) (2.67574,0.048984) (2.71598,0.042817) (2.75621,0.037306) (2.79645,0.032398) (2.83669,0.028046) (2.87692,0.024200) (2.91716,0.020813) (2.95740,0.017843) (2.99763,0.015247) (3.03787,0.012987) (3.07811,0.011026) (3.11834,0.009331) (3.15858,0.007871) (3.19882,0.006618) (3.23905,0.005546) (3.27929,0.004633) (3.31953,0.003858) (3.35976,0.003202) (3.40000,0.002649)};\draw[guide](0.14049,0.803290)--(1.54049,0.803290);\fill(0.74049,0.803290)circle(1.5pt)node[above right]{$x_*$};}
应联立$F'(x_*+a)+rF'(x_*-a)=0$和$F''(x_*+a)+rF''(x_*-a)=0$，解未知$x_*,a$。例如$F(x)=\ee^{-x^2}$，两次求导为$F'=-2x\ee^{-x^2}$、$F''=(-2+4x^2)\ee^{-x^2}$。记$u=x_*+a,v=x_*-a$，第一式给$r\ee^{-v^2}=-(u/v)\ee^{-u^2}$；代入第二式并乘$v$，得到
\[
 v(4u^2-2)-u(4v^2-2)=2(u-v)(2uv+1)=0.
\]
非零分离$a>0$对应$uv=-1/2$，即$a^2-x_*^2=1/2$。再代回第一式，
\[
 r=\frac{a+x_*}{a-x_*}\ee^{-4ax_*}.
\]
对$r=0.3$取弱峰一侧$x_*>0$，数值求这两个标量方程，得$a=1.023881$、$x_*=0.740495$。图\ref{fig:o3unequala}按该解画出峰、谷合并处的水平切线。等强$r=1$时对称性给$x_*=0$，才退回中点曲率条件。
\end{solution}

\section{孔径自相关、成像传递与边缘加权}
\subsection{从点像积分得到强度传递函数}
在一个方向上讨论无像差、空间平移不变的近轴成像。设像面点响应的场为
\[
 h(x)=C\int P(\xi)\exp\!\left(-\frac{2\pi\mathrm i\xi x}{\lambda f}\right)\dd\xi,
\]
其中$P$为有效光瞳函数，$f$为此式所用的有效焦距尺度。对互不相干物点，归一化强度点响应是$g(x)=|h(x)|^2/\int|h|^2\dd x$。物面强度先换算到几何像坐标，再与$g$作卷积：
\[
 I(x)=\int I_{\rm geom}(x')g(x-x')\dd x'.
\]
这一步使用了不同物点间的交叉项平均为零，不能用来处理激光照明下的任意相干物体。

记未归一化强度的傅里叶变换为$N(\nu)$。将$|h|^2$展开，并对$x$积分，
\begin{align*}
 N(\nu)&=\int |h(x)|^2\ee^{-2\pi i\nu x}\dd x\\
 &=|C|^2\iint P(\xi)P^*(\eta)
 \left\{\int\ee^{-2\pi ix[(\xi-\eta)/(\lambda f)+\nu]}\dd x\right\}
 \dd\xi\dd\eta\\
 &=|C|^2\lambda f\int
 P(\xi)P^*(\xi+\lambda f\nu)\dd\xi.
\end{align*}
最后一步在傅里叶分布意义下使用$\int\ee^{-2\pi ixt}\dd x=\delta(t)$；括号内的狄拉克函数为$\lambda f\,\delta(\eta-\xi-\lambda f\nu)$，其尺度因子不能漏去。以$N(0)=|C|^2\lambda f\int|P|^2\dd\xi$归一化，得到
\[
 G(\nu)=\int g(x)\ee^{-2\pi\mathrm i\nu x}\dd x
 =\frac{\int P(\xi)P^*(\xi+\lambda f\nu)\dd\xi}
 {\int |P(\xi)|^2\dd\xi}.
\]
分母使$G(0)=1$。对于实对称光瞳，$G$为实偶函数；带像差的复光瞳还可能产生传递相位，不能只保留模。

宽$a$的均匀一维光瞳取$P=1$于$|\xi|<a/2$。分子是两段错开$\lambda f\nu$的区间重叠长度，故
\[
 G(\nu)=\begin{cases}1-|\nu|/\nu_c,&|\nu|\leq\nu_c,\\0,&|\nu|>\nu_c,
 \end{cases}\qquad \nu_c=\frac a{\lambda f}.
\]
上式给出不相干强度成像的截止频率。相干成像传递的是复场，传递函数与光瞳本身对应，截止为$a/(2\lambda f)$。比较这两个数值时须保留各自的照明条件和被传递的量；仅凭截止频率的比值，还不足以比较实际物体的可分辨程度。
\ExpansionDiagram{fig:dep-o3-overlap}{错位光瞳的重叠长度决定一维不相干传递函数}{
\fill[gray!20](2.5,.1)rectangle(5,1.5);\draw[very thick](1,.5)--(5,.5);\draw[very thick](2.5,1.2)--(6.5,1.2);\node[below]at(3,.1){重叠$a-|\lambda f\nu|$};\draw[<->](1,1.8)--(2.5,1.8)node[midway,above]{$\lambda f\nu$};}
图\ref{fig:dep-o3-overlap}同时说明了截止的来由：光瞳错位超过整个宽度后，重叠积分为零，物体相应的强度频率分量不再出现在像中。

\begin{example}条纹能够成像但对比度已经降低
某一维不相干成像系统的均匀有效光瞳宽$a=2.0\,\mathrm{mm}$，有效焦距$f=100\,\mathrm{mm}$，波长$\lambda=500\,\mathrm{nm}$。几何像的强度为$I_0[1+0.8\cos(2\pi\nu x)]$，其中$\nu=20\,\mathrm{mm}^{-1}$。求实际像的调制度，并求保持调制度不低于0.4所允许的最高频率。
\ExpansionDiagram{fig:dep-o3-modulation-q}{平移不变成像使各正弦强度分量分别乘传递系数}{
\draw[->](0,0)--(3.2,0)node[right]{$x$};\draw[->](0,0)--(0,2.1)node[above]{$I/I_0$};\draw[dotted](0,1)--(3,1);\draw[domain=0:3,samples=90]plot(\x,{1+.8*cos(360*\x)});\node[above]at(1.5,2.2){几何像};\draw[->](3.8,1)--(4.7,1);\draw[->](5,0)--(8.2,0)node[right]{$x$};\draw[->](5,0)--(5,2.1)node[above]{$I/I_0$};\draw[dotted](5,1)--(8,1);\draw[domain=5:8,samples=90]plot(\x,{1+.4*cos(360*\x)});\node[above]at(6.5,2.2){实际像};}
\end{example}
\begin{solution}
统一长度单位后，$\lambda=0.0005\,\mathrm{mm}$，所以$\nu_c=2/(0.0005\times100)=40\,\mathrm{mm}^{-1}$。常数分量乘$G(0)=1$，正弦分量乘$G(20)=1-20/40=1/2$，因此
\[
 I(x)=I_0[1+0.4\cos(2\pi\nu x)].
\]
图\ref{fig:dep-o3-modulation-q}中的周期不变，变化的是振幅相对平均强度的比例。由$I_{\max}=1.4I_0$、$I_{\min}=0.6I_0$，可见度也是0.4。
\ExpansionDiagram{fig:dep-o3-modulation-a}{调制度要求选出比截止频率更低的可用频段}{
\draw[->](0,0)--(5,0)node[right]{$\nu/\nu_c$};\draw[->](0,0)--(0,2.4)node[above]{$G$};\draw[thick](0,2)--(4,0);\draw[dashed](0,1)--(2,1)--(2,0);\node[left]at(0,2){1};\node[left]at(0,1){$1/2$};\node[below]at(2,0){$1/2$};\node[below]at(4,0){1};}
一般输出调制度为$0.8(1-\nu/\nu_c)$，要求不低于0.4给$\nu\leq\nu_c/2=20\,\mathrm{mm}^{-1}$。图\ref{fig:dep-o3-modulation-a}中截止前仍有一段低对比度区。因此“频率未超过截止”与“达到给定对比度”是两个不同条件。这里假设探测器能连续记录；有限像素还会另加像素孔径与采样限制。
\end{solution}

\subsection{连续加权为何能压低旁瓣}
硬边孔径在边缘突然截止，其傅里叶场具有较慢的尾部衰减。让透过振幅在边缘连续降至零，可以改变尾部阶数。判断代价时须固定照明条件，不能把峰值归一化后的两条曲线直接当作相同通光量。

\begin{example}三角振幅孔径与均匀孔径
一维孔宽均为$a$、入射场振幅同为$U_0$。第一孔的振幅透过率为1，第二孔为$t(\xi)=1-2|\xi|/a$于$|\xi|\leq a/2$，外侧为零。求第二孔的远场场、第一零点及总透过功率，与第一孔比较。
\ExpansionDiagram{fig:dep-o3-triangle-q}{相同外宽的硬边光瞳与三角振幅光瞳}{
\draw[->](-2,0)--(2,0)node[right]{$\xi$};\draw[dashed](-1.5,0)--(-1.5,1.5)--(1.5,1.5)--(1.5,0);\draw[thick](-1.5,0)--(0,1.5)--(1.5,0);\node[above]at(0,1.6){$t$};\node[below]at(-1.5,0){$-a/2$};\node[below]at(1.5,0){$a/2$};}
\end{example}
\begin{solution}
图\ref{fig:dep-o3-triangle-q}中三角透过率是两个宽$a/2$的单位矩形函数卷积后乘$2/a$。每个矩形的傅里叶变换为$(a/2)\operatorname{sinc}(qa/4)$，故卷积定理给
\[
 F_\triangle(q)=\frac{aU_0}{2}\operatorname{sinc}^2(qa/4),
 \quad \frac{I_\triangle(q)}{I_\triangle(0)}
 =\operatorname{sinc}^4(qa/4).
\]
也可直接积分核对。对$q\ne0$，分部积分取$u=1-2\xi/a$、$\dd v=\cos(q\xi)\dd\xi$，得
\begin{align*}
 F_\triangle(q)&=2U_0\int_0^{a/2}(1-2\xi/a)\cos(q\xi)\dd\xi\\
 &=2U_0\left[(1-2\xi/a)\frac{\sin(q\xi)}q\right]_0^{a/2}
 +\frac{4U_0}{aq}\int_0^{a/2}\sin(q\xi)\dd\xi\\
 &=\frac{4U_0}{aq}\left[-\frac{\cos(q\xi)}q\right]_0^{a/2}
 =\frac{4U_0}{aq^2}[1-\cos(qa/2)].
\end{align*}
边界项为零，最后用$1-\cos(qa/2)=2\sin^2(qa/4)$，重现卷积法结果。首零点为$|q|=4\pi/a$，均匀孔首零为$2\pi/a$，所以三角加权的中央瓣宽了一倍。
\ExpansionDiagram{fig:dep-o3-triangle-a}{归一化场的零点移动；实线为三角加权，虚线为均匀孔}{
\draw[->](-.2,0)--(7.2,0)node[right]{$qa/(2\pi)$};\draw[->](0,0)--(0,2.3)node[above]{归一化场};
\draw[thick,domain=.01:6.6,samples=180]plot(\x,{2*(sin(90*\x)/(pi*\x/2))^2});\draw[dashed,domain=.01:6.6,samples=180]plot(\x,{2*sin(180*\x)/(pi*\x)});\node[below=12pt]at(1,0){1};\node[below=12pt]at(2,0){2};}
图\ref{fig:dep-o3-triangle-a}画场而非强度，均匀孔场会反号，三角孔场为非负的平方。三角孔强度尾部按$q^{-4}$包络衰减，均匀孔按$q^{-2}$衰减，但主瓣展宽且透过功率降低。

功率需在孔面算$|t|^2$：
\[
 P_\triangle\propto2|U_0|^2\int_0^{a/2}(1-2\xi/a)^2\dd\xi
 =2|U_0|^2\left[\xi-\frac{2\xi^2}{a}+\frac{4\xi^3}{3a^2}\right]_0^{a/2}
 =\frac{a|U_0|^2}{3}.
\]
均匀孔为$a|U_0|^2$，故总功率比为$1/3$，轴上强度比则由场比$1/2$得$1/4$。这两个比值不同，原因是整个角分布的宽度和旁瓣都发生了改变。
\end{solution}

\section{孔面相位误差与相干峰的损失}
\subsection{误差的相关性先于误差的大小}
有$N$个等振幅、等间距单元时，远场场为$A(\delta)=a\sum_{j=0}^{N-1}\ee^{\mathrm i(j\delta+\epsilon_j)}$。相位误差全相同时，只有一个共同相位因子，强度完全不变；误差为$\epsilon_j=j\epsilon$时，分布整体移到$\delta=-\epsilon$；独立随机误差才会把部分窄峰能量转成背景。只报“均方相位误差相同”，不足以确定实际图样。

设各$\epsilon_j$均值为零，服从联合高斯分布，协方差为$C_{jk}=\langle\epsilon_j\epsilon_k\rangle$。场强平方含成对项，使用高斯特征函数有
\[
 \left\langle\ee^{\mathrm i(\epsilon_j-\epsilon_k)}\right\rangle
 =\exp[-\tfrac12(C_{jj}+C_{kk}-2C_{jk})].
\]
因而平均强度直接由每一对单元的相位差方差决定。共同误差使这个方差为零；远距离误差不相关时，其交叉项受到抑制。此式把完全相干阵列与部分相干孔径联系起来，也说明局部加工误差、整体倾斜和共同振动的影响不同。

若误差独立、方差同为$\sigma_\phi^2$，对角项始终为1，非对角项均为$\ee^{-\sigma_\phi^2}$，于是
\[
 \langle I(\delta)\rangle=|a|^2\left[
 \ee^{-\sigma_\phi^2}\left|\sum_j\ee^{\mathrm ij\delta}\right|^2
 +N(1-\ee^{-\sigma_\phi^2})\right].
\]
平均强度由等间距阵列的窄峰和角度均匀的背景组成，实际图样还要乘单元包络。在一个相位周期内积分，非对角指数项均为零，总积分仍是$2\pi N|a|^2$。这个纯相位模型没有吸收能量，峰值下降的同时，其他方向的强度增加。

\begin{example}相位误差对峰高和暗点的不同影响
取$N=100$，各单元独立高斯相位误差标准差为$\sigma_\phi=0.5$弧度。求平均主峰强度相对理想峰的比例、理想暗点处的平均背景，并求保持主峰至少为理想峰90\%所允许的最大$\sigma_\phi$。若误差来自正常入射反射面的独立高度误差，波长为$500\,\mathrm{nm}$，换算高度标准差。
\ExpansionDiagram{fig:dep-o3-phase-q}{单元位置不变，但各单元出射场的相位独立偏离}{
\draw[->](0,0)--(7,0)node[right]{单元};\foreach\x/\a in{.5/80,1.5/105,2.5/60,3.5/95,4.5/115,5.5/70}{\fill(\x,0)circle(1.2pt);\draw[->](\x,.3)--++(\a:.8);}\node[above]at(3,1.7){$\epsilon_j$改变相量方向};}
\end{example}
\begin{solution}
图\ref{fig:dep-o3-phase-q}中主峰的理想相位差为零，故
\[
 \frac{\langle I_{\max}\rangle}{N^2|a|^2}
 =\ee^{-\sigma_\phi^2}+
 \frac{1-\ee^{-\sigma_\phi^2}}N.
\]
代入$\ee^{-0.25}\simeq0.778801$，得到主峰比例0.781013。理想零点处的相干和为零，但背景为$N(1-\ee^{-0.25})|a|^2\simeq22.1199|a|^2$，不再严格消光。
\ExpansionDiagram{fig:dep-o3-phase-a}{独立相位误差降低平均主峰并填高理想零点；示意取10单元}{
\draw[->](0,0)--(6.8,0)node[right]{$\delta$};\draw[->](0,0)--(0,2.6)node[above]{强度};\draw[dashed,domain=.003:6.28,samples=400]plot(\x,{.022*(sin(5*\x r)/sin(.5*\x r))^2});\draw[thick,domain=.003:6.28,samples=400]plot(\x,{.022*(exp(-.25)*(sin(5*\x r)/sin(.5*\x r))^2+10*(1-exp(-.25)))});}
图\ref{fig:dep-o3-phase-a}用10单元显示曲线形状，本题数值计算使用100单元。要求主峰比例不低于0.9，等价于
\[
 \ee^{-\sigma_\phi^2}\geq\frac{0.9N-1}{N-1},\qquad
 \sigma_\phi\leq\sqrt{-\ln(89/99)}\simeq0.3263.
\]
正常入射反射面高度差造成往返程差$2\delta h$，故$\sigma_\phi=4\pi\sigma_h/\lambda$，允许$\sigma_h\lesssim13.0\,\mathrm{nm}$。若表面误差在各单元间强相关，不能继续用独立误差的背景公式；例如整面同步平移只改变共同相位，相干峰不因此降低。
\end{solution}

\subsection{连续孔径的平均相位与相位方差}
对均匀孔径，目标轴上场相对理想场为$\langle\ee^{\mathrm i\phi(\xi)}\rangle_P$，其中平均按孔径面积归一化。小相位误差展开至二阶：
\[
 \langle\ee^{\mathrm i\phi}\rangle_P
 \simeq1+\mathrm i\langle\phi\rangle_P-\frac12\langle\phi^2\rangle_P,
\]
所以轴上强度比
\[
 \mathcal S\simeq1-[\langle\phi^2\rangle_P-\langle\phi\rangle_P^2].
\]
共同相位从方差中抵消。这是围绕指定理想像点的二阶近似，须让去掉平均值后的相位足够小。若波前含线性倾斜，固定轴上点的强度会降低，但把观察点移到新峰中心后即可补偿；若含二次曲率，则需移动像面或改变聚焦条件。

\begin{example}一维光瞳的小离焦容许量
宽$a$的均匀一维孔径经焦距$f$的理想透镜聚焦。观察平面偏离后焦面$\delta$，$|\delta|\ll f$，忽略振幅传播因子的相对小变化。求轴上峰强相对该面理想相位补偿后峰强的二阶损失，并给出损失不超过5\%的近似离焦范围。
\ExpansionDiagram{fig:dep-o3-defocus-q}{同一孔径在偏离焦面的屏上留下二次相位}{
\draw[->](0,0)--(6,0)node[right]{$z$};\draw[<->,thick](.5,-1.4)--(.5,1.4);\draw(.5,1.2)--(4,0)--(5,-.34);\draw(.5,-1.2)--(4,0)--(5,.34);\draw[dashed](4,-1)--(4,1)node[above]{$f$};\draw[thick](5,-1)--(5,1)node[above]{$f+\delta$};}
\end{example}
\begin{solution}
图\ref{fig:dep-o3-defocus-q}中孔面到屏的近轴传播保留$kx^2/[2(f+\delta)]$，透镜原有相位为$-kx^2/(2f)$，故残余$\phi(x)=bx^2$，其中$b\simeq-k\delta/(2f^2)$。均匀区间$[-a/2,a/2]$的矩为
\[
 \langle x^2\rangle=\frac1a\int_{-a/2}^{a/2}x^2\dd x
 =\frac1a\left[\frac{x^3}{3}\right]_{-a/2}^{a/2}=\frac{a^2}{12},
\]
\[
 \langle x^4\rangle=\frac1a\int_{-a/2}^{a/2}x^4\dd x
 =\frac1a\left[\frac{x^5}{5}\right]_{-a/2}^{a/2}=\frac{a^4}{80},
 \quad \operatorname{Var}(x^2)=a^4\left(\frac1{80}-\frac1{144}\right)=\frac{a^4}{180}.
\]
于是
\[
 1-\mathcal S\simeq\frac{b^2a^4}{180}
 =\frac{k^2\delta^2a^4}{720f^4}.
\]
\ExpansionDiagram{fig:dep-o3-defocus-a}{去掉平均相位后的孔面相位仍呈弯曲，不能靠共同相移消除}{
\draw[->](-2.5,0)--(2.5,0)node[right]{$x$};\draw[->](0,-1.8)--(0,1.8)node[above]{$\phi-\langle\phi\rangle$};\draw[thick,domain=-2:2,samples=70]plot(\x,{-.5*((\x)^2-4/3)});\node[above]at(-2,0){$-a/2$};\node[above]at(2,0){$a/2$};}
图\ref{fig:dep-o3-defocus-a}表明影响峰强的是相位离散程度。要求上式不超过0.05，给
\[
 |\delta|\lesssim\frac{6f^2}{ka^2}=
 \frac{3\lambda f^2}{\pi a^2}
\]
（空气中$k=2\pi/\lambda$）。边界对应的均方根相位约$\sqrt{0.05}$弧度，仍属小误差估计；若要求更精确的阈值，应对完整菲涅耳积分求强度。几何弥散斑和轴上波动峰强是不同描述，不能仅把几何斑半径为零作为所有波前误差消失的条件。
\end{solution}

\section{习题与解答}
\begin{exercise}\label{ex:o3fresnelcircle}
均匀平面波正照半径$a$的圆孔，求距孔$L$的轴上菲涅耳场和强度。说明轴上为何可随孔径大小交替增强、减弱，不能一直用远场轴上场正比面积解释。
\Diagram{fig:o3fresnelcircleq}{孔内同心环带到轴上点的二次相位不同}{\draw[curve](0,0)ellipse(.5 and 1.6);\draw[guide](0,0)ellipse(.25 and .8);\fill(5,0)circle(1.5pt)node[right]{$P$};\draw[curve](0,1.6)--(5,0);\draw[curve](0,-1.6)--(5,0);\draw[guide](0,0)--(5,0);\node[above]at(2.5,.2){$L$};}
\end{exercise}
\begin{solution}
图\ref{fig:o3fresnelcircleq}轴上近轴传播积分为
\[
 U(P)=\frac{U_0\ee^{ikL}}{i\lambda L}
 2\pi\int_0^a r\ee^{ikr^2/(2L)}\dd r.
\]
令$q=kr^2/(2L)$，$r\dd r=(L/k)\dd q$，上下限变为$0,\phi_a$，其中$\phi_a=\pi a^2/(\lambda L)$。因此
\[
 \int_0^a r\ee^{ikr^2/(2L)}\dd r
 =\frac Lk\int_0^{\phi_a}\ee^{iq}\dd q
 =\frac Lk\left[\frac{\ee^{iq}}i\right]_0^{\phi_a}
 =\frac{L}{ik}(\ee^{i\phi_a}-1).
\]
代回传播前因子，并用$k\lambda=2\pi$，得$U(P)=U_0\ee^{ikL}[1-\ee^{i\phi_a}]$。所以相对无遮挡平面波强度
\[
 I(P)/I_0=4\sin^2(\phi_a/2).
\]
\Diagram{fig:o3fresnelcirclea}{轴上强度按边缘二次相位振荡}{\draw[axis](0,0)--(6.8,0)node[right]{$\phi_a$};\draw[axis](0,0)--(0,2.8)node[above]{$I/I_0$};\draw[curve,domain=0:6.283,samples=110]plot(\x,{2.3*sin((\x/2) r)^2});\node[left]at(0,2.3){$4$};\node[below]at(3.1416,0){$\pi$};\node[below]at(6.283,0){$2\pi$};}
图\ref{fig:o3fresnelcirclea}在$\phi_a\ll1$时$I/I_0\simeq\phi_a^2\propto a^4$，恢复远场轴上强度正比面积平方。较大孔径加入的环带可能与原场反相，使轴上减弱，能量分到轴外。积分仍要求近轴、孔径远大于波长；不能把$a\to\infty$的振荡直接当无限无遮挡平面波的极限，因为该极限不再一致满足所用截断条件。
\end{solution}

\begin{exercise}\label{ex:o3poisson}
把上一题圆孔换为同半径不透明圆盘，仍以平面波正照。用互补原理求圆盘后轴上菲涅耳场，解释几何阴影中心的亮点及其模型条件。
\Diagram{fig:o3poissonq}{圆盘遮挡几何直达光，边缘绕射仍到达轴上}{\draw[curve](0,0)ellipse(.45 and 1.5);\draw[very thick](0,-1.5)--(0,1.5);\draw[guide](0,1.5)--(5,0);\draw[guide](0,-1.5)--(5,0);\fill(5,0)circle(1.5pt)node[right]{$P$};\node[above]at(2.5,1.2){几何阴影内};}
\end{exercise}
\begin{solution}
图\ref{fig:o3poissonq}圆盘与圆孔互补，标量场线性给$U_{\rm disk}+U_{\rm hole}=U_{\rm free}$。上一题$U_{\rm hole}=U_{\rm free}(1-\ee^{i\phi_a})$，所以$U_{\rm disk}=U_{\rm free}\ee^{i\phi_a}$，轴上强度$I_{\rm disk}=I_0$。
\Diagram{fig:o3poissona}{轴上边缘环带的相位相同，形成阴影中央亮点}{\draw[curve](0,0)circle(1.5);\fill(0,0)circle(3pt);\foreach\a in{0,60,120,180,240,300}{\draw[guide]({1.5*cos(\a)},{1.5*sin(\a)})--(0,0);}\node[below]at(0,-1.9){观察屏中心示意};}
图\ref{fig:o3poissona}的中心亮点由圆周边缘到轴上等相位的衍射波叠加形成。实际有限尺寸光源、盘边误差、非轴对准和有限照明会降低或展宽亮点；理想强度结论限于轴上；阴影区其余位置仍可很暗，总透过功率也因圆盘遮挡而减小。
\end{solution}

\begin{exercise}\label{ex:o3babinet}
同一平面内两互补孔径的透过函数满足$t_1(\xi,\eta)+t_2(\xi,\eta)=1$。推导两者远场的关系，说明细丝与同宽狭缝何时具有相同衍射强度。
\Diagram{fig:o3babinetq}{细丝遮挡与同宽狭缝互补}{\begin{scope}[xshift=-2cm]\draw[curve](-1,-1.2)rectangle(1,1.2);\draw[very thick](0,-1.2)--(0,1.2);\node[below]at(0,-1.5){细丝};\end{scope}\begin{scope}[xshift=2cm]\draw[curve](-1,-1.2)--(-.15,-1.2)--(-.15,1.2)--(-1,1.2);\draw[curve](1,-1.2)--(.15,-1.2)--(.15,1.2)--(1,1.2);\node[below]at(0,-1.5){狭缝};\end{scope}}
\end{exercise}
\begin{solution}
图\ref{fig:o3babinetq}由衍射积分线性，$U_1+U_2=U_{\rm free}$。理想无限平面波的无遮挡远场只在前向方向有集中分量，因此离开前向时$U_2=-U_1$、$|U_2|^2=|U_1|^2$。
\Diagram{fig:o3babineta}{轴外两复场等大反相，强度相同}{\draw[axis](0,0)--(2.4,0)node[right]{$U_1$};\draw[axis](0,0)--(-2.4,0)node[left]{$U_2$};\node[above]at(0,1){轴外$U_{\rm free}=0$};}
所以细丝轴外暗纹可按同宽单缝公式测其宽度，如图\ref{fig:o3babineta}。但中央透过背景、有限照明包络和物体附近矢量边界场不相同，不能说两者整个屏幕总光强或总透过功率都相同。原理适用于相同入射波和标量孔径模型。
\end{solution}

\begin{exercise}\label{ex:o3phasehalves}
宽$a$狭缝的上、下半缝照明振幅相同，但相位相差$\pi$。以缝中心为原点，取场$+U_0$于$0<\xi<a/2$、$-U_0$于$-a/2<\xi<0$。求远场强度和零点。
\Diagram{fig:o3halvesq}{同一孔径两半具有反号复振幅}{\draw[curve](0,-1.5)rectangle(.7,1.5);\draw[curve](0,0)--(.7,0);\node[right]at(1,.8){$+U_0$};\node[right]at(1,-.8){$-U_0$};\node[left]at(0,0){$a$};}
\end{exercise}
\begin{solution}
图\ref{fig:o3halvesq}中上下半缝的场符号相反。以$q=k\sin\theta$，分别作两个定积分：
\begin{align*}
 F(q)/U_0&=\int_0^{a/2}\ee^{-iq\xi}\dd\xi
 -\int_{-a/2}^0\ee^{-iq\xi}\dd\xi\\
 &=\left[\frac{\ee^{-iq\xi}}{-iq}\right]_0^{a/2}
 -\left[\frac{\ee^{-iq\xi}}{-iq}\right]_{-a/2}^0\\
 &=\frac{1-\ee^{-iqa/2}}{iq}
 -\frac{\ee^{iqa/2}-1}{iq}\\
 &=\frac{2-2\cos(qa/2)}{iq}
 =-\frac{2i}{q}[1-\cos(qa/2)].
\end{align*}
故$I\propto4U_0^2[1-\cos(qa/2)]^2/q^2$。由于孔面场是实奇函数，远场为纯虚奇函数，也可由积分对称性预先判断。
\Diagram{fig:o3halvesa}{轴上因两半相消为暗，光能分到两侧瓣}{\draw[axis](-3.5,0)--(3.7,0)node[right]{$qa/2$};\draw[axis](0,0)--(0,2.7)node[above]{相对强度};\draw[curve,domain=-8:-.01,samples=180]plot({\x/2.5},{3.8*(1-cos(\x r))^2/(\x*\x)});\draw[curve,domain=.01:8,samples=180]plot({\x/2.5},{3.8*(1-cos(\x r))^2/(\x*\x)});\node[below]at(0,0){$0$};\node[below]at(2.51327,0){$2\pi$};\node[below]at(-2.51327,0){$-2\pi$};}
图\ref{fig:o3halvesa}的轴上必须取极限。展开$1-\cos(qa/2)=q^2a^2/8+O(q^4)$，得到$F\simeq-iU_0qa^2/4\to0$。对$q\ne0$，其他零点要求$qa/2=2m\pi$，即$a\sin\theta=2m\lambda$，$m\ne0$。

令$u=qa/4$，可写成$F=-iU_0a\sin^2u/u$。首个正旁瓣在$0<u<\pi$内，场模正比$\sin^2u/u$。用商法则，
\[
 \frac{\dd}{\dd u}\frac{\sin^2u}{u}
 =\frac{2u\sin u\cos u-\sin^2u}{u^2}.
\]
内点$\sin u\ne0$，驻值方程为$2u\cos u=\sin u$，即$\tan u=2u$；取该区间的非零根。相位片没有改变孔面强度，按帕塞瓦尔关系总远场功率保持不变，原中央能量转移到两侧。
\end{solution}

\begin{exercise}\label{ex:o3sinusoidalphase}
在近轴薄元件的标量模型中，无限宽相位光栅透过率为$t(x)=\exp[i\phi\cos(2\pi x/b)]$。设正入射，$\lambda/b\ll1$，忽略反射及各低级衍射束方向余弦的差别。定义第$m$级的傅里叶系数$c_m$，在$|\phi|\ll1$时求零级、正负一级振幅到所需阶，并验证这一模型中的功率守恒到二阶。
\Diagram{fig:o3sinephaseq}{正弦相位调制不改变孔面振幅}{\draw[axis](0,0)--(5,0)node[right]{$x$};\draw[axis](0,-1.4)--(0,1.7)node[above]{相位};\draw[curve,domain=0:4.5,samples=100]plot(\x,{cos(240*\x)});\node[below]at(2.7,-1.4){周期$b$};}
\end{exercise}
\begin{solution}
图\ref{fig:o3sinephaseq}的周期函数展开为$t=\sum_mc_m\ee^{i2\pi mx/b}$，系数为
\[
 c_m=\frac1b\int_0^bt(x)\ee^{-i2\pi mx/b}\dd x.
\]
令$u=2\pi x/b$，弱调制展开为$t\simeq1+i\phi\cos u-\phi^2\cos^2u/2$。
\Diagram{fig:o3sinephasea}{弱相位光栅把少量零级功率转入两侧一级}{\draw[axis](-2.5,0)--(2.7,0)node[right]{$m$};\draw[very thick](0,0)--(0,2.4);\draw[very thick](-1.5,0)--(-1.5,.6);\draw[very thick](1.5,0)--(1.5,.6);\node[below]at(-1.5,0){$-1$};\node[below]at(1.5,0){$1$};}
用$\cos u=(\ee^{iu}+\ee^{-iu})/2$和$\cos^2u=(1+\cos2u)/2$，逐项收集指数函数：
\[
 t=1-\frac{\phi^2}{4}
 +\frac{i\phi}{2}(\ee^{iu}+\ee^{-iu})
 -\frac{\phi^2}{8}(\ee^{2iu}+\ee^{-2iu})+O(\phi^3).
\]
因此$c_0=1-\phi^2/4+O(\phi^4)$，$c_{\pm1}=i\phi/2+O(\phi^3)$；
二级振幅虽为$O(\phi^2)$，其功率仅从四阶起算。在题设的近轴归一化下，图\ref{fig:o3sinephasea}中三束的功率份额分别为
\[
 \eta_0=|c_0|^2=1-\frac{\phi^2}{2}+O(\phi^4),\qquad
 \eta_{\pm1}=|c_{\pm1}|^2=\frac{\phi^2}{4}+O(\phi^4).
\]
相加得到$\eta_0+\eta_{+1}+\eta_{-1}=1+O(\phi^4)$，即守恒到二阶。
图中竖线只标示三束分配，不按小参数的真实比例绘制；弱调制时一级功率远小于零级。

传播角满足$\sin\theta_m=m\lambda/b$，题设保证所保留的一级为近轴传播波。
若衍射角较大，以电场幅值计量时法向能流还含方向余弦因子；若某级不传播，它也不能作为远场输出功率相加。
大角度情形的能流分配应由实际边界条件和各级传播方向确定。孔面恒模条件$|t|=1$给出傅里叶系数的平方和，却不足以单独确定这些法向能流。相位调制较大时，还须保留更多衍射级。
\end{solution}

\begin{exercise}\label{ex:o3camera}
直径$D=50$ mm、焦距$f=200$ mm的理想圆孔成像系统在550 nm工作。求后焦面点像首暗环半径与瑞利角间隔，并与宽度同为50 mm的长直缝比较。
\Diagram{fig:o3cameraq}{圆孔点像落在后焦面}{\OptAxis\draw[<->](1,-1.3)--(1,1.3);\draw[curve](5,-1.6)--(5,1.6);\draw[curve](1,1.2)--(5,0);\draw[curve](1,-1.2)--(5,0);\node[above]at(1,1.3){$D$};\node[above]at(5,1.6){焦面};}
\end{exercise}
\begin{solution}
图\ref{fig:o3cameraq}圆孔首暗角$\theta_0\simeq1.22\lambda/D=1.342\times10^{-5}$弧度，焦面半径$\rho_0=f\theta_0\simeq2.684\,\mu\mathrm m$。
\Diagram{fig:o3cameraa}{圆孔圆斑半径与长缝单方向宽度使用不同系数}{\draw[curve](-1.5,0)circle(.7);\draw[curve](1.5,-1)--(1.5,1);\draw[curve](2.1,-1)--(2.1,1);\node[below]at(-1.5,-1.3){圆孔};\node[below]at(1.8,-1.3){长缝};}
同宽长缝首暗角$\lambda/a=1.10\times10^{-5}$弧度，焦面半宽为$2.20\,\mu\mathrm m$。图\ref{fig:o3cameraa}比较的都是中心到首零，圆孔有1.22因子，长缝没有。两者通光面积、二维像形和旁瓣分布不同，不能仅以一个更小数值判定所有成像指标更好。
\end{solution}

\begin{exercise}\label{ex:o3missinggeneral}
光栅周期与缝宽之比$b/a=p/q$，其中$p,q$为互素正整数且$p>q$。求缺级的全部整数主极大级次。
\Diagram{fig:o3missinggeneralq}{缺级要求光栅主峰和单缝暗纹重合}{\draw[axis](-.3,0)--(5,0)node[right]{级次};\foreach\x in{0,1,2,3,4}{\draw[curve](\x,0)--(\x,1);}\draw[guide](0,1.8)..controls(2,2.5)and(3,.1)..(4.5,0);\node[right]at(4.5,1.6){包络零点};}
\end{exercise}
\begin{solution}
图\ref{fig:o3missinggeneralq}重合条件$b\sin\theta=m\lambda$、$a\sin\theta=s\lambda$，$s\ne0$。相除$m/s=p/q$，所以$q m=p s$。由$p,q$互素，$m=kp,s=kq$，$k$为非零整数。
\Diagram{fig:o3missinggenerala}{最先缺掉的正级次为$p$，随后为其整数倍}{\draw[axis](0,0)--(6.7,0)node[right]{$m$};\foreach\x in{1,2,3,4,5,6}{\draw[curve](\x,0)--(\x,.8);}\foreach\x in{3,6}{\draw[thick](\x-.2,.5)--(\x+.2,.9);\draw[thick](\x-.2,.9)--(\x+.2,.5);}\node[above]at(3,1.2){例$p=3$};}
故缺级为$\pm p,\pm2p,\ldots$，前提这些级次本来可传播。图\ref{fig:o3missinggenerala}显示$p=3$例。若$b/a$为无理数，理想精确模型没有完全重合的非零级次，但可能有非常接近的弱主峰，实验上是否可见还取决于包络和动态范围。
\end{solution}

\begin{exercise}\label{ex:o3exposure}
双缝中一缝被振动镜调制相位$\phi(t)=\phi_a\cos\Omega t$。给出长曝光干涉可见度的积分表达式，并在$\phi_a\ll1$时求二阶近似。
\Diagram{fig:o3exposureq}{一臂相位作周期调制}{\draw[->](0,.8)--(4.5,.8);\draw[->](0,-.3)--(4.5,-.3);\draw[very thick](2,-.8)--(2,.2);\draw[<->](1.7,-1.3)--(2.3,-1.3);\node[below]at(2,-1.5){振动元件};}
\end{exercise}
\begin{solution}
图\ref{fig:o3exposureq}中两束等强，固定屏点的瞬时强度为$2I_1[1+\cos(\delta_0+\phi_a\cos\Omega t)]$。长曝光覆盖许多周期，可将时间平均改成对$u=\Omega t$的周期平均，
\[
 G=\frac1{2\pi}\int_0^{2\pi}\ee^{i\phi_a\cos u}\dd u=J_0(\phi_a).
\]
虚部为零，因为$u$与$u+\pi$的正弦项相反；实部保留。平均强度为$2I_1[1+J_0(\phi_a)\cos\delta_0]$，可见度是$|J_0(\phi_a)|$。
\Diagram{fig:o3exposurea}{小幅调制的平均相量比单位长度缩短}{\draw[curve](0,0)circle(1.4);\draw[axis](0,0)--(1.4,0);\draw[guide](0,0)--(1.2,.72);\draw[guide](0,0)--(1.2,-.72);\draw[very thick,->](0,0)--(1.15,0);\node[below]at(0,-1.8){平均相量沿实轴};}
图\ref{fig:o3exposurea}表示相量绕静态相位来回摆动后的平均。使用例\ref{ex:o3circular}已从角积分推导的$J_0$级数，分别取$j=0,1,2$项，得到
\[
 J_0(\phi_a)=1-\frac{\phi_a^2}{4}+\frac{\phi_a^4}{64}+\cdots.
\]
到二阶$\mathcal V\simeq1-\phi_a^2/4$。若振幅增大使$J_0$越过零点，条纹先消失再出现，且亮暗相位翻转$\pi$；因此大振幅下不能把二阶式外推成负可见度。相位方差虽为$\phi_a^2/2$，其分布并非高斯，直接套高斯相位噪声公式只会在最低非平凡阶巧合一致。短曝光还应保留实际起始相位和曝光区间。
\end{solution}

\begin{exercise}\label{ex:o3gratingorderoverlap}
光栅中波长$\lambda_1$的$m$级与波长$\lambda_2$的$(m+1)$级方向重合。求两波长关系，说明增加$N$能否分开这种重级。
\Diagram{fig:o3overlapq}{不同波长、不同级次可以具有同一出射角}{\draw[curve](0,-1.2)--(0,1.2);\draw[axis](0,0)--(3.7,1.6)node[right]{共同方向};\draw[guide](0,0)--(4,0);\node[above]at(1.9,1.8){$m\lambda_1=(m+1)\lambda_2$};}
\end{exercise}
\begin{solution}
图\ref{fig:o3overlapq}同角要求$b\sin\theta=m\lambda_1=(m+1)\lambda_2$，所以$\lambda_2=m\lambda_1/(m+1)$。两个理想峰中心完全重合，即使$N$增大、各峰变窄，也不会产生角位置分离。
\Diagram{fig:o3overlapa}{同中心的两峰变窄后仍重叠，需要级次筛选或额外色散}{\draw[axis](-2.3,0)--(2.5,0)node[right]{$\theta-\theta_0$};\draw[curve,domain=-2:2,samples=80]plot(\x,{2.2*exp(-4*\x*\x)});\draw[guide,domain=-2:2,samples=80]plot(\x,{1.5*exp(-8*\x*\x)});}
图\ref{fig:o3overlapa}说明同级分辨率$mN$不能处理不同级完全重叠。可用滤光片限制输入谱段，或用棱镜等另一色散方向交叉分散。若两波长的单孔包络不同，强度可不同，但仅凭同一位置的强度不能唯一拆出两个未知谱线。
\end{solution}

\begin{exercise}\label{ex:o3illumination}
总周期数$N$的光栅只有中间$N/2$条被均匀照亮，其余无光。求实际分辨率和主峰宽度相对全部照明的变化，讨论保持总入射功率与保持每缝照明两种条件。
\Diagram{fig:o3illumq}{实体缝数与实际照明缝数不同}{\draw[curve](0,-.6)rectangle(5,.8);\foreach\x in{.25,.75,1.25,1.75,2.25,2.75,3.25,3.75,4.25,4.75}{\draw[curve](\x,-.6)--(\x,.8);}\draw[guide](1.25,-1)--(1.25,1.2);\draw[guide](3.75,-1)--(3.75,1.2);\node[above]at(2.5,1.2){照明区域};}
\end{exercise}
\begin{solution}
图\ref{fig:o3illumq}实际相干求和只有$N_{\rm eff}=N/2$项，首零角半宽变为原来2倍，瑞利分辨率$mN/2$。计算应采用实际受到相干照明的条数。
\Diagram{fig:o3illuma}{缩小照明宽度加宽主峰；纵轴高度取决于功率约束}{\draw[axis](-2.5,0)--(2.7,0)node[right]{角偏移};\draw[curve,domain=-2.3:2.3,samples=80]plot(\x,{2.3*exp(-4*\x*\x)});\draw[guide,domain=-2.3:2.3,samples=80]plot(\x,{1.15*exp(-\x*\x)});}
若每缝场不变，主峰强度按项数平方降为1/4，总透过功率降为1/2；若把原总功率集中到半数缝，每缝强度加倍，则主峰降为1/2，总功率不变。图\ref{fig:o3illuma}采用后一种归一化示意。角宽结论不受这两种整体振幅标定的影响。
\end{solution}

\begin{exercise}\label{ex:o3tiltblaze}
反射闪耀光栅的入射方向相对宏观外法线与出射方向分居两侧，角度大小为$i,\theta$。小面法线相对宏观法线倾角$\alpha$。写出小面镜面条件和光栅相长条件，说明为何必须先画方向再定正负。
\Diagram{fig:o3tiltblazeq}{入射与出射分居法线两侧，采用角度大小约定}{\draw[curve](-2.5,0)--(2.5,0);\draw[guide](0,0)--(0,2.7);\draw[->](-1.5,2.1)--(0,0);\draw[->](0,0)--(2,1.7);\draw[guide](0,0)--(.6,2.6);\node[above]at(.6,2.6){小面法线};\node[left]at(-.7,1.3){$i$};\node[right]at(1,1){$\theta$};}
\end{exercise}
\begin{solution}
按图\ref{fig:o3tiltblazeq}，从表面指向光源的方向角为$-i$，出射为$+\theta$。镜面法线平分二者，故$\theta-i=2\alpha$。入射实际波矢沿向下向右，切向分量为$+k\sin i$；出射切向为$+k\sin\theta$，所以相长条件可取$b(\sin\theta-\sin i)=m\lambda$。
\Diagram{fig:o3tiltblazea}{以实际波矢切向分量之差确定级次，法线平分反向入射与出射}{\draw[axis](-2.2,0)--(2.8,0)node[right]{切向};\draw[axis](0,0)--(1.4,2)node[right]{$\vct k_o$};\draw[axis](-1.6,2)--(0,0)node[midway,left]{$\vct k_i$};\draw[guide](0,0)--(.4,2.6);}
联立$\theta=i+2\alpha$，得$\lambda_B=b[\sin(i+2\alpha)-\sin i]/m$。正常入射恢复$b\sin2\alpha/m$。图\ref{fig:o3tiltblazea}的角度约定不同于有些把入射与出射都按波矢同侧有向角定义的写法；公式应从图中切向波矢推得，不能把两个正弦的加减号脱离图使用。
\end{solution}

\begin{exercise}\label{ex:o3farfield}
宽$a=0.2$ mm的缝受$\lambda=500$ nm平面波照射，屏距$L=1$ m。估计菲涅耳数，求首暗纹位置，并说明把宽度加倍对近似条件与主瓣宽度的影响。
\Diagram{fig:o3farq}{有限宽缝到远屏的几何尺度}{\draw[curve](0,-1.2)--(0,-.2);\draw[curve](0,.2)--(0,1.2);\draw[curve](5,-1.5)--(5,1.8);\draw[guide](0,0)--(5,0);\node[above]at(2.5,0){$L$};\node[left]at(0,0){$a$};}
\end{exercise}
\begin{solution}
用全宽定义估计数$F_a=a^2/(\lambda L)=.08$；若以半宽定义，数值为.02。图\ref{fig:o3farq}孔边相对中心的二次相位是$\pi a^2/(4\lambda L)=.0628$弧度，足够小，夫琅禾费近似合理。
\Diagram{fig:o3fara}{宽缝使主瓣缩窄，同时加大被省略的二次相位}{\draw[axis](-2.7,0)--(2.9,0)node[right]{$x$};\draw[curve](-2,0)..controls(-.7,2.8)and(.7,2.8)..(2,0);\draw[guide](-1,0)..controls(-.4,2.8)and(.4,2.8)..(1,0);\node[above]at(0,2.4){$a$与$2a$示意};}
首暗纹$\sin\theta=\lambda/a=.0025$，屏坐标$x=L\tan\theta\simeq2.5$ mm。宽度加倍后首暗纹约1.25 mm，但菲涅耳数增为4倍，二次相位.251弧度，远场近似精度下降。图\ref{fig:o3fara}只比较主瓣宽度；若孔面照明强度固定，轴上强度还会增加，不能把归一化曲线当绝对光强不变。
\end{solution}

\begin{exercise}\label{ex:o3energy}
均匀平面波照亮单缝，缝宽由$a$增到$2a$。为何轴上强度增为4倍，总透过功率却只增为2倍？用积分说明。
\Diagram{fig:o3energyq}{缝宽变化，孔面单位面积照明保持不变}{\draw[curve](0,-1)rectangle(.6,1);\draw[curve](3,-1)rectangle(4.2,1);\node[below]at(.3,-1){$a$};\node[below]at(3.6,-1){$2a$};\draw[->](1,0)--(2.6,0);}
\end{exercise}
\begin{solution}
图\ref{fig:o3energyq}中，孔面场强固定，轴上各点同相，所以$U(0)\propto a$、$I(0)\propto a^2$。设一维孔函数$g(\xi)=1$于$|\xi|<a/2$，其他为零，其傅里叶变换$F(q)=a\operatorname{sinc}(qa/2)$。帕塞瓦尔关系给
\[
 \int_{-\infty}^{\infty}|F(q)|^2\frac{\dd q}{2\pi}
 =\int_{-\infty}^{\infty}|g(\xi)|^2\dd\xi=a.
\]
令$u=qa/2$，得到$\int_{-\infty}^{\infty}\operatorname{sinc}^2u\dd u=\pi$，也明确了所用积分常数。
\Diagram{fig:o3energya}{同一照明下的准确单缝强度；宽度加倍使轴上峰高增为4倍}{
\draw[axis](-3.1,0)--(3.3,0)node[right]{$u$};
\draw[axis](0,0)--(0,2.65)node[above]{$I/I_a(0)$};
\draw[curve]plot coordinates{(-2.800000,0.0051091) (-2.786667,0.0047656) (-2.773333,0.0044207) (-2.760000,0.0040761) (-2.746667,0.0037336) (-2.733333,0.0033951) (-2.720000,0.0030624) (-2.706667,0.0027374) (-2.693333,0.0024218) (-2.680000,0.0021177) (-2.666667,0.0018269) (-2.653333,0.0015511) (-2.640000,0.0012923) (-2.626667,0.0010523) (-2.613333,0.0008327) (-2.600000,0.0006353) (-2.586667,0.0004617) (-2.573333,0.0003136) (-2.560000,0.0001923) (-2.546667,0.0000995) (-2.533333,0.0000363) (-2.520000,0.0000041) (-2.506667,0.0000040) (-2.493333,0.0000371) (-2.480000,0.0001042) (-2.466667,0.0002061) (-2.453333,0.0003436) (-2.440000,0.0005172) (-2.426667,0.0007273) (-2.413333,0.0009741) (-2.400000,0.0012578) (-2.386667,0.0015784) (-2.373333,0.0019356) (-2.360000,0.0023291) (-2.346667,0.0027584) (-2.333333,0.0032228) (-2.320000,0.0037216) (-2.306667,0.0042538) (-2.293333,0.0048182) (-2.280000,0.0054136) (-2.266667,0.0060385) (-2.253333,0.0066912) (-2.240000,0.0073702) (-2.226667,0.0080734) (-2.213333,0.0087988) (-2.200000,0.0095443) (-2.186667,0.0103076) (-2.173333,0.0110863) (-2.160000,0.0118778) (-2.146667,0.0126794) (-2.133333,0.0134885) (-2.120000,0.0143022) (-2.106667,0.0151175) (-2.093333,0.0159316) (-2.080000,0.0167413) (-2.066667,0.0175436) (-2.053333,0.0183353) (-2.040000,0.0191133) (-2.026667,0.0198744) (-2.013333,0.0206154) (-2.000000,0.0213332) (-1.986667,0.0220247) (-1.973333,0.0226867) (-1.960000,0.0233163) (-1.946667,0.0239104) (-1.933333,0.0244662) (-1.920000,0.0249809) (-1.906667,0.0254518) (-1.893333,0.0258764) (-1.880000,0.0262522) (-1.866667,0.0265770) (-1.853333,0.0268487) (-1.840000,0.0270654) (-1.826667,0.0272254) (-1.813333,0.0273271) (-1.800000,0.0273693) (-1.786667,0.0273508) (-1.773333,0.0272709) (-1.760000,0.0271290) (-1.746667,0.0269247) (-1.733333,0.0266581) (-1.720000,0.0263293) (-1.706667,0.0259389) (-1.693333,0.0254878) (-1.680000,0.0249770) (-1.666667,0.0244080) (-1.653333,0.0237825) (-1.640000,0.0231026) (-1.626667,0.0223707) (-1.613333,0.0215896) (-1.600000,0.0207622) (-1.586667,0.0198919) (-1.573333,0.0189824) (-1.560000,0.0180377) (-1.546667,0.0170621) (-1.533333,0.0160602) (-1.520000,0.0150370) (-1.506667,0.0139978) (-1.493333,0.0129480) (-1.480000,0.0118934) (-1.466667,0.0108403) (-1.453333,0.0097949) (-1.440000,0.0087637) (-1.426667,0.0077538) (-1.413333,0.0067721) (-1.400000,0.0058260) (-1.386667,0.0049228) (-1.373333,0.0040704) (-1.360000,0.0032764) (-1.346667,0.0025488) (-1.333333,0.0018957) (-1.320000,0.0013253) (-1.306667,0.0008458) (-1.293333,0.0004656) (-1.280000,0.0001930) (-1.266667,0.0000364) (-1.253333,0.0000040) (-1.240000,0.0001043) (-1.226667,0.0003456) (-1.213333,0.0007359) (-1.200000,0.0012834) (-1.186667,0.0019960) (-1.173333,0.0028816) (-1.160000,0.0039476) (-1.146667,0.0052016) (-1.133333,0.0066507) (-1.120000,0.0083018) (-1.106667,0.0101615) (-1.093333,0.0122363) (-1.080000,0.0145321) (-1.066667,0.0170546) (-1.053333,0.0198090) (-1.040000,0.0228003) (-1.026667,0.0260330) (-1.013333,0.0295110) (-1.000000,0.0332381) (-0.986667,0.0372172) (-0.973333,0.0414511) (-0.960000,0.0459419) (-0.946667,0.0506911) (-0.933333,0.0556999) (-0.920000,0.0609687) (-0.906667,0.0664974) (-0.893333,0.0722856) (-0.880000,0.0783319) (-0.866667,0.0846346) (-0.853333,0.0911913) (-0.840000,0.0979990) (-0.826667,0.1050541) (-0.813333,0.1123524) (-0.800000,0.1198892) (-0.786667,0.1276589) (-0.773333,0.1356556) (-0.760000,0.1438728) (-0.746667,0.1523031) (-0.733333,0.1609388) (-0.720000,0.1697716) (-0.706667,0.1787925) (-0.693333,0.1879921) (-0.680000,0.1973604) (-0.666667,0.2068868) (-0.653333,0.2165603) (-0.640000,0.2263693) (-0.626667,0.2363020) (-0.613333,0.2463458) (-0.600000,0.2564879) (-0.586667,0.2667151) (-0.573333,0.2770136) (-0.560000,0.2873696) (-0.546667,0.2977687) (-0.533333,0.3081964) (-0.520000,0.3186377) (-0.506667,0.3290776) (-0.493333,0.3395008) (-0.480000,0.3498918) (-0.466667,0.3602350) (-0.453333,0.3705147) (-0.440000,0.3807151) (-0.426667,0.3908205) (-0.413333,0.4008149) (-0.400000,0.4106826) (-0.386667,0.4204078) (-0.373333,0.4299750) (-0.360000,0.4393687) (-0.346667,0.4485734) (-0.333333,0.4575742) (-0.320000,0.4663560) (-0.306667,0.4749044) (-0.293333,0.4832051) (-0.280000,0.4912439) (-0.266667,0.4990075) (-0.253333,0.5064825) (-0.240000,0.5136563) (-0.226667,0.5205164) (-0.213333,0.5270512) (-0.200000,0.5332493) (-0.186667,0.5391000) (-0.173333,0.5445931) (-0.160000,0.5497191) (-0.146667,0.5544689) (-0.133333,0.5588342) (-0.120000,0.5628075) (-0.106667,0.5663815) (-0.093333,0.5695502) (-0.080000,0.5723078) (-0.066667,0.5746495) (-0.053333,0.5765711) (-0.040000,0.5780692) (-0.026667,0.5791412) (-0.013333,0.5797852) (0.000000,0.5800000) (0.013333,0.5797852) (0.026667,0.5791412) (0.040000,0.5780692) (0.053333,0.5765711) (0.066667,0.5746495) (0.080000,0.5723078) (0.093333,0.5695502) (0.106667,0.5663815) (0.120000,0.5628075) (0.133333,0.5588342) (0.146667,0.5544689) (0.160000,0.5497191) (0.173333,0.5445931) (0.186667,0.5391000) (0.200000,0.5332493) (0.213333,0.5270512) (0.226667,0.5205164) (0.240000,0.5136563) (0.253333,0.5064825) (0.266667,0.4990075) (0.280000,0.4912439) (0.293333,0.4832051) (0.306667,0.4749044) (0.320000,0.4663560) (0.333333,0.4575742) (0.346667,0.4485734) (0.360000,0.4393687) (0.373333,0.4299750) (0.386667,0.4204078) (0.400000,0.4106826) (0.413333,0.4008149) (0.426667,0.3908205) (0.440000,0.3807151) (0.453333,0.3705147) (0.466667,0.3602350) (0.480000,0.3498918) (0.493333,0.3395008) (0.506667,0.3290776) (0.520000,0.3186377) (0.533333,0.3081964) (0.546667,0.2977687) (0.560000,0.2873696) (0.573333,0.2770136) (0.586667,0.2667151) (0.600000,0.2564879) (0.613333,0.2463458) (0.626667,0.2363020) (0.640000,0.2263693) (0.653333,0.2165603) (0.666667,0.2068868) (0.680000,0.1973604) (0.693333,0.1879921) (0.706667,0.1787925) (0.720000,0.1697716) (0.733333,0.1609388) (0.746667,0.1523031) (0.760000,0.1438728) (0.773333,0.1356556) (0.786667,0.1276589) (0.800000,0.1198892) (0.813333,0.1123524) (0.826667,0.1050541) (0.840000,0.0979990) (0.853333,0.0911913) (0.866667,0.0846346) (0.880000,0.0783319) (0.893333,0.0722856) (0.906667,0.0664974) (0.920000,0.0609687) (0.933333,0.0556999) (0.946667,0.0506911) (0.960000,0.0459419) (0.973333,0.0414511) (0.986667,0.0372172) (1.000000,0.0332381) (1.013333,0.0295110) (1.026667,0.0260330) (1.040000,0.0228003) (1.053333,0.0198090) (1.066667,0.0170546) (1.080000,0.0145321) (1.093333,0.0122363) (1.106667,0.0101615) (1.120000,0.0083018) (1.133333,0.0066507) (1.146667,0.0052016) (1.160000,0.0039476) (1.173333,0.0028816) (1.186667,0.0019960) (1.200000,0.0012834) (1.213333,0.0007359) (1.226667,0.0003456) (1.240000,0.0001043) (1.253333,0.0000040) (1.266667,0.0000364) (1.280000,0.0001930) (1.293333,0.0004656) (1.306667,0.0008458) (1.320000,0.0013253) (1.333333,0.0018957) (1.346667,0.0025488) (1.360000,0.0032764) (1.373333,0.0040704) (1.386667,0.0049228) (1.400000,0.0058260) (1.413333,0.0067721) (1.426667,0.0077538) (1.440000,0.0087637) (1.453333,0.0097949) (1.466667,0.0108403) (1.480000,0.0118934) (1.493333,0.0129480) (1.506667,0.0139978) (1.520000,0.0150370) (1.533333,0.0160602) (1.546667,0.0170621) (1.560000,0.0180377) (1.573333,0.0189824) (1.586667,0.0198919) (1.600000,0.0207622) (1.613333,0.0215896) (1.626667,0.0223707) (1.640000,0.0231026) (1.653333,0.0237825) (1.666667,0.0244080) (1.680000,0.0249770) (1.693333,0.0254878) (1.706667,0.0259389) (1.720000,0.0263293) (1.733333,0.0266581) (1.746667,0.0269247) (1.760000,0.0271290) (1.773333,0.0272709) (1.786667,0.0273508) (1.800000,0.0273693) (1.813333,0.0273271) (1.826667,0.0272254) (1.840000,0.0270654) (1.853333,0.0268487) (1.866667,0.0265770) (1.880000,0.0262522) (1.893333,0.0258764) (1.906667,0.0254518) (1.920000,0.0249809) (1.933333,0.0244662) (1.946667,0.0239104) (1.960000,0.0233163) (1.973333,0.0226867) (1.986667,0.0220247) (2.000000,0.0213332) (2.013333,0.0206154) (2.026667,0.0198744) (2.040000,0.0191133) (2.053333,0.0183353) (2.066667,0.0175436) (2.080000,0.0167413) (2.093333,0.0159316) (2.106667,0.0151175) (2.120000,0.0143022) (2.133333,0.0134885) (2.146667,0.0126794) (2.160000,0.0118778) (2.173333,0.0110863) (2.186667,0.0103076) (2.200000,0.0095443) (2.213333,0.0087988) (2.226667,0.0080734) (2.240000,0.0073702) (2.253333,0.0066912) (2.266667,0.0060385) (2.280000,0.0054136) (2.293333,0.0048182) (2.306667,0.0042538) (2.320000,0.0037216) (2.333333,0.0032228) (2.346667,0.0027584) (2.360000,0.0023291) (2.373333,0.0019356) (2.386667,0.0015784) (2.400000,0.0012578) (2.413333,0.0009741) (2.426667,0.0007273) (2.440000,0.0005172) (2.453333,0.0003436) (2.466667,0.0002061) (2.480000,0.0001042) (2.493333,0.0000371) (2.506667,0.0000040) (2.520000,0.0000041) (2.533333,0.0000363) (2.546667,0.0000995) (2.560000,0.0001923) (2.573333,0.0003136) (2.586667,0.0004617) (2.600000,0.0006353) (2.613333,0.0008327) (2.626667,0.0010523) (2.640000,0.0012923) (2.653333,0.0015511) (2.666667,0.0018269) (2.680000,0.0021177) (2.693333,0.0024218) (2.706667,0.0027374) (2.720000,0.0030624) (2.733333,0.0033951) (2.746667,0.0037336) (2.760000,0.0040761) (2.773333,0.0044207) (2.786667,0.0047656) (2.800000,0.0051091)};
\draw[very thick]plot coordinates{(-2.800000,0.0116154) (-2.786667,0.0114606) (-2.773333,0.0112039) (-2.760000,0.0108491) (-2.746667,0.0104015) (-2.733333,0.0098685) (-2.720000,0.0092589) (-2.706667,0.0085833) (-2.693333,0.0078534) (-2.680000,0.0070824) (-2.666667,0.0062845) (-2.653333,0.0054744) (-2.640000,0.0046676) (-2.626667,0.0038798) (-2.613333,0.0031266) (-2.600000,0.0024235) (-2.586667,0.0017854) (-2.573333,0.0012262) (-2.560000,0.0007589) (-2.546667,0.0003952) (-2.533333,0.0001450) (-2.520000,0.0000165) (-2.506667,0.0000161) (-2.493333,0.0001479) (-2.480000,0.0004138) (-2.466667,0.0008134) (-2.453333,0.0013439) (-2.440000,0.0020003) (-2.426667,0.0027750) (-2.413333,0.0036583) (-2.400000,0.0046386) (-2.386667,0.0057018) (-2.373333,0.0068327) (-2.360000,0.0080140) (-2.346667,0.0092275) (-2.333333,0.0104539) (-2.320000,0.0116732) (-2.306667,0.0128654) (-2.293333,0.0140101) (-2.280000,0.0150876) (-2.266667,0.0160789) (-2.253333,0.0169661) (-2.240000,0.0177327) (-2.226667,0.0183641) (-2.213333,0.0188477) (-2.200000,0.0191732) (-2.186667,0.0193330) (-2.173333,0.0193223) (-2.160000,0.0191391) (-2.146667,0.0187847) (-2.133333,0.0182632) (-2.120000,0.0175820) (-2.106667,0.0167516) (-2.093333,0.0157853) (-2.080000,0.0146994) (-2.066667,0.0135126) (-2.053333,0.0122460) (-2.040000,0.0109227) (-2.026667,0.0095674) (-2.013333,0.0082063) (-2.000000,0.0068662) (-1.986667,0.0055745) (-1.973333,0.0043583) (-1.960000,0.0032444) (-1.946667,0.0022584) (-1.933333,0.0014246) (-1.920000,0.0007650) (-1.906667,0.0002996) (-1.893333,0.0000454) (-1.880000,0.0000161) (-1.866667,0.0002221) (-1.853333,0.0006698) (-1.840000,0.0013617) (-1.826667,0.0022962) (-1.813333,0.0034672) (-1.800000,0.0048646) (-1.786667,0.0064738) (-1.773333,0.0082763) (-1.760000,0.0102497) (-1.746667,0.0123678) (-1.733333,0.0146014) (-1.720000,0.0169182) (-1.706667,0.0192836) (-1.693333,0.0216612) (-1.680000,0.0240134) (-1.666667,0.0263018) (-1.653333,0.0284881) (-1.640000,0.0305345) (-1.626667,0.0324048) (-1.613333,0.0340647) (-1.600000,0.0354826) (-1.586667,0.0366302) (-1.573333,0.0374832) (-1.560000,0.0380218) (-1.546667,0.0382312) (-1.533333,0.0381019) (-1.520000,0.0376305) (-1.506667,0.0368192) (-1.493333,0.0356769) (-1.480000,0.0342186) (-1.466667,0.0324654) (-1.453333,0.0304449) (-1.440000,0.0281904) (-1.426667,0.0257407) (-1.413333,0.0231399) (-1.400000,0.0204364) (-1.386667,0.0176828) (-1.373333,0.0149346) (-1.360000,0.0122496) (-1.346667,0.0096873) (-1.333333,0.0073074) (-1.320000,0.0051693) (-1.306667,0.0033307) (-1.293333,0.0018469) (-1.280000,0.0007694) (-1.266667,0.0001453) (-1.253333,0.0000161) (-1.240000,0.0004167) (-1.226667,0.0013746) (-1.213333,0.0029092) (-1.200000,0.0050314) (-1.186667,0.0077423) (-1.173333,0.0110335) (-1.160000,0.0148866) (-1.146667,0.0192729) (-1.133333,0.0241538) (-1.120000,0.0294807) (-1.106667,0.0351953) (-1.093333,0.0412306) (-1.080000,0.0475111) (-1.066667,0.0539540) (-1.053333,0.0604702) (-1.040000,0.0669653) (-1.026667,0.0733413) (-1.013333,0.0794976) (-1.000000,0.0853329) (-0.986667,0.0907469) (-0.973333,0.0956416) (-0.960000,0.0999235) (-0.946667,0.1035055) (-0.933333,0.1063081) (-0.920000,0.1082617) (-0.906667,0.1093085) (-0.893333,0.1094032) (-0.880000,0.1085159) (-0.866667,0.1066323) (-0.853333,0.1037557) (-0.840000,0.0999079) (-0.826667,0.0951299) (-0.813333,0.0894830) (-0.800000,0.0830488) (-0.786667,0.0759295) (-0.773333,0.0682483) (-0.760000,0.0601481) (-0.746667,0.0517919) (-0.733333,0.0433611) (-0.720000,0.0350550) (-0.706667,0.0270885) (-0.693333,0.0196914) (-0.680000,0.0131054) (-0.666667,0.0075828) (-0.653333,0.0033833) (-0.640000,0.0007720) (-0.626667,0.0000161) (-0.613333,0.0013823) (-0.600000,0.0051336) (-0.586667,0.0115262) (-0.573333,0.0208063) (-0.560000,0.0332071) (-0.546667,0.0489453) (-0.533333,0.0682183) (-0.520000,0.0912013) (-0.506667,0.1180441) (-0.493333,0.1488689) (-0.480000,0.1837676) (-0.466667,0.2227995) (-0.453333,0.2659898) (-0.440000,0.3133277) (-0.426667,0.3647653) (-0.413333,0.4202165) (-0.400000,0.4795567) (-0.386667,0.5426225) (-0.373333,0.6092123) (-0.360000,0.6790863) (-0.346667,0.7519684) (-0.333333,0.8275471) (-0.320000,0.9054773) (-0.306667,0.9853832) (-0.293333,1.0668603) (-0.280000,1.1494785) (-0.266667,1.2327856) (-0.253333,1.3163104) (-0.240000,1.3995672) (-0.226667,1.4820588) (-0.213333,1.5632819) (-0.200000,1.6427303) (-0.186667,1.7199001) (-0.173333,1.7942936) (-0.160000,1.8654241) (-0.146667,1.9328202) (-0.133333,1.9960299) (-0.120000,2.0546250) (-0.106667,2.1082048) (-0.093333,2.1564001) (-0.080000,2.1988764) (-0.066667,2.2353370) (-0.053333,2.2655262) (-0.040000,2.2892312) (-0.026667,2.3062844) (-0.013333,2.3165650) (0.000000,2.3200000) (0.013333,2.3165650) (0.026667,2.3062844) (0.040000,2.2892312) (0.053333,2.2655262) (0.066667,2.2353370) (0.080000,2.1988764) (0.093333,2.1564001) (0.106667,2.1082048) (0.120000,2.0546250) (0.133333,1.9960299) (0.146667,1.9328202) (0.160000,1.8654241) (0.173333,1.7942936) (0.186667,1.7199001) (0.200000,1.6427303) (0.213333,1.5632819) (0.226667,1.4820588) (0.240000,1.3995672) (0.253333,1.3163104) (0.266667,1.2327856) (0.280000,1.1494785) (0.293333,1.0668603) (0.306667,0.9853832) (0.320000,0.9054773) (0.333333,0.8275471) (0.346667,0.7519684) (0.360000,0.6790863) (0.373333,0.6092123) (0.386667,0.5426225) (0.400000,0.4795567) (0.413333,0.4202165) (0.426667,0.3647653) (0.440000,0.3133277) (0.453333,0.2659898) (0.466667,0.2227995) (0.480000,0.1837676) (0.493333,0.1488689) (0.506667,0.1180441) (0.520000,0.0912013) (0.533333,0.0682183) (0.546667,0.0489453) (0.560000,0.0332071) (0.573333,0.0208063) (0.586667,0.0115262) (0.600000,0.0051336) (0.613333,0.0013823) (0.626667,0.0000161) (0.640000,0.0007720) (0.653333,0.0033833) (0.666667,0.0075828) (0.680000,0.0131054) (0.693333,0.0196914) (0.706667,0.0270885) (0.720000,0.0350550) (0.733333,0.0433611) (0.746667,0.0517919) (0.760000,0.0601481) (0.773333,0.0682483) (0.786667,0.0759295) (0.800000,0.0830488) (0.813333,0.0894830) (0.826667,0.0951299) (0.840000,0.0999079) (0.853333,0.1037557) (0.866667,0.1066323) (0.880000,0.1085159) (0.893333,0.1094032) (0.906667,0.1093085) (0.920000,0.1082617) (0.933333,0.1063081) (0.946667,0.1035055) (0.960000,0.0999235) (0.973333,0.0956416) (0.986667,0.0907469) (1.000000,0.0853329) (1.013333,0.0794976) (1.026667,0.0733413) (1.040000,0.0669653) (1.053333,0.0604702) (1.066667,0.0539540) (1.080000,0.0475111) (1.093333,0.0412306) (1.106667,0.0351953) (1.120000,0.0294807) (1.133333,0.0241538) (1.146667,0.0192729) (1.160000,0.0148866) (1.173333,0.0110335) (1.186667,0.0077423) (1.200000,0.0050314) (1.213333,0.0029092) (1.226667,0.0013746) (1.240000,0.0004167) (1.253333,0.0000161) (1.266667,0.0001453) (1.280000,0.0007694) (1.293333,0.0018469) (1.306667,0.0033307) (1.320000,0.0051693) (1.333333,0.0073074) (1.346667,0.0096873) (1.360000,0.0122496) (1.373333,0.0149346) (1.386667,0.0176828) (1.400000,0.0204364) (1.413333,0.0231399) (1.426667,0.0257407) (1.440000,0.0281904) (1.453333,0.0304449) (1.466667,0.0324654) (1.480000,0.0342186) (1.493333,0.0356769) (1.506667,0.0368192) (1.520000,0.0376305) (1.533333,0.0381019) (1.546667,0.0382312) (1.560000,0.0380218) (1.573333,0.0374832) (1.586667,0.0366302) (1.600000,0.0354826) (1.613333,0.0340647) (1.626667,0.0324048) (1.640000,0.0305345) (1.653333,0.0284881) (1.666667,0.0263018) (1.680000,0.0240134) (1.693333,0.0216612) (1.706667,0.0192836) (1.720000,0.0169182) (1.733333,0.0146014) (1.746667,0.0123678) (1.760000,0.0102497) (1.773333,0.0082763) (1.786667,0.0064738) (1.800000,0.0048646) (1.813333,0.0034672) (1.826667,0.0022962) (1.840000,0.0013617) (1.853333,0.0006698) (1.866667,0.0002221) (1.880000,0.0000161) (1.893333,0.0000454) (1.906667,0.0002996) (1.920000,0.0007650) (1.933333,0.0014246) (1.946667,0.0022584) (1.960000,0.0032444) (1.973333,0.0043583) (1.986667,0.0055745) (2.000000,0.0068662) (2.013333,0.0082063) (2.026667,0.0095674) (2.040000,0.0109227) (2.053333,0.0122460) (2.066667,0.0135126) (2.080000,0.0146994) (2.093333,0.0157853) (2.106667,0.0167516) (2.120000,0.0175820) (2.133333,0.0182632) (2.146667,0.0187847) (2.160000,0.0191391) (2.173333,0.0193223) (2.186667,0.0193330) (2.200000,0.0191732) (2.213333,0.0188477) (2.226667,0.0183641) (2.240000,0.0177327) (2.253333,0.0169661) (2.266667,0.0160789) (2.280000,0.0150876) (2.293333,0.0140101) (2.306667,0.0128654) (2.320000,0.0116732) (2.333333,0.0104539) (2.346667,0.0092275) (2.360000,0.0080140) (2.373333,0.0068327) (2.386667,0.0057018) (2.400000,0.0046386) (2.413333,0.0036583) (2.426667,0.0027750) (2.440000,0.0020003) (2.453333,0.0013439) (2.466667,0.0008134) (2.480000,0.0004138) (2.493333,0.0001479) (2.506667,0.0000161) (2.520000,0.0000165) (2.533333,0.0001450) (2.546667,0.0003952) (2.560000,0.0007589) (2.573333,0.0012262) (2.586667,0.0017854) (2.600000,0.0024235) (2.613333,0.0031266) (2.626667,0.0038798) (2.640000,0.0046676) (2.653333,0.0054744) (2.666667,0.0062845) (2.680000,0.0070824) (2.693333,0.0078534) (2.706667,0.0085833) (2.720000,0.0092589) (2.733333,0.0098685) (2.746667,0.0104015) (2.760000,0.0108491) (2.773333,0.0112039) (2.786667,0.0114606) (2.800000,0.0116154)};
\node[left]at(0,.58){$1$};\node[left]at(0,2.32){$4$};
\node[right]at(1.4,1.5){细线：$a$};\node[right]at(1.4,1.1){粗线：$2a$};
\node[below]at(0,-.5){$u=\pi ax/(\lambda L)$};
}
图\ref{fig:o3energya}比较同一照明强度下的两个孔宽。屏上近轴强度写成$I(x)=Ca^2\operatorname{sinc}^2[\pi ax/(\lambda L)]$，于是
\begin{align*}
 \int I(x)\dd x
 &=Ca^2\frac{\lambda L}{\pi a}
 \int\operatorname{sinc}^2u\dd u\\
 &=C\lambda La.
\end{align*}
因此$a\to2a$时峰高乘4、横向尺度减半、积分功率乘2。图中两条曲线按sinc平方绘制，旁瓣也参与上述全平面积分。对有限高度的矩形孔，还要沿另一屏坐标积分，最终功率正比真实孔面积。此守恒关系以同一入射场强、同一无损传播条件为前提；若固定的是总入射功率，孔面振幅就应随孔宽重新归一化。
\end{solution}

\begin{exercise}\label{ex:o3oblique}
周期$b$的透射光栅受入射角$i$的平面波照射，入射和出射横向波矢分别为$k\sin i$、$k\sin\theta$。推导光栅方程和同级角色散，并讨论入射角是否直接改变$mN$。
\Diagram{fig:o3obliqueq}{斜入射与出射的横向波矢投影}{\draw[curve](-2,0)--(2,0);\draw[guide](0,-1.7)--(0,2.2);\draw[->](-1.2,-1.7)--(0,0);\draw[->](0,0)--(1.7,2.1);\node[left]at(-.6,-1){$i$};\node[right]at(.9,1.1){$\theta$};}
\end{exercise}
\begin{solution}
图\ref{fig:o3obliqueq}沿孔径横坐标相邻周期的净相位差为$kb(\sin i-\sin\theta)$。相长条件取整数级次约定为$b(\sin\theta-\sin i)=m\lambda$。固定$i$，微分得到$\dd\theta/\dd\lambda=m/(b\cos\theta)$。
\Diagram{fig:o3obliquea}{首零角宽与谱线位移共享同一投影因子}{\draw[axis](0,0)--(4.6,0)node[right]{$\theta$};\draw[curve](1,0)..controls(1.6,2.6)and(1.8,2.6)..(2.4,0);\draw[guide](2,0)..controls(2.6,2.6)and(2.8,2.6)..(3.4,0);\node[below]at(2.3,-.3){$\Delta\theta$};}
首零半宽仍为$\lambda/(Nb\cos\theta)$，两式之比使投影因子消去，故按瑞利判据仍有$\lambda/\Delta\lambda=mN$。图\ref{fig:o3obliquea}表示在固定级次、固定实际被照亮缝数时的比较。斜入射会改变足迹宽度、可传播级次和效率，所以在固定入射束截面条件下，$N$本身可能改变。
\end{solution}

\begin{exercise}\label{ex:o3defocus}
在夫琅禾费装置中，观察屏偏离透镜后焦面。为什么衍射场重新出现孔径二次相位？以一维自由传播说明保留哪些项。
\Diagram{fig:o3defocusq}{焦面前后不同观察平面}{\OptAxis\draw[<->](1.5,-1.5)--(1.5,1.5);\draw[guide](4,-1.7)--(4,1.7);\draw[curve](5,-1.7)--(5,1.7);\node[above]at(4,1.7){焦面};\node[above]at(5,2.2){偏移屏};}
\end{exercise}
\begin{solution}
图\ref{fig:o3defocusq}中，令透镜前孔面场为$U_A(\xi)$。薄透镜相位为$\exp[-ik\xi^2/(2f)]$，向屏传播$L$的近轴核为$\exp[ik(x-\xi)^2/(2L)]$，两者相乘并展开平方，得到
\[
 U(x,L)=C_L\ee^{ikx^2/(2L)}\int U_A(\xi)
 \ee^{\frac{ik\xi^2}{2}(1/L-1/f)}\ee^{-ikx\xi/L}\dd\xi.
\]
观察点公共相位$\exp[ikx^2/(2L)]$不改变单点强度，孔径内的二次相位却改变各面元的相干叠加，不能略去。
\Diagram{fig:o3defocusa}{只有$L=f$时二次相位恰好抵消}{\begin{scope}[yscale=.85]\draw[axis](-2.2,0)--(2.4,0)node[right]{$\xi$};\draw[axis](0,-1.5)--(0,2)node[above]{残余相位};\draw[curve,domain=-1.8:1.8,samples=60]plot(\x,{.4*\x*\x});\draw[guide,domain=-1.8:1.8,samples=60]plot(\x,{-.3*\x*\x});\node[right]at(1.8,1.3){$L<f$};\node[right]at(1.8,-1){$L>f$};\end{scope}}
图\ref{fig:o3defocusa}中$L=f$时二次项准确抵消，远场成为孔函数在$q=kx/f$处的傅里叶变换。若$L=f+\Delta L$，则$1/L-1/f=-\Delta L/[f(f+\Delta L)]\simeq-\Delta L/f^2$。以孔径全宽$a$、边缘$\xi=a/2$计，最大残余相位大小近似为
\[
 |\Delta\phi_{\rm edge}|\simeq\frac{k a^2|\Delta L|}{8f^2}.
\]
要求它远小于1，得到$|\Delta L|\ll8f^2/(ka^2)$。若改用半宽$a$记号，分母系数随之改变，不能混用。由此可见，判断离焦是否可略，须同时考虑波长、孔径与焦距；单独的$|\Delta L|/f\ll1$只保证几何距离接近，并不保证孔内相位误差小。
\end{solution}

\begin{exercise}\label{ex:o3aperturephase}
两条相同狭缝中心间隔$b$，在一条缝前加入均匀相位片$\phi_0$。求衍射分布；若$\phi_0=\pi$，轴上是否仍有中央主极大？
\Diagram{fig:o3phaseq}{相位片只改变一缝整体相位，不改变其宽度}{\begin{scope}[yscale=.75]\draw[curve](0,-1.2)rectangle(.4,-.4);\draw[curve](0,.4)rectangle(.4,1.2);\draw[very thick](-.2,.4)--(-.2,1.2);\node[left]at(-.2,.8){$\phi_0$};\draw[axis](.8,0)--(4,0)node[right]{远场};\end{scope}}
\end{exercise}
\begin{solution}
取图\ref{fig:o3phaseq}两缝中心坐标为0和$b$，则$U=U_a[1+\ee^{i(\phi_0-kb\sin\theta)}]$，故$I=4I_a(0)\operatorname{sinc}^2\beta\cos^2[(\phi_0-kb\sin\theta)/2]$。
\Diagram{fig:o3phasea}{附加$\pi$相位把轴上两缝相量变为反向}{\draw[->,thick](0,0)--(2,0)node[right]{$A$};\draw[->,thick](0,0)--(-2,0)node[left]{$-A$};\node[above]at(0,.35){$\theta=0$};}
在$\phi_0=\pi$时，图\ref{fig:o3phasea}轴上相消，强度为零。单缝包络仍以轴为中心，但干涉细纹平移半个周期；不能把包络中心自动称为总强度的主极大。若相位片有吸收，振幅失配会留下轴上残光。
\end{solution}


\chapter{偏振、晶体与界面电磁波}
\section{Jones矢量与偏振椭圆}\label{th:o4jones}
在均匀各向同性介质中，平面波为横波。取传播方向$+z$，固定相位约定$\exp[i(kz-\omega t)]$，省略公共相位后的Jones矢量为
\[
 \vct e=\begin{pmatrix}a\\b\ee^{i\delta}\end{pmatrix},\quad a,b\geq0.
\]
某固定位置以$\omega t$为时间变量，可写$E_x=a\cos\omega t$、$E_y=b\cos(\omega t-\delta)$。Jones矢量只适用于确定的完全偏振分量；自然光的随机偏振不能用一个恒定Jones矢量替代。

若$\delta=0$或$\pi$，两分量同相或反相，末端沿直线运动。若$a=b$且$\delta=\pm\pi/2$，轨迹为圆。一般情形为椭圆。取$a=b,\delta=\pi/2$时，在$t=0$电场沿$+x$，随后转向$+y$。从$+z$端朝光源看、且把$x$向右$y$向上画在纸上，此旋转为逆时针；换观察方向会改变名称，故以下均以这个时序说明旋向。

\Diagram{fig:o4states}{同一观察约定下的线、圆和椭圆轨迹}{\begin{scope}[xshift=-3cm]\draw[axis](-1.1,0)--(1.2,0);\draw[axis](0,-1.1)--(0,1.2);\draw[very thick](-.8,-.8)--(.8,.8);\node[below]at(0,-1.3){线};\end{scope}\draw[axis](-1.1,0)--(1.2,0);\draw[axis](0,-1.1)--(0,1.2);\draw[curve](0,0)circle(.9);\draw[->](.9,0)arc(0:50:.9);\node[below]at(0,-1.3){圆};\begin{scope}[xshift=3cm]\draw[axis](-1.1,0)--(1.2,0);\draw[axis](0,-1.1)--(0,1.2);\draw[curve,rotate=25](0,0)ellipse(1 and .5);\node[below]at(0,-1.3){椭圆};\end{scope}}

\subsection{消元方程与主轴}\label{th:o4ellipse}
由$E_x/a=\cos\omega t$和$E_y/b=\cos\omega t\cos\delta+\sin\omega t\sin\delta$消去时间，得
\begin{equation}\label{eq:o4ellipse}
 \frac{E_x^2}{a^2}+\frac{E_y^2}{b^2}
 -\frac{2E_xE_y\cos\delta}{ab}=\sin^2\delta.
\end{equation}
若$a$或$b$为零，应直接回到参数式，避免除零。设主轴相对$x$转角为$\psi$，令$E_\xi=E_x\cos\psi+E_y\sin\psi$、$E_\eta=-E_x\sin\psi+E_y\cos\psi$。对一个周期取平均，$\langle E_x^2\rangle=a^2/2$、$\langle E_y^2\rangle=b^2/2$、$\langle E_xE_y\rangle=ab\cos\delta/2$，于是
\[
 2\langle E_\xi E_\eta\rangle
 =ab\cos\delta\cos2\psi+\frac{b^2-a^2}{2}\sin2\psi.
\]
非退化椭圆的两主轴分量相差$\pm\pi/2$，其相关平均为零；令右边为零得到
\[
 \tan2\psi=\frac{2ab\cos\delta}{a^2-b^2}.
\]
实际应由$\operatorname{atan2}(2ab\cos\delta,a^2-b^2)$确定$2\psi$。半轴平方是矩阵$2\langle\vct E\vct E^{\mathsf T}\rangle$的两个本征值$\Lambda$，其方程为
\[
 (a^2-\Lambda)(b^2-\Lambda)-a^2b^2\cos^2\delta=0,
 \qquad \Lambda^2-(a^2+b^2)\Lambda+a^2b^2\sin^2\delta=0.
\]
解这个二次方程，长短半轴平方分别为
\[
 A_\pm^2=\frac{a^2+b^2\pm\sqrt{(a^2-b^2)^2+4a^2b^2\cos^2\delta}}2.
\]
这是时间参数轨迹的最大、最小场强平方，不是分别把$x,y$分量振幅当成半轴。只有$\cos\delta=0$或特殊对称情形，坐标轴才直接就是椭圆主轴。

\begin{example}\label{ex:o4ellipse}
Jones矢量为$(2,\ee^{i\pi/3})^{\mathsf T}E_0$。求椭圆方程、主轴方向、半轴及旋转时序。
\Diagram{fig:o4ellipseq}{两正交分量的振幅和相位差}{\draw[axis](-2.5,0)--(2.7,0)node[right]{$E_x/E_0$};\draw[axis](0,-1.5)--(0,1.7)node[above]{$E_y/E_0$};\draw[guide](-2,-1)rectangle(2,1);\node[above]at(1.5,2.3){$a=2,\ b=1,\ \delta=\pi/3$};}
\end{example}
\begin{solution}
采用实场$\operatorname{Re}(\vct E\ee^{-i\omega t})$。图\ref{fig:o4ellipseq}中$X=E_x/E_0=2\cos\omega t$、$Y=E_y/E_0=\cos(\omega t-\pi/3)$。由$Y=X/4+(\sqrt3/2)\sin\omega t$消去时间，
\[
 (Y-X/4)^2=\tfrac34(1-X^2/4)
 \quad\Longrightarrow\quad X^2/4+Y^2-XY/2=3/4.
\]
二次型矩阵为$Q=\bigl(\begin{smallmatrix}1/4&-1/4\\-1/4&1\end{smallmatrix}\bigr)$，其特征方程$\lambda^2-(5/4)\lambda+3/16=0$，得$\lambda_\pm=(5\pm\sqrt{13})/8$。
\Diagram{fig:o4ellipsea}{倾斜椭圆的长短轴；箭头表示由参数方程确定的运动方向}{\draw[axis](-2.5,0)--(2.7,0)node[right]{$X$};\draw[axis](0,-1.6)--(0,1.8)node[above]{$Y$};\draw[curve,domain=0:360,samples=100]plot({2*cos(\x)},{cos(\x-60)});\draw[guide](-2.1,-.636)--(2.1,.636);\draw[guide](-.4,1.32)--(.4,-1.32);\draw[->,thick](2,.5)--(1.97,.643);}
图\ref{fig:o4ellipsea}的半轴平方为右侧常数$3/4$除以相应特征值，故长短半轴分别满足
\[
 A^2/E_0^2=(5+\sqrt{13})/2,\qquad
 B^2/E_0^2=(5-\sqrt{13})/2.
\]
长轴对应较小的二次型特征值。旋转坐标消去交叉项给$\tan2\psi=2/3$；结合较小特征值的特征向量，取$\psi=\tfrac12\arctan(2/3)=16.85^\circ$，而非相差$90^\circ$的短轴方向。

判断旋向须回到时间参数：$t=0$处$(X,Y)=(2,1/2)$，$\dot X=0$、$\dot Y=(\sqrt3/2)\omega>0$，所以电场从右侧向上运动，与图中箭头一致。几何方程只含$\cos\delta$和$\sin^2\delta$，对$\delta$与$-\delta$给出同一几何轨迹，旋向须由时间参数确定。
\end{solution}

\section{自然光、部分偏振与检偏}\label{th:o4partial}
理想线偏振片透过轴单位矢量$\vct p$，其Jones矩阵为$P=\vct p\vct p^{\mathsf T}$。线偏振入射与轴成$\theta$时，振幅乘$\cos\theta$，光强满足马吕斯定律$I=I_0\cos^2\theta$。自然光在任意正交方向平均等分能量，故理想偏振片透过$I_0/2$，与透过轴角无关。

圆偏振经过线偏振片也透过一半，与转角无关，但出射变成线偏振。椭圆偏振转动检偏器一般有非零最大和最小值，分别正比$A_+^2,A_-^2$。自然光叠加线偏振光时，$I(\theta)=I_u/2+I_p\cos^2(\theta-\theta_0)$，所以$I_{\max}=I_u/2+I_p$、$I_{\min}=I_u/2$，偏振度$I_p/(I_u+I_p)=(I_{\max}-I_{\min})/(I_{\max}+I_{\min})$。

若光中含圆偏振分量，旋转单个线偏振片得到的只是线偏振调制度。完整偏振度还需测量圆偏振信息。定义强度单位一致的斯托克斯量$S_0=\langle|E_x|^2+|E_y|^2\rangle$、$S_1=\langle|E_x|^2-|E_y|^2\rangle$、$S_2=2\operatorname{Re}\langle E_x^*E_y\rangle$、$S_3=2\operatorname{Im}\langle E_x^*E_y\rangle$，则完全偏振$S_1^2+S_2^2+S_3^2=S_0^2$；一般偏振度为$\sqrt{S_1^2+S_2^2+S_3^2}/S_0$。

\begin{example}\label{ex:o4malus}
两理想偏振片透过轴正交，自然光先经过第一片，再在中间插入轴与第一片成$\theta$角的第三片。求透过强度及最优角度。
\Diagram{fig:o4malusq}{三片顺序为$0,\theta,\pi/2$}{\foreach\x/\t in{0/0,2/\theta,4/{\pi/2}}{\draw[curve](\x,-1)--(\x,1);\node[above]at(\x,1){$\t$};}\draw[axis](-1,0)--(5,0)node[right]{光};}
\end{example}
\begin{solution}
图\ref{fig:o4malusq}中，自然光的横向偏振方向均匀分布，第一片透过强度为$I_0\langle\cos^2\varphi\rangle=I_0/2$。它之后的光沿第一片轴线偏振，不再是自然光。中间片与该轴成$\theta$角，透过场投影乘$\cos\theta$，强度乘$\cos^2\theta$；最后片与中间轴夹角为$\pi/2-\theta$，再乘$\sin^2\theta$，故
\[
 I=\frac{I_0}{2}\cos^2\theta\sin^2\theta
 =\frac{I_0}{8}\sin^22\theta.
\]
\Diagram{fig:o4malusa}{逐次投影后偏振方向改变，不能把中间片视为单纯削弱光强}{\draw[axis](0,0)--(2.2,0)node[right]{$P_1$};\draw[axis](0,0)--(1.6,1.6)node[right]{$P_2$};\draw[axis](0,0)--(0,2.2)node[above]{$P_3$};\draw[guide](2,0)--(1,1)--(0,1);}
图\ref{fig:o4malusa}按传播次序标为$P_1,P_2,P_3$，中间插入片为$P_2$，末片为$P_3$。由$\sin^22\theta\leq1$，最大透过强度为$I_0/8$，取$\theta=\pi/4$模$\pi/2$；在$0\leq\theta\leq\pi/2$内唯一最优角为45度。

逐段能量也可核对：最优时$I_1=I_0/2$、$I_2=I_0/4$、$I_3=I_0/8$，每片都吸收部分功率。没有中间片时两轴正交，最后为零；插入后中间片把输出偏振改为45度，使它在最后轴上有非零投影。透过增加来自偏振态改变，整个过程仍有$I_3<I_2<I_1<I_0$，并未由吸收元件产生能量。
\end{solution}

\section{单轴晶体中的普通光与非普通光}\label{th:o4crystal}
各向异性介质的线性关系为$\vct D=\varepsilon_0\boldsymbol\varepsilon\vct E$，主轴系中介电张量可对角化。一般$\vct D$与$\vct E$不平行。无自由电荷仍要求$\vct k\cdot\vct D=0$，却不能直接推为$\vct k\cdot\vct E=0$。

单轴晶体取光轴为$z$，$\boldsymbol\varepsilon=\operatorname{diag}(n_o^2,n_o^2,n_e^2)$，这里$n_e$是传播方向垂直光轴时的非普通主折射率。波矢与光轴夹角为$\vartheta$。普通光电场垂直由$\vct k$和光轴张成的主截面，折射率恒为$n_o$。非普通光的$\vct D$在主截面内并垂直$\vct k$，其色散曲面为
\[
 \frac{k_x^2}{n_e^2}+\frac{k_z^2}{n_o^2}=\frac{\omega^2}{c^2},
\]
所以有效折射率
\[
 \frac1{n_{e,\rm eff}^2(\vartheta)}=
 \frac{\cos^2\vartheta}{n_o^2}+\frac{\sin^2\vartheta}{n_e^2}.
\]
沿光轴两种波简并，都为$n_o$；垂直光轴时分别为$n_o,n_e$。非普通光的能流方向由群速度或坡印廷矢量决定，一般不平行于波矢，产生走离。相位匹配用波矢，光线方向与光束分离则要用能流，二者不能混写。

\Diagram{fig:o4crystal}{波矢与位移矢量垂直，电场和能流可偏斜}{\draw[axis](-.3,0)--(3.4,0)node[right]{$x$};\draw[axis](0,-.3)--(0,3.3)node[above]{光轴$z$};\draw[axis](0,0)--(2,2.4)node[right]{$\vct k$};\draw[axis](0,0)--(-1.8,1.5)node[left]{$\vct D_e$};\draw[axis](0,0)--(-1,1.9)node[above]{$\vct E_e$};\draw[guide](0,0)--(2.85,1.5)node[right]{$\vct S_e$};}

\section{波片及偏振态变换}\label{th:o4plates}
波片正常入射，两个本征偏振轴在片面内，取快轴折射率$n_f$、慢轴$n_s>n_f$，厚度$h$。删去共同相位$\exp(ik_0n_fh)$后，Jones矩阵为$J=\operatorname{diag}(1,\ee^{i\delta})$，$\delta=2\pi(n_s-n_f)h/\lambda_0$。四分之一波片$\delta=\pi/2$模$2\pi$；半波片$\delta=\pi$模$2\pi$。共同相位不影响单束偏振，但在干涉仪中仍可能影响两臂相位。

如果波片快轴相对实验室$x$转角$\alpha$，先用$R(-\alpha)$把分量投到快慢轴，再乘对角矩阵，最后转回：$J(\alpha)=R(\alpha)\operatorname{diag}(1,\ee^{i\delta})R(-\alpha)$，其中$R(\alpha)=\bigl(\begin{smallmatrix}\cos\alpha&-\sin\alpha\\\sin\alpha&\cos\alpha\end{smallmatrix}\bigr)$。矩阵顺序代表实际操作，交换两片的顺序一般改变结果。

\begin{example}\label{ex:o4quarter}
线偏振入射与四分之一波片快轴成$\theta$，求输出Jones矢量、椭圆轴及旋转方向，特别讨论$\theta=45^\circ$。
\Diagram{fig:o4quarterq}{输入线偏振在快慢轴上的投影}{\draw[axis](0,0)--(3,0)node[right]{快$x$};\draw[axis](0,0)--(0,2.8)node[above]{慢$y$};\draw[axis](0,0)--(2,2)node[right]{输入};\node at(.9,.35){$\theta$};}
\end{example}
\begin{solution}
图\ref{fig:o4quarterq}输入为$E_0(\cos\theta,\sin\theta)^{\mathsf T}$，输出为$E_0(\cos\theta,i\sin\theta)^{\mathsf T}$。两个分量正交相差$\pi/2$，所以主轴就是快慢轴，半轴绝对值$E_0|\cos\theta|$和$E_0|\sin\theta|$。
\Diagram{fig:o4quartera}{45度输入产生圆偏振，固定位置由$+x$转向$+y$}{\draw[axis](-1.5,0)--(1.7,0)node[right]{$x$};\draw[axis](0,-1.5)--(0,1.7)node[above]{$y$};\draw[curve](0,0)circle(1.1);\draw[->,thick](1.1,0)arc(0:65:1.1);\node[right]at(2,.5){$(1,i)/\sqrt2$};}
若$0<\theta<\pi/2$，固定位置$t=0$电场沿$+x$，随时间向$+y$转，图\ref{fig:o4quartera}按此方向画箭头。$\theta=45^\circ$为圆；$\theta=0,90^\circ$仍是线偏振。把片转90度等效交换快慢轴，会反转圆偏振旋向，不能只说“经过四分之一波片必变圆”。
\end{solution}

\begin{example}\label{ex:o4half}
入射线偏振方向相对实验室$x$为$\theta$，半波片快轴方向为$\alpha$。求出射偏振方向，并求将偏振方向由10度转到70度的快轴位置。
\Diagram{fig:o4halfq}{半波片快轴与输入偏振的夹角}{\draw[axis](0,0)--(3,0)node[right]{$x$};\draw[axis](0,0)--(2.5,1)node[right]{$\theta$};\draw[guide](0,0)--(1.6,2.5)node[above]{$\alpha$};}
\end{example}
\begin{solution}
在图\ref{fig:o4halfq}的快慢轴基底，输入为$(\cos(\theta-\alpha),\sin(\theta-\alpha))$。半波片把慢轴分量反号，成为$(\cos(\theta-\alpha),-\sin(\theta-\alpha))$，这仍是线偏振，方向在快轴另一侧。
\Diagram{fig:o4halfa}{半波片把输入方向关于快轴作镜面对称}{\draw[axis](0,0)--(3,0)node[right]{$x$};\draw[axis](0,0)--(2.8,.49)node[right]{10度};\draw[guide](0,0)--(2.3,1.93)node[right]{快轴40度};\draw[axis](0,0)--(1.03,2.82)node[above]{70度};}
转回实验室，输出方向$\theta_o=2\alpha-\theta$模$\pi$。故图\ref{fig:o4halfa}可取$\alpha=40^\circ$；快轴再转90度得到同一线偏振方向但共同相位可能不同。半波片旋转$\Delta\alpha$，输出线偏振方向转$2\Delta\alpha$。
\end{solution}

\begin{example}\label{ex:o4identify}
未知光可能为完全线、圆、椭圆偏振，或自然光。给出理想线偏振片和可转四分之一波片的鉴别程序，说明部分偏振光需要怎样补充。
\Diagram{fig:o4identifyq}{先测线检偏响应，再加入可转波片}{\draw[axis](-1,0)--(6,0);\draw[curve](0,-1)--(0,1);\draw[curve](2,-1)--(2,1);\draw[curve](4,-1)--(4,1);\node[above]at(0,1){波片};\node[above]at(2,1){检偏器};\node[above]at(4,1){功率计};}
\end{example}
\begin{solution}
先移去图\ref{fig:o4identifyq}的波片，转检偏器：周期变化且能消光为完全线偏振；变化但不能消光可能为完全椭圆偏振，也可能为部分偏振；不变化可能是圆偏振或自然光，不能立即下结论。
\Diagram{fig:o4identifya}{波片把相位信息转成可由线检偏器测出的强度差}{\draw[curve](0,0)rectangle(2.2,1.2);\node at(1.1,.6){相位差};\draw[->](2.3,.6)--(3.3,.6);\draw[curve](3.5,0)rectangle(6,1.2);\node at(4.75,.6){旋转后的强度差};}
加入四分之一波片，转片和检偏器：圆偏振可变成线偏振而消光，自然光经过理想无损波片仍不随检偏角变化。完全椭圆偏振可把波片轴对准椭圆主轴，抵消$\pm\pi/2$相差而变线。图\ref{fig:o4identifya}说明波片提供的是原来缺少的相位信息。一般部分偏振不能靠“是否消光”完整确定，应测$0,90,45,135$度线偏振强度，以及两种圆偏振分析强度，构成$S_0,S_1,S_2,S_3$。若入射光含非零自然成分，仅转动这块无损波片和随后的一片检偏器，不能使总透射光完全消失。
\end{solution}

\section{界面边界条件与菲涅耳公式}\label{th:o4fresnel}
取界面$z=0$，介质1在$z<0$、介质2在$z>0$，无自由面电荷和面电流。入射面为$xz$面，入射、反射、透射波矢单位方向依次为$(\sin i,0,\cos i)$、$(\sin i,0,-\cos i)$、$(\sin r,0,\cos r)$。沿界面相位匹配给$n_1\sin i=n_2\sin r$，反射角等于入射角。

以下取无损非磁性介质。$s$偏振统一取电场单位矢量$\hat{\vct y}$；每条光线的$p$单位矢量取$\hat{\vct s}\times\hat{\vct k}$，使$\hat{\vct k}\times\hat{\vct p}=\hat{\vct s}$。于是入射$p$电场为$(\cos i,0,-\sin i)$，反射为$(-\cos i,0,-\sin i)$，两者在正入射极限的切向基方向相反。

\subsection{$s$偏振}
切向电场连续给$1+r_s=t_s$。磁场$\vct H=(n/\mu_0c)\hat{\vct k}\times\vct E$，切向磁场连续给$n_1\cos i(1-r_s)=n_2\cos r\,t_s$。消元得
\[
 r_s=\frac{n_1\cos i-n_2\cos r}{n_1\cos i+n_2\cos r},\quad
 t_s=\frac{2n_1\cos i}{n_1\cos i+n_2\cos r}.
\]
\subsection{$p$偏振}
在上述局部基约定下，切向电场给$\cos i(1-r_p)=\cos r\,t_p$，磁场给$n_1(1+r_p)=n_2t_p$。由后式取$t_p=n_1(1+r_p)/n_2$代入前式，整理为
\[
 n_2\cos i-n_1\cos r=(n_2\cos i+n_1\cos r)r_p,
\]
再回代求$t_p$，得
\[
 r_p=\frac{n_2\cos i-n_1\cos r}{n_2\cos i+n_1\cos r},\quad
 t_p=\frac{2n_1\cos i}{n_2\cos i+n_1\cos r}.
\]
正入射时$r_s=-r_p$，只是反射$p$基反向的结果；用同一固定切向电场比较，两物理反射系数相同。若$\mu_r\ne1$，磁场系数$n$应改用相应波导纳，不能把非磁性公式直接照搬。

\section{能流、互易与布儒斯特角}\label{th:o4energy}
反射率是反射法向能流与入射法向能流之比，故$R=|r|^2$。透射率为
\[
 T=\frac{n_2\cos r}{n_1\cos i}|t|^2.
\]
无损界面在存在传播透射波时$R+T=1$。若只取沿各自光线定义的强度比$I_t/I_i=(n_2/n_1)|t|^2$，斜入射时它与$R$之和一般不为1，因为比较的截面面积不同，不能据此宣称斜入射能量不守恒。

\begin{example}\label{ex:o4stokes}
对$s$偏振，正向界面振幅系数为$r,t$，反向入射为$r',t'$。证明$r'=-r$、$tt'=1-r^2$，并解释斯托克斯关系中各束光的相位方向。
\Diagram{fig:o4stokesq}{同一界面的正反向散射}{\draw[curve](-2,0)--(2,0);\draw[->](-1.5,-1.5)--(0,0);\draw[->](0,0)--(1.5,-1.5)node[right]{$r$};\draw[->](0,0)--(1.1,1.5)node[right]{$t$};\node[left]at(-2,-.6){$n_1$};\node[left]at(-2,.6){$n_2$};}
\end{example}
\begin{solution}
图\ref{fig:o4stokesq}取$s$偏振的共同切向电场方向作为振幅基底，令$a=n_1\cos i$、$b=n_2\cos r$。电场、磁场切向连续分别给$1+r=t$、$a(1-r)=bt$。将首式代入次式，解得
\[
 r=\frac{a-b}{a+b},\qquad t=\frac{2a}{a+b}.
\]
反向入射时交换$a,b$，因此$r'=(b-a)/(a+b)=-r$、$t'=2b/(a+b)$。
\Diagram{fig:o4stokesa}{时间反演的两束入射共同恢复原来的入射光}{\draw[curve](-2,0)--(2,0);\draw[->](-1.5,-1.5)--(0,0)node[midway,left]{$r$};\draw[->](1.1,1.5)--(0,0)node[midway,right]{$t$};\draw[->](0,0)--(1.5,-1.5)node[midway,right]{$1$};\node[above]at(-1.6,1.5){另一输出相消};}
图\ref{fig:o4stokesa}的两束反向入射按原反、透射振幅$r,t$取值。在恢复原入射光的输出端，合振幅为$r\,r+t\,t'$；另一端为$r\,t+t\,r'$。代入上述系数，
\[
 r^2+tt'=\frac{(a-b)^2+4ab}{(a+b)^2}=1,
 \qquad rt+tr'=t(r+r')=0.
\]
所以原波被恢复，另一输出相消。$r'=-r$表示反向反射相位相差$\pi$，不是反射率变成负数。能量透射率应为$T=(b/a)|t|^2$，由此$R+T=1$；不能直接令$T=t^2$。这里的时间反演图解采用无损传播及实系数的两侧入射端口。全反射的倏逝支、复系数或吸收介质须连同共轭、通道与能流条件重新处理。
\end{solution}

\subsection{布儒斯特角}\label{th:o4brewster}
$p$反射为零需$n_2\cos i=n_1\cos r$，再与$n_1\sin i=n_2\sin r$联立，非平凡折射率差时得到$i+r=\pi/2$，所以$\tan i_B=n_2/n_1$。这时反射光只含$s$偏振。若$n_1=n_2$，所有角度都没有界面反射，不能把它误解成只有45度一个布儒斯特角。

地面或水面的近布儒斯特反射主要为平行表面的$s$偏振，所以让眼镜透过轴接近竖直可削弱这部分眩光。实际粗糙表面、金属和吸收介质不能全部用理想介质界面的单一布儒斯特角描述。

\section{全反射相位与倏逝波}\label{th:o4tir}
设$n_1>n_2$，入射角超过$i_c=\arcsin(n_2/n_1)$。切向波矢$k_x=k_0n_1\sin i$不变，介质2的法向分量满足$k_{z2}^2=k_0^2n_2^2-k_x^2<0$。介质2位于$z>0$，为使场远离界面不发散，取$k_{z2}=i\kappa$，其中
\[
 \kappa=k_0\sqrt{n_1^2\sin^2i-n_2^2},\quad
 E_t\propto\ee^{ik_xx-\kappa z-i\omega t}.
\]
振幅衰减长度为$1/\kappa$，强度衰减长度为$1/(2\kappa)$。令$q=\sqrt{n_1^2\sin^2i-n_2^2}$，在菲涅耳系数中取$n_2\cos r=iq$，得
\[
 r_s=\frac{n_1\cos i-iq}{n_1\cos i+iq},\quad
 r_p=\frac{n_2^2\cos i-in_1q}{n_2^2\cos i+in_1q}.
\]
分子分母互为共轭，模均为1，但相位依角度和偏振而不同。全反射可改变偏振态；“反射率1”不等于“相位不变”。单界面介质2中存在倏逝场，时间平均法向能流为零，但沿界面可有非零能流。

\begin{example}\label{ex:o4tir}
玻璃$n_1=1.5$向空气$n_2=1$以45度入射，波长$\lambda_0$。求$s,p$反射相位及倏逝振幅衰减长度，并判断能否把一次反射等同于四分之一波片。
\Diagram{fig:o4tirq}{超过临界角后无传播折射光，但有衰减场}{\draw[curve](-2.5,0)--(2.5,0);\draw[guide](0,-2)--(0,2);\draw[->](-1.8,-1.8)--(0,0);\draw[->](0,0)--(1.8,-1.8);\foreach\y/\l in{.3/2,.7/1.3,1.1/.8,1.5/.5}{\draw[curve](-\l,\y)--(\l,\y);}\node[right]at(2,.7){倏逝场};}
\end{example}
\begin{solution}
图\ref{fig:o4tirq}中$q=\sqrt{2.25/2-1}=1/(2\sqrt2)$，$n_1\cos i=3/(2\sqrt2)$，所以$\phi_s=-2\arctan(1/3)\simeq-36.87^\circ$。$p$系数相位为$\phi_p=-2\arctan[n_1q/(n_2^2\cos i)]=-2\arctan(3/4)\simeq-73.74^\circ$。
\Diagram{fig:o4tira}{两偏振的单位模反射系数具有不同相位}{\draw[axis](-.3,0)--(2.6,0)node[right]{实部};\draw[axis](0,-2.5)--(0,.5)node[above]{虚部};\draw[curve](2,0)arc(0:-90:2);\draw[->,thick](0,0)--(1.6,-1.2)node[right]{$r_s$};\draw[->,thick](0,0)--(.56,-1.92)node[below]{$r_p$};}
图\ref{fig:o4tira}的相对相差为$-36.87^\circ$，不是$-90^\circ$，所以一次反射不等于四分之一波片。衰减常数$\kappa=2\pi q/\lambda_0=\pi/(\sqrt2\lambda_0)$，振幅长度$\ell=\sqrt2\lambda_0/\pi$。在给定基向量约定下比较两分量相位，还需同步追踪反射$p$基与实验室固定横向轴的关系。
\end{solution}

\begin{example}\label{ex:o4frustrated}
两块相同高折射率玻璃间隔空气层厚度$d$，入射角在单个玻璃—空气界面满足全反射。说明为何第二块玻璃足够靠近时可有透射，并给出厚间隙极限下的主要衰减尺度。
\Diagram{fig:o4frustratedq}{有限低折射率层中的倏逝耦合；间隙内不画传播光线}{
\draw[curve](-2.2,0)--(2.4,0);\draw[curve](-2.2,1.1)--(2.4,1.1);
\node[left]at(-2.2,-.6){$n_1$};\node[left]at(-2.2,.55){$n_2$};\node[left]at(-2.2,1.7){$n_1$};
\draw[->](-1.6,-1.6)--(0,0);\draw[->](0,0)--(1.6,-1.6);
\draw[->](.8,1.1)--(2.1,2.4);
\foreach\y/\l in{.23/1.1,.51/.8,.82/.55}{\draw[densely dotted](-\l,\y)--(\l,\y);}
\node[right]at(1.1,.55){倏逝场};
\draw[<->](2.8,0)--(2.8,1.1)node[midway,right]{$d$};
}
\end{example}
\begin{solution}
图\ref{fig:o4frustratedq}间隙内仍有相同实切向波矢和虚法向波矢$\pm i\kappa$。单界面半无限介质只许衰减支；有限间隙必须同时满足两端边界，场通常含$A\ee^{-\kappa z}+B\ee^{\kappa z}$。增长支在有限区间不发散，不能提前删去。
\Diagram{fig:o4frustrateda}{有限层两指数支叠加，可把能流耦合到右侧传播波}{\draw[axis](0,0)--(5.4,0)node[right]{$z$};\draw[guide](1,0)--(1,2.8);\draw[guide](4,0)--(4,2.8);\draw[curve,domain=1:4,samples=60]plot(\x,{2.2*exp(-.65*(\x-1))});\node[below]at(1,0){$0$};\node[below]at(4,0){$d$};\draw[->](4.2,.6)--(5,.6);}
先对$s$偏振完成匹配。令$k_0=2\pi/\lambda_0$、$a=n_1\cos i$、$b=\sqrt{n_1^2\sin^2i-n_2^2}$，于是$\kappa=k_0b$。把右界面场幅记为$t$，右侧只有出射波，故$E(d)=t$、$E'(d)=ik_0at$。间隙方程$E''=\kappa^2E$给
\[
 E(z)=t\left[\cosh\kappa(d-z)
 -i\frac ab\sinh\kappa(d-z)\right].
\]
它满足右界面的函数、导数两个条件。左侧入射幅归一为1，有$E(0)=1+r$、$E'(0)=ik_0a(1-r)$。代入上式，把导数条件除以$ik_0a$，再与场条件相加，得到
\[
 2=t\left[2\cosh\kappa d
 +i\left(\frac ba-\frac ab\right)\sinh\kappa d\right].
\]
左右传播介质相同，透射率为$|t|^2$。使用$\cosh^2u=1+\sinh^2u$，整理为
\[
 T_s=\left[1+\frac{(a^2+b^2)^2}{4a^2b^2}
 \sinh^2\kappa d\right]^{-1}.
\]
对$p$偏振，以磁场幅为未知量，界面连续的是$H$和$H'/n^2$。相同消元给上式的$a=\cos i/n_1$、$b=\sqrt{n_1^2\sin^2i-n_2^2}/n_2^2$，指数中的$\kappa$仍取原值。

当$d\to0$，两块同种玻璃接合，$T\to1$。当$\kappa d\gg1$，$\sinh\kappa d\simeq\ee^{\kappa d}/2$，所以
\[
 T\simeq\frac{16a^2b^2}{(a^2+b^2)^2}\ee^{-2\kappa d}.
\]
图\ref{fig:o4frustrateda}示意厚间隙中的主要衰减包络。单界面零法向平均能流的结论针对单一衰减支；有限层中两支场的交叉项可以传递实的法向平均能流，第二块玻璃因而接收到传播波。
\end{solution}
\section{偏振椭圆的代数结构}
\subsection{相位差怎样进入实轨迹}
令$X=E_x/a$、$Y=E_y/b$，则$X=\cos\omega t$、$Y=X\cos\delta+\sin\omega t\sin\delta$。将$Y-X\cos\delta$平方，等于$(1-X^2)\sin^2\delta$，移项得到$X^2+Y^2-2XY\cos\delta=\sin^2\delta$。这一步保留了轨迹几何，却把$\sin\delta$的正负平方掉，所以旋向必须另从时间参数恢复。

场强平方$E_x^2+E_y^2$可用二倍角展开为
\[
 \frac{a^2+b^2}2+
 \frac{a^2+b^2\cos2\delta}2\cos2\omega t
 +\frac{b^2\sin2\delta}2\sin2\omega t.
\]
后两项组成单一正弦，其振幅是$\tfrac12\sqrt{a^4+b^4+2a^2b^2\cos2\delta}$。因此最大、最小场强平方为$(a^2+b^2\pm\sqrt{a^4+b^4+2a^2b^2\cos2\delta})/2$，与前述主轴半径平方一致。椭圆长短半轴因而可由场模随时间变化的极大、极小值确定。

\subsection{退化情形的处理}
$\delta=0,\pi$时短半轴为零，轨迹退为直线；若$a=b$且$\delta=\pm\pi/2$，长短半轴相等，轨迹为圆，主轴方向任意。圆偏振时$\tan2\psi$的表达式成为$0/0$，正是因为圆上不存在唯一的主轴方向。

圆偏振的Jones矢量乘任意共同相位，固定位置上电场的起始角会改变，但旋转方向不变。同一椭圆旋向也不由相位$\delta$的主值名称决定，需同时说明时间指数约定、光传播方向和观察方向。所有图中的箭头均按本书$\ee^{-i\omega t}$与从$+z$端观察的约定作出。

\section{从Jones矢量到相干矩阵}\label{th:o4coherence}
对一般光束定义$2\times2$相干矩阵$C=\langle\vct e\vct e^\dagger\rangle$。它的对角元是两个正交分量的平均强度，非对角元保留它们的相关相位。对任何分析矢量$\vct a$，透射强度为$\vct a^\dagger C\vct a\geq0$，所以$C$是半正定厄米矩阵。

用本书的斯托克斯约定，
\[
 C=\begin{pmatrix}\langle|E_x|^2\rangle&\langle E_xE_y^*\rangle\\
 \langle E_yE_x^*\rangle&\langle|E_y|^2\rangle\end{pmatrix}
 =\frac12\begin{pmatrix}
 S_0+S_1&S_2-iS_3\\S_2+iS_3&S_0-S_1
 \end{pmatrix}.
\]
记$s^2=S_1^2+S_2^2+S_3^2$。将本征方程的行列式展开，
\begin{align*}
 0=\det(C-\lambda I)
 &=\left(\frac{S_0+S_1}{2}-\lambda\right)
 \left(\frac{S_0-S_1}{2}-\lambda\right)-\frac{S_2^2+S_3^2}{4}\\
 &=\lambda^2-S_0\lambda+\frac{S_0^2-s^2}{4}.
\end{align*}
二次方程给$\lambda_\pm=(S_0\pm s)/2$。半正定等价于$S_0\geq s\geq0$，故$S_1^2+S_2^2+S_3^2\leq S_0^2$。对非零光束定义$P=s/S_0$，则
\[
 P=\frac{\lambda_+-\lambda_-}{\lambda_++\lambda_-},\qquad
 P^2=1-\frac{4\det C}{(\operatorname{tr}C)^2}.
\]
完全偏振时$\lambda_-=0$、矩阵秩为1；自然光$C=(S_0/2)I$，两个本征值相同。对一般部分偏振，还可写$C=\lambda_-I+(\lambda_+-\lambda_-)\vct u_+\vct u_+^\dagger$：第一项是自然光，第二项是沿本征态$\vct u_+$的完全偏振分量。

\Diagram{fig:o4coherencematrix}{纯偏振是秩1相干矩阵，自然光平均后没有固定交叉相位}{\draw[curve](0,0)rectangle(2.6,1.2);\node at(1.3,.6){$C=\vct e\vct e^\dagger$};\draw[curve](3.7,0)rectangle(6.3,1.2);\node at(5,.6){$C=(S_0/2)I$};\node[above]at(1.3,1.3){完全偏振};\node[above]at(5,1.3){自然光};}
图\ref{fig:o4coherencematrix}也解释为何自然光不能选一个Jones矢量表示。若把两个正交确定场以固定相位直接相加，就得到某一种完全偏振态，而非自然光；自然光需要对随机相位或随机方向作统计平均。

无损波片的Jones矩阵$J$为酉矩阵，输出相干矩阵$C'=JCJ^\dagger$。本征值与迹不变，所以总强度和完整偏振度不变。波片可以在斯托克斯空间旋转偏振方向和椭圆率，却不能凭空把自然光变成完全偏振光。偏振片为投影矩阵，会丢弃一部分能量，才可以提高透射束偏振度。

\section{单轴晶体色散式的推导}
取$\vct k=(k_x,0,k_z)$、介电张量$\boldsymbol\varepsilon=\operatorname{diag}(n_o^2,n_o^2,n_e^2)$，假定非磁性。麦克斯韦方程给
\[
 \vct k\times(\vct k\times\vct E)+k_0^2\boldsymbol\varepsilon\vct E=0.
\]
若$\vct E$沿$y$，得到$k_x^2+k_z^2=k_0^2n_o^2$，即普通光。若$\vct E$在$xz$面，则
\[
 \begin{pmatrix}k_0^2n_o^2-k_z^2&k_xk_z\\
 k_xk_z&k_0^2n_e^2-k_x^2\end{pmatrix}
 \begin{pmatrix}E_x\\E_z\end{pmatrix}=0.
\]
要有非零场，行列式为零。展开、消去$k_x^2k_z^2$项，再除以$k_0^2n_o^2n_e^2$，得到$k_x^2/n_e^2+k_z^2/n_o^2=k_0^2$，就是非普通色散椭圆。

高斯定律给$k_xn_o^2E_x+k_zn_e^2E_z=0$，所以$\vct D\perp\vct k$，但$\vct E$不一定垂直。无损非色散介质的群速度方向为等频曲面的法线，通常也就是能流方向；非普通光因此发生走离。沿光轴$k_x=0$时，两支都给折射率$n_o$，横向任意线偏振可作为等价基底，不会发生通常意义的双折射分离。

\section{菲涅耳能量关系的逐项证明}
对$s$偏振，令$a=n_1\cos i$、$b=n_2\cos r$，$r_s=(a-b)/(a+b)$、$t_s=2a/(a+b)$。法向透射率$T_s=(b/a)t_s^2=4ab/(a+b)^2$，故$R_s+T_s=[(a-b)^2+4ab]/(a+b)^2=1$。

对$p$偏振，令分母$D=n_2\cos i+n_1\cos r$，则
\[
 1-r_p^2=\frac{4n_1n_2\cos i\cos r}{D^2},
\]
而法向透射率
\[
 T_p=\frac{n_2\cos r}{n_1\cos i}
 \frac{4n_1^2\cos^2i}{D^2}
 =\frac{4n_1n_2\cos i\cos r}{D^2}.
\]
因此两偏振都严格满足$R+T=1$，无需限定正入射。若取沿透射方向的强度比$I_t/I_i$而未乘$\cos r/\cos i$，相当于比较两种不同截面，当然不能直接与$R$相加作界面能量收支。

\Diagram{fig:o4fluxarea}{同一界面面积$A$对入射和透射光束的投影面积分别为$A\cos i,A\cos r$}{\draw[curve](-2.5,0)--(2.5,0);\draw[->](-1.7,-1.7)--(0,0)--(1,1.8);\draw[guide](0,-2)--(0,2.3);\node[left]at(-1,-1.2){$I_i\cos i$};\node[right]at(.8,1.1){$I_t\cos r$};\node[below]at(1.8,-.2){界面$A$};}
图\ref{fig:o4fluxarea}以固定界面面积作为收支基准。对于吸收介质，透射角和波阻抗可能变复，能流方向及损耗需要用复坡印廷矢量重新求出；本讲的实角无损公式不能直接把$n$换成复数后仍以同一实余弦解释。

\section{全反射的解析支与相位连续性}
当入射角从临界角以下增至以上，透射法向波矢由正实变为正虚。正虚支使$z>0$侧的$\ee^{ik_{z2}z}$衰减；负虚支则使场在无限远增长，与无穷远无输入的边界条件不符。若低折射率介质改在另一半空间，应重新按远离界面的衰减条件选取虚部符号。

反射系数在临界角以上沿复平面的单位圆变化。其模恒1，能量全部反回，但$s,p$相位差不恒定。若入射为两分量都有的线偏振，一次全反射一般变椭圆；两次适当角度的反射可以累积为四分之一波延迟。多次反射还要追踪各入射面如何转动，局部$s,p$基底之间若有几何旋转，就不能只把标量相位差相加。

倏逝场中$\vct E,\vct H$某些分量相差90度，平均法向坡印廷流为零，同时有非零瞬时能量往返交换。有限空气隙含两支指数解，交叉项能产生实的法向平均流，因此受抑全反射不违反单界面结果；它改变的是边界问题本身。有限间隙的透射率应由两端边界共同求出。

\subsection{部分偏振的相关矩阵与可实现条件}
第\ref{th:o4coherence}节已由相干矩阵的本征方程证明$S_1^2+S_2^2+S_3^2\leq S_0^2$。这个不等式也给实验数据的可实现条件：固定非零总强度$S_0$后，归一化矢量$(S_1,S_2,S_3)/S_0$位于单位球内，距球心的长度就是$P$。球面对应完全偏振，球心对应自然光。若重建值落在球外，应回查分析器标定、背景扣除或光源漂移；它不能作为某种“超完全偏振”状态。$S_0=0$时无光，不定义偏振度。
\Diagram{fig:o4stokesball}{固定总强度下的归一化斯托克斯矢量；球面对应完全偏振，球内对应部分偏振}{\draw[axis](-2,0)--(2.4,0)node[right]{$S_1/S_0$};\draw[axis](0,-1.8)--(0,2)node[above]{$S_3/S_0$};\draw[curve](0,0)circle(1.5);\draw[guide](0,0)ellipse(1.5 and .38);\draw[->,thick](0,0)--(.8,.7)node[right]{$P<1$};\fill(-.9,1.2)circle(1.5pt)node[left]{$P=1$};\node[below]at(.2,-1.9){$S_2$轴的投影略去};}
图\ref{fig:o4stokesball}把总强度固定后，偏振态表示为球内或球面点。只有线检偏器时能测得$S_0,S_1,S_2$，无法确定$S_3$；必须加入波片形成圆分析器，才能恢复全部相关信息。线检偏器扫描的可见度不等于总体偏振度；圆偏振的线分析强度与方向无关，但$P=1$。

对多束互不相干光，交叉相关项平均为零，故$C_{\rm tot}=\sum_jC_j$，四个斯托克斯量也分别相加。偏振度含矢量和的模，应在相加后计算：相反偏振方向可以抵消，即使每束的$P$都为1，总光也可能为自然光。无损波片与理想检偏器对$C$的作用已在第\ref{th:o4coherence}节区分，不能用偏振度的算术平均代替这一矩阵运算。

\section{从测量强度还原偏振态}
\begin{example}\label{ex:o4stokesmeasure}
同一未知光束经六种理想分析器测得强度：水平6、竖直4、45度7、135度3、正圆偏振分析6、负圆偏振分析4。圆分析正态定义为Jones矢量$(1,i)^{\mathsf T}/\sqrt2$。求斯托克斯量、完整偏振度以及线偏振扫描的可见度。
\Diagram{fig:o4measureq}{同一光束依次采用六种偏振分析设置}{\draw[curve](0,0)rectangle(2.2,1);\node at(1.1,.5){未知光};\draw[->](2.4,.5)--(3.3,.5);\node[right]at(3.5,.5){$H,V,D,A,+,-$};}
\end{example}
\begin{solution}
图\ref{fig:o4measureq}三组互补强度之和都为10，彼此相容。故$S_0=10$，$S_1=I_H-I_V=2$，$S_2=I_D-I_A=4$，$S_3=I_+-I_-=2$。完整偏振度$\mathcal P=\sqrt{2^2+4^2+2^2}/10=\sqrt6/5\simeq.4899$。
\Diagram{fig:o4measurea}{完整偏振矢量的长度大于其线偏振平面投影}{\draw[axis](0,0)--(3.4,0)node[right]{$S_1$};\draw[axis](0,0)--(0,3)node[above]{$S_2$};\draw[axis](0,0)--(-1.5,-1.2)node[left]{$S_3$};\draw[->,thick](0,0)--(1.8,2.2);\draw[guide](1.8,0)--(1.8,2.2);\node[right]at(2.1,1){三维分量示意};}
旋转线检偏器强度$I(\theta)=\tfrac12[S_0+S_1\cos2\theta+S_2\sin2\theta]$，调制振幅$\sqrt{20}/2$，所以可见度$\sqrt{20}/10=1/\sqrt5\simeq.4472$。图\ref{fig:o4measurea}说明单线检偏只看$S_1,S_2$投影，忽略$S_3$会低估完整偏振度。此组数据含不可定相的自然成分，不能指定一个唯一Jones矢量代表全部光。
\end{solution}

\begin{example}\label{ex:o4quarteranalysis}
证明快轴45度、慢轴135度的理想四分之一波片接水平检偏器，可测得正圆分析强度$S_0/2+S_3/2$。说明转片或交换快慢轴怎样测另一旋向。
\Diagram{fig:o4analysisq}{四分之一波片把圆基分量转成正交线基分量}{\draw[axis](-1,0)--(5.5,0);\draw[curve](0,-1)--(0,1);\draw[curve](3,-1)--(3,1);\node[above]at(0,1){$Q(45^\circ)$};\node[above]at(3,1){水平检偏};}
\end{example}
\begin{solution}
图\ref{fig:o4analysisq}中，快轴沿45度、慢轴沿135度，令快轴相位为零、慢轴增相$\pi/2$，波片矩阵为
\[
 Q_{45}=\frac12\begin{pmatrix}1+i&1-i\\1-i&1+i\end{pmatrix}.
\]
后接水平检偏器只保留第一行，因此$E'_H=[(1+i)E_x+(1-i)E_y]/2=(1+i)(E_x-iE_y)/2$。公共因子模平方为$1/2$，故
\[
 I_H'=\tfrac12\langle|E_x-iE_y|^2\rangle.
\]
\Diagram{fig:o4analysisa}{两圆偏振基在该波片后分别成为水平、竖直线偏振}{\node at(0,1){$(1,i)$};\node at(0,-.3){$(1,-i)$};\draw[->](1,1)--(2.3,1);\draw[->](1,-.3)--(2.3,-.3);\node[right]at(2.5,1){水平};\node[right]at(2.5,-.3){竖直};}
按本讲$S_3=-2\operatorname{Im}\langle E_xE_y^*\rangle$的约定，展开模平方，
\begin{align*}
 \langle|E_x-iE_y|^2\rangle
 &=\langle|E_x|^2+|E_y|^2+iE_xE_y^*-iE_yE_x^*\rangle\\
 &=S_0-2\operatorname{Im}\langle E_xE_y^*\rangle=S_0+S_3.
\end{align*}
因此测得$(S_0+S_3)/2$。图\ref{fig:o4analysisa}可用圆基$(1,i)/\sqrt2$、$(1,-i)/\sqrt2$逐个代入：前者经过波片成为水平线偏振，后者成为竖直线偏振，水平片分别全透与阻断。

快轴改为135度，或交换快慢轴，可测得$(S_0-S_3)/2$。波片本身是无损相位元件，输出总强度仍为$S_0$；圆偏振的选择发生在随后的检偏器中。描述“正圆”时必须同时保留时间因子与$S_3$符号约定，避免只用旋向名称造成歧义。
\end{solution}

\begin{example}\label{ex:o4ellipseinverse}
完全偏振光的长短半轴为$A,B$，$A\geq B\geq0$，长轴与$x$成$\psi$，从$+x$向$+y$的旋向取正。求斯托克斯量与椭圆几何的关系，说明哪一个量携带旋向。
\Diagram{fig:o4inverseellipseq}{椭圆的尺寸、倾角和旋向是独立几何信息}{\draw[axis](-2.2,0)--(2.5,0)node[right]{$x$};\draw[axis](0,-1.5)--(0,1.8)node[above]{$y$};\draw[curve,rotate=25](0,0)ellipse(1.8 and .8);\draw[guide](-1.8,-.84)--(1.8,.84);\node at(1,.2){$\psi$};}
\end{example}
\begin{solution}
在图\ref{fig:o4inverseellipseq}的主轴基中，正旋Jones矢量可取$(A,iB)$，所以$S_0=A^2+B^2$、$S_1=A^2-B^2$、$S_2=0$、$S_3=2AB$。将坐标旋转$\psi$，线偏振分量在斯托克斯平面旋转双角，得到
\[
 S_1=(A^2-B^2)\cos2\psi,\quad
 S_2=(A^2-B^2)\sin2\psi,\quad S_3=\pm2AB.
\]
\Diagram{fig:o4inverseellipsea}{相同轨迹反向走动，只改变$S_3$符号}{\begin{scope}[xshift=-2cm]\draw[curve](0,0)ellipse(1.2 and .65);\draw[->](1.2,0)arc(0:65:1.2 and .65);\node[below]at(0,-1){$S_3>0$};\end{scope}\begin{scope}[xshift=2cm]\draw[curve](0,0)ellipse(1.2 and .65);\draw[->](1.2,0)arc(0:-65:1.2 and .65);\node[below]at(0,-1){$S_3<0$};\end{scope}}
图\ref{fig:o4inverseellipsea}相同长短轴和倾角下，旋向只由$S_3$的符号区分。验证$S_1^2+S_2^2+S_3^2=(A^2+B^2)^2$，正是完全偏振条件；圆偏振$A=B$时$S_1=S_2=0$，长轴方向不再可定义，此时$\tan2\psi$不再确定唯一倾角。
\end{solution}

\section{晶体中的波矢与能流}
\begin{example}\label{ex:o4walkoff}
单轴晶体$n_o=1.60,n_e=1.50$，非普通光波矢与光轴成45度。求有效相折射率、能流方向和走离角。忽略频率色散。
\Diagram{fig:o4walkoffq}{波矢与光轴成$45^\circ$，能流偏向远离光轴的一侧}{
\draw[axis](0,0)--(3.8,0)node[right]{$x$};\draw[axis](0,0)--(0,3.4)node[above]{光轴$z$};
\draw[axis](0,0)--(2.15,2.15)node[above left]{$\vct k$};
\draw[axis](0,0)--({2.15*2.56/2.25},2.15)node[right]{$\vct S$};
\draw(0,.85)arc(90:45:.85);\node at(.6,1.12){$\vartheta$};
\draw(45:2.55)arc(45:41.308:2.55);
\draw[thin](43.154:2.55)--(3.1,1.1)node[right]{$\rho$};}
\end{example}
\begin{solution}
图\ref{fig:o4walkoffq}取光轴为$z$、波矢位于$x,z$平面。非普通波色散为$k_x^2/n_e^2+k_z^2/n_o^2=(\omega/c)^2$。令$k=n_{\rm eff}\omega/c$、$k_x=k\sin\vartheta$、$k_z=k\cos\vartheta$，得
\[
 \frac1{n_{\rm eff}^2}=\frac{\sin^2\vartheta}{n_e^2}
 +\frac{\cos^2\vartheta}{n_o^2}.
\]
取$\vartheta=45^\circ$、$n_o=1.60,n_e=1.50$，得到$n_{\rm eff}=1.5476$。这决定沿$\vct k$的相位速度$c/n_{\rm eff}$。
\Diagram{fig:o4walkoffa}{等频椭圆的外法线给出能流方向，虚线延长波矢方向}{
\draw[axis](0,0)--(3.8,0)node[right]{$k_x$};\draw[axis](0,0)--(0,3.5)node[above]{$k_z$（光轴）};
\draw[curve](2.4,0)arc(0:90:2.4 and 2.56);
\pgfmathsetmacro{\p}{1/sqrt(1/2.4^2+1/2.56^2)}
\fill(\p,\p)circle(1.5pt)node[below right]{$P$};
\draw[axis](0,0)--(\p,\p)node[midway,below right]{$\vct k$};
\draw[guide](\p,\p)--({\p+1.05},{\p+1.05});
\draw[axis](\p,\p)--({\p+.95*2.56/2.25},{\p+.95})node[right]{$\vct S$};
\draw(0,.7)arc(90:45:.7);\node at(.4,.94){$45^\circ$};
\node[align=left]at(3.4,.65){$\vartheta_S=48.69^\circ$\\$\rho=3.69^\circ$};}
无色散时，群速度为$\vct v_g=\nabla_{\vct k}\omega$。固定$k_z$，对色散式的两边关于$k_x$求导；再固定$k_x$，关于$k_z$求导，分别有
\[
 \frac{2k_x}{n_e^2}=\frac{2\omega}{c^2}\frac{\partial\omega}{\partial k_x},
 \qquad
 \frac{2k_z}{n_o^2}=\frac{2\omega}{c^2}\frac{\partial\omega}{\partial k_z}.
\]
解出两个偏导数，得到
\[
 v_{gx}=\frac{c^2k_x}{\omega n_e^2},\qquad
 v_{gz}=\frac{c^2k_z}{\omega n_o^2},
\]
因此图\ref{fig:o4walkoffa}的能流方向满足
\[
 \tan\vartheta_S=\frac{v_{gx}}{v_{gz}}
 =\frac{n_o^2}{n_e^2}\tan\vartheta
 =\frac{2.56}{2.25},\quad \vartheta_S=48.69^\circ.
\]
走离角$\rho=\vartheta_S-\vartheta=3.69^\circ$，方向远离光轴。图中的箭头沿等频椭圆在$P$点的外法线，虚线则沿波矢方向延长，两者夹角为$\rho$。若入射点取原点、出射端面为$z=h$，光束中心横坐标为$h\tan\vartheta_S$；相位法线方向对应$h\tan\vartheta$，二者横向差为$h(\tan\vartheta_S-\tan\vartheta)$。因此光斑位置与波面相位不能用同一个折射角描述。
\end{solution}

\begin{example}\label{ex:o4reflectionpolar}
自然光射向无损介质界面。给定菲涅耳反射率$R_s,R_p$，求反射光的线偏振度和总强度，讨论正入射、布儒斯特角和掠入射三个极限。
\Diagram{fig:o4reflectionq}{自然光的两正交份额分别被不同反射率筛选}{\node at(0,.6){$I_0/2:s$};\node at(0,-.6){$I_0/2:p$};\draw[->](1.2,.6)--(2.6,.6)node[midway,above]{$R_s$};\draw[->](1.2,-.6)--(2.6,-.6)node[midway,below]{$R_p$};\node[right]at(2.8,0){反射束};}
\end{example}
\begin{solution}
图\ref{fig:o4reflectionq}反射强度$I_r=I_0(R_s+R_p)/2$，两正交分量不相干，所以线偏振度$\mathcal P=|R_s-R_p|/(R_s+R_p)$。若$R_s>R_p$，主要偏振沿$s$。
\Diagram{fig:o4reflectiona}{布儒斯特角处纯$s$，两端极限偏振度回落}{\draw[axis](0,0)--(5.2,0)node[right]{$i$};\draw[axis](0,0)--(0,2.8)node[above]{$\mathcal P$};\draw[curve](0,0)..controls(1.7,.2)and(2.6,2.3)..(3.1,2.3)..controls(3.8,2.3)and(4.4,.7)..(4.8,0);\node[below]at(3.1,0){$i_B$};\node[below]at(4.8,0){$\pi/2$};}
正入射$R_s=R_p$，偏振度0；布儒斯特角$R_p=0$，偏振度1；掠入射两者都趋1，偏振度又趋0，图\ref{fig:o4reflectiona}表示这种趋势。若总反射强度为零，例如完全无折射率差的界面，则比值为$0/0$，无反射光可谈偏振，不能随意给偏振度0或1。
\end{solution}

\section{全反射延迟与边界附近能流}
\begin{example}\label{ex:o4rhomb}
玻璃向空气全反射两次，两次入射面共面且偏振基连续对应。求实现总相对延迟$\pi/2$所需的单次相位差条件；对$n=1.5$推导可用入射角方程。
\Diagram{fig:o4rhombq}{两次全反射的相位差相加}{\draw[curve](0,0)--(4,0)--(5,2)--(1,2)--cycle;\draw[->](.4,.8)--(1.2,0)--(3.2,2)--(4.4,.8);\node[below]at(1.2,0){$R_1$};\node[above]at(3.2,2){$R_2$};}
\end{example}
\begin{solution}
设$m=n_2/n_1=2/3$，$g=\sqrt{\sin^2i-m^2}$。单次相位$\phi_s=-2\arctan(g/\cos i)$、$\phi_p=-2\arctan(g/(m^2\cos i))$。图\ref{fig:o4rhombq}两次基底对应相同，需$|\phi_p-\phi_s|=\pi/4$。
\Diagram{fig:o4rhomba}{两相量相位差每次增加45度，两次累积90度}{\draw[axis](0,0)--(2.8,0)node[right]{$s$参考};\draw[->](0,0)--(1.6,-1.6)node[right]{一次};\draw[->,thick](0,0)--(0,-2.3)node[below]{两次};\draw[guide](2.3,0)arc(0:-90:2.3);}
令$x=g/\cos i>0$，用反正切差公式，条件为
\[
 \frac{x(1/m^2-1)}{1+x^2/m^2}=\tan(\pi/8).
\]
记$u=\tan(\pi/8)$，化为二次式$ux^2-(1-m^2)x+m^2u=0$，两根为
\[
 x_\pm=\frac{1-m^2\pm\sqrt{(1-m^2)^2-4m^2u^2}}{2u}.
\]
代入$m=2/3$，得$x_-=0.59795$、$x_+=0.74328$。又由$x^2=(\sin^2i-m^2)/(1-\sin^2i)$，整理为$(1+x^2)\sin^2i=x^2+m^2$，故
\[
 i_\pm=\arcsin\sqrt{\frac{x_\pm^2+m^2}{1+x_\pm^2}}.
\]
两解约为50.23度和53.26度，均大于临界角41.81度。图\ref{fig:o4rhomba}画的是相对延迟的累积；实际棱镜的入射、出射面若还引入振幅失配或相位差，也应计入总Jones矩阵。
\end{solution}

\begin{example}\label{ex:o4evanescentflux}
$s$偏振倏逝波在$z>0$介质中为$E_y=E_t\ee^{ik_xx-\kappa z-i\omega t}$，$k_x,\kappa$均为正实数。直接由麦克斯韦方程求平均坡印廷矢量，验证法向能流为零而沿界面不为零。
\Diagram{fig:o4evanq}{倏逝振幅沿法向衰减，相位沿界面前进}{\draw[axis](-.3,0)--(4.5,0)node[right]{$x$};\draw[axis](0,0)--(0,3.2)node[above]{$z$};\foreach\y/\l in{.3/3.8,.8/2.6,1.3/1.8,1.8/1.2}{\draw[curve](.2,\y)--(\l,\y);}\node[right]at(2,2.5){$\ee^{-\kappa z}$};}
\end{example}
\begin{solution}
取复波矢$\vct k=(k_x,0,i\kappa)$，图\ref{fig:o4evanq}的$\vct H=(\vct k\times\vct E)/(\mu\omega)$，所以$H_x=-i\kappa E_y/(\mu\omega)$、$H_z=k_xE_y/(\mu\omega)$。复坡印廷乘积为
\[
 \vct E\times\vct H^*
 =\frac{|E_y|^2}{\mu\omega}(k_x\hat{\vct x}-i\kappa\hat{\vct z}).
\]
\Diagram{fig:o4evana}{时间平均实能流平行界面，法向虚部表示往返储能交换}{\draw[curve](-.5,0)--(4.5,0);\foreach\y in{.5,1,1.5}{\draw[axis](.3,\y)--({3.5*exp(-.5*\y)},\y);}\node[right]at(3,1.7){$\langle S_z\rangle=0$};}
取一半实部，得$\langle S_x\rangle=k_x|E_t|^2\ee^{-2\kappa z}/(2\mu\omega)$，$\langle S_z\rangle=0$，如图\ref{fig:o4evana}。场在低折射率侧非零与无净法向透射并不矛盾。若存在第二界面产生另一指数支，交叉项可给实法向能流，必须重新计算而不能保留单支结果。
\end{solution}

\section{偏振元件的损耗与色散}
\subsection{相位延迟和选择性衰减改变不同的量}
以下继续以$C$表示相干矩阵、$J$表示Jones矩阵，元件作用为$C'=JCJ^\dagger$。第\ref{th:o4coherence}节已证明偏振度由$C$的迹与行列式确定；无损波片保持这两个量。二向色元件则对两个本征方向作不同衰减，改变$C$的本征值比例，因此可能以损失通光量为代价提高偏振度。

设主轴方向上的能量透过率为$t_x,t_y$，无额外相位延迟，则$D=\operatorname{diag}(\sqrt{t_x},\sqrt{t_y})$。若自然光入射，$C=I_0\mathbb I/2$，则$C_D=DCD^\dagger=I_0\operatorname{diag}(t_x,t_y)/2$。取其迹与两个本征值之差，得到出射强度和偏振度
\[
 I=\frac{I_0(t_x+t_y)}2,\qquad
 P=\frac{|t_x-t_y|}{t_x+t_y}.
\]
以上偏振度公式要求$I_0>0$且$t_x+t_y>0$；若两方向都完全遮光，出射强度为零，偏振度无定义。随后接一块无损波片，强度和偏振度不变，偏振成分却可变为椭圆偏振。两元件的先后顺序须按$QD$或$DQ$计算，它们一般不对易。

\begin{example}部分圆偏振光的产生与检偏
自然光强度为$I_0$，先通过主轴能量透过率$t_x=1,t_y=1/4$的二向色元件，再通过快轴与$x$成$45^\circ$的理想四分之一波片。按本讲Jones矩阵约定，删去共同相位后，$x$线偏振变为$(1,-\mathrm i)^{\mathsf T}/\sqrt2$。求出射强度、偏振度及旋转线检偏器的强度范围。比较交换两元件顺序的结果。
\ExpansionDiagram{fig:dep-o4-loss-q}{先选择性衰减、后相位延迟的偏振光路}{
\draw[->](0,0)--(7,0);\draw[fill=white](1,-.7)rectangle(2,.7)node[midway]{$D$};\draw[fill=white](3,-.7)rectangle(4,.7)node[midway]{$Q$};\draw[fill=white](5,-.7)rectangle(6,.7)node[midway]{$A_\theta$};\node[above]at(1.5,.9){$t_x\ne t_y$};\node[above]at(3.5,.9){$45^\circ$};}
\end{example}
\begin{solution}
由图\ref{fig:dep-o4-loss-q}先经$D$，得$C_D=I_0\operatorname{diag}(1,1/4)/2$。按第\ref{th:o4coherence}节的本征值分解，将它分成自然光与完全偏振光：
\[
 C_D=\frac{I_0}{8}\mathbb I+\frac{3I_0}{8}
 \begin{pmatrix}1&0\\0&0\end{pmatrix}.
\]
第一项是总强度$I_0/4$的自然光，第二项是强度$3I_0/8$的$x$线偏振。通过$Q$后，第一项不变，第二项按题设变为圆偏振，所以
\[
 I_{\rm out}=\frac{5I_0}{8},\qquad
 P=\frac{3I_0/8}{5I_0/8}=\frac35.
\]
\ExpansionDiagram{fig:dep-o4-loss-a}{相干矩阵分解中的自然成分与完全偏振成分}{
\draw(0,0)rectangle(2,1.2);\node at(1,.8){自然光};\node at(1,.35){$I_0/4$};\node at(2.5,.6){$+$};\draw(3,0)rectangle(5,1.2);\node at(4,.8){圆偏振};\node at(4,.35){$3I_0/8$};\node at(5.5,.6){$\longrightarrow$};\node at(7,.6){$P=3/5$};}
图\ref{fig:dep-o4-loss-a}中两种成分通过任意方向的理想线检偏器，都各透过一半。因此$I(\theta)=5I_0/16$恒定，虽然$P=3/5$并不为零。线检偏扫描没有起伏不能单独证明自然光。

若交换为先$Q$后$D$，自然光先经无损波片仍是自然光，最后输出恰为$C_D$。其强度和$P$仍相同，但偏振成分是线偏振，检偏强度为
\[
 I(\theta)=\frac{I_0}{2}\cos^2\theta+\frac{I_0}{8}\sin^2\theta,
\]
最大$I_0/2$、最小$I_0/8$。顺序变化改变了相干矩阵的方向结构，而本题两个顺序恰好有相同特征值；这不意味着任意入射态都保持相同输出强度。
\end{solution}

\subsection{复合波片的一阶消色差条件}
单片延迟量$\delta=2\pi h\Delta n(\lambda)/\lambda$同时受到折射率色散和$1/\lambda$因子的影响。即使忽略材料色散，延迟量也不与波长无关。两材料波片主轴平行或互相垂直时，延迟量可以代数相加：
\[
 \delta(\lambda)=\frac{2\pi}{\lambda}
 [h_1\Delta n_1(\lambda)+s h_2\Delta n_2(\lambda)],\qquad s=\pm1.
\]
这里$s=-1$表示把第二片的快慢轴对调。若两片轴任意夹角，应乘Jones矩阵，不能再直接加延迟量。

令$L(\lambda)=h_1\Delta n_1+s h_2\Delta n_2$。厚度固定时，对$\delta=2\pi L/\lambda$用商法则求导，得
\[
 \delta'(\lambda)=2\pi\left[\frac{L'(\lambda)}{\lambda}-\frac{L(\lambda)}{\lambda^2}\right].
\]
在目标波长$\lambda_0$处设计零级四分之一波片，并令$\delta'(\lambda_0)=0$，需要
\[
 L(\lambda_0)=\frac{\lambda_0}{4},\qquad
 L'(\lambda_0)=\frac{L(\lambda_0)}{\lambda_0}=\frac14.
\]
第二式中的非零$L'$补偿了$1/\lambda$因子的变化。厚度误差、温度膨胀和材料高阶色散仍须计入延迟量；一阶消色差条件没有消除这些影响。

\begin{example}两材料波片的厚度设计
两材料在$\lambda_0=500\,\mathrm{nm}$附近满足$\Delta n_1=0.010-10^{-5}(\lambda-500)$、$\Delta n_2=0.020-5\times10^{-5}(\lambda-500)$，式中波长数值以纳米为单位。将两片主轴交叉，求使总延迟为$\pi/2$且一阶消色差的正厚度。
\ExpansionDiagram{fig:dep-o4-achrom-q}{第二片主轴对调，使两片延迟量相减}{
\draw(0,0)rectangle(2,1.8);\draw[<->](.3,.9)--(1.7,.9)node[midway,above]{快轴1};\draw(3,0)rectangle(5,1.8);\draw[<->](4,.2)--(4,1.6);\node[right]at(4,1){快轴2};\node[below]at(1,0){$h_1$};\node[below]at(4,0){$h_2$};}
\end{example}
\begin{solution}
按图\ref{fig:dep-o4-achrom-q}取$s=-1$。厚度也用纳米表示，中心延迟条件和导数条件分别为
\[
 0.010h_1-0.020h_2=125,
 \qquad -10^{-5}h_1+5\times10^{-5}h_2=\frac14.
\]
第一式给$h_1=12500+2h_2$，代入第二式为$-0.125+3\times10^{-5}h_2=0.25$，故
\[
 h_2=12500\,\mathrm{nm}=12.5\,\mu\mathrm m,
 \qquad h_1=37.5\,\mu\mathrm m.
\]
两片单独的光程差为375 nm和250 nm，其差125 nm给$\pi/2$；导数之差为$-0.375-(-0.625)=0.25$，同时满足要求。
\ExpansionDiagram{fig:dep-o4-achrom-a}{两片光程差的差值随波长保持所需的一阶斜率}{
\draw[->](0,0)--(6,0)node[right]{$\lambda$};\draw[->](0,0)--(0,3.1)node[above]{光程差};\draw[thick](.5,2.3)--(5.3,1.7)node[right]{$h_1\Delta n_1$};\draw[dashed](.5,1.6)--(5.3,.6)node[right]{$h_2\Delta n_2$};\draw[<->](3,1.08)--(3,1.99)node[midway,right]{$L(\lambda)$};}
图\ref{fig:dep-o4-achrom-a}中的两条线都下降，但第二条下降更快，所以差值$L$随波长增加。本题给定的严格线性模型恰使$L=\lambda/4$，在该模型适用区间内总延迟恒定；实际材料的二阶以上色散会留下残差。这里不能凭线性模型推断无限宽带性能。
\end{solution}

\section{闭合偏振路径中的干涉相位}
\subsection{单束可以删去的共同相位，在两臂间仍可测量}
两个非零Jones矢量若只相差共同相位，其单束强度和偏振椭圆相同，与同一参考光干涉时却会产生不同的条纹位置。依次投影到几个偏振态后，即使偏振方向回到初态，也可能留下可测的路径相位。因此比较两臂时，应保留各元件作用后的完整复振幅。

归一化偏振态$|a\rangle$的理想投影元件写$P_a=|a\rangle\langle a|$。它选择该态并保留投影的复数系数。初态$|a\rangle$依次经$P_b,P_c,P_a$后，输出为
\[
 |a\rangle\langle a|c\rangle\langle c|b\rangle\langle b|a\rangle.
\]
三个内积的积包含透过幅度与路径相位。把任一态改写为$\ee^{\mathrm i\alpha}|a\rangle$并不改变投影算符，积中的相位也成对抵消。因此所求相位不是任选Jones代表元造成的假象。实际元件还可能有传播相位，干涉装置中需事先平衡或把它们明确加入模型。

\begin{example}交换偏振投影顺序，最后方向相同而条纹移动
取$|H\rangle=(1,0)^{\mathsf T}$、$|D\rangle=(1,1)^{\mathsf T}/\sqrt2$、$|R\rangle=(1,\mathrm i)^{\mathsf T}/\sqrt2$。干涉仪两臂在进入元件前均为场$E_0|H\rangle$。参考臂附加相位为$\delta$，信号臂先经$P_D$、再$P_R$、最后$P_H$。求信号透过率与干涉强度；再交换$P_D,P_R$，比较条纹。投影元件按上述算符定义，其他光程相位已补偿。
\ExpansionDiagram{fig:dep-o4-loop-q}{两条信号路径最终都回到水平偏振，参考臂保留相位基准}{
\node[draw]at(0,1){$H$};\node[draw]at(2,1){$D$};\node[draw]at(4,1){$R$};\node[draw]at(6,1){$H$};\draw[->](.4,1)--(1.6,1);\draw[->](2.4,1)--(3.6,1);\draw[->](4.4,1)--(5.6,1);\draw[->](0,-.3)--(6,-.3)node[midway,below]{参考臂$\ee^{\mathrm i\delta}H$};}
\end{example}
\begin{solution}
图\ref{fig:dep-o4-loop-q}第一步幅度系数为$\langle D|H\rangle=1/\sqrt2$，第二步为$\langle R|D\rangle=(1-\mathrm i)/2$，第三步为$\langle H|R\rangle=1/\sqrt2$。保留三个复系数相乘，得
\[
 c=\frac{1-\mathrm i}{4}=\frac1{2\sqrt2}\ee^{-\mathrm i\pi/4}.
\]
信号透过率$|c|^2=1/8$，最后偏振仍为$H$。参考臂与信号臂叠加时
\[
 I=I_0|\ee^{\mathrm i\delta}+c|^2
 =I_0\left[\frac98+\frac1{\sqrt2}\cos(\delta+\pi/4)\right],
\]
其中$I_0=|E_0|^2$采用同一强度单位。
\ExpansionDiagram{fig:dep-o4-loop-a}{两种顺序的信号系数互为共轭，参考相位扫描看到相反偏移}{
\draw[->](-.2,0)--(3.5,0)node[right]{实部};\draw[->](0,-2)--(0,2)node[above]{虚部};\draw[->,thick](0,0)--(2,-2)node[right]{$(1-\mathrm i)/4$};\draw[->,thick](0,0)--(2,2)node[right]{$(1+\mathrm i)/4$};\draw[dashed](2,-2)--(2,2);}
交换后$c'=\langle H|D\rangle\langle D|R\rangle\langle R|H\rangle=(1+\mathrm i)/4$。图\ref{fig:dep-o4-loop-a}中两信号模相同、辐角相差$\pi/2$，新强度为$I_0[9/8+\cos(\delta-\pi/4)/\sqrt2]$。两条纹相差四分之一个周期，可见度同为$4\sqrt2/9$。

两种顺序得到的信号具有相同强度和偏振，与共同参考臂叠加后才显出路径相位差。仅计算单束偏振时可以删去共同相位；计算干涉时，须另行保留它相对于参考臂的值。逐片归一化并丢掉投影复系数，会同时丢失透过率和条纹位移。
\end{solution}

\section{习题与解答}
\begin{exercise}\label{ex:o4twohalf}
两块理想半波片的快轴方向依次为$\alpha_1,\alpha_2$。求合成Jones矩阵和线偏振转角，说明为何要保留本征轴方向而不能只给光场乘一个共同相位。
\Diagram{fig:o4twohalfq}{两半波片的先后轴方向决定合成旋转}{\draw[axis](-1,0)--(5.5,0);\draw[curve](0,-1)--(0,1);\draw[curve](3,-1)--(3,1);\node[above]at(0,1){$H(\alpha_1)$};\node[above]at(3,1){$H(\alpha_2)$};}
\end{exercise}
\begin{solution}
删去每片的共同相位，图\ref{fig:o4twohalfq}半波片矩阵为
\[
 H(\alpha)=\begin{pmatrix}\cos2\alpha&\sin2\alpha\\
 \sin2\alpha&-\cos2\alpha\end{pmatrix}.
\]
令$c_j=\cos2\alpha_j$、$s_j=\sin2\alpha_j$，按传播顺序相乘得到
\[
 H(\alpha_2)H(\alpha_1)=
 \begin{pmatrix}c_2c_1+s_2s_1&c_2s_1-s_2c_1\\
 s_2c_1-c_2s_1&s_2s_1+c_2c_1\end{pmatrix}
 =R[2(\alpha_2-\alpha_1)].
\]
其中$R(\theta)=\bigl(\begin{smallmatrix}\cos\theta&-\sin\theta\\\sin\theta&\cos\theta\end{smallmatrix}\bigr)$。因此任意线偏振方向转过$2(\alpha_2-\alpha_1)$，椭圆的长轴同样旋转，旋向经两次翻转后恢复。
\Diagram{fig:o4twohalfa}{图取$\alpha_1=20^\circ,\alpha_2=60^\circ$，合成旋转$80^\circ$}{\draw[axis](0,0)--(3,0)node[right]{原方向};\draw[guide](0,0)--(2.53717,.92345)node[right]{$\alpha_1$};\draw[guide](0,0)--(1.5,2.59808)node[above]{$\alpha_2$};\draw[axis](0,0)--(.45149,2.56050)node[above left]{输出};\draw[->](1,0)arc(0:80:1);\node[below]at(1,-.6){$2(\alpha_2-\alpha_1)$};}
图\ref{fig:o4twohalfa}反映两次镜像的几何关系。交换波片顺序会使旋转角反号；快轴平行时合成为单位矩阵。半波片的相对延迟为$\pi$，不是对所有输入都加同一个标量相位；忽略本征轴会丢失整个旋转效应。
\end{solution}

\begin{exercise}\label{ex:o4geometricphase}
干涉仪两臂输入同一正圆偏振态$(1,i)^{\mathsf T}/\sqrt2$。一臂无偏振元件，另一臂依次通过快轴0度和$\alpha$的两块半波片。普通传播相位已补偿，两臂等强。求输出偏振、相对相位和合束强度。
\Diagram{fig:o4geometricq}{两臂最终偏振相同，波片组合仍可改变干涉相位}{\draw[->](0,1)--(5.5,1);\draw[->](0,-.5)--(5.5,-.5);\draw[curve](1.5,-1)--(1.5,0);\draw[curve](3.5,-1)--(3.5,0);\node[below]at(1.5,-1){$H(0)$};\node[below]at(3.5,-1){$H(\alpha)$};\node[right]at(5.5,.3){合束};}
\end{exercise}
\begin{solution}
由上一题，图\ref{fig:o4geometricq}测试臂的矩阵为$R(2\alpha)$。对圆Jones矢量直接作用，
\[
 R(2\alpha)\frac1{\sqrt2}\begin{pmatrix}1\\i\end{pmatrix}
 =\ee^{-2i\alpha}\frac1{\sqrt2}\begin{pmatrix}1\\i\end{pmatrix}.
\]
所以输出仍是同一正圆偏振，公共相位相对参考臂却改变$-2\alpha$。若单臂强度为$I_1$，两臂合场比例$1+\ee^{-2i\alpha}$，强度为$4I_1\cos^2\alpha$。
\Diagram{fig:o4geometrica}{两臂振幅相等，圆态的公共相位可由干涉测出}{\draw[axis](0,0)--(2.5,0)node[right]{参考相量};\draw[axis](0,0)--(1,-2.291288)node[right]{测试相量};\draw[guide](2.5,0)arc(0:-66.42182:2.5);\node at(1.4,-.5){$-2\alpha$};}
图\ref{fig:o4geometrica}的相位随波片转角改变，$\alpha=0$相长，$\alpha=\pi/2$相消。单臂的偏振椭圆保持不变，两臂的干涉强度却随这个公共相位变化。计算中删去的两片材料本身公共传播相位必须确已补偿，本题只隔离随轴角变化的部分。
\end{solution}

\begin{exercise}\label{ex:o4compensator}
完全偏振光Jones矢量为$(1,2\ee^{i\pi/3})^{\mathsf T}$。只允许转动一块四分之一波片，求可使出射成为线偏振的轴方向，并给出线偏振输出与椭圆主轴的关系。
\Diagram{fig:o4compq}{先由椭圆主轴确定补偿波片的本征轴}{\draw[axis](-1.8,0)--(2,0)node[right]{$x$};\draw[axis](0,-2.2)--(0,2.4)node[above]{$y$};\draw[curve,domain=0:360,samples=100]plot({cos(\x)},{2*cos(\x-60)});\node[right]at(1.5,1.7){待补偿};}
\end{exercise}
\begin{solution}
图\ref{fig:o4compq}有$a=1,b=2,\delta=\pi/3$。主轴角满足$\tan2\psi=-2/3$，用$\operatorname{atan2}(2,-3)$选长轴，$\psi\simeq73.155^\circ$。长短半轴平方为$A^2=(5+\sqrt{13})/2$、$B^2=(5-\sqrt{13})/2$。
\Diagram{fig:o4compa}{快慢轴对准长短轴后，90度相差被补成0或180度}{\draw[axis](-2,0)--(2.4,0)node[right]{长轴$\xi$};\draw[axis](0,-1.6)--(0,2)node[above]{短轴$\eta$};\draw[curve](0,0)ellipse(1.8 and .725);\draw[very thick](-1.8,.725)--(1.8,-.725);\draw[guide](-1.8,-.725)--(1.8,.725);}
原态$S_3>0$，在主轴基删公共相位可写$(A,iB)$。若快轴沿长轴、慢轴沿短轴，波片后为$(A,-B)$，输出线方向为$\psi-\arctan(B/A)$；若快慢互换，成为$(iA,iB)$，方向为$\psi+\arctan(B/A)$。图\ref{fig:o4compa}画出两种补偿线。因此波片快轴可沿$73.155^\circ$或其正交方向，输出方向由补偿后的两个非零分量合成，一般偏离原椭圆长轴。
\end{solution}

\begin{exercise}\label{ex:o4ellipsereconstruct}
完全偏振光经旋转线检偏器的最大、最小强度比为9:1，最大方向与$x$成30度。求椭圆长短轴比及长轴方向；仅凭这些数据能否确定旋向？
\Diagram{fig:o4reconstructq}{检偏最大方向与椭圆长轴一致}{\draw[axis](-2,0)--(2.5,0)node[right]{$x$};\draw[curve,rotate=30](0,0)ellipse(1.8 and .6);\draw[guide](-1.8,-1.039)--(1.8,1.039);\node[right]at(1.8,1.039){30度};}
\end{exercise}
\begin{solution}
图\ref{fig:o4reconstructq}完全偏振条件下，沿主轴检偏得到强度分别正比$A^2,B^2$，所以$A/B=3$，长轴为30度模180度。强度比不是半轴比9。
\Diagram{fig:o4reconstructa}{同一椭圆轨迹允许两种相反旋向}{\draw[curve](0,0)ellipse(1.6 and .6);\draw[->](1.6,0)arc(0:60:1.6 and .6);\draw[->](1.6,0)arc(0:-60:1.6 and .6);\node[right]at(2,.3){需圆偏振分析};}
图\ref{fig:o4reconstructa}两种旋向有相同线检偏曲线，需波片测$S_3$符号。若没有“完全偏振”前提，9:1还可能来自自然光与线偏振混合，不能唯一恢复椭圆半轴。
\end{solution}

\begin{exercise}\label{ex:o4halfcircular}
证明理想半波片将任意完全偏振光的椭圆旋向反转，并求其长轴方向变化。输入长轴角$\psi$，快轴角$\alpha$。
\Diagram{fig:o4halfcircq}{半波片在快慢轴上把一个分量反号}{\draw[axis](-2,0)--(2.3,0)node[right]{快轴};\draw[axis](0,-1.7)--(0,2)node[above]{慢轴};\draw[curve,rotate=25](0,0)ellipse(1.6 and .7);}
\end{exercise}
\begin{solution}
图\ref{fig:o4halfcircq}中，在波片快、慢轴基底去掉共同相位，半波片为$H=\operatorname{diag}(1,-1)$。对任意时间的实电场，它使$(E_f(t),E_s(t))$变为$(E_f(t),-E_s(t))$，即关于快轴的镜像。镜像保留长短半轴，长轴相对快轴的夹角由$\psi-\alpha$变为$-(\psi-\alpha)$，所以非圆椭圆的输出方向为
\[
 \psi_o=2\alpha-\psi\pmod\pi.
\]
\Diagram{fig:o4halfcirca}{镜像变换保留椭圆形状，反转沿轨迹的有向旋转}{\draw[axis](-2.4,0)--(2.6,0)node[right]{快轴};\draw[curve,rotate=25](0,0)ellipse(1.8 and .65);\draw[guide,rotate=-25](0,0)ellipse(1.8 and .65);\node[below]at(0,-1.8){$\psi_o=2\alpha-\psi$};}
图\ref{fig:o4halfcirca}中的轨迹形状不变，有向面积却反号。代数上，轴基底的互相关$E_fE_s^*$变为$-E_fE_s^*$，故$S_3$变号；绕传播轴旋转实坐标基底不改变$S_3$，这一结论在实验室基底同样成立。完全偏振光的椭圆率大小保持，旋向反转。

有两个退化情形须分开：线偏振的$S_3=0$，本来没有椭圆旋向，只需使用线方向公式；圆偏振满足长短半轴相同，长轴方向$\psi$没有定义，不能再用长轴角描述输出。对圆Jones矢量直接作用$H$，$(1,i)$变为$(1,-i)$，能明确判断圆旋向反转。公共相位改变时间参数的起点，不改变上述几何结论。
\end{solution}

\begin{exercise}\label{ex:o4twoquarter}
两块相同四分之一波片快轴平行或正交。分别求合成Jones矩阵，并说明两片都存在时偏振是否一定改变。
\Diagram{fig:o4twoquarterq}{两片波片的相对轴角决定相位差相加或相消}{\draw[curve](0,-1)--(0,1);\draw[curve](3,-1)--(3,1);\draw[axis](-1,0)--(4.5,0);\node[above]at(0,1){$Q_1$};\node[above]at(3,1){$Q_2$};}
\end{exercise}
\begin{solution}
按图\ref{fig:o4twoquarterq}，快轴平行时两矩阵均为$\operatorname{diag}(1,i)$，乘积$\operatorname{diag}(1,-1)$，等效半波片。快轴正交时第二片在同一基底为$\operatorname{diag}(i,1)$，乘积$iI$。
\Diagram{fig:o4twoquartera}{正交轴相消后只剩共同相位，所有输入偏振态保持}{\node at(0,.5){平行：$Q^2=H$};\node at(0,-.6){正交：$Q_\perp Q=iI$};\draw[->](2,.5)--(3.3,.5);\draw[->](2,-.6)--(3.3,-.6);\node[right]at(3.5,.5){相差累加};\node[right]at(3.5,-.6){偏振不变};}
图\ref{fig:o4twoquartera}的正交轴组合在所用约化矩阵中为$iI$，所有输入偏振态保持。若恢复每片删去的快轴传播相位$\phi_{f1},\phi_{f2}$，完整矩阵为$\ee^{i(\phi_{f1}+\phi_{f2})}iI$；与参考臂干涉时须保留这整个共同相位，不能仅由$iI$断言实物插入后的总相移恰为$\pi/2$。平行轴组合虽等效半波片，对沿其本征轴的输入仍只改变共同相位。
\end{solution}

\begin{exercise}\label{ex:o4criticaldecay}
全反射入射角为$i=i_c+\epsilon$，$0<\epsilon\ll1$，$i_c=\arcsin(n_2/n_1)$。求倏逝振幅衰减长度随$\epsilon$的领先阶，说明临界角附近的尺度。
\Diagram{fig:o4criticalq}{刚超过临界角，低折射率侧场衰减很慢}{\draw[curve](-2.3,0)--(2.3,0);\draw[guide](0,-1.5)--(0,2);\draw[->](-1.6,-1.3)--(0,0);\draw[->](0,0)--(1.6,-1.3);\node[right]at(1.2,1){长衰减尾部};}
\end{exercise}
\begin{solution}
图\ref{fig:o4criticalq}的$\kappa=k_0\sqrt{n_1^2\sin^2i-n_2^2}$。记$F(i)=n_1^2\sin^2i-n_2^2$，对平方和三角函数逐层求导有$F'(i)=2n_1^2\sin i\cos i$。利用$F(i_c)=0$、$\sin i_c=n_2/n_1$，得到
\begin{align*}
 F(i_c+\epsilon)&=F(i_c)+F'(i_c)\epsilon+O(\epsilon^2)\\
 &=2n_2\sqrt{n_1^2-n_2^2}\,\epsilon+O(\epsilon^2).
\end{align*}
将这一领先项代入$\ell=1/(k_0\sqrt F)$，得
\[
 \ell=\frac1\kappa\simeq
 \frac{\lambda_0}{2\pi\sqrt{2n_2\sqrt{n_1^2-n_2^2}\,\epsilon}}.
\]
\Diagram{fig:o4criticala}{衰减长度在临界角右侧按$\epsilon^{-1/2}$增大}{\draw[axis](0,0)--(4.5,0)node[right]{$\epsilon$};\draw[axis](0,0)--(0,3.5)node[above]{$\ell$};\draw[curve,domain=.2:4,samples=70]plot(\x,{1.4/sqrt(\x)});}
图\ref{fig:o4criticala}的发散说明接近临界角时必须考虑有限光束角谱和第二界面距离。若光束有角度展宽，一部分成分可能低于临界角，单一平面波的无穷衰减长度不再完整描述实际束形。
\end{solution}

\begin{exercise}\label{ex:o4magneticinterface}
正常入射到无损磁性介质界面，介质波阻抗分别为$Z_1=\sqrt{\mu_1/\varepsilon_1}$、$Z_2=\sqrt{\mu_2/\varepsilon_2}$。推导共同切向电场基底的$r,t,R,T$，说明何时可只用折射率。
\Diagram{fig:o4magneticq}{正常入射时电场与磁场的比例由波阻抗决定}{\draw[curve](2,-1.3)--(2,1.3);\draw[->](0,.5)--(1.9,.5);\draw[->](1.9,-.4)--(0,-.4);\draw[->](2.1,.5)--(4.5,.5);\node[above]at(.8,1){$Z_1$};\node[above]at(3.3,1){$Z_2$};}
\end{exercise}
\begin{solution}
图\ref{fig:o4magneticq}中，统一取切向电场沿$x$，入射波沿$+z$、反射波沿$-z$。由$\vct H=\hat{\vct k}\times\vct E/Z$，反射磁场相对其电场多一个方向反号。将入射电场幅度取1，边界条件为
\[
 1+r=t,\qquad \frac{1-r}{Z_1}=\frac t{Z_2}.
\]
用$t=1+r$代入，得到$Z_2-Z_2r=Z_1+Z_1r$，因此$r=(Z_2-Z_1)/(Z_2+Z_1)$、$t=2Z_2/(Z_1+Z_2)$。
\Diagram{fig:o4magnetica}{阻抗匹配可无反射，即使相折射率不同}{\draw[curve](0,-1.5)--(0,1.7);\node[above]at(-1.7,1.5){$Z_1$};\node[above]at(1.8,1.5){$Z_2=Z_1$};\draw[axis](-3,0)--(3.2,0)node[right]{$z$};\draw[->,thick](-2.3,0)--(-2.3,1)node[above]{$E_i$};\draw[->,thick](1.5,0)--(1.5,1)node[above]{$E_t=E_i$};\node[below]at(-1.6,-.6){$E_r=0$};\node[below]at(1.5,-.6){$T=1$};}
无损介质$Z_1,Z_2>0$，平面波平均能流大小为$|E|^2/(2Z)$。图\ref{fig:o4magnetica}中反、入射处于同一介质，透射处于介质2，所以
\[
 R=|r|^2,\qquad T=\frac{Z_1}{Z_2}|t|^2
 =\frac{4Z_1Z_2}{(Z_1+Z_2)^2}.
\]
两者相加的分子为$(Z_2-Z_1)^2+4Z_1Z_2=(Z_1+Z_2)^2$，故$R+T=1$。当$Z_1=Z_2$时无反射，即使两介质的折射率不同也成立，因为阻抗依赖$\mu/\varepsilon$，折射率依赖$\mu\varepsilon$。两介质均非磁性、磁导率相同时，$Z\propto1/n$，才可把阻抗比换成反折射率比，恢复通常正常入射公式。
\end{solution}

\begin{exercise}\label{ex:o4uncertaintystate}
两束互不相干光分别为正、负圆偏振，强度为$3I$和$I$。求总斯托克斯量和偏振度；旋转线检偏器能否观察到强度变化？
\Diagram{fig:o4mixtureq}{两相反圆偏振分量不相干地叠加}{\begin{scope}[xshift=-2cm]\draw[curve](0,0)circle(.8);\draw[->](.8,0)arc(0:70:.8);\node[below]at(0,-1.1){$3I$};\end{scope}\begin{scope}[xshift=2cm]\draw[curve](0,0)circle(.8);\draw[->](.8,0)arc(0:-70:.8);\node[below]at(0,-1.1){$I$};\end{scope}}
\end{exercise}
\begin{solution}
图\ref{fig:o4mixtureq}两圆态各有$S_1=S_2=0$，$S_3$分别等于其强度的正、负值。直接加斯托克斯量，得$S_0=4I,S_1=S_2=0,S_3=2I$，偏振度$1/2$。
\Diagram{fig:o4mixturea}{线检偏透过恒为$2I$，波片分析后才显出圆偏振不平衡}{\draw[axis](0,0)--(4.6,0)node[right]{检偏角};\draw[axis](0,0)--(0,2.8)node[above]{透射强度};\draw[very thick](0,1.7)--(4,1.7);\node[left]at(0,1.7){$2I$};}
线检偏强度恒为$S_0/2=2I$，如图\ref{fig:o4mixturea}。虽然它的旋转曲线与自然光相同，光束仍有一半的圆偏振成分；加四分之一波片后两线方向透过强度成为$3I$与$I$，即可识别。独立光束按斯托克斯量相加；Jones矢量直接相加只适用于有确定相位关系的场。
\end{solution}

\begin{exercise}\label{ex:o4stack}
自然光从空气入射，穿过一叠相互平行的无损玻璃片后回到空气，每个界面均处于相应方向的布儒斯特角。设共经历$M$个界面（$M$为偶数），忽略反射光重新进入透射主束。若每界面$s$偏振反射率为$R_s$，求直接透射主束光强和线偏振度；另说明若在玻璃内部截取主束，光强应怎样换算。
\Diagram{fig:o4stackq}{玻片堆逐渐削弱$s$分量，保留$p$分量}{\foreach\x in{0,1,2,3}{\draw[curve](\x,-.9)--({\x+.5},.9);}\draw[axis](-1,0)--(4.5,0)node[right]{主束};\node[above]at(1.8,1.3){每面为布儒斯特角};}
\end{exercise}
\begin{solution}
图\ref{fig:o4stackq}的自然光初始$p,s$各占$I_0/2$。每界面的法向能流透射率为$T_p=1$、$T_s=1-R_s$，而沿光束方向定义的强度还须换算截面。若第$j$面前、后的光线与共同法线分别成$\theta_{j-1},\theta_j$，有
\[
 I_{\sigma,j}\cos\theta_j
 =T_{\sigma,j}I_{\sigma,j-1}\cos\theta_{j-1},
 \qquad \sigma=p,s.
\]
逐面相乘，所有中间余弦因子抵消。最终又回到空气，$\theta_M=\theta_0$，所以
\[
 I_p=\frac{I_0}{2},\qquad
 I_s=\frac{I_0}{2}(1-R_s)^M.
\]
\Diagram{fig:o4stacka}{两正交强度份额决定偏振度}{\draw[curve](0,0)rectangle(4,.6);\node at(2,.3){$I_p=I_0/2$};\draw[curve](0,-1)rectangle(1.5,-.4);\node[right]at(1.7,-.7){$I_s=(I_0/2)(1-R_s)^M$};}
总强度$I=I_0[1+(1-R_s)^M]/2$，线偏振度$\mathcal P=[1-(1-R_s)^M]/[1+(1-R_s)^M]$。图\ref{fig:o4stacka}在$0<R_s<1$、$M\to\infty$时趋于完全$p$偏振、主束强度$I_0/2$。

若在玻璃内的第$j$面后截取主束，上述两个强度须共同乘$\cos\theta_0/\cos\theta_j$，并将幂次$M$改为$j$；偏振度中共同因子抵消。这样便保留了法向通量与光束强度的区别。若多次反射旁束也进入探测孔径，须重新求和；有吸收或表面误差时，$p$分量也会损失。
\end{solution}

\begin{exercise}\label{ex:o4partialmeasure}
未知光由自然光和某方向线偏振光组成。旋转检偏器得到最大强度8、最小强度2，单位相同。求入射总强度、自然成分和线偏振成分。若去掉“线偏振”这一前提，还能唯一确定偏振度吗？
\Diagram{fig:o4partialq}{检偏器扫描得到非零的最低强度}{\draw[axis](0,0)--(5.2,0)node[right]{角度};\draw[axis](0,0)--(0,3)node[above]{透射强度};\draw[curve,domain=0:180,samples=80]plot({\x/40},{(2+6*cos(\x)^2)/3});\node[left]at(0,2.667){$8$};\node[left]at(0,.667){$2$};}
\end{exercise}
\begin{solution}
按图\ref{fig:o4partialq}，$I_u/2=2$，所以$I_u=4$；$I_p+I_u/2=8$给$I_p=6$。入射总强度$I_0=10$，偏振度$6/10=.6$。
\Diagram{fig:o4partiala}{线检偏扫描测得$S_1,S_2$，不能单独得到$S_3$}{\draw[curve](0,0)rectangle(2.5,1);\node at(1.25,.5){线检偏器};\draw[->](2.6,.5)--(3.5,.5);\node[right]at(3.5,.5){$S_0,S_1,S_2$};\node[below]at(2,-.6){圆偏振信息需波片};}
若完全偏振分量允许为椭圆或圆偏振，图\ref{fig:o4partiala}说明圆偏振信息$S_3$未测得。单次扫描只能给线性调制度.6，即完整偏振度的下限；它可以介于.6与1之间。附加适当四分之一波片测圆偏振分析强度才能完成判定。
\end{solution}

\begin{exercise}\label{ex:o4plateorder}
线偏振沿$x$，先过快轴45度的四分之一波片，再过快轴0度的半波片。求输出Jones矢量，并比较交换两片顺序的结果。
\Diagram{fig:o4orderq}{两块转角不同的波片，矩阵按传播次序从右向左相乘}{\draw[axis](-1,0)--(5.5,0);\draw[curve](0,-1)--(0,1);\draw[curve](2.8,-1)--(2.8,1);\node[above]at(0,1){$Q(45^\circ)$};\node[above]at(2.8,1){$H(0)$};\node[left]at(-1,.5){$x$线偏振};}
\end{exercise}
\begin{solution}
由旋转矩阵，$Q(45^\circ)=\tfrac12\bigl(\begin{smallmatrix}1+i&1-i\\1-i&1+i\end{smallmatrix}\bigr)$，$H(0)=\operatorname{diag}(1,-1)$。图\ref{fig:o4orderq}的输出为$\tfrac12(1+i,-1+i)^{\mathsf T}$，删去共同相位可写$(1,i)^{\mathsf T}/\sqrt2$。
\Diagram{fig:o4ordera}{交换顺序使最终圆偏振的旋向相反}{\begin{scope}[xshift=-2cm]\draw[curve](0,0)circle(.9);\draw[->](.9,0)arc(0:70:.9);\node[below]at(0,-1.2){先$Q$后$H$};\end{scope}\begin{scope}[xshift=2cm]\draw[curve](0,0)circle(.9);\draw[->](.9,0)arc(0:-70:.9);\node[below]at(0,-1.2){先$H$后$Q$};\end{scope}}
交换后$H(0)$先作用于$(1,0)$不改偏振，输出为$\tfrac12(1+i,1-i)^{\mathsf T}$，删共同相位为$(1,-i)^{\mathsf T}/\sqrt2$。图\ref{fig:o4ordera}的两箭头相反；不对易来自两波片的本征轴不同，而不是共同相位差。
\end{solution}

\begin{exercise}\label{ex:o4normalenergy}
空气向折射率$n=1.5$玻璃正常入射。求振幅反射、透射系数以及能量反射、透射率。再解释为何透射场振幅小于入射场并不表示透过能量只有振幅比的平方。
\Diagram{fig:o4normalq}{正常入射界面，采用共同切向电场基底}{\draw[curve](2,-1.4)--(2,1.4);\draw[->](0,.6)--(1.9,.6);\draw[->](1.9,0)--(0,0);\draw[->](2.1,.6)--(4.5,.6);\node[above]at(.8,1){$n_1=1$};\node[above]at(3.3,1){$n_2=1.5$};}
\end{exercise}
\begin{solution}
图\ref{fig:o4normalq}统一切向场有$r=(1-1.5)/(1+1.5)=-.2$、$t=2/(1+1.5)=.8$。所以$R=.04$，$T=(n_2/n_1)t^2=1.5\times.64=.96$。
\Diagram{fig:o4normala}{能量份额为4\%反射、96\%透射}{\draw[curve](0,0)rectangle(5,1);\draw[curve](.2,0)--(.2,1);\node[above]at(.1,1){4\%};\node at(2.6,.5){96\%};\node[below]at(2.5,0){$R+T=1$};}
图\ref{fig:o4normala}满足能量守恒。平面波光强为$n\varepsilon_0cE_0^2/2$，介质2具有不同波导纳，因此$t^2=.64$只是电场振幅平方比。若使用局部$p$基，振幅$r_p=+.2$，能量反射率仍相同，必须同时注明基向量改变。
\end{solution}

\begin{exercise}\label{ex:o4brewsternumerical}
自然光从空气射向$n=1.5$玻璃，入射角为布儒斯特角。求该角、折射角、反射光强相对入射光强的比例及反射偏振方向。
\Diagram{fig:o4brewsterq}{布儒斯特条件使反射与透射光线垂直}{\draw[curve](-2.5,0)--(2.5,0);\draw[guide](0,-1.7)--(0,2);\draw[->](-2,-1.333)--(0,0);\draw[->](0,0)--(2,-1.333);\draw[->](0,0)--(1.2,1.8);\node[right]at(2,-1.333){反射};\node[right]at(1.2,1.8){透射};}
\end{exercise}
\begin{solution}
图\ref{fig:o4brewsterq}满足$\tan i_B=1.5$，所以$i_B\simeq56.31^\circ$、$r_B\simeq33.69^\circ$。$p$反射为零，$s$振幅系数代入$\cos i_B=1/\sqrt{1+n^2}$、$\cos r_B=n/\sqrt{1+n^2}$得$r_s=(1-n^2)/(1+n^2)=-5/13$。
\Diagram{fig:o4brewstera}{自然光的两个偏振份额中仅$s$有反射}{\draw[curve](0,0)rectangle(2.8,.7);\node at(1.4,.35){$I_0/2$的$s$};\draw[->](2.9,.35)--(4,.35);\node[right]at(4,.35){乘$25/169$};\node[below]at(2,-.7){$p$反射为0};}
自然光的$s$份额只有$I_0/2$，故总反射率为$(1/2)(25/169)=25/338\simeq.07396$，如图\ref{fig:o4brewstera}。反射光完全为$s$偏振，电场垂直入射面；不能把$R_s$直接当自然光的总反射率。
\end{solution}

\begin{exercise}\label{ex:o4retardthickness}
某波片在目标波长$\lambda_0$处双折射差为$\Delta n$。求零级四分之一波片厚度及厚度误差$\Delta h$引起的相位误差。若忽略材料色散，偏离波长时会怎样？
\Diagram{fig:o4thicknessq}{波片中快慢分量走过相同厚度却积累不同相位}{\draw[curve](1,-1)rectangle(3,1);\draw[->](0,.4)--(4.5,.4)node[right]{快};\draw[->](0,-.4)--(4.5,-.4)node[right]{慢};\draw[<->](1,-1.4)--(3,-1.4)node[midway,below]{$h$};}
\end{exercise}
\begin{solution}
图\ref{fig:o4thicknessq}中令$\Delta n=n_s-n_f>0$，慢轴相对快轴的相位延迟为$\delta=2\pi\Delta n h/\lambda_0$。零级四分之一波片取最小正延迟$\pi/2$，故$h_0=\lambda_0/(4\Delta n)$。固定波长、材料时厚度误差给
\[
 \Delta\delta=\frac{2\pi\Delta n}{\lambda_0}\Delta h
 =\frac\pi2\frac{\Delta h}{h_0}.
\]
若快慢轴命名相反，延迟符号随之改变，最小厚度应使用$|\Delta n|$。
\Diagram{fig:o4thicknessa}{延迟偏离90度后，45度线偏振不再变成严格圆偏振}{\draw[axis](-1.5,0)--(1.7,0);\draw[axis](0,-1.5)--(0,1.7);\draw[guide](0,0)circle(1);\draw[curve,rotate=45](0,0)ellipse(1.2 and .75);\node[right]at(1.8,.7){圆变椭圆};}
图\ref{fig:o4thicknessa}中，45度线偏振入射后的Jones矢量正比$(1,\ee^{i\delta})$。若$\delta=\pi/2+\epsilon$，其归一化$S_3/S_0=\sin\delta\simeq1-\epsilon^2/2$，$S_2/S_0=\cos\delta\simeq-\epsilon$，不再严格为圆偏振。

固定厚度且忽略色散时，$\delta(\lambda)=(\pi/2)\lambda_0/\lambda$，从$1/(\lambda_0+\Delta\lambda)$一阶展开得$\Delta\delta\simeq-(\pi/2)\Delta\lambda/\lambda_0$。保留材料色散则应微分整个比值：
\begin{align*}
 \frac{\dd\delta}{\dd\lambda}
 &=2\pi h\left[\frac1\lambda\frac{\dd\Delta n}{\dd\lambda}-\frac{\Delta n}{\lambda^2}\right]\\
 &=\delta\left[\frac1{\Delta n}\frac{\dd\Delta n}{\dd\lambda}-\frac1\lambda\right].
\end{align*}
多级波片的设计延迟为$\pi/2+2m\pi$，其相对厚度误差或相对波长偏移对应更大的绝对相位误差。容差由完整延迟的导数确定；先模去$2\pi$会丢失多级波片的灵敏度。
\end{solution}


\end{document}
